\section{Classification: corrections and realization results}

\begin{claim}{D.0}{Refuted}{Desuspension-closed nullity does not characterize a localization}
Even for an accessible reflective class $\mathcal Z$ having all the closure properties of local spaces, desuspension closure of
$\{X:\Map_*(X,Z)\simeq *\text{ for every }Z\in\mathcal Z\}$
does not imply reversibility of its associated localizer. The requested classification needs an additional condition on maps, not only this condition on null objects.
\end{claim}
\begin{proof}
Let $\mathcal S$ be representatives of all countable connected integrally acyclic pointed CW complexes, and let $P$ be nullification at the set of maps $A\to *$, $A\in\mathcal S$. Existence is the usual set-generated nullification theorem, in the precise form of \citet[Theorem 1.1]{Bousfield1996ULP}; alternatively the set of category-of-simplices maps defines the proper localizer in \citet[Theorem 6.1.11]{FGHCisinski2006}. Let $\mathcal Z$ be its full class of local pointed spaces.

Every connected acyclic CW complex is a filtered pointed homotopy colimit of countable acyclic subcomplexes. To see this, start with any countable subcomplex. Its homology has countably many classes; choose in the ambient acyclic complex a finite cellular filling for each class, add the finitely many cells of each filling, and repeat countably. The union is countable and acyclic. Starting with the union of countably many prescribed subcomplexes proves filteredness. Thus $P$ kills every acyclic space. Conversely every nullification generator is an integral homology equivalence, so its generated localizer is contained in $W_{\mathbb Z}$ and it kills only nonempty acyclic spaces. Consequently the null class of $\mathcal Z$ is exactly the integral-acyclic class, which is closed under desuspension by the suspension isomorphism in reduced homology. As a set-generated local-object class, $\mathcal Z$ is reflective and closed under homotopy limits, retracts and cotensors, and under sufficiently filtered homotopy colimits.

Nevertheless its localizer is strictly smaller than $W_{\mathbb Z}$. An explicit detecting local object is $B\Sigma_3$. A map from a connected acyclic space to this $1$-type is classified by a homomorphism from its perfect fundamental group to $\Sigma_3$. Every subgroup of $\Sigma_3$ is solvable, so every such homomorphism is trivial. The based mapping space into a $1$-type from a connected source has contractible components; hence $B\Sigma_3$ is $P$-local.

Let $T$ be the presentation $2$-complex with one vertex, loops $a,b$, and one $2$-cell attached by $aba(bab)^{-1}$. Let $u:T\to S^1$ send both loops once around the circle; its attaching word has total exponent zero, so it extends. The cellular boundary is the column $(1,-1)^t$. Thus $H_1(T;\mathbb Z)=\mathbb Z$, $H_2(T;\mathbb Z)=0$, and $u$ is an integral homology equivalence. Its homotopy cofiber is connected and acyclic. But the homomorphism
\[
 a\longmapsto(12),\qquad b\longmapsto(23)
 \quad\text{from }\pi_1T\text{ to }\Sigma_3
\]
satisfies the braid relation and is nonabelian. It cannot factor through $\pi_1u:\pi_1T\to\mathbb Z$, even up to conjugacy. Hence $\Map(u,B\Sigma_3)$ is not surjective on components and $u$ is not a $P$-equivalence. Its cofiber is $P$-aspherical, so the cofiber criterion fails. This proves the assertion directly; it also exhibits the distinction underlying the proper-aspherical shadow of Cisinski's Theorem 6.1.11. The finite relation and boundary calculations are checked in \texttt{code/CD\_finite\_checks.py}.
\end{proof}

\begin{claim}{D.1.1}{Proved}{The exact single-target test and the incorrect infinite-loop lead}
For a pointed space $Z$, its null class is closed under desuspension if and only if
\[
 \{X:\Map_*(X,\Omega Z)\simeq *\}
 \subseteq
 \{X:\Map_*(X,Z)\simeq *\}.                         \tag{D.1.1}
\]
Equivalently there is no pointed $X$ whose mapping space into $Z$ has weakly contractible constant component and has another component. If $Z$ is in the resolving class generated by $\Omega Z$, (D.1.1) holds. The suggestion that every infinite loop space satisfies it is false.
\end{claim}
\begin{proof}
The adjunction $\Map_*(\Sigma X,Z)\simeq\Map_*(X,\Omega Z)$ proves the equivalence with desuspension closure. The latter mapping space is the loop space at the constant map of $\Map_*(X,Z)$, so weak contractibility means precisely that its constant component is weakly contractible. A pointed mapping-space functor preserves homotopy limits, proving the resolving-class sufficient condition. For the counterexample take $Z=K(\mathbb Z,n)$ and $X=S^n$, $n\geq1$. Then $\Map_*(\Sigma X,Z)$ is weakly contractible, while $\pi_0\Map_*(X,Z)=\mathbb Z$. These $Z$ are infinite loop spaces. The loss of components in passing to a loop space is the exact error in the proposed sufficient condition.
\end{proof}

\begin{claim}{D.1.2}{Refuted}{A single integral or localized sphere does not have the property}
For every $n\geq1$ and every prime $p$, the null classes of the individual targets $S^n$ and $S^n_{(p)}$ fail desuspension closure. This does not contradict the family statement A.1.
\end{claim}
\begin{proof}
For $n=1$, take $X=S^1$ and use $S^1_{(p)}=K(\mathbb Z_{(p)},1)$. Assume $n\geq2$ and put $d=n$ when $n$ is odd and $d=2n-1$ when $n$ is even. Serre's finiteness theorem, opened in \citet[Theorem 1.21]{HatcherSS}, says that $\pi_dS^n$ is $\mathbb Z$ plus a finite group, all groups above degree $d$ are finite, and the only possible additional infinite group is $\pi_nS^n=\mathbb Z$.

For the localized target put
\[
 X=\operatorname{Tel}\bigl(S^{d-1}\xrightarrow{p}S^{d-1}
                         \xrightarrow{p}S^{d-1}\longrightarrow\cdots\bigr).
\]
The mapping space is the homotopy limit of the tower $\Omega^{d-1}S^n_{(p)}$ with bonding maps induced by degree $p$. Its homotopy-group towers are multiplication-by-$p$ towers. On each finite $p$-group these towers are pro-zero, so both limit and derived limit vanish. On $\mathbb Z_{(p)}$, the limit vanishes and
\[
 \lim{}^1\bigl(\mathbb Z_{(p)}\xleftarrow{p}\mathbb Z_{(p)}\xleftarrow{p}\cdots\bigr)
       =\mathbb Z_p/\mathbb Z_{(p)}\neq0.
\]
For clarity, this formula follows from the difference map
$(x_i)\mapsto(x_i-px_{i+1})$ on the product: summing $p^ia_i$ in $\mathbb Z_p$ identifies its cokernel with the displayed quotient; changing a preimage changes the sum by an element of $\mathbb Z_{(p)}$.

The homotopy-limit exact sequence therefore makes every positive homotopy group of $\Map_*(X,S^n_{(p)})$ vanish, while its components contain, and in fact equal, this nonzero derived limit. The possible group in degree $d-1$ has zero ordinary inverse limit and does not change that conclusion. Here $d-1\geq2$, so the mapping tower consists of double-loop spaces and the component calculation is a group calculation. Consequently its loop space $\Map_*(\Sigma X,S^n_{(p)})$ is weakly contractible, but the original mapping space is not.

For the integral target use the telescope whose successive degrees are $2,3,4,\ldots$. Every finite homotopy-group tower is now pro-zero. The corresponding derived limit of the integral tower is $\widehat{\mathbb Z}/\mathbb Z\neq0$, since the products of the bonding degrees are cofinal among the positive integers under divisibility. The identical argument gives the integral counterexample. These are countable CW complexes and actual phantom-component witnesses, not merely formal failures to infer components from loops.
\end{proof}

\begin{claim}{D.1.3}{Proved}{The group targets}
For any group $G$, including every nilpotent group, the null class of $K(G,1)$ is desuspension-closed if and only if $G=1$.
\end{claim}
\begin{proof}
If $G\neq1$, take $X=S^1$. The based mapping space into $K(G,1)$ is the discrete group $G$, up to weak equivalence, whereas the mapping space from $\Sigma S^1=S^2$ is weakly contractible. For $G=1$ the null class consists of all pointed spaces, so the property holds.
\end{proof}

\begin{claim}{D.1.4}{Proved}{The target $BU$}
The null class of $BU$ is closed under desuspension.
\end{claim}
\begin{proof}
Bott periodicity gives $\Omega^2BU\simeq\mathbb Z\times BU$, with the base component identified with $BU$. The opened statement \cite[Theorem 2.11]{HatcherVB} says that external multiplication by the Bott class induces the double-suspension isomorphism in complex $K$-theory. Its representing Bott map is consequently an isomorphism on every homotopy group of the indicated representing spaces, and Whitehead's theorem gives the space equivalence just used. Suppose $\Map_*(\Sigma X,BU)$ is weakly contractible. Its loop space gives $\Map_*(X,\Omega^2BU)\simeq *$. If $X$ had a component other than its base component, the discrete factor $\mathbb Z$ would already give nonconstant based maps to $\mathbb Z\times BU$, a contradiction. Thus $X$ is connected. Every based map from $X$ to the discrete factor is zero, so $\Map_*(X,\Omega^2BU)\simeq\Map_*(X,BU)$. This proves the assertion. The use of periodicity, rather than infinite-loop structure alone, is essential.
\end{proof}

\begin{claim}{D.1.5}{Proved}{The Moore-space classification}
For prime $p$ and $n\geq1$, the single-target null class of the usual Moore space $M(\mathbb Z/p,n)$ is desuspension-closed except when $(p,n)=(2,1)$. The simply connected odd-primary and mod-$2$ cases are proved in D.1.7 and D.1.8, respectively; the fundamental-group cases, including the explicit $\RP^2$ counterexample, are proved in D.1.9. The finite-$H$-space question is completely decided in D.1.6.
\end{claim}
The telescope proof for spheres does not transfer to simply connected finite Moore targets: their homotopy groups are finite $p$-groups, so it produces no nonzero $\lim^1$ component. The affirmative Moore results below instead use explicit loop-space retractions, collection-level finite-wedge constructions, and a convergent simply connected resolution.

\begin{claim}{D.1.6}{Proved}{The complete finite-$H$-space test}
Let $Z$ be an $H$-space having the homotopy type of a finite CW
complex. Its single-target null class
\[
 \{X:\Map_*(X,Z)\simeq *\}
\]
is closed under desuspension if and only if $Z$ is weakly
contractible. Homotopy associativity or commutativity of the
multiplication is not assumed. The assertion includes disconnected
$H$-spaces, with the unit as basepoint.
\end{claim}
\begin{proof}
We first suppose that $Z$ is connected. A connected $H$-space is
simple: its fundamental group is abelian, and the action on its
higher homotopy groups is trivial. Both assertions follow by
comparing concatenation with the multiplication, using the unit.
The finite-type theorem for nilpotent spaces
\cite[Theorem~4.5.2]{I_MayPonto2012} then implies that every
$\pi_i Z$ and every integral homology group of the universal cover
$\widetilde Z$ is finitely generated. The rational $H$-space
theorem says that $Z$ has the rational homotopy type of a finite
product of odd-dimensional spheres
\cite[Theorem~9.1.1 and Corollary~9.1.2]{I_MayPonto2012}.
This statement applies to $H$-spaces without homotopy associativity:
the source explicitly uses quasi-Hopf algebras for that reason.

Suppose first that some rational homotopy group in degree at least
two is nonzero, and let $d$ be its highest degree. Then $d$ is odd
and $d\geq3$, and
\[
 \pi_d Z\cong\ZZ^a\oplus T,\qquad a>0,\quad T\text{ finite}.
\]
All $\pi_jZ$ with $j>d$, as well as $\pi_{d-1}Z$, are finite.
Define the countable pointed CW complex
\[
 X=\operatorname{Tel}\bigl(
 S^{d-1}\xrightarrow{2}S^{d-1}
 \xrightarrow{3}S^{d-1}
 \xrightarrow{4}\cdots\bigr).
\]
Mapping out of this homotopy colimit identifies $\Map_*(X,Z)$
with the homotopy limit of the tower $\Omega^{d-1}Z$ whose maps
on homotopy groups are successive multiplication by $2,3,4,\ldots$.
For a finite group in this tower the composites eventually vanish.
Its inverse limit and its first derived inverse limit are therefore
zero. For the free group $\ZZ^a$, the inverse limit is zero and
\[
 \lim{}^1\bigl(\ZZ^a\xleftarrow{2}\ZZ^a
                 \xleftarrow{3}\ZZ^a
                 \xleftarrow{4}\cdots\bigr)
       \cong(\widehat{\ZZ}/\ZZ)^a\ne0.                 \tag{D.1.6.1}
\]
Here $\widehat{\ZZ}$ is the profinite completion. To see the
formula, write $b_0=1$ and
$b_i=\prod_{j=0}^{i-1}(j+2)$, so the subgroups $b_i\ZZ$ are
cofinal among all finite-index subgroups of $\ZZ$.
For the difference map on the product,
$D(x)_i=x_i-(i+2)x_{i+1}$, the convergent profinite series
$\sum_{i\geq0}b_i y_i$ identifies its cokernel with
$\widehat{\ZZ}/\ZZ$. Telescoping shows that $D(x)$ maps to $x_0$
modulo $\ZZ$; conversely the divisibility of the successive
remainders constructs a preimage whenever the series is integral.
The same calculation in each coordinate gives (D.1.6.1).

The homotopy-limit exact sequences now give
\[
 \pi_j\Map_*(X,Z)=0\quad(j\geq1),\qquad
 \pi_0\Map_*(X,Z)\cong(\widehat{\ZZ}/\ZZ)^a.
\]
In the second formula the ordinary inverse limit of
$\pi_{d-1}Z$ is zero; the finite summand of $\pi_dZ$ contributes
zero derived limit. For $j\geq1$, the derived-limit term comes
from $\pi_{d+j}Z$, hence vanishes, and the ordinary-limit term
vanishes even when it comes from $\pi_dZ$. Since $d-1\geq2$,
these are the usual group-valued loop-space exact sequences.
Consequently $\Map_*(\Sigma X,Z)\simeq
\Omega\Map_*(X,Z)$ is weakly contractible, while
$\Map_*(X,Z)$ is not. This is the required counterexample in
the first case.

It remains to handle the possibility that the only rational
homotopy occurs in degree one. The universal cover $\widetilde Z$
inherits an $H$-space multiplication by lifting the multiplication
of $Z$; the unit homotopies lift as well. Its homotopy groups are
torsion, so it is rationally acyclic. Moreover its lifted CW
structure has the same finite dimension as a finite CW model of
$Z$. Combining this bound with the finite generation already
proved, its integral homology is finitely generated in total.

We use exactly Browder's Corollary~7.2
\cite[p.~47]{I_Browder1961}: for a connected $H$-space with each
integral homology group finitely generated and homology zero
above degree $m$, if $H_m$ is nonzero then $H_m\cong\ZZ$.
Its hypotheses require bounded finitely generated homology,
not a finite CW complex or a finite-sheeted covering.
If $\widetilde Z$ had any nonzero positive homology, its highest
such group would therefore contain nonzero rational homology.
This contradicts rational acyclicity. Thus $\widetilde Z$ is
integrally acyclic and simply connected, and is contractible by
the Hurewicz and Whitehead theorems.

It follows that $Z\simeq K(G,1)$ for the finitely generated
abelian group $G=\pi_1Z$. This group is torsion-free. Indeed,
if it contained a subgroup $C_p$, the corresponding covering
of the finite-dimensional CW model of $Z$ would be a
finite-dimensional $K(C_p,1)$. This contradicts
$H^j(C_p;\FF_p)\ne0$ in every degree $j\geq0$, obtained directly
from the two-periodic cyclic resolution: its two differentials
are zero with trivial $\FF_p$ coefficients. Hence
$G\cong\ZZ^r$ and $Z$ is a torus up to homotopy.
If $r>0$, the space $X=S^1$ has
$\Map_*(\Sigma X,Z)\simeq *$, whereas
$\Map_*(X,Z)\simeq G$ is a nontrivial discrete space.
If $r=0$, then $Z$ is contractible. This finishes the connected
case.

Finally let $Z$ have several components and let $Z_e$ be its
identity component. This is a finite connected $H$-space.
If it is noncontractible, the preceding counterexamples have
connected source $X$, so both mapping spaces into $Z$ are the
same as the corresponding based mapping spaces into $Z_e$.
If $Z_e$ is contractible but $Z$ is disconnected, take $X=S^0$:
then $\Map_*(\Sigma X,Z)=\Omega Z_e\simeq *$ and
$\Map_*(X,Z)=Z$ is not weakly contractible. In the remaining
case $Z$ itself is weakly contractible and every pointed mapping
space into it is weakly contractible, so the desired closure
holds.
\end{proof}

\begin{claim}{D.1.7}{Proved}{The simply connected odd-primary Moore targets}
Let $p$ be odd and let $n\geq2$. For the usual Moore space
$M(\ZZ/p,n)=S^n\cup_p e^{n+1}$, its single-target null class is
closed under desuspension. There is no finiteness hypothesis on the
source. More precisely, for any $d\geq3$ and any pointed space $X$,
\[
 \Map_*(X,\Omega P^d(p))\simeq *
 \quad\Longrightarrow\quad
 \Map_*(X,P^e(p))\simeq *\quad\hbox{for every }e\geq3,
 \tag{D.1.7.1}
\]
where $P^e(p)=S^{e-1}\cup_p e^e$ is indexed by its top cell.
\end{claim}
\begin{proof}
Fix $X$ and write
\[
 \mathcal R_X=\{Y:\Map_*(X,Y)\simeq *\}.
\]
This class is closed under homotopy limits, loops and homotopy
retracts. The proof has three steps: obtain all sufficiently high
looped Moore targets; use a convergent loop--suspension resolution
to remove their loops; and use finite cone cubes to lower their
dimensions. All the closure statements concern the whole based
mapping space, including its set of components.

\medskip\noindent
\textbf{The precise retraction input and the bottom-dimensional case.}
For odd $p$ the smash splitting is
\[
 P^a(p)\wedge P^b(p)
       \simeq P^{a+b}(p)\vee P^{a+b-1}(p),\qquad a,b\geq2.
 \tag{D.1.7.2}
\]
It is the splitting used in the proof of
\cite[Lemma~1.5, pp.~13--14]{I_CMN1987}; the hypotheses and formula
are also explicitly recorded in \cite[Lemma~4.5]{I_Boyde2021}.
We use two retraction theorems from the original CMN paper:
\[
 \Omega P^{2a+2}(p)
 \simeq S^{2a+1}\{p\}\times
       \Omega\bigvee_{j\geq0}P^{4a+2ja+3}(p),
       \qquad a\geq1,                              \tag{D.1.7.3}
\]
and the assertion that a mod-$p$ homology monomorphism between
loop spaces of bouquets of simply connected mod-$p$ Moore spaces
has a left homotopy inverse. These are respectively
\cite[Proposition~1.1, p.~9, and Proposition~1.7, p.~15]{I_CMN1987}.
The symbol $S^{2a+1}\{p\}$ denotes the fibre of the degree-$p$
map on the $p$-local sphere; only the Moore-bouquet factor will be
used. CMN's retractions are $p$-local statements. Every simply
connected finite mod-$p$ Moore space, and every finite wedge or
smash of these spaces, has finite $p$-primary homotopy groups by
the Serre-class theorem. These spaces are therefore already
$p$-local. The retractions used here consequently apply to their
ordinary homotopy types. The displayed range $a\geq1$ was checked
on the printed page of Proposition~1.1, including its endpoint.

For clarity we do not use CMN's odd-dimensional Corollary~1.9
outside its printed range $a\geq2$. In particular the case
$P^3(p)$ is established as follows. Let $d\geq3$ be odd, set
$A=P^{d-1}(p)$, and let
\[
 [\id,\id]:\Sigma(A\wedge A)\longrightarrow\Sigma A=P^d(p)
\]
be the generalized Whitehead product. The lower Moore summand in
(D.1.7.2), before suspension, gives a map
$j:P^{2d-3}(p)\to A\wedge A$. Put
\[
 h=[\id,\id]\circ\Sigma j:P^{2d-2}(p)\longrightarrow P^d(p).
\]
We check that $\Omega h$ is a mod-$p$ homology monomorphism.
Write
\[
 H_*(\Omega P^d(p);\FF_p)=T(u,v),\qquad
 |u|=d-2,\quad |v|=d-1,\quad\beta v=u.
\]
This tensor-algebra formula is the James theorem for the connected
space $A$. The adjoint Whitehead product has the graded commutator
as its homology map. Since $|u|$ is odd and $|v|$ is even, the
two generators of $H_*(\Omega P^{2d-2}(p);\FF_p)$ have images,
after choosing the compatible Moore orientations,
\[
                  2u^2,\qquad vu-uv.                \tag{D.1.7.4}
\]
Here is a check that the second formula is independent of the
chosen complementary summand. Normalize the lower generator of
$j$ to $u\otimes u$. Its upper generator then maps to
$c(v\otimes u)+b(u\otimes v)$ with $c-b=1$, since
$\beta(v\otimes u)=u\otimes u$ and
$\beta(u\otimes v)=-u\otimes u$. The commutator sends this
combination to $(c-b)(vu-uv)$. It sends $u\otimes u$ to $2u^2$.
Thus the two images also satisfy the required Bockstein relation.

These images freely generate a tensor subalgebra. Indeed, with
the lexicographic order $v>u$, their leading words are $uu$ and
$vu$, respectively, with nonzero coefficients. Both words have
length two. A concatenation of them has a unique division into
two-letter blocks, so distinct monomials in the two proposed
generators have distinct leading words. This proves injectivity.
CMN Proposition~1.7 therefore makes $\Omega P^{2d-2}(p)$ a
homotopy retract of $\Omega P^d(p)$. Its source and target are
loop spaces of simply connected Moore spaces even when $d=3$;
the proposition has no additional connectivity restriction here.
Consequently an even-dimensional looped Moore space belongs to
$\mathcal R_X$ whenever the original $\Omega P^d(p)$ does.

\medskip\noindent
\textbf{All sufficiently high looped Moore spaces.}
Apply (D.1.7.3) to that even-dimensional Moore space, and retain
two bouquet summands. For some integers $a_1,a_2\geq3$ we obtain
\[
 \Omega(P^{a_1}(p)\vee P^{a_2}(p))\in\mathcal R_X.
 \tag{D.1.7.5}
\]
For example the first two summands of (D.1.7.3) have top
dimensions $4a+3$ and $6a+3$.
The Hilton--Milnor factor indexed by the basic commutator with
$k$ occurrences of the first letter and one of the second is
\[
 \Omega\Sigma\bigl(
      P^{a_1-1}(p)^{\wedge k}\wedge P^{a_2-1}(p)\bigr),
      \qquad k\geq1.
\]
It is a homotopy retract of (D.1.7.5). Repeated use of
(D.1.7.2) makes every $\Omega P^D(p)$ with
\[
        k(a_1-2)+a_2\ \leq D\leq\ k(a_1-1)+a_2
 \tag{D.1.7.6}
\]
a retract of this factor. Each integer in the interval occurs:
the smash of $k+1$ Moore spaces has a Moore summand with each of
the $k+1$ possible dimension reductions $0,1,\ldots,k$.
The lower endpoint for $k+1$ minus the upper endpoint for $k$
is $a_1-2-k$. Thus the integer intervals overlap or are adjacent
for all sufficiently large $k$, and their union contains every
integer above some bound $N$. We have proved
\[
                  \Omega P^D(p)\in\mathcal R_X
                  \quad(D\geq N).                  \tag{D.1.7.7}
\]
The Hilton--Milnor theorem and its Moore-smash application used
here are also given in \cite[Proposition~2.2 and its proof,
pp.~18--19]{I_CMN1987}; the interval calculation is the extra
step which gives every sufficiently large dimension.

We will need a slightly stronger consequence. If $W$ is a finite
wedge of Moore spaces of top dimension at least $N$, every
Hilton--Milnor factor of $\Omega W$ splits into looped Moore
spaces of top dimension at least $N$. A word of length $q$ has
minimum possible top dimension
$\sum_{i=1}^q e_i-2q+2\geq q(N-2)+2\geq N$.
The factors' connectivities increase with word length, so the
weak product supplied by Hilton--Milnor is weakly equivalent to
the ordinary product. Products and retracts are allowed in
$\mathcal R_X$, proving $\Omega W\in\mathcal R_X$.
The same conclusion holds for a finite-type wedge: truncate it
by top dimension, use the projections that collapse the tails,
and take the homotopy limit. In each fixed homotopy degree this
tower is eventually constant, so its homotopy limit is the
original wedge. This argument uses target connectivity, not
compactness of $X$.

\medskip\noindent
\textbf{Removing loops and lowering dimensions.}
Let $K$ be any finite wedge of Moore spaces of top dimension at
least $N$, and put $T=\Omega\Sigma$. The iterated James splitting
from C.9 gives
\[
 \Sigma T^rK\simeq
       \bigvee_{m\geq1}\bigvee^{c_{r,m}}\Sigma K^{\wedge m},
       \qquad r\geq0,                              \tag{D.1.7.8}
\]
where each multiplicity $c_{r,m}$ is finite. To recall why,
$\Sigma TV\simeq\bigvee_{q\geq1}\Sigma V^{\wedge q}$;
distributing smash products and moving the single suspension
between factors gives the induction. At a fixed weight $m$
there are only finitely many ordered decompositions of $m$ into
positive weights, so all the multiplicities remain finite.
By (D.1.7.2), the right side of (D.1.7.8) is a finite-type wedge
of Moore spaces of top dimension at least $N+1$. Its loop space
therefore belongs to $\mathcal R_X$ by (D.1.7.7) and the preceding
paragraph. Thus $T^{r+1}K\in\mathcal R_X$ for every $r\geq0$.

For every simply connected pointed space the derived
loop--suspension resolution converges:
\[
             K\simeq\holim_{[q]\in\Delta}T^{q+1}K.
 \tag{D.1.7.9}
\]
This is exactly \cite[Theorem~1.8(b)]{FGHBlomquistHarper2017},
with $r=1$ in that paper, or
\cite[Sections~7.9 and~9.9]{FGHBousfield2003}. There is no
finite-type hypothesis on the mapping source. Closure under
homotopy limits proves that $\mathcal R_X$ contains every such
finite wedge $K$.

Finally we spell out the desuspension mechanism, so that this
step does not assume the property being proved. If $V$ is simply
connected and every finite wedge of $\Sigma V$ lies in a
homotopy-limit-closed class $\mathcal R$, form the strongly
homotopy cocartesian $s$-cube generated by $s$ copies of
$V\to CV$. Each noninitial vertex is a finite wedge of
$\Sigma V$, including the empty wedge. Let $Y_s$ be the homotopy
limit of the punctured cube. Then $Y_s\in\mathcal R$.
The higher Blakers--Massey theorem makes the comparison
$V\to Y_s$ increasingly connected as $s$ increases. The cone
cubes have compatible projection maps, and hence
$V\simeq\holim_sY_s\in\mathcal R$. This is the Barratt
argument in \cite[Theorem~8, p.~7]{I_Strom2012}.

Apply this argument first to every finite wedge of $P^{N-1}(p)$,
whose suspension is a finite wedge of $P^N(p)$. It puts all
these wedges in $\mathcal R_X$. Repeat, one dimension at a time,
down to top dimension three. Together with the already proved
high-dimensional assertion, this proves (D.1.7.1). Taking
$e=d=n+1$ and using the suspension--loop adjunction proves the
claim.
\end{proof}

The odd-prime hypothesis is used both in (D.1.7.2) and in the
nonzero coefficient $2$ in (D.1.7.4). Neither step gives the
mod-$2$ Moore assertion by substitution.

\begin{claim}{D.1.8}{Proved}{The simply connected mod-two Moore targets}
For every $n\geq2$, the single-target null class of
$M(\ZZ/2,n)$ is closed under desuspension, with no restriction on
the source space. More strongly, nullity into any
$\Omega P^d(2)$, $d\geq3$, implies nullity into every
$P^e(2)$, $e\geq3$. The additional geometric input is the explicit
homotopy retraction
\[
 \Omega\bigl(P^{11d-15}(2)\vee P^{11d-15}(2)\bigr)
       \quad\hbox{of}\quad \Omega P^d(2),\qquad d\geq4.
 \tag{D.1.8.1}
\]
For the bottom case one can first use the retraction of
$\Omega P^6(2)$ from $\Omega P^3(2)$.
\end{claim}
\begin{proof}
We give the two-summand construction, including a reproducible
finite certificate for its only larger calculation. It does not
use an odd-primary splitting of Moore smash products.

\medskip\noindent
\textbf{The geometric maps.}
Fix $d\geq4$ and write
\[
 A=P^{d-1}(2),\qquad |u|=d-2,\quad |v|=d-1,\qquad
 s=|u|+|v|=2d-3,
\]
where $u,v$ generate $\widetilde H_*(A;\FF_2)$.
The space $A$ is a simply connected finite suspension and is
already $2$-local. For odd $m$, let
$\widetilde L_m(A)$ be the telescope of the
Dynkin--Specht--Wever map $\frac1m\beta_m$ on
$\Sigma A^{\wedge m}$. There are maps
\[
 p_m:\Sigma A^{\wedge m}\longrightarrow\widetilde L_m(A),
 \qquad
 i_m:\widetilde L_m(A)\longrightarrow\Sigma A^{\wedge m}
\]
such that $p_mi_m\simeq\id$ and, on mod-$2$ homology,
\[
 (i_mp_m)_*(x_1\cdots x_m)
       =\beta_m(x_1\cdots x_m)
       =[[\cdots[x_1,x_2],x_3],\ldots,x_m].
 \tag{D.1.8.2}
\]
Suspension symbols on homology are suppressed here and below.
This splitting is constructed by the two complementary
telescopes and suspension comultiplication in
\cite[Section~2, pp.~3248--3249]{I_ChenWu2013}.
It does not require the representing map to be a co-H map.
Since $A$ is a suspension, \cite[Theorem~1.6]{I_Wu1998}
makes $\Omega\widetilde L_{11}(A)$ a homotopy retract of
$\Omega\Sigma A$. The sequence in that theorem can be the odd
primes; its conditions exclude divisibility by $2$ and divisibility
of one chosen index by another. Thus the index $11$ is permitted.

There is a map
\[
       \alpha:S^s\longrightarrow A\wedge A,
       \qquad \alpha_*(\iota_s)=uv+vu,
 \tag{D.1.8.3}
\]
by \cite[Proposition~2.1, pp.~3249--3250]{I_ChenWu2013}.
The map
$e_3=p_3\Sigma(\id_A\wedge\alpha):
\Sigma^{s+1}A\to\widetilde L_3(A)$
is a homotopy equivalence. To verify it in this case, the image
of $\beta_3$ on $\widetilde H_*(\Sigma A^{\wedge3};\FF_2)$
has the two basis elements $uuv+vuu$ and $uvv+vvu$;
$e_3$ takes the two generators to these elements. Thus it is a
mod-$2$ homology equivalence between simply connected spaces
whose integral homology is finite $2$-primary torsion in bounded
degrees. It is an integral homology equivalence and hence a
homotopy equivalence. Equivalently this is the first step of
\cite[proof of Proposition~3.1, p.~3251]{I_ChenWu2013}.
Define $r_3=e_3^{-1}p_3$. Its homology map has the complete table
\[
\begin{array}{c|cccccccc}
 w&uuu&uuv&uvu&uvv&vuu&vuv&vvu&vvv\\ \hline
 (r_3)_*(w)&0&0&u&v&u&v&0&0.
\end{array}                                      \tag{D.1.8.4}
\]

Iterate this actual projection, moving suspension factors through
smash products, to obtain
\[
 \varphi_k:\Sigma A^{\wedge(2k+1)}
                     \longrightarrow\Sigma^{ks+1}A,
 \qquad
 \varphi_1=r_3,\qquad
 \varphi_k=(\Sigma^{(k-1)s}r_3)
                (\varphi_{k-1}\wedge\id_A\wedge\id_A).
 \tag{D.1.8.5}
\]
Its homology map applies the table (D.1.8.4) successively to the
first three letters, to the result and the next two letters, and
so on. This also describes what happens when an intermediate
value is zero: the final value is zero.

Set $B=\Sigma^{5s+1}A=P^{11d-15}(2)$ and
$L=\widetilde L_{11}(A)$. Let $\sigma=(1\ 8)$ and
$\tau=(1\ 3)$ act on the eleven smash factors. With
$j_0=\Sigma(\id_A\wedge\alpha^{\wedge5}):
B\to\Sigma A^{\wedge11}$, define
\[
 j_1=p_{11}j_0,\quad j_2=p_{11}\sigma j_0,
 \qquad
 r_1=\varphi_5i_{11},\quad r_2=\varphi_5\tau i_{11}.
 \tag{D.1.8.6}
\]
All four maps have now been specified geometrically.

\medskip\noindent
\textbf{The exact finite calculation.}
For either generator of $\widetilde H_*(B;\FF_2)$ the matrix
of the four composites $r_i j_j$ is
\[
                     \begin{pmatrix}1&1\\1&0\end{pmatrix}.
 \tag{D.1.8.7}
\]
Here are explicit finite instructions and intermediate integer
counts that verify the matrix. They also fix all permutation and
commutator conventions, so that no representation-theoretic
decomposition is being assumed. Write $u=0$, $v=1$ and put
\[
 w(\epsilon,x)=
 (x,\epsilon_1,1-\epsilon_1,\ldots,
       \epsilon_5,1-\epsilon_5),
 \qquad\epsilon\in\{0,1\}^5,\quad x\in\{0,1\}.
\]
For a word $z=(z_1,\ldots,z_{11})$ and
$S\subseteq\{2,\ldots,11\}$, let $D_S(z)$ be the word obtained
by listing the letters with indices in $S$ in decreasing order,
then $z_1$, then the remaining letters in increasing order.
Expansion of the successive commutators gives the exact identity
\[
               \beta_{11}(z)=\sum_{S\subseteq\{2,\ldots,11\}}
                                  D_S(z)\quad\hbox{over }\FF_2.
 \tag{D.1.8.8}
\]
Indeed at each commutator step the new letter is put either on
the left or on the right; the former choices accumulate in
reverse order and the latter in their original order.

Let $f(z)$ be the repeated ternary reduction in (D.1.8.4), with
value $\bot$ if a zero is encountered. For
$(\rho_1,\rho_2)=(\id,\tau)$ and
$(\gamma_1,\gamma_2)=(\id,\sigma)$ define the integer counts
\[
 C_{ij}(x)=
 \sum_{\epsilon\in\{0,1\}^5}
 \sum_{S\subseteq\{2,\ldots,11\}}
  \mathbf1\{f(\rho_iD_S(\gamma_jw(\epsilon,x)))=x\}.
\]
Applying the eight-entry table five times gives
\[
             C(0)=C(1)=
                  \begin{pmatrix}3125&3125\\3039&2460\end{pmatrix}.
 \tag{D.1.8.9}
\]
Reduction modulo two is (D.1.8.7). Each entry has exactly
$2^5\cdot2^{10}$ candidate terms, so this is a bounded exhaustive
calculation. A second calculation cancels equal words after
each commutator in (D.1.8.2) before doing any projections; the
two embedding supports then have sizes $62$ and $126$, and
give the same matrix. The complete short verifier
\texttt{code/I\_lie\_moore.py} implements both formulas
independently, compares them on every input word, and checks
the displayed integer and mod-$2$ matrices. Its exact output
is \texttt{code/outputs/I\_lie\_moore.json}.

\medskip\noindent
\textbf{From the matrix to the wedge retraction.}
Being a retract of a suspension, $L$ has a comultiplication
\[
 \Delta_L=(p_{11}\vee p_{11})
                \Delta_{\Sigma A^{\wedge11}}i_{11}:L\to L\vee L
\]
whose two projections are homotopic to the identity. Thus the
two projections in (D.1.8.6) give one actual map to a wedge:
\[
 R=(r_1\vee r_2)\Delta_L:L\longrightarrow B\vee B,
 \qquad J=j_1\vee j_2:B\vee B\longrightarrow L.
\]
By (D.1.8.7), $RJ$ is a mod-$2$ homology equivalence. It is a
homotopy equivalence: its mapping cone is simply connected,
has finite $2$-primary integral homology, and has zero mod-$2$
homology, hence zero integral homology. Replacing $R$ by
$(RJ)^{-1}R$ proves that $B\vee B$ is a homotopy retract of
$L$. Looping, and using the Wu retraction above, proves
(D.1.8.1).

For $d=3$ use the explicit Moore-space retraction
\[
                  \Omega P^6(2)\quad\hbox{of}\quad\Omega P^3(2).
 \tag{D.1.8.10}
\]
It is the case $n=2$, $k=1$ of the formula
$\Omega P^{n+1+k(2n-1)}(2)\to\Omega P^{n+1}(2)$ stated
in \cite[proof of Corollary~1.4, p.~3258]{I_ChenWu2013},
where the printed range is $n\geq2$, $k\geq1$.
We only need this endpoint, not the broader prime-indexed product
formula in that paper's Theorem~1.3. Applying (D.1.8.1) with
$d=6$ then supplies the required two equal summands. This use of
the paper's explicit $P^3(2)$ application does not identify an
odd-degree telescope of $\RP^2$ with a CW complex having finitely
many ordinary cells in each degree; that identification would
be false, since the telescope has $\pi_2=\ZZ_{(2)}$.

\medskip\noindent
\textbf{Why two equal summands suffice.}
Suppose now $\Omega P^d(2)\in\mathcal R_X$, where $d\geq3$.
We have just obtained an integer $D>d$ with
\[
       \Omega(P^D(2)\vee P^D(2))\in\mathcal R_X.
\]
Put $K=P^{D-1}(2)$. The Hilton--Milnor word containing $q-1$
copies of the first letter and one copy of the second makes
$\Omega\Sigma K^{\wedge q}$ a retract for every $q\geq2$;
the case $q=1$ is either original wedge summand. Consequently
all these loop spaces belong to $\mathcal R_X$.

If $V$ is any finite wedge of copies of $K$, the same theorem
shows that the loops on every finite wedge of spaces
$\Sigma V^{\wedge m}$ belong to $\mathcal R_X$: all its factors
have the form $\Omega\Sigma K^{\wedge q}$. Factor connectivities
increase with word length, so the weak product is weakly
equivalent to the product. The iterated James splitting
(D.1.7.8), now used without any Moore smash splitting, expresses
$\Sigma(\Omega\Sigma)^rV$ as a finite-type wedge of precisely
these smash powers. Its truncation tower has increasingly
connected tails. Therefore all
$(\Omega\Sigma)^{r+1}V$ belong to $\mathcal R_X$.
The convergent simply connected resolution (D.1.7.9) proves
$V\in\mathcal R_X$.

We have proved nullity for every finite wedge of $P^{D-1}(2)$.
The finite cone-cube argument at the end of D.1.7 lowers this
whole family by one dimension at a time, down to top dimension
three. In particular $P^d(2)\in\mathcal R_X$, as $3\leq d<D$.
Moreover $P^{D-1}(2)$ and its loop space now belong to
$\mathcal R_X$. Repeating the same construction with $d=D-1$
increases the upper bound to $11(D-1)-15$, and iteration makes
these bounds unbounded. Thus every $P^e(2)$ with $e\geq3$
belongs to $\mathcal R_X$, proving the stronger assertion.
Adjunction proves the requested desuspension closure. The only
homotopy limits taken in this proof are limits in the targets;
the source $X$ has been arbitrary throughout.
\end{proof}

\begin{claim}{D.1.9}{Proved}{The Moore spaces with cyclic fundamental group}
For odd $p$, the single-target null class of
$M(\ZZ/p,1)=P^2(p)$ is closed under desuspension, for arbitrary
source spaces. For $p=2$ this assertion is false:
$P^2(2)=\RP^2$ admits an explicit simply connected
factorial-telescope counterexample. Consequently, among the
usual Moore targets $M(\ZZ/p,n)$ with $p$ prime and $n\geq1$,
the only failure is $M(\ZZ/2,1)$.
\end{claim}
\begin{proof}
Let $p$ be odd, put $Y=P^2(p)$, and suppose
$\Map_*(X,\Omega Y)\simeq *$. We may replace $X$ by a CW
model. This hypothesis forces $X$ to be connected: otherwise
$S^0$ is a pointed retract of $X$, and $\Omega Y$ would be a
retract of the contractible mapping space. That is impossible,
since $\pi_0\Omega Y=\ZZ/p$.

The universal cover $V$ of $Y$ satisfies
\[
                         V\simeq\bigvee^{p-1}S^2.
 \tag{D.1.9.1}
\]
Indeed $Y$ is the presentation complex of $\langle a\mid a^p\rangle$,
and its $p$-sheeted universal cover is a simply connected finite
two-dimensional complex of Euler characteristic $p$. Its only
nonzero reduced integral homology is free of rank $p-1$ in degree
two. Representatives of a homology basis under the Hurewicz
isomorphism give a homology equivalence from the displayed wedge;
the simply connected Whitehead theorem makes it a homotopy
equivalence.

Since $X$ is connected, its based maps to $\Omega Y$ land in
the base component $\Omega V$. Thus $\Omega V$ belongs to
$\mathcal R_X=\{Z:\Map_*(X,Z)\simeq *\}$.
There are at least two sphere summands in (D.1.9.1).
Hilton--Milnor on these two summands gives a retract
$\Omega S^{q+1}$ for every $q\geq1$: use a word of length
$q$ with $q-1$ occurrences of one letter and one of the other
when $q\geq2$, and a sphere summand when $q=1$.
Hence
\[
                      \Omega S^d\in\mathcal R_X
                      \qquad(d\geq2).              \tag{D.1.9.2}
\]

Here is the unlooping and finite-target step with its source
hypotheses made explicit. If $W$ is a finite wedge of spheres
of dimension at least two, iterated James splitting expresses
$\Sigma(\Omega\Sigma)^rW$ as a finite-type wedge of spheres.
Hilton--Milnor and (D.1.9.2) put the loops on every finite
truncation in $\mathcal R_X$. Increasing connectivity of the
tails and closure under homotopy limits put
$(\Omega\Sigma)^{r+1}W$ in $\mathcal R_X$ for every $r\geq0$.
The simply connected resolution (D.1.7.9) now gives
$W\in\mathcal R_X$. In particular every $S^d$, $d\geq2$,
belongs to this class.

The class $\mathcal R_X$ is a resolving kernel and is closed
under extensions by fibrations. We can therefore apply
\cite[Corollary~12(a), pp.~8--9]{I_Strom2012}, which puts every
simply connected finite complex in $\mathcal R_X$ as soon as
the sphere dimensions include the indicated infinite arithmetic
progression. Its proof uses the collection-level wedge theorem,
the cone-length theorem, and finite cone cubes. It has no
finite-type hypothesis on $X$; the finite-type source hypothesis
in that paper's Theorem~1 enters only for the stronger assertion
about arbitrary finite-dimensional targets. We only need the
two finite targets
\[
                           V,\ \Sigma Y\in\mathcal R_X.
 \tag{D.1.9.3}
\]
In particular $\Omega\Sigma Y\in\mathcal R_X$.

Consider any based map $f:X\to Y$ and the unit
$\eta_Y:Y\to\Omega\Sigma Y$. The composite $\eta_Yf$ is
nullhomotopic by (D.1.9.3). On fundamental groups the unit is
the abelianization map
\[
 \pi_1Y\longrightarrow H_1(Y;\ZZ)
     \cong H_2(\Sigma Y;\ZZ)
     \cong\pi_2\Sigma Y=\pi_1\Omega\Sigma Y.
\]
For $Y=P^2(p)$ it is an isomorphism $\ZZ/p\to\ZZ/p$.
Thus $\pi_1f=0$, so $f$ lifts to the universal cover $V$.
The lift is nullhomotopic by (D.1.9.3); hence so is $f$.
We have proved that $\Map_*(X,Y)$ has only its constant
component. Its loop space is the given contractible space
$\Map_*(X,\Omega Y)$, so all its positive homotopy groups
also vanish. Therefore $Y\in\mathcal R_X$, as required.
Notice that no loop--suspension convergence statement has been
applied to the nonsimply connected space $Y$ itself.

For $p=2$, put
\[
 X=\operatorname{Tel}\bigl(S^2\xrightarrow{2}S^2
              \xrightarrow{3}S^2\xrightarrow{4}\cdots\bigr).
\]
This is simply connected. The integral $S^2$ calculation in
D.1.2 gives
\[
 \Map_*(X,\Omega S^2)\simeq *,\qquad
 \pi_0\Map_*(X,S^2)
       =\lim{}^1\bigl(\ZZ\xleftarrow{2}\ZZ
                      \xleftarrow{3}\ZZ\xleftarrow{4}\cdots\bigr)
       =\widehat\ZZ/\ZZ\ne0.
\]
Every based map from $X$, or from a suspension of $X$, to
$\RP^2$ lifts uniquely through its universal cover $S^2$.
The corresponding based mapping-space comparison is a weak
equivalence. Thus
$\Map_*(X,\Omega\RP^2)\simeq *$ whereas
$\Map_*(X,\RP^2)\not\simeq *$. This proves the failure and,
together with D.1.7--D.1.8, the complete stated classification.
\end{proof}


\begin{claim}{D.2.1}{Proved}{Every finite thick subcategory occurs}
Let $\mathcal C\subseteq\Ho(\Sp)^\omega$ be a thick subcategory of finite spectra and set $\mathcal L=\loc(\mathcal C)$. Then
\[
 T R(\mathcal L)=\mathcal L,
 \qquad T R(\mathcal L)\cap\Ho(\Sp)^\omega=\mathcal C.
\]
Thus every Hopkins--Smith type function occurs for a reversible fundamental localizer. The classification input meant here is exactly \citet[Theorem 7]{HopkinsSmith1998}: a thick subcategory of finite $p$-local spectra is one of the chromatic subcategories $\mathcal C_n$, including zero. Our realization argument does not require that classification theorem.
\end{claim}
\begin{proof}
Every finite spectrum is, after a sufficiently large positive suspension, the suspension spectrum of a finite pointed CW complex. This follows by realizing its finite list of stable cell attachments after a common suspension; stable homotopies between the attaching maps are realized after increasing that suspension once more. Hence $\mathcal L$ is generated by the suspension spectra it contains, and the checked formula $TR(\mathcal L)=\loc(\mathcal L\cap\{\text{suspension spectra}\})$ from POLES F.5 gives the first equality.

We recall the compactness argument for the second equality. A localizing subcategory generated by a set of compact objects consists of filtered homotopy colimits of their finite iterated cofibers, shifts and retracts. One can verify this by cellular approximation: at stage zero take a coproduct of shifts of the generators mapping to a proposed object, and successively attach such cells to kill the kernel on maps out of every generator. The sequential comparison is an equivalence within the generated subcategory. At a fixed finite stage, compactness lets each finite collection of attaching maps factor through a finite collection of cells of preceding stages. Induction on the finite stage shows that the resulting object is a filtered homotopy colimit of finite constructions. Sequential colimits are themselves filtered, so this remains true at the limit. If a compact object lies in this subcategory, its identity factors through one of these finite constructions. It is a retract of that object and therefore belongs to the original thick subcategory. The reverse inclusion is immediate. Apply this to the finite generators $\mathcal C$.
\end{proof}

\begin{claim}{D.2.2}{Proved}{All nonzero torsion finite thick subcategories occur without $RT$-fixedness}
Every nonzero thick subcategory $\mathcal C$ of torsion finite spectra occurs as the finite part of $T(W)$ for an accessible reversible localizer $W$ which is not $RT$-fixed. In fact $T(W)=\loc(\mathcal C)$ can be imposed.
\end{claim}
\begin{proof}
Choose a prime $p$ for which $\mathcal C$ contains a nonzero $p$-primary finite spectrum. Use the sphere family of A.1 if $p=2$ and the Toda family of A.6 otherwise, and denote its reversible null class by $\Nul_p$.

We need one consequence of the unstable Kahn--Priddy argument. For every finite $p$-primary spectrum $F$, all sufficiently positive infinite-loop spaces $\Omega^\infty\Sigma^rF$ lie in $\Nul_p$. To prove it, the space form of Kahn--Priddy for $B\mathbb Z/p$, explicitly stated in the main theorem of \citet[p.~103]{KahnPriddy1978}, gives $(Q_0S^0)_{(p)}$ in the reversible closure of $B\mathbb Z/p$. The proof of C.3, with $2$ replaced by $p$, then gives $\Omega^\infty\Sigma^r(\mathbb S/p)$ for all $r\geq1$. Finite $p$-primary spectra form the thick subcategory generated by $\mathbb S/p$: alternatively, if $p^a\id_F=0$, then $F$ is a retract of $F/p^a$, which is built from $\mathbb S/p$ by finite cells and extensions. Such an exponent exists because $\operatorname{End}(F)$ is a finite $p$-group. A finite construction by shifts, cofibers and retracts can be suspended far enough that all its infinite-loop-space fibre sequences have connected bases. At a cofiber step the first and second spaces are already aspherical; the fibre rule in C.1 gives the third. Shifts are absorbed in the common suspension and retracts are preserved. This proves the consequence without asserting that arbitrary fibres of equivalences are aspherical.

Write a finite torsion spectrum as the finite wedge of its primary summands; the stable Chinese-remainder idempotents give this decomposition, and thickness puts each summand of an object of $\mathcal C$ back in $\mathcal C$. For each primary $F\in\mathcal C$ choose a finite simply connected $X_F$ and $r_F$ with $\Sigma^\infty X_F\simeq\Sigma^{r_F}F$. If $F$ is $p$-primary choose $r_F$ sufficiently large as above and use the aspherical generator $QX_F\to *$. For the other primes use $X_F\to *$. These latter spaces are null for the $p$-local target family because their $p$-localizations are contractible.

Let $W$ be the reversible closure of this set of object-to-point maps. It is accessible by B.4. All its generators lie in $W_{\Nul_p}$, so $W\subseteq W_{\Nul_p}$. They also lie in $R(\mathcal L)$ for $\mathcal L=\loc(\mathcal C)$: in the $QX_F$ case the Snaith splitting expresses its suspension spectrum as homotopy orbits of smash powers of $\Sigma^\infty X_F$, all in the tensor ideal $\mathcal L$. Thus $T(W)\subseteq\mathcal L$. Conversely the first summand in the same splitting retracts onto $\Sigma^\infty X_F$, and the other-prime generators already have that spectrum. Therefore $T(W)$ contains every generator of $\mathcal L$ and equals $\mathcal L$.

Finally choose a nonzero $p$-primary $F\in\mathcal C$. Its finite space $X_F$ is not null for the target family. Indeed the Spanier--Whitehead dual of its nonzero finite suspension spectrum is nonzero, so some stable cohomotopy group of $X_F$ localized at $p$ is nonzero. After suspending the domain and taking a sufficiently large odd-dimensional sphere as target, the stable class is represented unstably. Such odd spheres occur in both target families. Hence $X_F\notin\Nul_p$, while $\Sigma^\infty X_F\in\mathcal L=T(W)$. The map $X_F\to *$ lies in $RT(W)$ and not in $W$, as required.
\end{proof}

\begin{claim}{D.2.3}{Proved}{The full finite thick subcategory occurs without fixedness}
There is an accessible reversible localizer $W$ with $T(W)=\Ho(\Sp)$ and $W\neq RT(W)$. Thus the full finite thick subcategory also occurs.
\end{claim}
\begin{proof}
Use G.5: the smallest reversible closure of $B\mathbb Z/2\to *$ and $B\mathbb Z[1/2]\to *$ is contained in the non-total completed-sphere localizer, but its stable part contains $\mathbb S/2$ and $\mathbb S[1/2]$, which together generate the sphere. The completed EHP comparisons and nonmembership of $\RP^2$ are proved there.
\end{proof}

\begin{claim}{D.2.4}{Open}{The zero finite thick subcategory}
The remaining finite-part question is whether a non-$RT$-fixed reversible localizer can have $T(W)\cap\Ho(\Sp)^\omega=0$. If the stronger condition $T(W)=0$ is imposed, the answer is no: $W\subseteq R(0)=W_{\mathbb Z}$ and every reversible localizer contains $W_{\mathbb Z}$, so $W=W_{\mathbb Z}$. This does not settle a nonzero stable part having no nonzero finite objects.
\end{claim}

\begin{claim}{D.2.5}{Proved}{A rational stable part cannot give the missing separation}
If a reversible fundamental localizer $W$ satisfies
\[
 0\ne T(W)\subseteq\loc(H\mathbb Q),
\]
then $T(W)=\loc(H\mathbb Q)$ and
$W=R(\loc(H\mathbb Q))$. In particular it is $RT$-fixed.
\end{claim}
\begin{proof}
Write $\mathcal A_W$ for its pointed aspherical spaces. The stable
category $\loc(H\mathbb Q)$ is precisely the rational stable category:
the rationalized sphere is $H\mathbb Q$, and its shifts generate that
category. Since $T(W)$ is nonzero, some $X\in\mathcal A_W$ has nonzero
reduced suspension spectrum. It is connective and rational, so its
first nonzero stable homotopy group is its first nonzero integral
homology group, a nonzero rational vector space. All its reduced
integral homology groups are rational. Also $X$ is connected, since
otherwise its reduced $H_0$ is a nonzero free abelian group.

Suspend $X$ enough that $Z=\Sigma^rX$ is simply connected and its first
nonzero reduced homology occurs in an odd degree $d\geq3$.
The rationalization criterion for simply connected spaces
\cite[Theorems~6.1.1--6.1.2]{FGHMayPonto2012} makes $Z$ rational.
Choose $a\in H_dZ$ and a rational-linear functional
$\ell:H_dZ\to\mathbb Q$ with $\ell(a)=1$.
Hurewicz realizes $a$ by $S^d\to Z$, which extends through
$S^d_{\mathbb Q}$ since $Z$ is rational. The universal coefficient
theorem, using injectivity of $\mathbb Q$, realizes $\ell$ by
$Z\to K(\mathbb Q,d)$. For odd $d$ the latter target is
$S^d_{\mathbb Q}$, by \cite[Corollary~6.7.3]{FGHMayPonto2012}.
The composite is the identity on the sole nonzero homotopy group, so
$S^d_{\mathbb Q}$ is a homotopy retract of $Z$. It belongs to
$\mathcal A_W$.

For $n\geq1$, the natural rational homology comparison and simple
connectivity give
$\Sigma S^n_{\mathbb Q}\simeq S^{n+1}_{\mathbb Q}$; at $n=1$ one
may use the telescope model of $K(\mathbb Q,1)$.
Suspension and desuspension closure therefore put all positive
rational spheres in $\mathcal A_W$.

We give the needed cellular step explicitly. Let $Y$ be a simply
connected rational space, with a CW model having one vertex and no
one-cells. Inductively rationalize the boundary spheres of its cells.
At a stage already constructed as a simply connected rational space,
each attaching map extends from its ordinary sphere to the rational
sphere. Attach its rationalized cone by a homotopy cofiber, doing all
cells in the dimension by a wedge. The new space is simply connected
and has rational integral homology, hence remains rational; comparison
with the ordinary skeleton is a rational homology equivalence by the
cofiber exact sequence. At the sequential homotopy colimit the result
is $Y_{\mathbb Q}\simeq Y$. All boundaries, wedges, cofibers and
colimits used are in $\mathcal A_W$, so $Y\in\mathcal A_W$.
For the two-cells the boundary is a wedge of $K(\mathbb Q,1)$'s and
its cofiber is simply connected. No nilpotence assertion about that
wedge of circles is needed.

Finally, if $U$ has rational reduced suspension spectrum, reduced
$H_0$ makes it connected, and $\Sigma U$ is simply connected with
rational homology. The preceding argument and desuspension give
$U\in\mathcal A_W$. This is exactly the aspherical class of
$R(\loc(H\mathbb Q))$; the opposite inclusion already follows from
$W\subseteq R(T(W))$. The reversible cofiber criterion identifies
the localizers. Rational spheres in their null class generate
$\loc(H\mathbb Q)$, giving the asserted stable part.
\end{proof}

\begin{claim}{D.2.6}{Proved}{A concrete nonzero stable part with no finite spectra}
For every prime $p$, put
\[
 \mathcal L_p=\loc\bigl(\Sigma^\infty K(\mathbb Z/p,2)\bigr).
\]
Then $\mathcal L_p\ne0$ and
$\mathcal L_p\cap\Ho(\Sp)^\omega=0$.
The smallest reversible localizer containing
$K(\mathbb Z/p,2)\to *$ is accessible and has stable part
exactly $\mathcal L_p$.
\end{claim}
\begin{proof}
The exact form of Lee's theorem used here is the opened
\cite[Corollary~1.4.1]{JYuan2024}: for a simply connected $\pi$-finite
space $A$, the collapse gives an equivalence
$\mathbb S\simeq F(\Sigma^\infty A_+,\mathbb S)$.
Here $\pi$-finite means finitely many components and finitely many
nonzero homotopy groups, all finite. Taking $A=K(\mathbb Z/p,2)$
and splitting off the basepoint gives
\[
                  F(\Sigma^\infty A,\mathbb S)=0.
\]
The annihilator $\{E:F(E,\mathbb S)=0\}$ is localizing, since
mapping spectra take coproducts to products, cofiber sequences to
fiber sequences, and shifts to inverse shifts. It contains
$\mathcal L_p$. A finite spectrum in this annihilator has zero
Spanier--Whitehead dual, hence is zero by finite duality.
On the other hand $H_2(A;\mathbb Z)=\mathbb Z/p$, so the generator
is nonzero.

Let $W$ be the specified smallest reversible localizer. It is
accessible by B.4, and $W\subseteq R(\mathcal L_p)$ implies
$T(W)\subseteq\mathcal L_p$. Its aspherical generator gives the
opposite inclusion. This proves all assertions without asserting
that $W$ is nonfixed.
\end{proof}

These results narrow D.2.4 but do not settle it. Neither a nonzero
finite space nor $B\mathbb Z/p$ can serve as a missing aspherical
object for this stable part. Replacing the generator by
$QK(\mathbb Z/p,2)$ also gives no separator: the original
Eilenberg--Mac Lane space is an infinite-loop retract of its $Q$.
The stable-cohomotopy claims in
\cite[Section~10.12 and Lemma~10.13]{EH5Bousfield1997} additionally
rule out $B(\mathbb Z/p^\infty)$ and $\mathbb CP^\infty$ as naive
replacement candidates. A genuinely unstable separator is still
required.

\paragraph{Current endpoint of the following obstruction calculation.}
The successive finite calculations below culminate in D.2.61:
the actual Moore composite $h_Ml_4$ is zero, the related
class $\alpha\sigma'$ has exact order two, and the detecting
integral class is stably spherical.  D.2.59 also determines
all seven primary degree-fifteen classes and their integral lifts.
The unresolved stages mentioned in the earlier reductions are
settled to that extent by these later results.  What remains
open in D.2.4 is the full-source extension and its consequences
through all later nullification stages; no finite calculation
below is asserted to settle that question on its own.
\subsection*{A finite-suspension test for the smallest reversible class}

\begin{claim}{D.2.7}{Proved}{Ordinary nullification after a bounded suspension}
\label{D27nullification}
Let $A$ be a connected pointed space whose reduced integral homology is
$p$-primary torsion. Suppose that, for some integer $r\geq1$, $A$ kills
every $r$-connected generalized Eilenberg--Mac Lane space (GEM) with
$p$-primary torsion homotopy groups. Then the aspherical class of the
smallest reversible fundamental localizer containing $A\to *$ is
\[
 \mathcal A_{\min\mathrm{Rev}(A)}
   =\{X\in\cS_*:P_A(\Sigma^rX)\simeq *\}.
 \tag{D.2.7a}
\]
Here $P_A$ is ordinary unstable $A$-nullification. In particular:
\begin{enumerate}
\item For $A=K(\mathbb F_p,n)$, $n\geq1$, one may take
      $r=\max\{1,n-1\}$.
\item For $A=K(C_{p^\infty},n)$, $n\geq1$, where
      $C_{p^\infty}=\mathbb Z[1/p]/\mathbb Z$, one may take $r=n$.
\item If $B\mathbb Z/p$ is $A$-acyclic, one may take $r=1$.
      Thus for every connected $V$ with $p$-primary reduced integral
      homology,
\[
 \mathcal A_{\min\mathrm{Rev}(B\mathbb Z/p\vee V)}
  =\{X:P_{B\mathbb Z/p\vee V}(\Sigma X)\simeq *\}.
 \tag{D.2.7b}
\]
\end{enumerate}
The statements impose no bound on cardinalities or torsion exponents
and use no $p$-completion.
\end{claim}

\begin{proof}
Put $\mathcal C_A=\{Y:P_AY\simeq *\}$ and
$\mathcal N_r=\{X:\Sigma^rX\in\mathcal C_A\}$.
Ordinary nullification exists for every space $A$
\citep[Theorem~1.1]{Bousfield1996ULP}. Its mapping-space universal
property shows that $\mathcal C_A$ is closed under weak equivalences,
retracts and pointed homotopy colimits. In a homotopy cofiber sequence,
membership of the first and second terms, or of the first and third
terms, implies membership of the remaining term: apply the derived
mapping-space functor into every $A$-null space. These same properties
hold for $\mathcal N_r$, since suspension preserves pointed homotopy
colimits and cofiber sequences.

Every $A$-acyclic space has $p$-primary reduced integral homology.
Indeed, $A\to *$ is an $H\mathbb Z[1/p]$-equivalence; every
$H\mathbb Z[1/p]$-local space is therefore $A$-null. The universal
property of $P_A$ implies that every $A$-acyclic space is
$H\mathbb Z[1/p]$-acyclic. Since $\mathbb Z[1/p]$ is flat, its
vanishing reduced homology means exactly that every element of each
reduced integral homology group is killed by a power of $p$.

We prove the essential desuspension implication. Suppose
$\Sigma X\in\mathcal N_r$, so that
$\Sigma^{r+1}X\in\mathcal C_A$. The preceding paragraph and the
suspension homology isomorphism imply that every
$\widetilde H_i(X;\mathbb Z)$ is $p$-primary torsion. In particular
$X$ is connected, because $\widetilde H_0(X;\mathbb Z)$ is free
abelian. Hence $\Sigma^rX$ is $r$-connected. By Dold--Thom,
$SP^\infty\Sigma^rX$ is an $r$-connected GEM with $p$-primary
torsion homotopy groups, and the hypothesis gives
$SP^\infty\Sigma^rX\in\mathcal C_A$.

The exact classical input is \citet[Theorem~3.3]{Bousfield1996ULP}:
for a connected space $U$ and an integer $k\geq1$,
\[
 \langle\Sigma U\rangle
  =\langle\Sigma^kU\vee SP^\infty\Sigma U\rangle.
\]
The brackets here denote ordinary nullity classes: equality says that
the two displayed spaces kill exactly the same spaces under ordinary
nullification. Apply this identity to the connected space
$U=\Sigma^{r-1}X$, with $k=2$. Both summands
$\Sigma^{r+1}X$ and $SP^\infty\Sigma^rX$ are $A$-acyclic, so their
wedge is $A$-acyclic. The identity gives
$\Sigma^rX\in\mathcal C_A$, proving $X\in\mathcal N_r$.

For the third cofiber implication, if $Y,C\in\mathcal N_r$ in a
cofiber sequence $X\to Y\to C$, use the next cofiber sequence
$Y\to C\to\Sigma X$ and the already proved first implication.
It gives $\Sigma X\in\mathcal N_r$, and desuspension gives
$X\in\mathcal N_r$. Thus A.3 constructs a reversible fundamental
localizer with aspherical class $\mathcal N_r$. It contains $A\to *$
because every positive suspension of $A$ is $A$-acyclic.

For minimality, let $W$ be any reversible fundamental localizer
containing $A\to *$. The ordinary localizer generated by this one map
is contained in $W$, so every space in $\mathcal C_A$ is
$W$-aspherical. If $X\in\mathcal N_r$, then
$\Sigma^rX\in\mathcal C_A$; reversible desuspension, applied $r$
times, puts $X$ in the aspherical class of $W$. The reversible cofiber
criterion identifies the localizers themselves. This proves
\textup{(D.2.7a)}.

For the first listed application, the opened statement of
\citet[Lemma~8.13]{EH5Bousfield1997} says that $K(G,n)$, for
$n\geq2$ and an arbitrary nonzero $p$-torsion abelian group $G$,
kills each $(n-1)$-connected generalized polyGEM with $p$-torsion
homotopy groups, subject to the additional condition
$\pi_nX/p\pi_nX=0$ only when $G/pG=0$. A GEM is a generalized
polyGEM. Taking $G=\mathbb F_p$ removes the additional condition,
and gives $r=n-1$ when $n\geq2$.

For $n=1$, $B\mathbb Z/p$ kills $K(\mathbb F_p,2)$ by the
simplicial bar construction: its levels are finite powers of
$B\mathbb Z/p$, nullification preserves finite products
\citep[proof of Theorem~1.7]{Bousfield1996ULP}, and its acyclic
class is closed under geometric realization. Applying Lemma~8.13 to
$K(\mathbb F_p,2)$ gives all simply connected $p$-torsion GEMs.
This proves $r=1$ and also proves the third listed application by
transitivity of ordinary acyclicity.

For $G=C_{p^\infty}$ and $n\geq2$, the extra condition in
Lemma~8.13 cannot be discarded: $G/pG=0$. It is automatically
satisfied on $n$-connected GEMs, giving $r=n$. To include $n=1$,
use the universal coefficient injection
\[
 \operatorname{Ext}_{\mathbb Z}(C_{p^\infty},\mathbb F_p)
  \lhook\joinrel\longrightarrow
 H^2(BC_{p^\infty};\mathbb F_p).
\]
The group on the left is nonzero: the extension
$0\to\mathbb F_p\to C_{p^\infty}\xrightarrow{p}C_{p^\infty}\to0$
does not split, since a direct summand of a divisible abelian group is
divisible. Thus $K(\mathbb F_p,2)$ is not
$BC_{p^\infty}$-null. By
\citet[Lemma~9.13]{EH5Bousfield1997}, an Eilenberg--Mac Lane
space $K(\mathbb F_p,m)$ is either acyclic or local for any specified
localization. It is therefore $BC_{p^\infty}$-acyclic. Lemma~8.13
again supplies all simply connected $p$-torsion GEMs, giving $r=1$.

Finally these Eilenberg--Mac Lane generators have $p$-primary
reduced integral homology. One can use the nilpotent localization
criterion away from $p$, or start with the finite cyclic groups and
use the bar construction, followed by filtered colimits for
$C_{p^\infty}$. This verifies the homology premise in every case.
\end{proof}

\begin{remark}
For the zero-finite-part example in D.2.6, the exact remaining unstable
test becomes
\[
 \Sigma^\infty X\in\loc(\Sigma^\infty K(\mathbb F_p,2)),
 \qquad P_{K(\mathbb F_p,2)}(\Sigma X)\not\simeq *.
\]
Likewise \textup{(D.2.7b)} turns the interval tests generated by
$B\mathbb Z/p\vee V$ into ordinary-nullification calculations.
Neither formula supplies the required separating space by itself;
D.2.4 and the strict-intermediate-element problem in D.3 retain
their stated status. The initial suspension is essential:
Theorem~3.3 makes no assertion that the unsuspended ordinary
$B\mathbb Z/p$-acyclic class is desuspension-closed.
\end{remark}

\begin{claim}{D.2.8}{Proved}{Finite primary space generators are $RT$-fixed}
Let $V$ be a nonzero finite connected pointed complex with $p$-primary
reduced integral homology, where ``nonzero'' means
$\Sigma^\infty V\not\simeq0$. Then
\[
 \min\mathrm{Rev}(V\to *)
   =R\bigl(\loc(\Sigma^\infty V)\bigr).
 \tag{D.2.8}
\]
This equality concerns all pointed source spaces, without a finiteness
hypothesis on them. Thus the finite primary generators themselves
give fixed localizers; the use of infinite-loop generators in D.2.2
is material.
\end{claim}

\begin{proof}
Put $W=\min\mathrm{Rev}(V\to *)$ and
$\mathcal L=\loc(\Sigma^\infty V)$. The inclusion
$W\subseteq R(\mathcal L)$ follows from the defining minimum.
It remains to prove the reverse inclusion on aspherical spaces.

Let $\Sigma^\infty X\in\mathcal L$. The spectrum $H\mathbb Z[1/p]$
annihilates $\Sigma^\infty V$, so it annihilates $\mathcal L$.
Consequently $X$ has $p$-primary reduced integral homology. Its
reduced $H_0$, being free abelian, is zero; hence $X$ is connected.
The space $Y=\Sigma X$ is simply connected and has $p$-primary
integral homology. The nilpotent homological localization criterion
therefore gives $Y\simeq Y_{(p)}$.

The nonzero finite $p$-primary spectrum $\Sigma^\infty V$ has some
finite type $n+1\geq1$. For each $0\leq i\leq n$, the telescope
$T(i)$ annihilates $\Sigma^\infty V$, and thus every spectrum in
$\mathcal L$. In particular $Y\to *$ is a $T(i)$-homology
equivalence for all these $i$.

We use the exact finite-generator theorem stated and proved in
\citet[Theorem~4.1]{FGHBarkanHeutsShi2026}, attributed there to
Bousfield's 1994 work. It says that, for an arbitrary map of spaces
$f$, being a $T(i)$-homology equivalence for every $i\leq n$ is
equivalent to
\[
 P_{\Sigma V}(\Sigma^k f_{(p)})
 \quad\hbox{being an equivalence for all sufficiently large }k.
\]
Apply it to $Y\to *$, using $Y\simeq Y_{(p)}$. We obtain
$P_{\Sigma V}(\Sigma^kY)\simeq *$ for some $k$.
The ordinary localizer generated by $\Sigma V\to *$ lies in $W$,
since $V$ and its suspension are $W$-aspherical. Thus
$\Sigma^{k+1}X$ is $W$-aspherical. Repeated reversible
desuspension puts $X$ in the same aspherical class. The reversible
cofiber criterion now gives the equality of localizers in
\textup{(D.2.8)}.

Only the acyclicity of finite type-$(n+1)$ spectra for $T(i)$,
$i\leq n$, and the displayed primary theorem are used. No
identification of telescopic and Morava-$K$ localizations is assumed.
\end{proof}

\begin{claim}{D.2.9}{Proved}{The connective stable obstruction and its limit}
Let $C\in\Ho(\Sp)_{\geq n}$, $n\geq1$, have $p$-primary integral
homology and $H_n(C;\mathbb F_p)\neq0$. Every
$E\in\loc(C)\cap\Ho(\Sp)_{\geq n}$ satisfies
\[
 P_C^{\Sp}E\simeq0,
 \tag{D.2.9a}
\]
where $P_C^{\Sp}$ is spectral nullification using the mapping
\emph{space} out of $C$, equivalently all its nonnegative
suspensions. Thus if $C=\Sigma^\infty A$, the infinite-loop space
$\Omega^\infty E$ is ordinarily $A$-acyclic.

In particular put $B=K(\mathbb F_p,2)$,
$\mathcal I_p=\loc(\Sigma^\infty B)$ and
$A_r=Q\Sigma^rB$, for $r\geq0$. Every
$E\in\mathcal I_p\cap\Ho(\Sp)_{\geq r+2}$ gives an ordinarily
$A_r$-acyclic infinite-loop space. The localizers
\[
 W_r=\min\mathrm{Rev}(A_r\to *)
\]
form a descending chain with common stable part $\mathcal I_p$,
which has zero finite part. The exact remaining test is
\[
 B\in\mathcal A_{W_r}
 \quad\Longleftrightarrow\quad
 P_{A_r}(\Sigma^{r+1}B)\simeq *.
 \tag{D.2.9b}
\]
Moreover, for $r\geq1$,
\[
 \begin{aligned}
 P_{A_r}A_{r-1}&\simeq K(\mathbb F_p,r+1),\\
 W_{r-1}&=W_r\vee
   \min\mathrm{Rev}(K(\mathbb F_p,r+1)\to *).
 \end{aligned}
 \tag{D.2.9c}
\]
The connective assertion does not decide this test, and neither a
strict inclusion in this chain nor its collapse is asserted.
\end{claim}

\begin{proof}
Write $C_i=\Sigma^{-i}C$ and $Y_i=P_{C_i}^{\Sp}E$, for
$i\geq0$. Their natural maps point from $Y_i$ to $Y_{i+1}$.
The precise input is \citet[Theorem~2.6]{Bousfield1996ULP}:
because $C$ satisfies the $n$-supported $p$-torsion condition,
\[
 \fib(Y_i\longrightarrow Y_{i+1})
   \simeq\Sigma^{n-i-1}HG_i
 \tag{D.2.9d}
\]
for some $p$-primary torsion abelian group $G_i$. Hence, for
$j>i$, the fiber of $Y_i\to Y_j$ lies in
$\Ho(\Sp)_{\leq n-i-1}$, by extension closure. Set
$Y_\infty=\colim_iY_i$. Filtered colimits in spectra are exact
and preserve these homotopy-group bounds, so
\[
 \fib(Y_i\to Y_\infty)\in\Ho(\Sp)_{\leq n-i-1}.
 \tag{D.2.9e}
\]

Fix $i$, and use the map $Y_{i+1}\to Y_\infty$. Its fiber lies
in $\Ho(\Sp)_{\leq n-i-2}$, and its suspended fiber lies in
$\Ho(\Sp)_{\leq n-i-1}$. On the other hand
$C_i\in\Ho(\Sp)_{\geq n-i}$. The $t$-structure orthogonality
therefore shows that
\[
 \Map_{\Sp}(C_i,Y_{i+1})
   \xrightarrow{\simeq}\Map_{\Sp}(C_i,Y_\infty).
\]
The first mapping space is contractible because $Y_{i+1}$ is
$C_{i+1}$-null and $C_i=\Sigma C_{i+1}$. Thus $Y_\infty$ is
null for every positive and negative shift of $C$. The shift from
$Y_i$ to $Y_{i+1}$ in this paragraph controls components as well
as positive homotopy groups; a vanishing fiber alone would not
have sufficed.

The cone-attachment construction of spectral nullification puts
each $\cofib(E\to Y_i)$ in $\loc(C)$. Since $E\in\loc(C)$,
also $Y_i$ and $Y_\infty$ belong to $\loc(C)$. A member of
$\loc(C)$ that is right orthogonal to all shifts of $C$ is zero:
its mapping spectrum out of every member of $\loc(C)$ vanishes,
including its own identity. Hence $Y_\infty=0$.
Equation~\textup{(D.2.9e)} with $i=0$ now gives
$Y_0\in\Ho(\Sp)_{\leq n-1}$. But nullification by
$n$-connective $C$ preserves $n$-connectivity, since attaching
cones on nonnegative suspensions of $C$ stays in
$\Ho(\Sp)_{\geq n}$. Thus $Y_0=0$, proving
\textup{(D.2.9a)}. This argument neither assumes $C$ compact
nor interchanges its mapping-space functor with a filtered
colimit.

The passage to infinite-loop spaces is exactly
\citet[Theorem~2.10 and Proposition~3.1]{Bousfield1996ULP}:
\[
 P_A\Omega^\infty E
   \simeq\Omega^\infty P_{\Sigma^\infty A}^{\Sp}E.
\]
For $A_r$, Snaith splitting identifies
$\loc(\Sigma^\infty A_r)=\mathcal I_p$: its linear summand is
$\Sigma^r\Sigma^\infty B$, and all higher extended powers
belong to its localizing tensor ideal. The first homology group
of $A_r$ is $H_{r+2}(A_r;\mathbb Z)=\mathbb F_p$, and
$\Sigma^\infty A_r$ is $(r+2)$-connective. The already proved
assertion therefore applies with $n=r+2$.

The connected infinite-loop space $A_{r+1}$ is the delooping of
$A_r$. Its simplicial bar construction, preservation of finite
products by ordinary nullification, and closure of its acyclic
class under realization show that $A_r$ kills $A_{r+1}$.
Consequently $W_{r+1}\subseteq W_r$. Also
$W_0=\min\mathrm{Rev}(B\to *)$: the infinite-loop structure
on $B$ makes it a retract of $QB$, while every space kills its
free infinite-loop space by Proposition~3.1 of the same source.
Snaith splitting gives $T(W_r)=\mathcal I_p$ for every $r$,
and D.2.6 gives its zero finite part.

We compute the adjacent ordinary nullification explicitly. First
$H\mathbb F_p\in\mathcal I_p$: smash its generator with
$H\mathbb F_p$ and split the resulting nonzero
$H\mathbb F_p$-module into a wedge of shifts of
$H\mathbb F_p$. Its degree-two homology supplies a nonzero
summand, and $\mathcal I_p$ is closed under retracts and shifts.
For $r\geq1$, set $E=\Sigma^{r-1}\Sigma^\infty B$. Its lower
Postnikov section is
\[
 \tau_{\leq r+1}E\simeq\Sigma^{r+1}H\mathbb F_p.
\]
Both $E$ and this section belong to $\mathcal I_p$, so their
fiber $\tau_{\geq r+2}E$ does also. Apply
\textup{(D.2.9a)} with $C=\Sigma^\infty A_r$ and $n=r+2$ to
kill that fiber. The displayed lower section is already
$C$-null by $t$-structure orthogonality. Spectral nullification
preserves this cofiber sequence because its third term is
$C$-null \citep[Corollary~2.3]{Bousfield1996ULP}. Hence
\[
 P_{\Sigma^\infty A_r}^{\Sp}E
   \simeq\Sigma^{r+1}H\mathbb F_p.
\]
Theorem~2.10 then gives the first equality in
\textup{(D.2.9c)}, including its natural nullification unit
$A_{r-1}\to K(\mathbb F_p,r+1)$. This unit is a
$W_r$-equivalence. Consequently adjoining the collapse of its
target to $W_r$ is equivalent to adjoining the collapse of
$A_{r-1}$. Since $W_r\subseteq W_{r-1}$, the second equality
in \textup{(D.2.9c)} follows. Thus an adjacent strictness question
is reduced to the membership of one specified
Eilenberg--Mac Lane space, rather than settled by a connectivity
comparison.

Finally $A_r$ kills its infinite symmetric product, which has
$K(\mathbb F_p,r+2)$ as a homotopy retract. Lemma~8.13 of
\citet{EH5Bousfield1997} shows that it kills every
$(r+1)$-connected $p$-torsion GEM. Apply D.2.7 with the bound
$r+1$ to obtain \textup{(D.2.9b)}. That target is a suspension
of $B$; the spectral argument does not supply an infinite-loop
description of it. This is the exact unproved step separating
the reduction from a zero-finite-part counterexample.
\end{proof}

\begin{claim}{D.2.10}{Proved}{Stable sphere annihilation implies finite-target Miller nullity}
Put
\[
 \mathcal A_{\mathbb S}
 =\{E\in\Ho(\Sp):\operatorname{F}(E,\mathbb S)\simeq0\},
\]
where $\operatorname{F}$ denotes the derived function spectrum, so the
condition includes every shift. If a pointed space $Y$ satisfies
$\Sigma^\infty Y\in\mathcal A_{\mathbb S}$, then, for every finite simply
connected pointed complex $K$ and every prime $p$,
\[
 \Map_*(Y,K)\simeq *,\qquad \Map_*(Y,K_{(p)})\simeq *.
\]
Consequently
\[
 R(\mathcal A_{\mathbb S})\subseteq W_S,
 \qquad R(\mathcal I_p)\subseteq W_S,
 \quad \mathcal I_p=\loc\bigl(\Sigma^\infty K(\FF_p,2)\bigr).
\]
The second inclusion uses D.2.6. Thus finite simply connected targets,
including the sphere tests defining $W_S$, cannot separate the
zero-finite-part candidate from its stable closure. This does not decide
whether that candidate is fixed by $RT$.
\end{claim}
\begin{proof}
Write $\mathcal C=\operatorname{coloc}(\mathbb S)$ for the smallest full
stable subcategory containing $\mathbb S$ and closed under retracts and
small homotopy limits. For each fixed $E$, the kernel of
$\operatorname{F}(E,-)$ has these closure properties. We will show that
the spectra underlying the Goodwillie layers of $K$ and $K_{(p)}$ lie
in $\mathcal C$.

\emph{Finite-group classifying spectra.}
For a finite $p$-group $P$, let $\mathbb S_P$ be the genuine equivariant
sphere. The $p$-completed Segal theorem gives
\begin{equation}\label{eq:D-stable-Miller-Segal}
 ((\mathbb S_P)^P)^{\wedge}_p
 \simeq \operatorname{F}(\Sigma^\infty BP_+,\mathbb S^{\wedge}_p).
\end{equation}
The precise comparison is \cite[Theorem~1.1]{FGHBergsakerRognes2023}
with $B=\mathbb S$; its introduction explicitly identifies $B^{\wedge P}$
with $\mathbb S_P$. Completion commutes here with the function spectrum:
use $X^{\wedge}_p=\holim_a X/p^a$, exactness, and dualizability of the
finite Moore spectra. The tom Dieck splitting
\cite[Corollary~2.3 and Section~3]{FGHCarlsson1992} is
\[
 (\mathbb S_P)^P\simeq
 \bigvee_{(H\leq P)}\Sigma^\infty B(N_P(H)/H)_+.
\]
The $H=1$ summand and \eqref{eq:D-stable-Miller-Segal} show that
$(\Sigma^\infty BP_+)^{\wedge}_p$ is a retract of a homotopy limit of
$\mathbb S^{\wedge}_p$. Since
$\mathbb S^{\wedge}_p=\holim_a\mathbb S/p^a$ belongs to $\mathcal C$,
so does this retract.

Let $G$ be any finite group and choose a Sylow $p$-subgroup $P$.
The stable transfer of the finite covering $BP\to BG$ gives
\[
 \Sigma^\infty BG_+\xrightarrow{\operatorname{tr}}
 \Sigma^\infty BP_+\xrightarrow{j}\Sigma^\infty BG_+.
\]
On ordinary homology, $j\operatorname{tr}$ is multiplication by
$[G:P]$: a singular simplex has exactly that many lifts. Thus this
composite becomes an equivalence after $p$-completion, and
$(\Sigma^\infty BG_+)^{\wedge}_p$ is a retract of
$(\Sigma^\infty BP_+)^{\wedge}_p$. This uses the homology formula for
the transfer: its cofiber is bounded below, and vanishing mod-$p$
homology implies that its smash product with $\mathbb S/p$ vanishes
by the stable Hurewicz theorem. It therefore has zero $p$-completion.
No equality of the stable composite with an integer multiple of the
identity is needed. The stable transfer construction is recalled
in \cite[Section~2(B)]{FGHCarlsson1992}.

The reduced spectrum $\Sigma^\infty BG$ has finite
$|G|$-primary homotopy groups in each degree. Indeed its connective
Atiyah--Hirzebruch spectral sequence has terms
$\widetilde H_s(BG;\pi_t\mathbb S)$, with $s\geq1$, $t\geq0$,
and finitely many terms in each total degree. These groups are finite
and annihilated by $|G|$: use the bar resolution, which is finite in
each degree, and the ordinary transfer to the trivial group.
Consequently the natural map
\[
 \Sigma^\infty BG\longrightarrow
 \prod_{p\mid |G|}(\Sigma^\infty BG)^{\wedge}_p
\]
is an equivalence on every homotopy group. This finite product, and
then $\Sigma^\infty BG_+\simeq\mathbb S\vee\Sigma^\infty BG$, belong
to $\mathcal C$.

\emph{Localization of the orbit cells.}
We also need $(\Sigma^\infty BG_+)_{(p)}\in\mathcal C$.
First the cofiber of
\[
 \mathbb S\longrightarrow\prod_q\mathbb S^{\wedge}_q
\]
is $HV$, where $V=(\prod_q\ZZ_q)/\ZZ$. This follows on homotopy
groups: the map is an isomorphism in positive degrees and injective
in degree zero. The group $V$ is a nonzero rational vector space.
For a prime $q$, multiplication by $q$ is surjective by subtracting
an integer with the prescribed residue in the $q$-coordinate; it is
injective because a rational number with denominator $q$ lying in
every $\ZZ_r$ must be an integer. Hence $H\QQ$ is a retract of $HV$
and lies in $\mathcal C$. Every rational vector space embeds as a
direct summand in a product of copies of $\QQ$, by choosing a basis
and splitting the inclusion of its direct sum into the product.
Applying $H$ proves that $HU\in\mathcal C$ for every rational vector
space $U$. This uses products and retracts, with no coproduct-closure
assumption on $\mathcal C$.

The arithmetic pullback
\[
 \mathbb S_{(p)}\simeq
 \mathbb S^{\wedge}_p\mathop{\times}^{h}_{H\QQ_p}H\QQ
\]
therefore lies in $\mathcal C$. One can check this pullback directly:
in positive degrees it has the $p$-primary stable stems, while
$\ZZ_p\cap\QQ=\ZZ_{(p)}$ and $\ZZ_p+\QQ=\QQ_p$ give degree zero
and the vanishing of degree $-1$. Since the reduced classifying
spectrum has degreewise finite homotopy groups,
$(\Sigma^\infty BG)_{(p)}\simeq(\Sigma^\infty BG)^{\wedge}_p$.
Adding $\mathbb S_{(p)}$ proves the required localized assertion.

It follows that if $F$ is a genuine finite $G$-spectrum, then both
$F_{hG}$ and $(F_{hG})_{(p)}$ lie in $\mathcal C$. Indeed finite
genuine spectra are obtained by finite cofiber constructions and
retracts from the cells $G/H_+\wedge\mathbb S^d$, whose homotopy
orbits are $\Sigma^d\Sigma^\infty BH_+$. This assertion is only
needed for genuinely finite equivariant models; it does not assume
that an arbitrary action on an underlying finite spectrum has such
a model.

\emph{The identity tower.}
The identity functor has layers
\[
 D_k I(X)\simeq
 \Omega^\infty\bigl(\partial_k I\wedge
                    (\Sigma^\infty X)^{\wedge k}\bigr)_{h\Sigma_k}.
\]
By \cite[Theorem~2.1]{EH5AroneMahowald1999}, $\partial_k I$ is the
Spanier--Whitehead dual of a finite partition complex with its
specified $\Sigma_k$-action. This is the restriction of a genuine
finite $\Sigma_k$-spectrum. A finite simplicial model of $K$ gives a
finite equivariant CW model of $K^{\wedge k}$ after subdivision;
duality and smash products preserve genuine finiteness. Thus the
spectrum in the displayed layer for $X=K$ belongs to $\mathcal C$.
For $X=K_{(p)}$, use the natural equivalence
$\Sigma^\infty K_{(p)}\simeq(\Sigma^\infty K)_{(p)}$ and the
localized orbit assertion just proved. The layer formula extends
from finite complexes to all spaces because the identity and its
layers preserve filtered homotopy colimits.

Put $E=\Sigma^\infty Y$. Our hypothesis implies
$\operatorname{F}(E,V)=0$ for every $V\in\mathcal C$.
Adjunction therefore makes $\Map_*(Y,D_kI(X))$ contractible in both
cases. Starting from $P_0I(X)=*$, the fibre sequences of the tower
give $\Map_*(Y,P_kI(X))\simeq *$ for all $k$. The identity is
$1$-analytic \cite[Example~4.3]{FGHGoodwillie1992}; hence its tower
converges on the simply connected spaces $K$ and $K_{(p)}$
\cite[Theorem~1.13]{FGHGoodwillie2003}. Mapping out of $Y$ preserves
this homotopy limit, proving both mapping-space assertions without
any compactness assumption on $Y$.

Finally, for $f$ in $R(\mathcal A_{\mathbb S})$, apply the assertion
to its pointed homotopy cofiber $C(f)$ and to $K=S^n$, $n\geq2$,
at $p=2$. This is exactly the cofiber condition defining $W_S$.
There is no component exception: a nonbasepoint component of a
pointed space makes $\mathbb S$ a retract of its reduced suspension
spectrum, so $\Sigma^\infty C(f)\in\mathcal A_{\mathbb S}$ already
forces $C(f)$ to be connected. The class $\mathcal A_{\mathbb S}$
is localizing, since $\operatorname{F}(-,\mathbb S)$ takes
homotopy colimits to homotopy limits. D.2.6 puts its generator
$\Sigma^\infty K(\FF_p,2)$ in this class, so
$\mathcal I_p\subseteq\mathcal A_{\mathbb S}$ and the last inclusion
follows.
\end{proof}

\begin{claim}{D.2.11}{Proved}{Trivial periodic Lie algebras jointly detect every target}
Fix a prime $p$ and a finite positive height $n$. Let
$\mathcal V=\Sp_{T(n)}$ or $\Sp_{K(n)}$, and put
$\mathfrak L=\operatorname{Lie}(\mathcal V)$.
The trivial Lie algebras generate $\mathfrak L$ under small
colimits and retracts. In fact a countable set of trivial Lie
algebras suffices. In particular,
\[
 \Map_{\mathfrak L}(\operatorname{triv}U,Z)\simeq *
 \quad\hbox{for every }U\in\mathcal V
 \quad\Longrightarrow\quad Z\simeq0.
 \tag{D.2.11a}
\]

For the ordinary-space consequence, use Heuts's fully faithful
realization of the unstable $K(2)$-local category
$\mathcal M_{K(2)}$ inside pointed $p$-local spaces. For
\[
 B=K(\FF_p,2),\qquad A_r=Q\Sigma^rB,\qquad r\geq0,
\]
one has
\[
 Z\in\mathcal M_{K(2)},\quad
 \Map_*(A_r,Z)\simeq *
 \quad\Longrightarrow\quad Z\simeq *.
 \tag{D.2.11b}
\]
Thus no noncontractible target in that unstable periodic category
can separate $B$ from the reversible class generated by $A_r$.
This does not compute the ordinary-space nullification in
\textup{(D.2.9b)}, or compare it with periodic localization.
\end{claim}
\begin{proof}
We first prove the categorical assertion, including its set-size
qualification, and then establish the ordinary mapping-space
interface needed for \textup{(D.2.11b)}.

\medskip\noindent
\textbf{Products and the generated class.}
Corollary~B of \citet{I_BlansHeuts2026} proves that the cartesian
product in $\operatorname{Lie}(\Sp)$ preserves weakly contractible
colimits separately in each variable. This property passes to
$\mathfrak L$. Indeed stable Bousfield localization induces a
reflection on spectral Lie algebras whose underlying spectrum is
the localized underlying spectrum. Its operations are obtained
by localizing the original operations and using
\[
 L_E\bigl(\mathbb L(k)\wedge X^{\wedge k}\bigr)
 \simeq
 L_E\bigl(\mathbb L(k)\wedge(L_EX)^{\wedge k}\bigr).
\]
Here $E=T(n)$ or $K(n)$, and the equivalence follows because
$E$-acyclic spectra form a tensor ideal. Operadic composition and
its coherences localize as well. This is the construction made
explicit for the $K(n)$-local reflection in
\cite[Proposition~4.26]{I_Heuts2021}.
The reflection is left exact: limits of algebras are created on
underlying spectra, and stable localization is exact. Apply it
to the product-colimit comparison in ordinary spectral Lie
algebras. Local colimits are reflected ordinary colimits, and
finite products are preserved, proving the assertion in
$\mathfrak L$.

Suppose for the moment that $\mathcal G$ is a collection of
trivial algebras closed under finite products and containing
zero. Let $\mathcal C$ be its closure under small colimits and
retracts. Then $\mathcal C$ is closed under finite products.
Here is the size-independent closure argument. First fix
$A\in\mathcal G$ and consider the class of $X$ for which
$X\times A\in\mathcal C$. It contains $\mathcal G$, since
products of trivial algebras are trivial algebras on direct sums.
It is closed under retracts. Given a diagram $(X_i)_{i\in I}$
in this class, adjoin an initial vertex to $I$ and put zero
there. The extended diagram $X'$ has the same colimit as the
original diagram, while its indexing category $I^{\triangleleft}$
is weakly contractible. Consequently
\[
 \left(\colim_{i\in I}X_i\right)\times A
 \simeq
 \colim_{i\in I^{\triangleleft}}(X'_i\times A).
 \tag{D.2.11c}
\]
The additional value on the right is $0\times A=A$, which is in
$\mathcal C$, and every other value is in $\mathcal C$ by
assumption. The same argument includes the empty diagram.
Thus $X\times A\in\mathcal C$ for every $X\in\mathcal C$.
Now fix any $Y\in\mathcal C$ and repeat the argument for the
class of $X$ with $X\times Y\in\mathcal C$. It contains all
generators by the first step, and the added initial value is
now $Y\in\mathcal C$. This proves binary-product closure, hence
finite-product closure. We have not asserted preservation of
arbitrary colimits by products: the extra initial value in
\textup{(D.2.11c)} is essential.

\medskip\noindent
\textbf{A fibre sequence as an action-bar realization.}
For a fibre sequence $F\to E\to D$ in $\mathfrak L$ there is an
augmented simplicial object with terms
\[
 [q]\longmapsto F\times(\Omega D)^q
\]
whose realization is $E$. Its face maps encode the canonical
$\Omega D$-action on $F$; they need not describe a trivial action.

To construct and check it, form the \v{C}ech nerve of $0\to D$.
This is the bar construction of the group object $\Omega D$.
Its realization is $D$: the conservative forgetful functor
$U:\mathfrak L\to\mathcal V$ preserves finite limits and
geometric realizations, and in the stable category $\mathcal V$
the corresponding realization is
$\Sigma\Omega U(D)\simeq U(D)$.
Preservation of geometric realizations follows from preservation
of sifted colimits by the spectral Lie monad, whose tensor
products preserve colimits separately. Equivalently, this is
the bar calculation in the proof of
\cite[Proposition~4.11]{I_BlansHeuts2026}.

Pull this augmented nerve back along $E\to D$. Its $q$th term
is $F\times(\Omega D)^q$, by the finite-limit description of
the \v{C}ech action. After applying $U$, realization commutes
with this pullback: finite limits in a stable category are
finite colimits, and geometric realization is exact. The
realized augmentation therefore becomes
\[
 U(E)\mathop{\times}_{U(D)}
       |U(\check C(0\to D))|\simeq U(E).
\]
Conservativity proves the assertion in $\mathfrak L$. In
particular, for the product-closed class constructed above,
\[
 F\in\mathcal C,\quad \Omega D\in\mathcal C
 \quad\Longrightarrow\quad E\in\mathcal C.
 \tag{D.2.11d}
\]
Every simplicial term lies in $\mathcal C$, as does its
realization. The premise concerns $\Omega D$, not merely $D$.
No join axiom or general fibre/base-to-total closure of a
nullification kernel is used.

\medskip\noindent
\textbf{Finite weights and an explicit countable family.}
Write $\mathbb S_{\mathcal V}$ for the local unit. For every
$\ell\in\ZZ$, let
\[
 G_\ell=\operatorname{Free}_{\mathbb L}
                (\Sigma^\ell\mathbb S_{\mathcal V}),
 \qquad J_{\ell,j}=(G_\ell)_j
\]
denote the free algebra and its weight-$j$ spectrum. By
\cite[Theorem~4.22 and Corollary~4.23]{I_Heuts2021},
$J_{\ell,j}=0$ for $j>2p^n$, and already for $j>p^n$ when
$\ell$ is odd. The stated range explicitly includes every
integer $\ell$, including negative integers; the proof treats
negative shifts by partition-complex duality. For $K(n)$-local
spectra apply the monoidal reflection to the $T(n)$-local
formula.

Take $\mathcal G$ to consist of the algebras
\[
 \operatorname{triv}
 \left(\bigoplus_{a=1}^t
       \Sigma^{-k_a}J_{\ell_a,j_a}\right),
 \qquad
 \begin{gathered}
 t\geq0,\quad k_a\geq0,\quad \ell_a\in\ZZ,\\
 1\leq j_a\leq2p^n.
 \end{gathered}
 \tag{D.2.11e}
\]
This is a countable set of specified objects. It contains zero
and is closed under finite products and loops. Let $\mathcal C$
again be its colimit and retract closure. The preceding two
arguments apply to this actual set, without supposing the
existence of a nullification generated by a proper class.

For $m\geq0$, form the positively graded algebra
$(G_\ell)_{\leq m}$ by setting weights above $m$ equal to zero.
Operations add input weights, so these truncations and their
projections are well-defined operadic algebras and algebra maps.
There are fibre sequences
\[
 \operatorname{triv}J_{\ell,m}\longrightarrow
 (G_\ell)_{\leq m}\longrightarrow(G_\ell)_{\leq m-1}
 \qquad(m\geq1).
 \tag{D.2.11f}
\]
The fibre has only weight $m$. Its operations of arity at least
two land in weight at least $2m>m$ and vanish. This verification
is in the category of graded operadic algebras, so it includes
the inherited higher coherences, not just the binary bracket.

Induct on $m$, proving simultaneously that
\[
 \Omega^k(G_\ell)_{\leq m}\in\mathcal C
                   \quad\hbox{for every }k\geq0.
 \tag{D.2.11g}
\]
For $m=0$ the object is zero. Apply $\Omega^k$ to
\textup{(D.2.11f)}. Its fibre is
$\operatorname{triv}(\Sigma^{-k}J_{\ell,m})$, a generator in
\textup{(D.2.11e)}. The loop of its base is
$\Omega^{k+1}(G_\ell)_{\leq m-1}$, which belongs to
$\mathcal C$ by the induction hypothesis for this value of
$k+1$. Formula~\textup{(D.2.11d)} proves the induction step.
At $m=2p^n$ this gives $G_\ell\in\mathcal C$ for every
integer $\ell$. This induction does not assume that
$\mathcal C$ is generally closed under loops.

Local sphere shifts generate $\mathcal V$ under colimits:
present an ordinary spectrum by sphere cells and apply the
localization. The free algebra functor preserves colimits, so
every free Lie algebra lies in $\mathcal C$. Every algebra is
the geometric realization of its free-forgetful monadic bar
resolution. Hence $\mathcal C=\mathfrak L$.

The class of objects $X$ with $\Map_{\mathfrak L}(X,Z)\simeq *$
is closed under colimits and retracts. If it contains all the
trivial algebras, or just the countable family
\textup{(D.2.11e)}, it therefore contains every algebra. Taking
$X=Z$ makes the identity nullhomotopic and gives $Z\simeq0$.
This proves \textup{(D.2.11a)} and the stronger countable
assertion. In particular nullification at that countable family,
if expressed as an accessible localization, is the zero
localization. The argument itself only used the explicit
closure class and mapping-space orthogonality.

\medskip\noindent
\textbf{The precise ordinary-to-periodic interface.}
In the remainder $n=2$. Use the fully faithful models of
$\mathcal S_{v_2}$ and $\mathcal M_{K(2)}$ in
\cite[Definition~3.11, Corollary~2.8 and Proposition~4.26]
{I_Heuts2021}. Let $d_3$ be the connectivity parameter there,
choosing the lower parameter with $d_2\leq d_3$.
Here $L_i^f$ means finite-space nullification at the chosen
type-$(i+1)$ suspension, as in that source.

Suppose a spectrum $H$ is $(d_3+1)$-connective and is both
rationally and $T(1)$-acyclic. Put $X=\Omega^\infty H$.
The identity $\Phi_i\Omega^\infty=L_{T(i)}$, recorded on
page~4 of \citet{I_Heuts2021}, makes its lower periodic
homotopy vanish. Theorem~3.7 of that paper, at height one,
gives $L_1^fX\simeq *$. Proposition~3.6 and the identity
$L_1^fL_2^f=L_1^f$ show that $L_2^fX$ is still
$d_3$-connected and $L_1^f$-acyclic. Thus it belongs to the
specified full subcategory $\mathcal S_{v_2}$ and the
monochromatic functor satisfies $MX\simeq L_2^fX$.

Every realized $Z\in\mathcal M_{K(2)}$ is $L_2^f$-local.
If $F_Z\in\operatorname{Lie}(\Sp_{K(2)})$ is its Lie
algebra, nullification adjunction, fullness, and the
identification of the periodic infinite-loop functor with
the trivial Lie algebra functor give
\[
 \Map_*(\Omega^\infty H,Z)
 \simeq
 \Map_{\operatorname{Lie}(\Sp_{K(2)})}
       (\operatorname{triv}L_{K(2)}H,F_Z).
 \tag{D.2.11h}
\]
For the last identification use Theorem~2.6 and the discussion
on page~8 of \citet{I_Heuts2021}, followed by its
$K(2)$-local reflection. The connective hypotheses justify
the ordinary-space mapping adjunction; we have not treated
$M$ as a reflection from all pointed spaces.

\medskip\noindent
\textbf{Obtaining every trivial source from $A_r$-acyclic spaces.}
Put $C_0=\Sigma^\infty B$ and $E_0=L_{K(2)}C_0$.
The spectrum $C_0$ is rationally and $T(1)$-acyclic. For
the second assertion without an implicit higher telescope
conjecture, first use
\cite[Section~10.11]{EH5Bousfield1997} to get
$K(1)_*B=0$ in reduced homology. Snaith splitting makes
$K(1)_*QB=0$, since each extended power of $C_0$ is
$K(1)$-acyclic. To make the $K(1)$ versus complex $K/p$
interface explicit, Theorem~11.3 makes $QB\to *$ a virtual
$K(1)$-equivalence; immediately before Theorem~11.11,
Bousfield identifies virtual $K(1)$-equivalences with virtual
$K/p$-equivalences. Theorem~11.11 then gives zero
$v_1$-periodic homotopy for the connected $H$-space $QB$
(the same implication summarized in Corollary~11.12).
Finally $\Phi_1QB\simeq L_{T(1)}C_0$ gives the assertion.
It follows that every object of $\loc(C_0)$ has the same
two lower stable acyclicity properties.

We need a unit retract in $\Sp_{K(2)}$. The object $E_0$ is
dualizable by \cite[Corollary~8.8]{I_HoveyStrickland1999}.
The explicit calculation of
\cite[Theorem~11.1(a)]{I_RavenelWilson1980}, with
$n=q=2$ and cyclic group of order $p$, gives unreduced
$K(2)$-rank $p$, concentrated in even degrees. The extension
to $p=2$ is explicitly recorded in
\cite[Remark~4.4]{I_HopkinsRavenelWilson1998}; see also its
Corollary~1.3. Therefore $K(2)_*E_0$ has even reduced rank
$p-1$.

For odd $p$, apply the symmetric graded K\"unneth isomorphism
to the categorical trace of $\id_{E_0}$. It gives multiplication
by $p-1$ on $K(2)$-homology. Thus
\[
 \mathbb S_{K(2)}\xrightarrow{\operatorname{coev}}
 E_0\otimes DE_0\xrightarrow{\operatorname{ev}\circ\tau}
 \mathbb S_{K(2)}
\]
is an equivalence, and rescaling the second map gives a
retraction onto the unit. At $p=2$ no homotopy-commutativity
of Morava $K$ is assumed. The $K(2)$-ranks of $E_0$ and
$DE_0$ are both one. Coevaluation is nonzero on
$K(2)$-homology, since otherwise a duality triangle would
induce the zero map on $K(2)_*E_0$. It is therefore an
isomorphism between rank-one graded modules, and hence an
equivalence of local spectra. This proves in every case that
\[
 \mathbb S_{K(2)}\text{ is a retract of }E_0\otimes DE_0.
 \tag{D.2.11i}
\]
Here $D$ is duality in $\Sp_{K(2)}$ and $\otimes$ its
localized tensor product.

Fix $r\geq0$ and an arbitrary $U\in\Sp_{K(2)}$. In ordinary
spectra form
\[
 V=\tau_{\geq N}(DE_0\otimes U),\qquad H=C_0\wedge V,
 \tag{D.2.11j}
\]
where $N$ is large enough that $H$ is at least
$\max\{r+2,d_3+1\}$-connective. A bounded-above spectrum
is $T(2)$-acyclic: if $T(2)$ is the telescope of a positive
$v_2$ self-map on a finite spectrum $T$, then the terms
$\Sigma^{-kd}(T\wedge W)$ have zero homotopy in any fixed
degree once $k$ is large. The positive self-map is supplied by
\cite[Theorem~9]{HopkinsSmith1998}. Thus their telescope
vanishes for bounded-above $W$, and $W$ is also
$K(2)$-acyclic. Applying this to the complementary
Postnikov section in \textup{(D.2.11j)} gives
\[
 L_{K(2)}V\simeq DE_0\otimes U,\qquad
 L_{K(2)}H\simeq E_0\otimes DE_0\otimes U.
 \tag{D.2.11k}
\]

The spectrum $H$ belongs to $\mathcal I_p=\loc(C_0)$:
this localizing subcategory is a tensor ideal in spectra.
It has $p$-primary integral homology, and its chosen
connectivity allows D.2.9 to be applied. Consequently
$\Omega^\infty H$ is ordinarily $A_r$-acyclic. Its
rational and $T(1)$-acyclicity also permits
\textup{(D.2.11h)}. If $Z$ satisfies the premise of
\textup{(D.2.11b)}, these facts give
\[
 \Map_{\operatorname{Lie}(\Sp_{K(2)})}
 \bigl(\operatorname{triv}(E_0\otimes DE_0\otimes U),F_Z\bigr)
 \simeq *.
\]
By \textup{(D.2.11i)}, $\operatorname{triv}U$ is a retract
of this source. Hence $F_Z$ is orthogonal to every trivial
Lie algebra. The already proved categorical assertion makes
$F_Z=0$, and fullness of the realization gives $Z\simeq *$.
The auxiliary $N$ can be increased for every fixed $r$, so
this proves \textup{(D.2.11b)} for all $r\geq0$, including
those below the periodic model's connectivity cutoff.

The proof covers every periodic target, even those with an
unbounded Lie weight filtration. The finite filtration was
used only for a set of free generators of the category.
It does not imply that ordinary $A_r$-nullity is preserved
by periodic localization, and therefore does not settle
\textup{(D.2.9b)}.
\end{proof}

\begin{claim}{D.2.12}{Proved}{The suspended stabilization unit has no retraction}
Let $B=K(\FF_2,2)$ and $A_r=Q\Sigma^rB$, where
$Q=\Omega^\infty\Sigma^\infty$. For every $r\geq1$ and every
pointed map
\[
 g:\Sigma A_r\longrightarrow\Sigma^{r+1}B,
\]
the homomorphism
\[
 \widetilde H_{r+3}(g;\FF_2)=0.
 \tag{D.2.12a}
\]
Consequently the suspended stabilization unit
\[
 \Sigma\eta:\Sigma^{r+1}B\longrightarrow\Sigma A_r
 \tag{D.2.12b}
\]
has no homotopy retraction. Its canonical stable retraction,
obtained from the suspension-spectrum/infinite-loop-space
counit, cannot be represented by a map of these spaces.
This conclusion does not assert that the entire mapping space
is contractible or determine $P_{A_r}\Sigma^{r+1}B$.
\end{claim}
\begin{proof}
All homology and cohomology in this proof have coefficients in
$\FF_2$. Put $X=\Sigma^rB$, $n=r+2\geq3$, and
$V=\widetilde H_*(X)$. Write $\operatorname{Sq}_*^i$ for the
homology operation dual to $\operatorname{Sq}^i$.

\medskip\noindent
\textbf{The four low-degree cohomology groups.}
Serre's calculation, including its explicit $K(\FF_2,2)$
example, gives
\[
 H^2(B)=\FF_2\{u\},\quad
 H^3(B)=\FF_2\{v\},\quad
 H^4(B)=\FF_2\{u^2\},\quad
 H^5(B)=\FF_2\{uv,w\},
 \tag{D.2.12c}
\]
where $v=\operatorname{Sq}^1u$ and
$w=\operatorname{Sq}^2v$; see
\cite[\S2, no.~8, Theorem~2 and no.~9, pp.~203--204]{I_Serre1953}.
Let $a,b,d\in V$ be the suspended homology classes dual to
$u,v,w$, respectively, with $d$ zero on the suspended class
$uv$. Their degrees are $n,n+1,n+3$, and
\[
 \operatorname{Sq}_*^1b=a,\qquad
 \operatorname{Sq}_*^2b=0,\qquad
 \operatorname{Sq}_*^1d=0,\qquad
 \operatorname{Sq}_*^2d=b.
 \tag{D.2.12d}
\]
For the third equality use
$\operatorname{Sq}^1(u^2)=0$; the other equalities follow
directly from \textup{(D.2.12c)}, naturality, and compatibility
of Steenrod operations with suspension. These compatibility and
Cartan formulas are recorded in
\cite[\S1, nos.~2.1--2.3, p.~199]{I_Serre1953}.
Moreover
\[
 \operatorname{Sq}_*^1c=0\quad(c\in V_{n+2}).
 \tag{D.2.12e}
\]
Indeed its value is a multiple of $b$, and applying
$\operatorname{Sq}_*^1$ again gives that multiple of $a$,
whereas $(\operatorname{Sq}_*^1)^2=0$.

\medskip\noindent
\textbf{The tensor algebra and the source classes.}
The James construction identifies
\[
 H_*(\Omega\Sigma X)\cong T(V)
       =\bigoplus_{k\geq0}V^{\otimes k}
 \tag{D.2.12f}
\]
as a Pontryagin algebra. The generators are the images of
$X\to\Omega\Sigma X$. This follows from the reduced-product
homology calculation and the canonical equivalence in
\cite[\S3, especially (3.3)--(3.4); Theorem~5.6,
Corollary~5.7 and Lemma~5.5]{I_James1955}.
Since $r\geq1$, the reduced diagonal of $X$ is zero in homology;
all of $V$ is therefore primitive in this Hopf algebra. The
diagonal on $T(V)$ is the multiplicative diagonal determined by
these primitives. Steenrod operations preserve tensor length:
they preserve $V$ by naturality of the one-letter inclusion,
and their effect on words is given by the Cartan formula.
Let $\ell_k:T(V)\to V^{\otimes k}$ be the projection onto
tensor length $k$.

Denote the images of $a,b,d$ under $X\to A_r$ by
$\bar a,\bar b,\bar d$. They are primitive and satisfy the
same Steenrod relations. The Pontryagin algebra $H_*(A_r)$
is commutative, since $A_r$ is an infinite loop space. Hence
the Cartan formula gives
\[
 \operatorname{Sq}_*^2(\bar b\bar d)=\bar b^2,
 \qquad
 \operatorname{Sq}_*^1(\bar b^2)
   =\bar a\bar b+\bar b\bar a=0.
 \tag{D.2.12g}
\]

Suppose that the adjunct $f:A_r\to\Omega\Sigma X$ of $g$
is nonzero on bottom homology. Then $f_*(\bar a)=a$ and
$f_*(\bar b)=b$: the groups in degrees $n$ and $n+1$ of
the target are one-dimensional, and
$\operatorname{Sq}_*^1b=a$. We will use only naturality of
the diagonal and Steenrod operations for $f$; multiplicativity
of $f_*$ is not required.

\medskip\noindent
\textbf{The coefficient forced by the diagonal.}
Let $\lambda:V\to\FF_2$ extract the coefficient of $b$,
and set $\mu=\lambda\operatorname{Sq}_*^2$.
Thus $\lambda$ is supported in degree $n+1$ and $\mu$ in
degree $n+3$. Relation \textup{(D.2.12e)} says
$\lambda\operatorname{Sq}_*^1=0$.
Put
\[
 t=\ell_1f_*(\bar d),\qquad
 z=\ell_2f_*(\bar b\bar d)\in V\otimes V.
\]
Naturality and preservation of tensor length imply
$\mu(t)=\lambda\ell_1f_*(\bar b)=1$.
Primitivity of $\bar b,\bar d$ gives
\[
 \overline\Delta(\bar b\bar d)
    =\bar b\otimes\bar d+\bar d\otimes\bar b.
\]
Projecting naturality of this equation onto tensor length
$(1,1)$ in the target gives
\[
 z+\tau z=b\otimes t+t\otimes b,
 \tag{D.2.12h}
\]
where $\tau$ interchanges the factors. Here only length-two
words contribute to the left side, because every letter is
primitive. Applying $\lambda\otimes\mu$ yields
\[
 (\lambda\otimes\mu+\mu\otimes\lambda)(z)=1.
 \tag{D.2.12i}
\]
On the other hand, the coefficient of the word $bb$ in
$\operatorname{Sq}_*^2z$ is exactly this expression: the
middle Cartan term vanishes because
$\lambda\operatorname{Sq}_*^1=0$. By
\textup{(D.2.12g)} and naturality, we have proved
\[
 \operatorname{coeff}_{bb}\bigl(\ell_2f_*(\bar b^2)\bigr)=1.
 \tag{D.2.12j}
\]
This argument also covers $n=3$. In that case $f_*(\bar d)$
may contain a length-two term $aa$; the projection defining
$t$ removes it, and it contributes nothing to
\textup{(D.2.12h)}.

\medskip\noindent
\textbf{The incompatible Bockstein.}
The length-two part of total degree $2n+2$ can contain only
$bb$ and terms $ac,ca$ with $c\in V_{n+2}$. Equations
\textup{(D.2.12d)--(D.2.12e)} give
\[
 \operatorname{Sq}_*^1(bb)=ab+ba,\qquad
 \operatorname{Sq}_*^1(ac)=
 \operatorname{Sq}_*^1(ca)=0.
\]
Thus \textup{(D.2.12j)} forces the length-two part of
$\operatorname{Sq}_*^1f_*(\bar b^2)$ to be the nonzero word
sum $ab+ba$. Other tensor lengths cannot cancel it. This
contradicts \textup{(D.2.12g)} and naturality, proving that
$f$ is zero on bottom homology.

The unit $X\to QX$ and the homology suspension
$H_n(\Omega\Sigma X)\to\widetilde H_{n+1}(\Sigma X)$
are isomorphisms in these bottom degrees, by connectivity and
the Hurewicz theorem. Adjunction therefore identifies the
bottom homology map of $f$ with that of $g$. This proves
\textup{(D.2.12a)}. The map \textup{(D.2.12b)} is a bottom
homology isomorphism, so it cannot have a homotopy retraction.
The stable counit is its stable retraction by the adjunction
triangle identity; any space-level representative would also
be a bottom homology isomorphism, which is impossible.

The proof above treats maps from one suspension of $A_r$.
The corollary below extends the low-degree homology conclusion
to every suspension, but does not control successive attachments
in nullification. No conclusion about the value of
$P_{A_r}\Sigma^{r+1}B$ is inferred.
\end{proof}

\paragraph*{Corollary to D.2.12: first cone attachments.}
\begingroup\itshape
Keep $r\geq1$, put $n=r+2$, and let $Y=\Sigma^{r+1}B$.
For every $k\geq0$ and every map $h:\Sigma^kA_r\to Y$,
\[
 \widetilde H_i(h;\FF_2)=0
 \qquad(n+1\leq i\leq n+3).
 \tag{D.2.12k}
\]
In degree $n+4$, its image is contained in the line dual to
the suspended class $uv\in H^5(B;\FF_2)$.
Thus $n+4=r+6$ is the first degree not excluded by these
computations; existence of a map hitting that line is not asserted.
If $C_h$ is the homotopy cofiber of $h$, then
\[
 \widetilde H_i(Y;\FF_2)\longrightarrow
 \widetilde H_i(C_h;\FF_2)
 \quad\hbox{is injective for }n+1\leq i\leq n+3.
 \tag{D.2.12l}
\]
The bottom integral homology homomorphism is also injective.
The same injections hold for the cofiber of any small wedge
of maps of the displayed form, including the first simultaneous
attachment of cones on all maps $\Sigma^kA_r\to Y$.
These statements concern the first attachment to $Y$; no
iteration through subsequent cone attachments is claimed.
\par\endgroup
\begin{proof}
First consider the bottom target degree $n+1$. For $k\geq2$
the source has no homology in that degree, and $k=1$ is
\textup{(D.2.12a)}. For $k=0$, restrict $h$ along the unit
$\eta:X=\Sigma^rB\to QX=A_r$. Since $X$ is
$(n-1)$-connected and $n\geq3$, the Freudenthal suspension
theorem \cite[Corollary~4.24]{FGHHatcher2002} makes this unit an isomorphism on homology in degree
$n+1$. If $h$ were nonzero in that degree, the restricted map
$j=h\eta:X\to Y$ would therefore satisfy
\[
 j^*(\sigma^{r+1}u)=\sigma^rv.
\]
Applying $\operatorname{Sq}^2$ gives
$j^*(\sigma^{r+1}u^2)=\sigma^rw$.
But $\operatorname{Sq}^1(u^2)=0$, whereas the Adem and
top-square identities give
\[
 \operatorname{Sq}^1w
 =\operatorname{Sq}^1\operatorname{Sq}^2v
 =\operatorname{Sq}^3v=v^2\ne0.
\]
Here the nonvanishing is part of Serre's polynomial-algebra
calculation, and the Adem identity is recorded in
\cite[\S4, no.~32.2, p.~224]{I_Serre1953}.
Naturality of $\operatorname{Sq}^1$ is a contradiction.

Write $a_Y,b_Y,c_Y$ for the classes in target degrees
$n+1,n+2,n+3$ dual to the suspended $u,v,u^2$.
They satisfy
\[
 \operatorname{Sq}_*^1b_Y=a_Y,
 \qquad\operatorname{Sq}_*^2c_Y=a_Y.
\]
The already proved vanishing in degree $n+1$ and naturality
force the image in each of the next two degrees to be zero.
In degree $n+4$ there are two basis classes, dual to $uv,w$.
The latter maps to $b_Y$ under $\operatorname{Sq}_*^2$ and
the former maps to zero. Hence the image can involve only
the $uv$-dual class. The homology exact sequence of the
cofiber proves \textup{(D.2.12l)}.

Finally $\widetilde H_{n+1}(Y;\mathbb Z)=\mathbb Z/2$.
The integral homology image of $h$ in that group is zero:
its reduction modulo $2$ injects into mod-$2$ homology by
the universal coefficient sequence, and
\textup{(D.2.12k)} annihilates it. The integral cofiber
sequence gives the asserted bottom injection. Homology takes
the indicated small wedge to a direct sum, so the same
vanishing and cofiber-sequence arguments prove the simultaneous
version. No integral
injection in degree $n+3$ is inferred from mod-$2$
injectivity, and the altered cofiber has not been shown to
satisfy the same Steenrod and loop-Hopf-algebra hypotheses
as $Y$.
\end{proof}

\begin{claim}{D.2.13}{Proved}{A finite obstruction to the canonical Hopf extension}
Put $B=K(\FF_2,2)$, $X=\Sigma B$, $Y=\Sigma^2B$,
$C_mX=\Omega^m\Sigma^mX$, and $S_m=\Sigma^2C_mX$.
Let $j:X\to QX$ be the stabilization unit and set
\[
 h_0=\eta_2\wedge\id_B:\Sigma^2X=\Sigma^3B\longrightarrow Y,
 \qquad \eta_2:S^3\longrightarrow S^2.
 \tag{D.2.13a}
\]
This map is stably essential, although its ordinary homology map is zero.
The suspension--loop counits give compatible extensions $h_m:S_m\to Y$
for $m=1,2$. The canonical $h_2$ extends through relative dimension
eleven in a relative CW approximation of $S_2\to\Sigma^2QX$,
but it does not extend across $S_2\to S_4$.

The coefficient calculation below would force the double adjunct of
any pointed extension $h:\Sigma^2QX\to Y$ of $h_0$ to fail to
preserve the Pontryagin square of the degree-five class $C$ specified
below. The additional multiplication-deviation argument of D.2.16
proves that no ordinary extension exists, already at $S_4$.
The general finite-detection argument remains useful: an extension
exists if and only if $h_0$ extends to every $S_m$, and a negative
answer is detected at a finite stage and on a finite relative subcomplex.
\end{claim}
\begin{proof}
Throughout the homology calculation coefficients are $\FF_2$;
write $\beta=\operatorname{Sq}_*^1$ and $T=\Omega^2Y$.
The argument uses Steenrod naturality for arbitrary maps. It uses
Pontryagin multiplicativity only for the specified double-loop map
$\Omega^2h_0$.

\medskip\noindent
\textbf{The Hopf attachment and its first extensions.}
The cofiber of $h_0$ is $\mathbb CP^2\wedge B$. Let $t$ be the
degree-two class of $\mathbb CP^2$, and let $u\in H^2B$ and
$v=\operatorname{Sq}^1u\in H^3B$. Serre's computation gives
$v\ne0$ and
\[
 \operatorname{Sq}^3(tu)=t^2v\ne0.
 \tag{D.2.13b}
\]
If the stable map were zero, its cofiber would split, compatibly with
the inclusion of $Y$, as $\Sigma^2\Sigma^\infty B\vee
\Sigma^4\Sigma^\infty B$. The unique degree-four class would then
have zero $\operatorname{Sq}^3$, since $\operatorname{Sq}^3u=0$.
This contradicts \textup{(D.2.13b)}. The homology map itself is zero
by the K\"unneth theorem applied to $\eta_2\wedge\id_B$.
The cohomology and operation conventions here are those of
\cite[\S1, nos.~2.1--2.4; \S2, nos.~8--9, pp.~199, 203--204]{I_Serre1953}.

For $0\le m\le2$ let
\[
 r_m=\Sigma^{2-m}\varepsilon^{(m)}:
 \Sigma^2\Omega^m\Sigma^mX\longrightarrow\Sigma^2X,
 \qquad h_m=h_0r_m.
 \tag{D.2.13c}
\]
The adjunction triangle identities show that these extend $h_0$;
naturality of the counit shows their compatibility under stabilization.
In particular the double adjunct of $h_2$ is
$g_2=\Omega^2h_0:C_2X\to T$.

\medskip\noindent
\textbf{The one prescribed label value.}
Use the following labels for classes of $\widetilde H_*B$, and for
their images in $H_*T$ under the unit:
\[
\begin{array}{c|rrrrrrr}
\text{label}&a&b&c&d&e&f&x\\
\text{degree}&2&3&4&5&5&6&6\\
\text{dual monomial}&u&v&u^2&w&uv&u^3&v^2
\end{array}
 \qquad w=\operatorname{Sq}^2v.
 \tag{D.2.13d}
\]
These are the complete positive-degree homology bases through degree
six. Serre's polynomial description and Cartan give
\[
 \beta b=a,\quad \beta c=\beta d=\beta e=\beta f=0,
 \quad\beta x=d+e,\quad \operatorname{Sq}_*^2c=a,
 \quad\overline\Delta c=a\otimes a.
 \tag{D.2.13e}
\]
For example $\operatorname{Sq}^1w=v^2=\operatorname{Sq}^1(uv)$,
which proves the formula for $\beta x$.

Write $\xi=\xi_1$ for the top operation on the homology of an
$E_2$-space, of degree $|\xi z|=2|z|+1$, and
$\{z,z'\}=\lambda_1(z,z')$ for its degree-one Browder bracket.
Let $g_0:X\to T$ be the double adjunct of $h_0$. Then
\[
 C=\sigma c\in H_5X,\qquad g_{0*}(C)=\xi a.
 \tag{D.2.13f}
\]
Here is a verification that does not identify reduced Hopf rotation
with an unreduced BV operation. The once-adjunct
$\varphi_0:\Sigma^2B\to\Omega\Sigma^2B$ has homology image
\[
 \varphi_{0*}(\sigma^2c)=(\sigma a)^2
       \in T(\widetilde H_*\Sigma B).
 \tag{D.2.13g}
\]
Indeed the quadratic part of a minimal Adams--Hilton attaching
differential is the reduced diagonal of the cone. In
$\mathbb CP^2\wedge B$, the product of the two old classes $tu$
is $t^2u^2$, giving precisely the displayed coefficient. The linear
part vanishes because $h_0$ is homology-zero, and a word of length
three has degree at least nine, whereas the input has degree six.
This is the relative suspension-cone form of the Adams--Hilton
calculation in \cite[\S3.1, p.~300; proof of Proposition~7.4,
p.~309]{FGHHalperinLemaire1995}. The tensor-algebra identification
is the James calculation \cite[\S3; Theorem~5.6 and Corollary~5.7]{I_James1955}.

Homology suspension $H_5T\to H_6\Omega Y$ sends $\xi a$ to
$(\sigma a)^2$. The free $E_2$ basis in degree five is
$d,e,ab,\xi a$. The first two have independent one-letter
suspensions, the last has the displayed two-letter suspension,
and $ab$ suspends to zero. Since $a,b$ are primitive,
$\overline\Delta(ab)=a\otimes b+b\otimes a\ne0$.
Thus the primitive kernel of homology suspension in this degree
is zero. The class $g_{0*}(C)$ is primitive because $X$ is a
suspension. This proves \textup{(D.2.13f)}.

\medskip\noindent
\textbf{Two coefficient identities in the target.}
The free $E_2$ homology theorem describes $H_*T$ as the polynomial
algebra on a basis of the free restricted, degree-one Lie algebra
on the labels of $\widetilde H_*B$; use
\cite[Chapter III, Theorems~3.1--3.2 and Lemma~3.8]{I_CohenLadaMay1976}.
In particular pure restricted powers of a label are independent
of the Lie monomials containing a bracket. The exact top-operation
relations, including their Browder terms, are
\[
 \beta\xi z=(|z|\bmod2)z^2+\{\beta z,z\},\qquad
 \operatorname{Sq}_*^2\xi z=\xi(\beta z)+\{z,\operatorname{Sq}_*^2z\}.
 \tag{D.2.13h}
\]
The bracket has the Steenrod Cartan formulas and $\{z,z\}=0$.
These are \cite[Chapter III, Theorems~1.2--1.4,
pp.~216--219]{I_CohenLadaMay1976}; their homology-suspension formula
also supplies the operation used above.

Let $\lambda:H_{10}T\to\FF_2$ select the coefficient of
$(\xi a)^2$, and let $\mu:H_{11}T\to\FF_2$ select the coefficient
of $\xi(\xi a)$. We claim
\[
 \lambda\beta=\mu,\qquad
 \mu\operatorname{Sq}_*^2=0:H_{13}T\longrightarrow\FF_2.
 \tag{D.2.13i}
\]
All basis possibilities can be checked in these degrees as follows.
The degree-five Lie generators are $d,e,\xi a$; hence the pure
top generators in degree eleven are $\xi d,\xi e,\xi(\xi a)$.
Their Bocksteins are respectively $d^2,e^2,(\xi a)^2$ by
\textup{(D.2.13e),(D.2.13h)}. Bocksteins of bracket expressions
retain a bracket factor after the Cartan formula is expanded, so
cannot contribute the specified pure monomial. A decomposable
class could contribute only through a degree-six Bockstein
preimage of $\xi a$. There is none: the complete degree-six
basis is
\[
 f, x, ac, b^2, a^3, \{a,b\},
 \tag{D.2.13j}
\]
whose Bocksteins are respectively $0,d+e,0,0,0,0$.
This proves the first identity.

For the second, products have decomposable Steenrod images,
labels remain labels, and brackets remain bracket expressions.
A pure top generator of degree thirteen has the form $\xi z$
with $z$ a degree-six Lie generator, hence $z=f,x,$ or $\{a,b\}$.
Its square-operation image modulo the bracket ideal is
$\xi(\beta z)$. The displayed degree-six list shows that
$\beta z$ has no $\xi a$ component; it is a combination of
$d,e$. Thus $\xi(\beta z)$ has no $\xi(\xi a)$ component.
This proves the second identity, and in particular
$\lambda\beta\operatorname{Sq}_*^2=0$ on $H_{13}T$.

\medskip\noindent
\textbf{The obstruction already at the fourth stage.}
Use $C$ also for its image in $H_5C_mX$, and put
$A_0=\sigma a\in H_3C_mX$. They satisfy
$\beta C=\beta A_0=0$ and $\operatorname{Sq}_*^2C=A_0$.
At the fourth stage $\xi_3C=Q^8C$ is the \emph{top} operation.
Its exact Nishida formula is
\[
 \operatorname{Sq}_*^2\xi_3C
       =Q^6C+\lambda_3(C,A_0),\qquad
 \beta Q^6C=C^2.
 \tag{D.2.13k}
\]
The omitted possible ordinary term is $Q^7\beta C=0$.
The coefficient of $Q^6C$ is odd. These are the top and non-top
formulas of \cite[Chapter III, Theorems~1.1--1.3,
especially pp.~217--218]{I_CohenLadaMay1976}.
The Bockstein of the Browder correction is zero because both
its inputs are Bockstein-zero. Consequently
\[
 \beta\operatorname{Sq}_*^2\xi_3C=C^2
       \quad\hbox{in }H_{10}C_4X.
 \tag{D.2.13l}
\]
This step retains the top-operation correction; the uncorrected
Nishida formula alone would not justify a fourth-stage assertion.

Suppose $h_2$ extended to $S_4$. Taking double adjuncts would give
an ordinary pointed map $G:C_4X\to T$ restricting to
$g_2=\Omega^2h_0$. This specified restriction is multiplicative,
so \textup{(D.2.13f)} implies
$\lambda G_*(C^2)=\lambda(\xi a)^2=1$.
But Steenrod naturality and \textup{(D.2.13i),(D.2.13l)} give
\[
 \lambda G_*(C^2)
   =\lambda\beta\operatorname{Sq}_*^2G_*(\xi_3C)=0,
 \tag{D.2.13m}
\]
a contradiction. No multiplicativity was assumed for $G$.
The relation detecting this obstruction lies in homology degree
thirteen before the two suspensions.

In the full infinite-loop source the ordinary Nishida identities
give $\operatorname{Sq}_*^2Q^8C=Q^6C$ and $\beta Q^6C=C^2$.
Thus every ordinary extension $g:QX\to T$ would satisfy
\[
 \lambda g_*(C^2)=0,
 \qquad \lambda\bigl(g_*(C)^2\bigr)=1.
 \tag{D.2.13n}
\]
In particular it cannot be an H-map. For an arbitrary map the
first value is not constrained to equal the second; this coefficient
calculation alone therefore does not decide the full extension
question. D.2.16 uses an actual multiplication-deviation map and
additional cohomology identities to exclude every ordinary extension.

\medskip\noindent
\textbf{The positive range through relative dimension eleven.}
Freudenthal makes $B\to\Omega\Sigma B$ three-connected and
$\Omega\Sigma B\to\Omega^2\Sigma^2B$ four-connected.
Their composite is three-connected. Since $\pi_3B=0$, it follows
that $\pi_5Y=\pi_3T=0$; see
\cite[Corollary~4.24]{FGHHatcher2002} for the connectivity form.
Choose a bottom cell $i:S^3\to X$ and set $X'=X/S^3$.
The quotient is three-connected, since $i$ surjects onto
$H_3X=\ZZ/2$ and the quotient is simply connected.
The map $h_0\Sigma^2i$ is null, so the cofiber property gives
$\bar h:\Sigma^2X'\to Y$ with $\bar h\Sigma^2q\simeq h_0$,
where $q:X\to X'$.

The stabilization
$t':\Sigma^2\Omega^2\Sigma^2X'\to\Sigma^2QX'$
is eleven-connected: $\Sigma^2X'$ is five-connected,
Freudenthal gives eleven-connectivity before looping, and
looping twice and suspending twice restore that same bound.
The homotopy pullback
\[
 P=(\Sigma^2QX)\times^h_{\Sigma^2QX'}
                 (\Sigma^2\Omega^2\Sigma^2X')
 \tag{D.2.13o}
\]
therefore maps eleven-connectedly to $\Sigma^2QX$.
Naturality supplies a lift of $S_2$ to $P$. The other projection,
followed by the counit and $\bar h$, gives $P\to Y$ whose
restriction is $h_2$, by counit naturality and
$\bar h\Sigma^2q\simeq h_0$.
Replace $P\to\Sigma^2QX$ by a fibration. Its fiber is
ten-connected, so the specified lift extends over every relative
cell of dimension at most eleven. Composition with $P\to Y$
gives the asserted positive range. This argument uses an actual
pullback; it does not suspend a fiber sequence as a fiber sequence.

\medskip\noindent
\textbf{Finite detection for arbitrary choices of extensions.}
Serre's class theorem gives finite $2$-primary homology for $B$
in every positive degree and finite $2$-primary homotopy groups
for the simply connected space $Y$. The exact interfaces are
\cite[Chapter I, \S6.4; Chapter II, Proposition~8; Chapter III,
Theorem~1 and Corollaries~1--2, pp.~265, 271--274]{I_SerreClasses1953}.
All the sources above admit countable CW models. For a finite
CW complex $K$, Postnikov obstruction theory makes $[K,Y]_*$
finite and makes $\pi_1\Map_*(K,Y)$ finite at every component.
For an exhaustion $K=\bigcup K_s$ by finite subcomplexes,
the mapping-tower argument consequently gives
\[
 [K,Y]_*\cong\lim_s[K_s,Y]_*.
 \tag{D.2.13p}
\]
The obstruction to injectivity is a $\lim^1$ of finite groups;
its tower is Mittag--Leffler, so it vanishes. Surjectivity follows
by realizing compatible maps and homotopies on the telescope.
The right side gives $[K,Y]_*$ a compact Hausdorff profinite
topology. Restriction along any map of countable CW sources is
continuous, since its restriction to a finite subcomplex factors
through a finite subcomplex after cellular approximation.

Let $E_m\subseteq[S_m,Y]_*$ consist of the classes restricting
to $[h_0]$ on $S_0=\Sigma^2X$. These are closed compact sets, with
continuous restriction maps $E_{m+1}\to E_m$. If every $E_m$
is nonempty, the compatibility conditions in $\prod_mE_m$ have
the finite intersection property: any finite list is satisfied
by choosing a class at its largest index and restricting it.
Hence $\lim_mE_m$ is nonempty, without a surjectivity assumption
on the bonding maps. The mapping-telescope surjection realizes
such a compatible family on
$\operatorname{hocolim}_mS_m\simeq\Sigma^2QX$.
Conversely an extension clearly restricts to all finite stages.
At a fixed stage, apply the same finite Postnikov and compactness
argument to an exhaustion of a mapping-cylinder CW pair for
$S_0\to S_m$. Failure to extend is detected on a finite subpair.
This equivalence concerns \emph{all} choices in $E_m$. The
obstructed canonical choice above alone does not prove $E_4$ empty;
that stronger conclusion is proved separately in D.2.16.
\end{proof}

\begin{claim}{D.2.14}{Proved}{The odd-primary Pr\"ufer ideal is the full primary ideal}
Let
\[
 D=\Sigma^\infty\mathbb{C}P^\infty,
 \qquad J_p=\operatorname{loc}(D\wedge S/p),
\]
where the suspension spectrum is reduced. The following statements hold.
\begin{enumerate}
\item If an $E_1$-ring $R$ satisfies $F(D,R)\simeq0$ and the image
  $\eta_R\in\pi_1R$ of the stable Hopf class is zero, then $R\simeq0$.
\item For every prime $p$, the inclusion
  $J_p\subseteq\operatorname{loc}(S/p)$ is proper if and only if there
  is a nonzero derived $p$-complete $E_\infty$-ring $R$ such that
  $F(D,R)\simeq0$.
\item For every odd prime $p$,
  \[
   \operatorname{loc}\bigl(\Sigma^\infty B C_{p^\infty}\bigr)
      =J_p=\operatorname{loc}(S/p).
   \tag{D.2.14a}
  \]
  In particular this ideal contains the finite Moore spectrum $S/p$.
\end{enumerate}
The first statement uses an actual trivialization of the coefficient
line bundle. It does not infer Eilenberg--Moore convergence from the
finiteness of $S^1$.
\end{claim}

\begin{proof}
\emph{The Euler sequences and the Thom fiber.}
Put $G=S^1$, and let $V=\mathbb C$ be its standard complex
representation. The equivariant Euler cofiber sequence is
\[
 \Sigma^\infty G_+\longrightarrow S\longrightarrow S^V.
 \tag{D.2.14b}
\]
Taking homotopy orbits gives
\[
 S\longrightarrow\Sigma^\infty BG_+
       \longrightarrow\operatorname{Th}(L),
\]
where $L$ is the universal complex line over $BG=\mathbb CP^\infty$.
The first map comes from $EG\times_GG=EG\to BG$; since $EG$ is
contractible, it is the inclusion of a basepoint up to homotopy.
Consequently $\operatorname{Th}(L)\simeq D$.

Tensoring \textup{(D.2.14b)} with $S^{-V}$ before taking homotopy
orbits instead gives
\[
 S^{-2}\xrightarrow{j}\operatorname{Th}(-L)
   \longrightarrow\Sigma^\infty BG_+
   \xrightarrow{\delta}S^{-1}.
 \tag{D.2.14c}
\]
Indeed $G_+\wedge S^{-V}$ is induced from the underlying sphere
$S^{-2}$. The map $j$ is a Thom-fiber inclusion: explicitly its source
before collapsing $EG$ is
$EG_+\wedge_G(G_+\wedge S^{-V})\simeq EG_+\wedge S^{-2}$,
and restriction to any point of $EG$ includes the corresponding
fiber in the Thom colimit. This also recovers the usual
circle-transfer sequence; see the first displayed fibration in
\cite[\S1]{EH5MadsenSchlichtkrull2000}.

\emph{The unit extends when $\eta_R=0$.}
Let $u:S^{-2}\to\Sigma^{-2}R$ be the shifted unit. Its obstruction to
extension over $j$ is the composite
\[
 \Sigma^\infty BG_+\xrightarrow{\delta}S^{-1}
      \xrightarrow{\Sigma u}\Sigma^{-1}R.
 \tag{D.2.14d}
\]
Since $\Sigma^\infty BG_+\simeq S\vee D$ and $F(D,R)=0$,
restriction to the basepoint gives an isomorphism
\[
 [\Sigma^\infty BG_+,\Sigma^{-1}R]
       \xrightarrow{\ \cong\ }[S,\Sigma^{-1}R].
\]
The restriction of \textup{(D.2.14d)} is the image under the unit of
an element of $\pi_1S\cong\mathbb Z/2$, generated by $\eta$.
It is therefore zero when $\eta_R=0$. Exactness of mapping out of
\textup{(D.2.14c)} supplies
\[
 t:\operatorname{Th}(-L)\longrightarrow\Sigma^{-2}R,
 \qquad t j=u.
 \tag{D.2.14e}
\]

\emph{The extension trivializes the full local system.}
Extend $t$ $R$-linearly to a map of left $R$-modules
$R\wedge\operatorname{Th}(-L)\to\Sigma^{-2}R$.
The Thom object is the colimit of the local system
$R\wedge S^{-L}$ of left $R$-modules over $BG$.
By the colimit adjunction, this map is the same as a natural
transformation
\[
 R\wedge S^{-L}\longrightarrow
          \operatorname{const}_{BG}(\Sigma^{-2}R).
 \tag{D.2.14f}
\]
Its value at the selected basepoint is the identity of
$\Sigma^{-2}R$, by \textup{(D.2.14e)} and the identification of $j$.
Every point of $BG$ can be joined to that basepoint. Naturality along
a path conjugates the corresponding fiber maps by equivalences.
Thus every component of \textup{(D.2.14f)} is an equivalence, and the
natural transformation itself is an equivalence of local systems.
This is the Thom-orientation criterion of
\cite[Definitions 2.20 and 2.22, Theorem 2.24]{EH5ABGHR2014}; it
retains the entire coherent sphere-bundle action.

For the stated $E_1$ generality, dualize pointwise using
$F_R(-,R)$ from left to right $R$-modules. Sphere lines are
dualizable, so \textup{(D.2.14f)} gives an equivalence of right
$R$-module local systems
\[
 S^L\wedge R\simeq
          \operatorname{const}_{BG}(\Sigma^2R).
\]
This step needs no commutative multiplication on $R$.
Taking colimits and then the augmentation $BG_+\to S$ produces a
map $\operatorname{Th}(L)\wedge R\to\Sigma^2R$ whose restriction
to the selected fiber is the identity. Precomposition with the unit
therefore gives
\[
 D\simeq\operatorname{Th}(L)\longrightarrow\Sigma^2R
\]
restricting to the unit on the bottom Thom sphere $S^2$.
But $F(D,R)=0$ makes this map null. The unit of $R$ is consequently
null, so unitality gives $\operatorname{id}_R=0$ and hence $R=0$.
This proves the first statement.

\emph{The exact local-unit criterion.}
Every localizing subcategory of spectra is a tensor ideal: for a
fixed object $A$ in the subcategory, the spectra $X$ with
$A\wedge X$ in that subcategory form a localizing subcategory
containing $S$, hence all spectra. Since $J_p$ is generated by one
object, it is the kernel of an accessible exact localization $L$.
Its equivalences are preserved by tensoring with arbitrary spectra.
The localized category therefore has its localized symmetric
monoidal structure, and its inclusion in spectra is lax symmetric
monoidal. Its unit gives an $E_\infty$-ring $R_0=L(S)$; this does
not assert that $L$ is smashing. These are precisely the monoidal
localization statements in
\cite[Example 2.2.1.7, Proposition 2.2.1.9]{FGHLurieHA}.

If $J_p$ is proper in $\operatorname{loc}(S/p)$, then
\[
 R_0/p\simeq L(S/p)\ne0.
\]
Let $C=\operatorname{fib}(S\to S[1/p])$. For any spectrum $E$,
the fiber sequence
\[
 F(S[1/p],E)\longrightarrow E\longrightarrow F(C,E)
\]
exhibits $F(C,E)$ as its derived $p$-completion: the first term is
$p$-invertible, and the last is right orthogonal to $S[1/p]$ because
$C\wedge S[1/p]=0$. This is orthogonality to all $p$-invertible
spectra, which are generated by $S[1/p]$.
In particular $F(C,E)/p\simeq E/p$.
Derived $p$-completion is a monoidal localization, since
$S/p$-equivalences are tensor-stable. Hence
\[
 R=\widehat{R_0}_p\simeq F(C,R_0)
\]
is an $E_\infty$-ring, with $R/p=R_0/p\ne0$.
It is still $J_p$-local, since
\[
 F(D\wedge S/p,F(C,R_0))
   \simeq F(D\wedge S/p\wedge C,R_0)=0.
\]
It follows that $p$ acts invertibly on $F(D,R)$.
This function spectrum is itself derived $p$-complete: if $X$ is
$p$-invertible, then
$F(X,F(D,R))\simeq F(X\wedge D,R)=0$.
Being both $p$-invertible and derived $p$-complete, it is zero.
Thus a proper inclusion produces the required nonzero ring $R$.

Conversely, if such a ring $R$ exists, then
$F(D\wedge S/p,R)=0$, so $R$ is $J_p$-local.
Equality $J_p=\operatorname{loc}(S/p)$ would make multiplication
by $p$ invertible on $R$. Its derived $p$-completeness would then
force $R=0$, a contradiction. This proves the second statement.
No multiplicative structure on the spectral cofiber $R_0/p$ has
been assumed.

\emph{Odd primes and the Pr\"ufer comparison.}
If $p$ is odd, multiplication by $2$ is invertible on every derived
$p$-complete spectrum. Hence $\eta_R=0$, and the first statement
excludes the ring in the second. Thus $J_p=\operatorname{loc}(S/p)$.

For completeness, write $P=\Sigma^\infty BC_{p^\infty}$.
The abelian-group extension
$0\to\mathbb Z\to\mathbb Z[1/p]\to C_{p^\infty}\to0$
gives a fiber sequence
\[
 K(\mathbb Z[1/p],1)\longrightarrow BC_{p^\infty}
       \longrightarrow K(\mathbb Z,2).
\]
The first space is the telescope of the degree-$p$ maps on $S^1$,
so it is mod-$p$ acyclic. The map $P\to D$ is therefore an
$H\mathbb F_p$-equivalence. Its cofiber $Q$ is connective.
The groups $H_*(Q;\mathbb Z)$ have multiplication by $p$ invertible,
by the mod-$p$ coefficient sequence. Consequently $Q/p$ has zero
integral homology, and stable Hurewicz for a bounded-below spectrum
gives $Q/p=0$. Thus $P/p\simeq D/p$.
This is also the Pr\"ufer-to-circle comparison used in
\cite[proof of Theorem 1.5]{EH5Levy2025}.

The reduced homology of $BC_{p^\infty}$ is $p$-primary, as follows
either from its filtered union of finite cyclic groups or from the
same extension. Hence $P[1/p]=0$, by stable Hurewicz.
The cofiber of $P\to P[1/p]$ is the filtered homotopy colimit of
the cofibers $P/p^r$ and is $\Sigma P$.
Each $P/p^r$ is a finite extension of copies of $P/p$; conversely
$P/p$ lies in $\operatorname{loc}(P)$.
Therefore
\[
 \operatorname{loc}(P)=\operatorname{loc}(P/p)
     =\operatorname{loc}(D/p)=J_p,
\]
which completes \textup{(D.2.14a)}.
\end{proof}

At $p=2$, the second statement is still an exact criterion, and
the first identifies the possible obstruction $\eta_R\ne0$.
D.2.15 removes this obstruction using the reduced circle transfer
and the endomorphism ring of $R\wedge C\eta$, proving the equality
at $2$ as well. The action of $\eta$ on $C\eta$ itself is nonzero;
the transfer is what makes its coefficient version vanish.

\begin{claim}{D.2.15}{Proved}{Reduced complex projective space generates spectra; the Pr\"ufer ideal at every prime}
Let $D=\Sigma^\infty\mathbb CP^\infty$ be the reduced suspension
spectrum. For every $E_1$-ring spectrum $R$,
\[
 F(D,R)\simeq0\quad\Longrightarrow\quad R\simeq0.
 \tag{D.2.15a}
\]
Consequently
\[
 \begin{gathered}
 \operatorname{loc}(D)=\mathrm{Sp},\\
 \operatorname{loc}\bigl(\Sigma^\infty BC_{p^\infty}\bigr)
   =\operatorname{loc}(D\wedge S/p)=\operatorname{loc}(S/p)
 \end{gathered}
 \tag{D.2.15b}
\]
for every prime $p$, including $2$. In particular the degree-one
Pr\"ufer generator has the full finite $p$-primary thick part.
\end{claim}

\begin{proof}
The first part of D.2.14 proves \textup{(D.2.15a)} whenever the image
of $\eta$ in $\pi_1R$ is zero. We first remove this hypothesis for
derived $2$-complete rings, and then remove the completeness
hypothesis.

\emph{The bottom reduced circle transfer is an odd multiple of $\nu$.}
Use the Euler--Thom sequence \textup{(D.2.14c)} and split off the
basepoint of $\mathbb CP^\infty_+$:
\[
 S^{-2}\longrightarrow\operatorname{Th}(-L)
  \longrightarrow S\vee D\xrightarrow{\delta}S^{-1}.
\]
Let $t:D\to S^{-1}$ be the restriction of $\delta$, let
$i:S^2\to D$ include the bottom cell, and put
\[
 b=\Sigma(ti):S^3\longrightarrow S.
 \tag{D.2.15c}
\]
We will show that $b=k\nu$ in
$\pi_3S_{(2)}\cong\mathbb Z/8$ with $k$ odd. Here $\nu$ is the
stable quaternionic Hopf class.

There are natural maps of homotopy fibers
\[
 F_{ti}\longrightarrow F_t\longrightarrow F_\delta
       \simeq\operatorname{Th}(-L),
 \qquad F_{ti}\simeq\Sigma^{-2}Cb,
 \tag{D.2.15d}
\]
where $Cb=\operatorname{cofib}(b:S^3\to S)$.
Each of the two maps in \textup{(D.2.15d)} induces an isomorphism
on mod-$2$ cohomology in degrees $-2$ and $2$. This follows directly
from the fiber-sequence cohomology exact sequences: the maps
$S^2\to D\to S\vee D$ are isomorphisms in degree $2$, all their
sources and targets have zero cohomology in degrees $-2,-1$, and
the common target $S^{-1}$ supplies the degree-$-2$ fiber class.
The extra basepoint summand changes only degree $0$ here.

Let $u\in H^{-2}(\operatorname{Th}(-L);\mathbb F_2)$ be the
virtual Thom class and let $x=c_1(L)\bmod2$. The Thom identity and
the Whitney product formula give
\[
 w(-L)=(1+x)^{-1}=1+x+x^2+\cdots,
 \qquad \operatorname{Sq}^4u=x^2u\ne0.
 \tag{D.2.15e}
\]
The ordinary Thom identity is given in
\cite[\S8, pp.~91--93]{EH5MilnorStasheff1974}.
Its use for $-L$ requires no additional orientation assertion:
over $\mathbb CP^2$, choose a finite-dimensional complement $L^\perp$
to $L$ in a trivial complex bundle, apply the ordinary identity to
$L^\perp$, and desuspend by the dimension of that trivial bundle.
The Whitney product formula gives exactly the inverse total class
in \textup{(D.2.15e)}.

Pulling \textup{(D.2.15e)} through \textup{(D.2.15d)} shows that
$\operatorname{Sq}^4$ is nonzero from the bottom to the top class
of $Cb$. Thus $b$ has mod-$2$ Hopf invariant one. Since
$\pi_3S_{(2)}=\mathbb Z/8\{\nu\}$, this says precisely that its
coefficient $k$ is odd; see
\cite[Example 7.17, pp.~45--46, and equation (8.1)]{EH5Adams1966}.
Only this parity is needed, so no choice of sign or normalization
of the circle transfer is involved.

\emph{A derived $2$-complete ring.}
Suppose now that $R$ is derived $2$-complete and $F(D,R)=0$.
The composite of $t$ with the shifted unit
$S^{-1}\to\Sigma^{-1}R$ is null. Its bottom restriction makes
$b_R=0$. All odd integers act invertibly on a derived $2$-complete
spectrum, so \textup{(D.2.15c)--(D.2.15e)} imply
\[
 \nu_R=0.
 \tag{D.2.15f}
\]

Let $C\eta=\operatorname{cofib}(\eta:S^1\to S)$ and set
$M=R\wedge C\eta$, viewed as a left $R$-module. The stable action
of $\eta$ on $C\eta$ factors as
\[
 \eta_{C\eta}:\Sigma C\eta\longrightarrow S^3
       \xrightarrow{\nu}S\longrightarrow C\eta;
 \tag{D.2.15g}
\]
this is \cite[equation (8.4), p.~1279]{EH5AngeltveitRognes2005}.
For the implication needed here one can also derive a possibly
non-normalized version directly. The composite with the suspended
bottom inclusion is zero, so $\eta_{C\eta}$ factors through $S^3$.
The map $\pi_3S\to\pi_3C\eta$ is surjective, since multiplication
by $\eta$ from $\pi_1S$ to $\pi_2S$ is an isomorphism. Thus, at
$2$, the remaining factor comes from a multiple of $\nu$.
Either description and \textup{(D.2.15f)} give
$\eta_M=0$ as an $R$-linear map.

Consider the unital $E_1$-ring
\[
 A=\operatorname{End}_R(M).
\]
Because $C\eta$ is finite, its underlying spectrum is
\[
 A\simeq(C\eta)^\vee\wedge R\wedge C\eta.
 \tag{D.2.15h}
\]
The class of spectra $Z$ with $F(D,Z)=0$ is stable and closed
under finite colimits and retracts. It therefore contains
$R\wedge K$ for every finite spectrum $K$. In particular
\textup{(D.2.15h)} gives $F(D,A)=0$.
The image of $\eta$ under the unit of $A$ is exactly the
$R$-linear self-map $\eta_M$, so $\eta_A=0$.
D.2.14(1), which applies to $E_1$-rings, now gives $A=0$.
Its unit represents $\operatorname{id}_M$, hence $M=0$.

It follows that $\eta:\Sigma R\to R$ is an equivalence, since its
cofiber is $M$. But $\eta^4=0$ in the stable sphere; this relation
is recalled explicitly in
\cite[proof of Lemma 8.11, p.~1285]{EH5AngeltveitRognes2005}.
An invertible nilpotent action forces $R=0$. This proves
\textup{(D.2.15a)} for derived $2$-complete $E_1$-rings.

\emph{An arbitrary ring.}
Let $R$ now be any $E_1$-ring with $F(D,R)=0$, and put
$C=\operatorname{fib}(S\to S[1/2])$.
As established in the local-unit argument of D.2.14,
\[
 \widehat R_2\simeq F(C,R),\qquad
 \widehat R_2/2\simeq R/2,
 \tag{D.2.15i}
\]
and derived $2$-completion is a monoidal localization, so
$\widehat R_2$ inherits an $E_1$-ring structure. Also
\[
 F(D,\widehat R_2)\simeq F(D\wedge C,R)=0.
\]
Indeed $\operatorname{loc}(D)$ is a tensor ideal, and $R$ is right
orthogonal to that whole localizing subcategory. The completed
case makes $\widehat R_2=0$. By \textup{(D.2.15i)}, $R/2=0$,
so multiplication by $2$ is invertible on $R$. Since $2\eta=0$,
we have $\eta_R=0$; D.2.14(1) gives $R=0$.

\emph{From ring detection to actual localizing generation.}
The subcategory $\mathcal I=\operatorname{loc}(D)$ is generated by
one object and is a stable tensor ideal. Its accessible exact
localization therefore has a localized symmetric monoidal
structure; the inclusion of local objects is lax symmetric
monoidal. In particular the underlying spectrum of the localized
unit, $L_{\mathcal I}S$, is a unital $E_\infty$-ring and is
$D$-null; see
\cite[Example 2.2.1.7, Proposition 2.2.1.9]{FGHLurieHA}.
The detector \textup{(D.2.15a)} gives $L_{\mathcal I}S=0$.
Thus $S\in\mathcal I$, and $\mathcal I=\mathrm{Sp}$.
This argument does not require the localization to be smashing.

Finally smash a localizing construction of $S$ from $D$ with
$S/p$. The functor $-\wedge S/p$ is exact and preserves all
colimits, so it proves
$S/p\in\operatorname{loc}(D\wedge S/p)$.
The reverse inclusion follows from the tensor-ideal property of
$\operatorname{loc}(S/p)$. The Pr\"ufer comparison
\[
 \operatorname{loc}(\Sigma^\infty BC_{p^\infty})
       =\operatorname{loc}(D\wedge S/p)
\]
was proved for every prime in D.2.14. This completes
\textup{(D.2.15b)}.
\end{proof}

For historical comparison, the equality of smash Bousfield classes
$\langle\mathbb CP^\infty\rangle=\langle S\rangle$ was already
reported as a theorem of Hopkins in the discussion preceding
\cite[Conjecture 8.7]{JHoveyPalmieri1999}.
The proof above establishes localizing generation directly; it
does not assume the general comparison between Bousfield order
and localizing membership discussed in that paper's Section~9.

\begin{claim}{D.2.16}{Proved}{A fourth-stage obstruction to the Hopf coefficient}
Let $B=K(\FF_2,2)$, $X=\Sigma B$, $Y=\Sigma^2B$, and
$C_mX=\Omega^m\Sigma^mX$. Write $a\in H_2(B;\FF_2)$ for
the bottom class, $c\in H_4(B;\FF_2)$ for the class dual to
the square of the fundamental cohomology class, and $C=\sigma c$.
For every ordinary pointed map
\[
 G:C_4X\longrightarrow T=\Omega^2Y,
\]
the coefficient of the top $E_2$ generator $\xi_1a$ in
$G_*(C)$ is zero. In particular, the map
\[
 h_0=\eta_2\wedge\id_B:\Sigma^3B\longrightarrow\Sigma^2B
 \tag{D.2.16a}
\]
does not extend across
$\Sigma^2X\to\Sigma^2C_4X$, for any choice of ordinary pointed
extension. Consequently it does not extend across
$\Sigma^2X\to\Sigma^2QX$.

This excludes the particular homology-zero Hopf attachment whose
cofiber is $\mathbb CP^2\wedge B$. It does not assert survival
through all subsequent $Q\Sigma B$-nullification attachments.
\end{claim}
\begin{proof}
All homology and cohomology in this proof have coefficients $\FF_2$.
Put $T=\Omega^2Y$, and let $g_0:X\to T$ be the double adjunct of
$h_0$. We use the homology labels $a,b,c,d,e,f,x$ of D.2.13,
dual respectively to $u,v,u^2,w,uv,u^3,v^2$ in
$H^*B$, where
\[
 |u|=2,\qquad v=\operatorname{Sq}^1u,\qquad
 w=\operatorname{Sq}^2v,\qquad t=\operatorname{Sq}^4w.
 \tag{D.2.16b}
\]
Their degrees are $2,3,5,9$ respectively. Serre's calculation
makes these independent polynomial generators; in particular
$t\ne0$ \cite[\S2, nos.~8--9, pp.~203--204]{I_Serre1953}.
Write $\beta=\operatorname{Sq}_*^1$ on homology,
$\xi=\xi_1$ for the top $E_2$ homology operation, and
$\{\ ,\ \}=\lambda_1$ for the Browder bracket.

\medskip\noindent
\textbf{The forced failure of multiplicativity.}
Set $C=\sigma c\in H_5X$, and use the same letter for its image
in the loop stages. The calculation of D.2.13 gives
\[
 g_{0*}(C)=\xi a.
 \tag{D.2.16c}
\]
Let $L\in H^{10}T$ be the polynomial-basis functional dual to
$(\xi a)^2$. The fourth-stage calculation in D.2.13 proves, for
\emph{every} ordinary map $G:C_4X\to T$,
\[
 \langle L,G_*(C^2)\rangle=0.
 \tag{D.2.16d}
\]
Here $C^2$ is the Pontryagin square. Explicitly,
$\beta\operatorname{Sq}_*^2\xi_3C=C^2$ in the fourth stage,
where the top Nishida relation's Browder term has zero
Bockstein. On the target, the coefficient of $(\xi a)^2$ in
$\beta\operatorname{Sq}_*^2H_{13}T$ is zero. Thus
\textup{(D.2.16d)} uses only Steenrod naturality for $G$.

\medskip\noindent
\textbf{Exact cohomology identities in the target.}
In addition to $L$, define polynomial-basis dual classes
\[
\begin{array}{c|ccccc}
\text{class}&U&V&W&L&L_b\\
\text{degree}&5&7&9&10&14\\
\text{dual monomial}&\xi a&\xi b&\xi c&(\xi a)^2&(\xi b)^2.
\end{array}
 \tag{D.2.16e}
\]
We will use the following four identities:
\[
 \operatorname{Sq}^2L=UV,\qquad
 \operatorname{Sq}^4L=L_b+UW,\qquad
 \operatorname{Sq}^2L_b=0,\qquad
 \operatorname{Sq}^2\operatorname{Sq}^1L=0.
 \tag{D.2.16f}
\]
The last is the target coefficient identity of D.2.13. We check
all possibilities for the other three.

The free $E_2$ homology theorem describes $H_*T$ as a polynomial
algebra on a free restricted degree-one Lie basis on the labels
of $\widetilde H_*B$; see
\cite[Chapter III, Theorems~3.1--3.2 and Lemma~3.8]{I_CohenLadaMay1976}.
Every pure top generator $\xi z$ has odd degree. Steenrod
operations send labels to labels, and send bracket expressions
to the bracket ideal. They cannot send a positive-degree
factor to a nonzero scalar. Thus a pure two-factor polynomial
coefficient in an even-degree Steenrod image must come from an
initial product of two Lie generators. No top generator exists
in that even input degree.

The exact top Nishida formula gives
\[
 \operatorname{Sq}_*^2\xi b=\xi a,\qquad
 \operatorname{Sq}_*^2\xi c=\{c,a\},\qquad
 \operatorname{Sq}_*^4\xi c=\xi a.
 \tag{D.2.16g}
\]
For the last identity, use
$\beta c=0$, $\operatorname{Sq}_*^2c=a$, and
$\operatorname{Sq}_*^3c=\operatorname{Sq}_*^4c=0$.
The surviving ordinary term is $Q^3a=\xi a$ and all the
Browder terms vanish. These are the original formulas of
\cite[Chapter III, Theorems~1.2--1.4, pp.~216--219,
especially pp.~217--218]{I_CohenLadaMay1976}.

There is no degree-six Bockstein preimage of the indecomposable
$\xi a$. The full degree-six basis is
\[
 f,x,ac,b^2,a^3,\{a,b\},
\]
with Bockstein images $0,d+e,0,0,0,0$. There is also no
degree-eight Bockstein preimage of the indecomposable $\xi b$:
even-degree indecomposables are labels or brackets, and their
Bocksteins cannot become a pure top generator; a decomposable
class has decomposable Bockstein.

The coefficient of $(\xi a)^2$ in
$\operatorname{Sq}_*^2H_{12}T$ is consequently exactly the
coefficient of $\xi a\cdot\xi b$ in its input. The partition
$1+1$ would require the absent degree-six Bockstein preimage.
Likewise the coefficient of $(\xi a)^2$ in
$\operatorname{Sq}_*^4H_{14}T$ is the sum of the input
coefficients of $(\xi b)^2$ and $\xi a\cdot\xi c$.
The partitions $1+3$ and $3+1$ again require the absent
Bockstein preimage. The partitions $2+2$ and $0+4$ are
precisely those displayed in \textup{(D.2.16g)}.

These two product-coefficient functionals are exactly the cup
products $UV$ and $UW$. This point requires the coalgebra, since
$c$ is not primitive. Coproducts of labels contain only labels;
coproducts of Browder generators retain a bracket in at least
one factor; and there are no even-degree top generators in
degrees twelve or fourteen. For a product to have the two
specified indecomposable top generators in its tensor factors,
its factors must therefore be exactly the indicated two pure
top generators. A product of lower label components cannot
equal an indecomposable $\xi a$, $\xi b$, or $\xi c$.
This proves the first two identities of \textup{(D.2.16f)}
as exact cohomology identities, rather than just as projections
on a polynomial subspace.

Finally, a degree-sixteen class could contribute $(\xi b)^2$
after $\operatorname{Sq}_*^2$ only through a degree-nine
square-operation preimage of $\xi b$ or a degree-eight
Bockstein preimage. The latter was excluded above. The only
pure top generator in degree nine is $\xi c$, whose image in
\textup{(D.2.16g)} is a bracket. Products and the other
indecomposables cannot supply the required preimage. This gives
$\operatorname{Sq}^2L_b=0$ and completes the check of
\textup{(D.2.16f)}.

\medskip\noindent
\textbf{An actual map representing the multiplication deviation.}
Suppose that $G:C_4X\to T$ has coefficient one on $\xi a$
in $y=G_*(C)$. Set $g=Gj_4$, where
$j_4:X\to C_4X$ is the unit. On $X\times X$, form the two maps
\[
 F=G\mu_4(j_4\times j_4),\qquad
 H=\mu_T(g\times g),
 \tag{D.2.16h}
\]
where $\mu_4$ and $\mu_T$ are the loop multiplications.
They agree up to a pointed homotopy on $X\vee X$.
Their pointwise difference, formed by multiplying $F$ with the
pointwise inverse of $H$ in the H-group $T$, is therefore null
on that wedge. Choosing the nullhomotopy and using the pointed
cofiber universal property produces an ordinary map
\[
 \delta:X\wedge X\simeq\Sigma^2(B\wedge B)\longrightarrow T.
 \tag{D.2.16i}
\]
One may use a loop-group model for strict inverses. No
multiplicativity of $G$, coherent symmetry of $\delta$, or
factorization through a homotopy-orbit square is assumed.

The classes $C$ and $y$ are primitive. On $C\otimes C$, the homology of
the pointwise difference is computed by convolution. Its two
cross terms are both $y^2$ and cancel in characteristic two.
The antipode fixes the primitive $y$ and its square. In the
polynomial algebra $H_*T$, the coefficient of $(\xi a)^2$
in $y^2$ is exactly the coefficient of $\xi a$ in $y$, namely one.
Equation \textup{(D.2.16d)} therefore gives
\[
 \langle L,\delta_*(C\otimes C)\rangle
  =\langle L,G_*(C^2)\rangle+\langle L,y^2\rangle=1.
 \tag{D.2.16j}
\]
The quotient $X\times X\to X\wedge X$ identifies the K\"unneth
cross summand with the reduced homology of the smash product.
Thus the chosen nullhomotopy on the wedge does not alter this
cross-homology calculation.

\medskip\noindent
\textbf{The finite contradiction on the smash source.}
Let $z\in\widetilde H^8(B\wedge B)$ be the double desuspension
of $\delta^*L$. By \textup{(D.2.16j)}, its coefficient of
$u^2\otimes u^2$ is one. All positive cup products on
$X\wedge X$ vanish, since it is a double suspension.
Pulling back \textup{(D.2.16f)} and using compatibility of
Steenrod operations with suspension yields
\[
 \operatorname{Sq}^2z=0,\qquad
 \operatorname{Sq}^2\operatorname{Sq}^1z=0,\qquad
 \operatorname{Sq}^2\operatorname{Sq}^4z=0.
 \tag{D.2.16k}
\]
For the third equality, first pull back
$\operatorname{Sq}^4L=L_b+UW$, and then use
$\operatorname{Sq}^2L_b=0$.

The complete degree-eight basis of
$\widetilde H^*(B\wedge B)$ is
\[
\begin{array}{lll}
 u\otimes u^3,&u\otimes v^2,&v\otimes w,\\
 v\otimes uv,&u^2\otimes u^2,&w\otimes v,\\
 uv\otimes v,&u^3\otimes u,&v^2\otimes u.
\end{array}
 \tag{D.2.16l}
\]
Cartan gives
\[
 \ker(\operatorname{Sq}^2\text{ in degree }8)
       =\langle z_0,z_1\rangle,
 \quad
 \begin{cases}
 z_0=u^2\otimes u^2+u\otimes v^2+v^2\otimes u,\\
 z_1=v\otimes w+w\otimes v.
 \end{cases}
 \tag{D.2.16m}
\]
Here is an exhaustive coefficient check. The terms
$u\otimes u^4$ and $u^4\otimes u$ force the coefficients of
$u\otimes u^3$ and $u^3\otimes u$ to vanish. The unique terms
$w\otimes uv$ and $uv\otimes w$ force the coefficients of
$v\otimes uv$ and $uv\otimes v$ to vanish. The remaining
images force the three coefficients in $z_0$ to agree and the
two in $z_1$ to agree. This proves \textup{(D.2.16m)}.

Furthermore,
\[
 \operatorname{Sq}^1z_0=\operatorname{Sq}^1z_1
       =v\otimes v^2+v^2\otimes v,
\]
\[
 \operatorname{Sq}^2\operatorname{Sq}^1z_0
 =\operatorname{Sq}^2\operatorname{Sq}^1z_1
       =w\otimes v^2+v^2\otimes w\ne0.
 \tag{D.2.16n}
\]
The first two conditions in \textup{(D.2.16k)}, with the required
coefficient one, therefore force
\[
 z=z_0+z_1,
 \qquad\operatorname{Sq}^1z=0.
 \tag{D.2.16o}
\]
The Adem relation
$\operatorname{Sq}^2\operatorname{Sq}^4
 =\operatorname{Sq}^6+\operatorname{Sq}^5\operatorname{Sq}^1$
now gives
\[
 \operatorname{Sq}^2\operatorname{Sq}^4z
       =\operatorname{Sq}^6z\ne0.
 \tag{D.2.16p}
\]
Indeed the coefficient of $w\otimes t$ is one: it comes from
$\operatorname{Sq}^2v\otimes\operatorname{Sq}^4w$ on the term
$v\otimes w$, and none of the other four terms of
\textup{(D.2.16o)} can contribute that basis element.
This contradicts the third condition in \textup{(D.2.16k)}.

Thus no $G$ can have the asserted nonzero $\xi a$ coefficient.
Equation \textup{(D.2.16c)} and double adjunction prove the
claimed fourth-stage nonextension of $h_0$. Restricting a
hypothetical full $QX$-extension to this stage gives the same
contradiction.
The finite coefficient calculation was independently verified
by exhaustive enumeration of the nine-dimensional source space;
the displayed proof does not depend on that computation.
\end{proof}

\begin{claim}{D.2.17}{Proved}{An actual realization on the six-skeleton}
Put $B=K(\FF_2,2)$, $X=\Sigma B$, $A=QX$, $Y=\Sigma^2B$, and
$T=\Omega^2Y$. We use the labels $a,b,c,d,e,f,x$ of
$\widetilde H_*(B;\FF_2)$ dual, respectively, to
\[
 u,\quad v=\operatorname{Sq}^1u,\quad u^2,\quad
 w=\operatorname{Sq}^2v,\quad uv,\quad u^3,\quad v^2.
 \tag{D.2.17a}
\]
The same letters denote their images in $H_*T$. Write
$\xi=\xi_1$ for the top double-loop operation, and set
\[
 k=\{a,b\},\qquad p_3=f+ac+a^3,\qquad
 r=b^2+k+p_3\in H_6(T;\FF_2).
 \tag{D.2.17b}
\]
There is an actual pointed map $\rho:S^6\to T$ with Hurewicz
image $r$.

There is a CW model of $A$ with finite six-skeleton $K$, and an
ordinary pointed map $G:K\to T$, with the following values on
the indicated lifts of the unit classes and the bottom product:
\[
\begin{array}{c|rrrrrr}
 z&A_0=\sigma a&B_0=\sigma b&C=\sigma c&D=\sigma d&E=\sigma e&A_0^2\\
 \hline
 G_*z&0&a^2&0&b^2+k&r&0.
\end{array}
 \tag{D.2.17c}
\]
Thus the remaining low-dimensional pattern with $s=1$, $\gamma=0$,
and $\epsilon=1$ has an actual six-dimensional realization.
This assertion does not give an extension over the seven-cells,
over $C_4X$, or over $A$.
\end{claim}
\begin{proof}
All unlabeled homology groups in the calculation have coefficients
$\FF_2$. We write $\beta=\operatorname{Sq}_*^1$; the notation
$\beta_2$ always means the differential on the second page of the
integral Bockstein spectral sequence.

\medskip\noindent
\textbf{The degree-six primitive and its second Bockstein.}
The free double-loop homology basis and the coproduct give
\[
 \operatorname{Prim}H_6T=\langle b^2,k,p_3\rangle,
 \qquad
 \operatorname{Sq}_*^2(b^2)=\operatorname{Sq}_*^2(p_3)=a^2,
 \qquad \operatorname{Sq}_*^2(k)=0.
 \tag{D.2.17d}
\]
For completeness, the entire degree-six basis is
$b^2,k,f,x,ac,a^3$. The class $x$ has the unique $b\otimes b$
term in its reduced coproduct. The remaining nonprimitive terms are
\[
 \overline\Delta f=a\otimes c+c\otimes a,
 \qquad
 \overline\Delta(ac)=\overline\Delta f+a^2\otimes a+a\otimes a^2,
 \qquad
 \overline\Delta(a^3)=a^2\otimes a+a\otimes a^2.
\]
This proves the asserted primitive basis. Cartan and the label
operations give the square calculation. These are the conventions
of \cite[Chapter~III, Theorems~1.2--1.4 and 3.1--3.2]{I_CohenLadaMay1976}.

We need the following integral refinement:
\[
 [b^2]=[k],\qquad
 \beta_2[b^2]=\beta_2[k]=[\xi a]\ne0,
 \qquad [p_3]=0
 \quad\hbox{on the second Bockstein page.}
 \tag{D.2.17e}
\]
Here is a chain calculation fixing its coefficient. Use the bottom
Moore map $M(\ZZ/2,2)\to B$, with integral reduced chains $a_2,b_3$
and $db=2a$. In the arity-two configuration circle take free
$\Sigma_2$-cells $e_0,e_1$ with $de_1=(\tau-1)e_0$.
The action on tensors includes the Koszul sign:
$\tau(a\otimes b)=b\otimes a$ and
$\tau(b\otimes b)=-b\otimes b$. In the coinvariant tensor complex put
\[
 z=e_1(a\otimes a),\qquad
 q=e_0(b\otimes b),\qquad h=e_1(a\otimes b).
\]
With the total-complex differential,
\[
\begin{split}
 dq&=2\bigl(e_0(a\otimes b)-e_0(b\otimes a)\bigr),\\
 dh&=e_0(b\otimes a)-e_0(a\otimes b)-2z,\qquad dz=0,\\
 d(q+2h)&=-4z.
\end{split}
 \tag{D.2.17f}
\]
The reductions of $q$ and $z$ represent $b^2$ and $\xi a$.
Thus the corrected integral representative $q+2h$ proves
$\beta_2[b^2]=[\xi a]$. In particular, we are not asserting that
$dq$ is divisible by four. The same model gives
$\beta(\xi b)=b^2+k$. The configuration-chain conventions are those
in the proof of \cite[Chapter~III, Lemma~3.10, p.~232]{I_CohenLadaMay1976};
naturality along the bottom Moore map gives these identities in $T$.
There is no degree-six first-Bockstein preimage of $\xi a$:
on the complete degree-six basis its only nonzero value is
$\beta x=d+e$. Finally, with $g$ the label dual to $u^2v$,
\[
 \beta(g+ae+bc+a^2b)=p_3.
\]
This proves all parts of \textup{(D.2.17e)}.

It follows that a spherical degree-six class with nonzero double
homology suspension must be $r$. Indeed a spherical class is primitive
and is killed by $\operatorname{Sq}_*^2$ and $\beta_2$. The square
condition in \textup{(D.2.17d)} makes its coefficients on $b^2$ and
$p_3$ equal, and the Bockstein condition makes its coefficients on
$b^2$ and $k$ equal. Double homology suspension kills $b^2$ and $k$
and sends $p_3$ to $\sigma^2f$.

\medskip\noindent
\textbf{The required class is actually spherical.}
Let $M=\mathbb RP^3\times\mathbb RP^3$.
The proof of \cite[Theorem~1.1, p.~9]{EH5AkhmetievRepovsSkopenkov2001}
constructs a smooth embedding
\[
 M\times I^2\hookrightarrow\mathbb R^8.
 \tag{D.2.17g}
\]
Restricting to the middle slice gives an embedding of $M$ with a
trivialized rank-two normal bundle: the two interval directions
project to a frame of its normal quotient. Write
$t_1,t_2\in H^1(M;\FF_2)$ for the two projective-space generators,
and choose $j:M\to B$ classifying $t_1t_2$. The framed
Pontryagin--Thom collapse gives an actual map
\[
 \theta:S^8\longrightarrow
 \operatorname{Th}(\varepsilon_M^2)
 \cong S^2\wedge M_+
 \xrightarrow{S^2\wedge j_+} S^2\wedge B=Y,
 \tag{D.2.17h}
\]
where $j_+$ sends the disjoint basepoint to the basepoint of $B$.
Its Hurewicz image is $\sigma^2j_*[M]$. Serre's calculation gives
the degree-six cohomology basis $u^3,v^2$
\cite[\S2, nos.~8--9, pp.~203--204]{I_Serre1953}, and
\[
 \langle j^*(u^3),[M]\rangle
   =\langle t_1^3t_2^3,[M]\rangle=1,
 \qquad
 j^*(v^2)=(t_1^2t_2+t_1t_2^2)^2
          =t_1^4t_2^2+t_1^2t_2^4=0.
 \tag{D.2.17i}
\]
Hence $j_*[M]=f$. The double adjunct $\rho:S^6\to T$ has double
homology suspension $\sigma^2f$, so the preceding uniqueness
calculation gives $\rho_*[S^6]=r$. This uses the actual rank-two
normal framing in \textup{(D.2.17g)}, rather than a framing in a
larger Euclidean space.

\medskip\noindent
\textbf{The canonical map before correction.}
Let
\[
 h_0=\eta_2\wedge\id_B:\Sigma^3B\to\Sigma^2B,
 \qquad
 g_2=\Omega^2h_0:C_2X\to T.
 \tag{D.2.17j}
\]
On the unit $X\to C_2X$ its restriction $g_0$ has values
\[
 g_0(A_0)=0,\quad g_0(B_0)=a^2,\quad g_0(C)=\xi a,
 \quad g_0(D)=b^2,\quad g_0(E)=k.
 \tag{D.2.17k}
\]
We include the check of the entry on $B_0$, which is not a formal
consequence of the entry on $C$. Write $X_7=\sigma x$.
The cofiber of $h_0$ is $\mathbb CP^2\wedge B$. Its reduced
diagonal gives the quadratic once-adjoint image of $X_7$ as the
square of the loop label $\sigma b$. The linear part is zero,
since $h_0$ is zero on homology; tensor length three begins in
degree nine, above this degree-eight calculation. This is the
relative Adams--Hilton computation used in D.2.13, with the
quadratic differential and cell-attachment conventions of
\cite[\S3.1, p.~300; proof of Proposition~7.4, p.~309]{FGHHalperinLemaire1995}.

The primitive degree-seven basis of $H_*T$ is
\[
 \xi b,\quad \{a,c\},\quad
 P=g+ae+bc+a^2b,\quad Q=h+ad,
 \tag{D.2.17l}
\]
where $g,h$ are the labels dual to $u^2v,uw$.
Their homology suspensions are, respectively, the square of
$\sigma b$, the commutator of $\sigma a$ and $\sigma c$,
and the two independent linear labels $\sigma g,\sigma h$.
Thus homology suspension is injective on this primitive subspace,
and $g_0(X_7)=\xi b$. The top operation formula gives
\[
 \operatorname{Sq}_*^3(\xi b)=(\beta b)^2=a^2,
 \qquad \operatorname{Sq}_*^3X_7=B_0.
\]
Naturality proves $g_0(B_0)=a^2$. The value on $C$ was proved in
D.2.13. The degree-three primitive basis is $\langle b\rangle$,
and $\beta b=a\ne0$, whereas $\beta A_0=0$; hence $g_0(A_0)=0$.
The same quadratic diagonal calculation has zero value on $D$
and the commutator value on $E$. By \textup{(D.2.17d)}, the
kernel of homology suspension on the degree-six primitives is
$\langle b^2\rangle$. Finally
$\operatorname{Sq}_*^2D=B_0$ and
$\operatorname{Sq}_*^2E=0$ determine the possible $b^2$
coefficients, giving the remaining two entries of
\textup{(D.2.17k)}. Since $g_2$ is a double-loop map, it also sends
the product $A_0^2$ to zero.

\medskip\noindent
\textbf{The quotient of the six-skeleton.}
The map $C_2X\to QX$ is seven-connected. Indeed $\Sigma^2X$ is
four-connected, its Freudenthal unit is nine-connected, and two
loops leave connectivity seven; successive stabilization maps
have increasing connectivity. A relative CW approximation thus
adds no cells below dimension eight. We choose finite skeleta of
$C_2X$ and the relative approximation, and take $K=A^{(6)}$ from
this model. Its four-skeleton may be taken to be
$K^{(4)}=M(\ZZ/2,3)$. The unit classes have cellular representatives
in $K$, and the bottom product is represented in the second
configuration stage.

After collapsing $K^{(4)}$, the quotient is four-connected and
has cells only in dimensions five and six. Moreover
\[
 H_5(K/K^{(4)};\ZZ)=H_5(A;\ZZ)=H_4(B;\ZZ)=\ZZ/4.
 \tag{D.2.17m}
\]
Here $H_4(K^{(4)};\ZZ)=0$, and cells of dimension at least seven
do not affect $H_5$. The stable splitting
\cite[Theorem~1.1 and \S2]{Snaith1974} identifies $H_5A$ with
the weight-one group: all weights at least two begin in dimension
six for $X=\Sigma B$. The last equality follows from the
order-four calculation
\cite[proof in \S6, p.~1340]{I_Mahowald1964}, which identifies
$H^5(B;\ZZ)=\ZZ/4$, and the integral universal coefficient theorem;
the relevant homology is finite two-primary torsion.

The six-cell attachments to a wedge of five-spheres are entirely
determined by their degree matrix. Smith normal form and cell
slides therefore give
\[
 K/K^{(4)}\simeq M(\ZZ/4,5)\vee\bigvee_j S^6,
 \tag{D.2.17n}
\]
after removal of contractible cell pairs. We retain the free
six-spheres. Their role is essential in the following choice.
The order-four Bockstein has
$\beta_2[D]=\beta_2[E]=[C]$. Equivalently, the integral
Pontryagin-square generator has reduction $w+uv$
\cite[\S2.3, p.~1336]{I_Mahowald1964}.
Thus, in an integral cellular choice lifting these mod-two
classes, the boundary degrees of the cells $D,E$ on the
five-cell are $4,-4$, respectively. Their sum is an integral
free six-cycle. The bottom product $A_0^2$ is an independent
free six-cycle, in the configuration weight-two summand.

We can consequently choose pointed maps
\[
 q:K/K^{(4)}\to M(\ZZ/4,5),\qquad
 \ell:K/K^{(4)}\to S^6
 \tag{D.2.17o}
\]
with the following values. The map $q$ has degree one on the
five-cell, sends $D,E$ to the mod-two top Moore class, and
kills $A_0^2$. For $\ell$, collapse the five-skeleton and project
to the $E$ six-cell. Then
\[
 \ell_*D=0,\qquad \ell_*E=[S^6],\qquad \ell_*(A_0^2)=0.
 \tag{D.2.17p}
\]
Any additional free six-spheres are sent to the basepoint.
These maps are choices on the actual quotient cell complex.
No equation involving a seven-dimensional source class is being
imposed on an arbitrary map from $K$.

\medskip\noindent
\textbf{The Moore correction and its top-cell choice.}
Let $u_4:S^4\to Y$ be the bottom map, of order two by Hurewicz,
and let $\nu_4:S^7\to S^4$ be the quaternionic Hopf map.
Set $\alpha=u_4\nu_4$. The Hopf--Hilton formula gives
\[
 0=(2u_4)\nu_4=2\alpha+[u_4,u_4],
 \qquad 2[u_4,u_4]=[2u_4,u_4]=0,
 \tag{D.2.17q}
\]
so $4\alpha=0$. The formula is obtained by pinching the target
of $\nu_4$ into two four-spheres: the only additional component
in this dimension is the Whitehead product, with coefficient
the Hopf invariant one. Its quadratic form and composition rule
are recorded in \cite[(8.2)(4), (8.4), and Lemma~8.5,
pp.~4756--4757]{I_Baues1996}.

The double adjunct $\overline\alpha:S^5\to T$ has Hurewicz
image $\xi a$. To verify the coefficient, first work with
$\nu_4$: its cofiber is $\mathbb HP^2$, whose bottom class has
nonzero square. Its once-adjoint Hurewicz image is therefore
the square of the degree-three loop generator. Homology
suspension in $\Omega^2S^4$ sends $\xi a$ to that square and
has zero primitive kernel in degree five. Naturality under
$u_4$ proves the assertion in $T$.

Since $4\overline\alpha=0$, it extends to a map
$\mu:M(\ZZ/4,5)\to T$. Its top mod-two class is primitive,
is killed by $\operatorname{Sq}_*^2$, and has second
Bockstein $\xi a$. Equations \textup{(D.2.17d)--(D.2.17e)}
therefore force its image to have the form
\[
 k+t r,\qquad t\in\FF_2.
 \tag{D.2.17r}
\]
The top-cell coaction
$M(\ZZ/4,5)\to M(\ZZ/4,5)\vee S^6$, followed by
$\mu\vee\rho$, changes $t$ by one while retaining the given
bottom map. We choose the extension with top image exactly $k$.

Let $p:K\to K/K^{(4)}$ be the quotient map and form the
pointwise sum in the H-group $T$:
\[
 G=g_2|_K+\mu q p+\rho\ell p:K\longrightarrow T.
 \tag{D.2.17s}
\]
Both corrections vanish on $K^{(4)}$. On $C$ the first
correction adds $\xi a$ and the second adds zero. On $D,E$,
the first adds $k$ to each, while the second adds zero to $D$
and $r$ to $E$. All these source classes are primitive, so
the homology of the pointwise sum has no cross terms. The
product class $A_0^2$ is also primitive in characteristic two;
$g_2$ kills it by multiplicativity and both corrections kill
its chosen free-sphere representative. Combining this with
\textup{(D.2.17k)} gives exactly \textup{(D.2.17c)}.

The separate free-sphere choice in \textup{(D.2.17p)} is what
sets the value on $E$ correctly. The construction does not
verify the seven-cell attaching relations or the required
integral orders of their images. Those are further extension
conditions, and neither a map on $A$ nor an ordinary
nullification calculation is inferred from the finite map.
The seven-cell conditions are computed and satisfied in D.2.18 below.
\end{proof}

\begin{claim}{D.2.18}{Proved}{The exact degree-eight group and the seven-skeleton extension}
Retain the notation of D.2.17, and let
\[
 \overline\alpha:S^5\longrightarrow T=\Omega^2\Sigma^2B
\]
be the double adjunct of $u_4\nu_4:S^7\to\Sigma^2B$, where
$u_4:S^4\to\Sigma^2B$ is the bottom class. Then
\[
 \pi_5T=\ZZ/4\{\overline\alpha\},\qquad
 \pi_6T=\ZZ/4\{\rho\},\qquad
 2\rho=\overline\alpha\eta_5\ne0.
 \tag{D.2.18a}
\]
Here $\rho$ can be \emph{any} element whose mod-two Hurewicz image is
$r=b^2+\{a,b\}+p_3$, including the actual Pontryagin--Thom element
of D.2.17. In particular, no element with that Hurewicz image has
order two. Equivalently,
\[
 \pi_8(\Sigma^2K(\FF_2,2))=\ZZ/4\{\widehat\rho\},\qquad
 2\widehat\rho=u_4\nu_4\eta_7.
 \tag{D.2.18b}
\]

There is a CW model of $A=Q\Sigma B$ with finite seven-skeleton
$K_7$ for which the map in D.2.17 can be chosen to extend to an
ordinary pointed map $G_7:K_7\to T$. Its stated six-dimensional
homology values remain unchanged. This assertion concerns the
seven-skeleton; the eight-cell and later extension conditions are
not settled by it.
\end{claim}
\begin{proof}
Throughout the proof of the homotopy-group calculation we work
two-locally. The groups of $Y=\Sigma^2B$, its connected loop spaces,
and the connected fibers used below are finite two-primary in
positive degrees, by the Serre-class theorem. Thus the final
finite-group conclusions are also integral conclusions. The
low cohomology of $B$ and the free double-loop homology operations
are those already fixed in D.2.17; see
\cite[\S2, nos.~8--9]{I_Serre1953} and
\cite[Chapter~III, Theorems~1.2--1.4 and 3.1--3.2]{I_CohenLadaMay1976}.

\medskip\noindent
\textbf{A three-connected fiber.}
The bottom class defines a map $Y\to K(\FF_2,4)$. Its double
loop is an $H$-map
\[
 q:T\longrightarrow B.
\]
The double-suspension unit $j:B\to T$ is a section of $q$:
$qj$ is the identity on the unique nonzero homotopy group of $B$.
Put $F=\operatorname{fib}(q)$. Multiplication in the grouplike
$H$-space $T$ gives a space equivalence
\[
 B\times F\longrightarrow T,\qquad (b,x)\longmapsto j(b)x.
 \tag{D.2.18c}
\]
One obtains an inverse using $q$ and translation by the inverse
of $j(q(t))$. A multiplicative section is not required.
Freudenthal makes $j$ three-connected, and the fiber sequence
then makes $F$ three-connected. Consequently
\[
 \pi_iF\cong\pi_iT\quad(i\ge3),\qquad
 \pi_6F\cong\pi_6\Sigma^\infty F.
 \tag{D.2.18d}
\]
The second isomorphism places the required calculation in the
stable range, despite the original group's unstable description.

\medskip\noindent
\textbf{All fiber homology through degree seven.}
To avoid confusing the label $e\in H_5B$ with a six-dimensional
fiber class, set
\[
 A_4=a^2,\quad z_5=e+ab,\quad u_5=\xi a,\quad
 d_6=b^2,\quad e_6=b^2+\{a,b\},\quad f_6=p_3.
\]
The complete table for $\widetilde H_*(F;\FF_2)$ in these degrees is
\[
\begin{array}{c|c|c|c}
 \text{degree}&\text{basis}&\beta=\operatorname{Sq}_*^1
                    &\operatorname{Sq}_*^2\\ \hline
 4&A_4&0&0\\
 5&z_5,u_5&A_4,0&0,0\\
 6&d_6,e_6,f_6&0,0,0&A_4,A_4,A_4\\
 7&V=\xi b,\ R=\{a,c\},\ P,\ Q&e_6,0,f_6,0&u_5,0,z_5,z_5
\end{array}
 \tag{D.2.18e}
\]
Here, as in D.2.17,
\[
 P=g+ae+bc+a^2b,\qquad Q=h+ad,
\]
where $g,h$ are the labels dual to $u^2v,uw$ in $H^*B$.
These are fiber classes, rather than merely elements of the
linear kernel of $q_*$. Indeed the splitting
\textup{(D.2.18c)} is an isomorphism of homology coalgebras.
Primitives in the tensor product of two connected coalgebras
are sums of primitives in its two factors. All displayed
classes are primitive and map to zero in $H_*B$.
For $u_5$ this also follows from its spherical realization by
$\overline\alpha$ in D.2.17. The degree-six assertion follows
from the divided-power multiplication in $H_*B$; in degree
seven $H_*B$ has no primitives, since $H^*B$ has no
indecomposable generator in that degree.

For completeness, the full dimensions of $H_iT$, for
$i=4,5,6,7$, are $2,4,6,9$. The product splitting subtracts
the base and mixed terms and leaves fiber dimensions $1,2,3,4$.
Thus no degree-seven attachment is omitted by restricting to
the displayed primitive list. The square on $V$ is the Nishida
identity $\operatorname{Sq}_*^2(\xi b)=\xi a$. For the other
possibly nonzero entries, Cartan gives
\[
 \operatorname{Sq}_*^2g=\operatorname{Sq}_*^2h=e,
 \qquad
 \operatorname{Sq}_*^2(bc)=\operatorname{Sq}_*^2(ad)=ab,
\]
which verifies the two entries $z_5$. The remaining entries
are the coproduct, Cartan, and Bockstein calculations of D.2.17.
Its explicit integral configuration chain further gives
\[
 \beta_2[d_6]=[u_5]\ne0.
 \tag{D.2.18f}
\]
The universal coefficient theorem and these first and second
Bocksteins determine the integral groups:
\[
 H_4F=\ZZ/2,\qquad H_5F=\ZZ/4,\qquad
 H_6F=(\ZZ/2)^2.
 \tag{D.2.18g}
\]
In the last degree, $e_6,f_6$ are the two independent first
Bockstein images; $d_6$ is the Tor class belonging to the
order-four group in degree five.

\medskip\noindent
\textbf{A finite stable cell model, with every attachment detected.}
Take a finite cellular approximation to
$\Sigma^\infty F_{(2)}$ through dimension eight. Such an
approximation exists because the low homotopy groups are finitely
generated and $F$ is three-connected. In this finite range,
cancel each adjacent cell pair with odd cellular incidence.
An odd incidence is an equivalence of two-local sphere cells;
ordinary row and column elimination in the attaching cofiber
diagram splits off that contractible pair. The remaining
cellular differential is zero modulo two. Its cell counts
through dimension seven therefore equal the dimensions in
\textup{(D.2.18e)}. Dimension eight was included so that any
extraneous seven-cell killed by an odd eight-cell incidence is
also removed. Now discard cells of dimension eight. This does
not change $\pi_5$ or $\pi_6$. We denote the resulting stable
skeleta by $L^{(i)}$; no assertion about an unstable minimal
CW model is needed.

Using \textup{(D.2.18e)--(D.2.18g)}, normalize the integral
cellular differentials through dimension six as
\[
 \partial z_5=2A_4,\qquad \partial d_6=4u_5,
 \qquad \partial e_6=\partial f_6=0.
 \tag{D.2.18h}
\]
The changes can preserve every indicated mod-two label:
the nonzero second Bockstein fixes the odd coefficient of
$\partial d_6/4$, and any other multiple of $4u_5$ is eliminated
by an even multiple of $d_6$. Consequently
$L^{(5)}=P^5(2)\vee S^5_{u_5}$, where
$P^5(2)=M(\ZZ/2,4)$. Its stable fifth homotopy group is
$\ZZ/2\{i\eta_4\}\oplus\ZZ_{(2)}\{u_5\}$.
The square column of \textup{(D.2.18e)} therefore specifies
all three six-cell attachments: they have the same
$i\eta_4$ component, and the degree components are $4,0,0$.

Change six-cell basis to $e_6,d_6-e_6,r_6=f_6-e_6$.
This gives the stable wedge
\[
 L^{(6)}\simeq C\vee M(\ZZ/4,5)\vee S^6_{r_6},\qquad
 C=\operatorname{cofib}\bigl((2,\eta_4):S^4\vee S^5\to S^4\bigr).
 \tag{D.2.18i}
\]
The mod-two reduction of $r_6$ is exactly the spherical
class $r=e_6+f_6$.

Here is a direct computation of every possible seven-cell
attaching component. The first two stable stems are
$\pi_1^S=\ZZ/2\{\eta\}$ and
$\pi_2^S=\ZZ/2\{\eta^2\}$, with $\eta^2\ne0$
\cite[\S2, equation~(2.1), p.~149]{I_Oka1986}.
The cofiber sequence defining $C$ gives an injection
\[
 \pi_6C\longrightarrow
 \pi_6(C/S^4)=\ZZ/2\{\eta_5\}\oplus\ZZ_{(2)}\{\iota_6\}
\]
whose image is $\ZZ/2\oplus2\ZZ_{(2)}$.
Indeed the preceding map onto $\pi_6S^4$ is surjective
because it contains multiplication by $\eta$ on
$\pi_1^S$, and the following map is reduction modulo two
on the second summand. Hence
\[
 \pi_6C=\ZZ_{(2)}\{\gamma\}\oplus\ZZ/2\{\delta\},
 \qquad \gamma\longmapsto(0,2),\quad
 \delta\longmapsto(\eta_5,0).
 \tag{D.2.18j}
\]
In particular $h(\gamma)=2e_6$, while the square from
the five-cell $z_5$ into a new seven-cell detects precisely
its $\delta$ coefficient. This last assertion follows by
collapsing the bottom $S^4$: the quotient is an isomorphism
on the cohomology degrees five and seven used by that square.
Similarly the degree-four Moore cofiber sequence gives
\[
 \pi_6M(\ZZ/4,5)=\ZZ/2\{t=i\eta_5\};
\]
the square from $u_5$ into the new seven-cell detects $t$.
Finally $\pi_6S^6=\ZZ_{(2)}\{\rho_0\}$.
Thus the integral degrees and these two square coefficients
determine each attachment completely. There is no additional
secondary component left unspecified.

The degree-six cycles in the original cell coordinates are
freely generated by $e_6,f_6$. Their boundary subgroup is
$2\ZZ_{(2)}\{e_6,f_6\}$, by \textup{(D.2.18g)}.
After division by two the degree-seven boundary matrix is a
surjection onto a free rank-two module; its reduction modulo
two sends $V,P$ to $e_6,f_6$ and kills $Q,R$.
Choose preimages of the two basis vectors congruent to $V,P$
modulo two. An initial error is twice a vector, so surjectivity
corrects it by twice a chain. The kernel also reduces onto
the kernel modulo two, because the quotient is free. Hence
the zero-boundary basis can be chosen congruent to $Q,R$.
All square coefficients are preserved, and the four integral
boundaries are exactly $2e_6,2f_6,0,0$.

In the coordinates of \textup{(D.2.18i)}, the complete
attaching-map table is consequently
\[
\begin{array}{c|c}
 \text{seven-cell}&\text{attaching map in }\pi_6L^{(6)}\\ \hline
 V&\gamma+t\\
 P&\gamma+2\rho_0+\delta\\
 Q&\delta\\
 R&0.
\end{array}
 \tag{D.2.18k}
\]
Taking its cokernel gives
\[
 \pi_6F=
 \frac{\ZZ_{(2)}\gamma\oplus\ZZ/2\delta
       \oplus\ZZ/2t\oplus\ZZ_{(2)}\rho_0}
      {\langle\gamma+t,\ \gamma+2\rho_0+\delta,\ \delta\rangle}
 =\ZZ/4\{\rho_0\},\qquad 2\rho_0=t\ne0.
 \tag{D.2.18l}
\]
Also $\pi_5C=0$, so \textup{(D.2.18i)} gives
$\pi_5F=\ZZ/4$ with Hurewicz generator $u_5$; seven-cells
and higher cells do not change it. The class
$\overline\alpha$ has exactly this Hurewicz image by D.2.17,
and is therefore a generator. The image of $t$ is
$\overline\alpha\eta_5$; changing the odd generator of
$\pi_5F$ does not change this order-two element.
The Hurewicz image of $\rho_0$ is $r$. Every element with
that image is an odd generator of $\ZZ/4$, so its double
is the same nonzero $t$. This proves
\textup{(D.2.18a)--(D.2.18b)} for the actual geometric
representative as well as every other such representative.

\medskip\noindent
\textbf{The source quotient and the actual seven-stage extension.}
The map $C_2X\to A$ is seven-connected, as in D.2.17,
so choose a finite $K_7=A^{(7)}$ from a CW model of $C_2X$.
The canonical map $g_2=\Omega^2(\eta_2\wedge\id_B)$ is
already defined on it. The correction is zero on
$K_4=M(\ZZ/2,3)$, so its extension is a problem on the
four-connected quotient $Q_7=K_7/K_4$.

We record all relevant source cells. The six-skeleton of
the quotient has the Moore bottom $c_5=C$, Moore top
$d_6=D$, and free six-cycles
$e_6^{\rm src}=D+E$ and $z_6^{\rm src}=A_0^2$.
Here the integral lifts of $D,E$ have opposite boundaries
$4c_5,-4c_5$, respectively. The Snaith splitting gives
\[
 H_6(A;\ZZ)=H_5(B;\ZZ)\oplus
 H_6(D_2X;\ZZ)=(\ZZ/2)^2.
 \tag{D.2.18m}
\]
For the first term, the low dimensions from Serre and
$H_4(B;\ZZ)=\ZZ/4$ make $H_5B$ cyclic two-primary;
$\beta x=d+e\ne0$ fixes its exponent as two. The
order-four input is the Pontryagin-square calculation in
\cite[\S2.3 and proof of Lemma~6.1]{I_Mahowald1964}.
The weight-two term is its bottom group
$H_0(\Sigma_2;H_3X\otimes H_3X)=\ZZ/2$; there is no
lower term to support a differential. Higher weights begin
in degree nine. These uses of the stable splitting are
exactly \cite[Theorem~1.1 and \S2]{Snaith1974}.

The full degree-seven mod-two source basis is
\[
 F_7=\sigma f,\quad X_7=\sigma x,\quad
 R_7=Q^4A_0,\quad A_0B_0.
\]
Its Bocksteins are $0,e_6^{\rm src},z_6^{\rm src},
z_6^{\rm src}$, respectively. The product $A_0B_0$
must be retained here, even though it is not primitive.
On the unique degree-five cohomology class of the quotient,
\[
 \operatorname{Sq}^2(C^\vee)=\sigma(v^2).
 \tag{D.2.18n}
\]
It evaluates to one on $X_7$ and zero on the other three
basis elements. This follows on unit classes from
$\operatorname{Sq}^2(u^2)=v^2$, on $R_7$ from the
Nishida identity, and on $A_0B_0$ from Cartan.

Integral Smith reduction cancels any unit six/seven pairs.
Equation \textup{(D.2.18m)} and the Bocksteins normalize the
two remaining nonzero boundaries as
$2e_6^{\rm src},2z_6^{\rm src}$, preserving the specified
mod-two classes $X_7,R_7$. The zero-boundary classes include
$F_7$ and $A_0B_0+R_7$. Any additional seven-cell class
killed by higher cells has zero value under
\textup{(D.2.18n)}, by naturality.
There is no extra attaching invariant: the quotient is
four-connected and
\[
 \pi_6\bigl(M(\ZZ/4,5)\vee S^6_e\vee S^6_z\bigr)
 =\ZZ/2\{i\eta_5\}\oplus\ZZ\{e\}\oplus\ZZ\{z\}.
\]
The degree-five square detects the first summand exactly.
Thus the actual quotient has the homotopy type
\[
 \Bigl(M(\ZZ/4,5)\vee S^6_e\vee S^6_z\Bigr)
 \mathop{\cup}_{i\eta_5+2e}e^7
 \mathop{\cup}_{2z}e^7
 \ \vee\bigvee S^7.
 \tag{D.2.18o}
\]
In this dimension and connectivity the stable attachment
calculation is the unstable one; all possible Whitehead
products have larger dimension. Extra zero-degree
attachments vanish because their only possible $\eta$
component has just been checked to be zero.

Use the Moore extension
$\mu:M(\ZZ/4,5)\to T$ chosen in D.2.17, with bottom
$\overline\alpha$ and top mod-two image $k=\{a,b\}$.
On the two free six-spheres use $\rho$ and zero,
respectively. The only potentially nonzero obstruction
in \textup{(D.2.18o)} is
\[
 \overline\alpha\eta_5+2\rho=0,
 \tag{D.2.18p}
\]
by \textup{(D.2.18a)}. The other displayed attachment and
all free seven-sphere attachments impose no condition.
Choose nullhomotopies to extend the correction across the
seven-cells, pull it back along $K_7\to Q_7$, and add it
to $g_2$ using the multiplication of $T$.
On $C,D,E,A_0^2$ the correction has values
$\xi a,k,k+r,0$, respectively. Adding the canonical
values of D.2.17 gives precisely its displayed values,
including $D\mapsto b^2+k$ and $E\mapsto r$.
This constructs $G_7$ and finishes the stated finite-stage
extension.
\end{proof}

\begin{claim}{D.2.19}{Proved}{The next target group and the eight-skeleton extension}
Retain the notation of D.2.17--D.2.18. Then
\[
 \pi_7T=\pi_9(\Sigma^2K(\FF_2,2))\cong\ZZ/2,
 \qquad \rho\eta_6=0.
 \tag{D.2.19a}
\]
The generator has mod-two Hurewicz image $R=\{a,c\}$ in
$H_7T$. If $p:M(\ZZ/4,5)\to S^6$ is the top pinch,
$\zeta:S^7\to M(\ZZ/4,5)$ satisfies $p\zeta=\eta_6$,
and $\mu:M(\ZZ/4,5)\to T$ is the Moore extension used in
D.2.17, then
\[
 \mu\zeta=0.
 \tag{D.2.19b}
\]
The ordinary map on the seven-skeleton in D.2.18 can be chosen
to extend over a finite eight-skeleton $K_8$ of a CW model of
$A=Q\Sigma K(\FF_2,2)$, retaining all the homology values
in D.2.17. This is a finite-stage extension; no extension
over the full source or nullification conclusion is asserted.
\end{claim}
\begin{proof}
\textbf{The target degree is still stable.}
Let $F=\fib(T\to B)$ be the three-connected fiber of D.2.18.
It is a double loop space, since its defining map is the
double loop of $Y\to K(\FF_2,4)$. Write $F=\Omega Z$.
The adjunction triangle identity gives a left inverse
\[
 \Omega\epsilon_Z:\Omega\Sigma F\longrightarrow F
\]
to the suspension--loop unit. Thus that unit is injective
on every homotopy group. Freudenthal makes it surjective
on $\pi_7$, and $\Sigma F$ is four-connected. Hence
\[
 \pi_7F\cong\pi_8\Sigma F\cong\pi_7\Sigma^\infty F.
 \tag{D.2.19c}
\]
This removes the possible first unstable kernel, including
its torsion, without assuming that a proposed list of
Whitehead products exhausts it. The same retraction is
described by \cite[Theorem~1.8, p.~173]{I_James1955}.
We compute two-locally, as in D.2.18; the resulting finite
groups are also the integral groups.

\medskip\noindent
\textbf{Every degree-eight fiber class and its operations.}
Use the fiber classes
\[
 A_4=a^2,\quad z_5=e+ab,\quad u_5=\xi a,\quad
 d_6=b^2,\quad e_6=b^2+k,\quad f_6=p_3,
 \qquad k=\{a,b\}.
\]
Let $i,j,l\in H_8B$ be dual to $u^4,uv^2,vw$.
The full free double-loop basis of $H_8T$ is
\[
\begin{gathered}
 i,j,l,\ af,ax,bd,be,c^2,a^2c,ab^2,a^4,\\
 b\xi a,ak,\{a,d\},\{a,e\},\{b,c\}.
\end{gathered}
 \tag{D.2.19d}
\]
There is no further top-operation term in even degree eight.
The restricted Browder identities give
$\{a,\xi a\}=0$ and zero brackets on identical inputs;
these are the basis conventions of
\cite[Chapter~III, Theorems~1.2--1.4]{I_CohenLadaMay1976}.
The splitting $T\simeq B\times F$ subtracts dimensions
$3,1,2,3$ from the base and the mixed terms, leaving
$\dim H_8F=7$. A complete basis is
\[
\begin{split}
 W_8&=a^4,\qquad C_8=c^2,\qquad L_8=l+bd,\\
 M_8&=j+ax+be+ab^2,\\
 J_1&=\{a,d\},\qquad
 J_2=\{a,e\}+ak,\qquad J_3=\{b,c\}+ak.
\end{split}
 \tag{D.2.19e}
\]
All except $C_8$ are primitive. The reduced diagonal of
$C_8$ is $A_4\otimes A_4$. Each primitive has zero image
in $B$, which has no degree-eight primitive homology.
For the remaining class, $q_*C_8=0$ and
$(q_*\otimes1)\Delta C_8=1\otimes C_8$ show that it is
in the fiber factor. Thus these are actual fiber classes.

Here is the complete operation table needed below:
\[
\begin{array}{c|c|c|c}
 &\beta=\operatorname{Sq}_*^1&\operatorname{Sq}_*^2
                         &\operatorname{Sq}_*^4\\ \hline
 W_8&0&0&0\\
 C_8&0&0&A_4\\
 L_8&Q&d_6&0\\
 M_8&Q&f_6&0\\
 J_1&0&d_6+e_6&0\\
 J_2&0&0&0\\
 J_3&R&d_6+e_6&0
\end{array}
 \tag{D.2.19f}
\]
We give checks fixing the less immediate coefficients.
Serre's cohomology calculation and Cartan give
\[
 \operatorname{Sq}^1(uw)=vw+uv^2,
 \qquad \operatorname{Sq}^2(u^3)=u^4+uv^2.
\]
Consequently
\[
\begin{split}
 \beta L_8&=h+ad=Q,\\
 \beta M_8&=h+a(d+e)+ae=Q,\\
 \operatorname{Sq}_*^2L_8&=b^2=d_6,\\
 \operatorname{Sq}_*^2M_8&=f+ac+a^3=f_6,\\
 \operatorname{Sq}_*^2\{a,d\}
 &=\operatorname{Sq}_*^2\{b,c\}=k=d_6+e_6.
\end{split}
\]
Cartan gives the other first- and second-square entries.
For the last column, instability identifies the fourth
square of the degree-four cohomology class with its cup
square. Primitives annihilate that product, whereas the
diagonal of $C_8$ evaluates to one. This checks the last
column independently of label-operation formulas.
The primary conventions are those of
\cite[\S2, nos.~8--9]{I_Serre1953} and
\cite[Chapter~III, Theorems~3.1--3.2]{I_CohenLadaMay1976}.

The universal coefficient theorem, $H_6F=(\ZZ/2)^2$,
and the two independent first-Bockstein images $Q,R$
in \textup{(D.2.19f)} give
\[
 H_7F=(\ZZ/2)^2.
 \tag{D.2.19g}
\]

\medskip\noindent
\textbf{The entire stable seven-skeleton group.}
Take the finite stable approximation used in D.2.18 through
dimension nine before cancelling odd-unit cell pairs. This
also removes extraneous eight-cells. Write $L_7$ for its
seven-skeleton. Its six-skeleton is
\[
 L_6=C\vee M(\ZZ/4,5)\vee S^6_{\rho_0},
 \qquad C=\cofib\bigl((2,\eta_4):S^4\vee S^5\to S^4\bigr),
\]
and its seven-cell attachments are
\[
 V:\gamma+t,\qquad
 P:\gamma+2\rho_0+\delta,\qquad Q:\delta,\qquad R:0.
 \tag{D.2.19h}
\]
The first three stable stems, including
$\pi_3^S{}_{(2)}=\ZZ/8\{\nu\}$, are recorded in
\cite[\S2, (2.1), p.~149]{I_Oka1986}.
The same Chang cofiber sequence now gives
\[
 0\longrightarrow\ZZ/2\{v_0=i\nu\}
 \longrightarrow\pi_7C
 \longrightarrow\ZZ/2\{\eta_5^2\}\longrightarrow0.
 \tag{D.2.19i}
\]
The class $\delta\eta_6$ maps nontrivially to the last
group and has order two, so this extension splits.
The class $\gamma\eta_6$ maps to zero and is therefore
zero or $v_0$. The fourth-square coefficient from the
bottom four-class into the cone of $v_0$ equals one,
by comparison with the Hopf-invariant-one map $\nu$.
It vanishes for $\gamma\eta_6$: this map factors through
$S^6$, and the induced map from the cone of $\eta_6$
kills the degree-four class while preserving the new
top degree-eight class. Hence
\[
 \gamma\eta_6=0.
 \tag{D.2.19j}
\]

For the Moore summand one has
\[
 \pi_7M(\ZZ/4,5)
 =\ZZ/2\{t\eta_6\}\oplus\ZZ/2\{b_0\},
 \qquad p_*b_0=\eta_6.
 \tag{D.2.19k}
\]
Here the order-two lift can be constructed from the
cofiber map $M(\ZZ/2,5)\to M(\ZZ/4,5)$ that has
degree two on bottom cells and degree one on top cells.
The relation $2\widetilde\eta=i\eta^2$ in the mod-two
Moore spectrum, given in
\cite[\S2, (2.2)(a)--(b), p.~149]{I_Oka1986}, maps to
$2i\eta^2=0$. The top pinch proves independence from
the bottom $\eta^2$ class.
Set $c_0=\rho_0\eta_6$ in the last sphere summand.
The preceding maps induced by \textup{(D.2.19h)} kill
exactly $t\eta_6$ and $\delta\eta_6$ in $\pi_7L_6$.
The kernel of the corresponding map on $\pi_6$ is free
on twice the $Q$ cell and the $R$ cell. The cofiber
sequence therefore gives
\[
 \pi_7L_7=
 \ZZ_{(2)}\{\kappa\}\oplus\ZZ_{(2)}\{\lambda\}
 \oplus(\ZZ/2)\{v_0,b_0,c_0\},
 \quad h(\kappa)=2Q,\quad h(\lambda)=R.
 \tag{D.2.19l}
\]
The extension by the free group splits; the choices of
its lifts are normalized next.

\medskip\noindent
\textbf{Complete detectors for the eight-cell attachments.}
For a map $S^7\to L_7$ with zero mod-two Hurewicz image,
its cone has a new top degree-eight cohomology class.
The second-square coefficients from degree six and the
fourth-square coefficient from degree four are additive
invariants of such stable maps. This follows from
naturality of cones under the stable pinch map.
On the torsion classes they are
\[
\begin{array}{c|c|c}
 &\operatorname{Sq}_*^2&\operatorname{Sq}_*^4\\\hline
 v_0&0&A_4\\
 b_0&d_6+e_6&0\\
 c_0&e_6+f_6&0.
\end{array}
 \tag{D.2.19m}
\]
For $b_0$, collapse the Moore bottom: its top
$\eta$-lift is detected by the six-class $d_6-e_6$.
For $c_0$, use the sphere class $f_6-e_6$. These
factorizations also make their fourth-square coefficients
zero. The coefficient for $v_0$ is the Hopf-invariant-one
coefficient; its second-square coefficients vanish by
the bottom-sphere comparison.

For $\kappa$, the top Bockstein in its cone is $Q$.
The Adem identity
$\operatorname{Sq}^2\operatorname{Sq}^2
=\operatorname{Sq}^3\operatorname{Sq}^1$, together with
$\operatorname{Sq}_*^2Q=z_5$ and $\beta z_5=A_4$,
makes the sum of its three second-square coordinates odd.
Adding $b_0,c_0$ makes this vector exactly $d_6$;
adding $v_0$ makes its fourth-square coefficient zero.
Fix this choice of $\kappa$.
The element $2\lambda$ has zero second- and
fourth-square coefficients: compare its cone with the
mod-two Moore cone of the degree-two map of $S^7$.
The comparison kills the relevant degree-six and
degree-four classes and preserves the top class.

Every eight-cell attachment has even integral degree
on $R$. Its coefficients on $\kappa$ and $\lambda$,
together with \textup{(D.2.19m)}, consequently determine
it completely. The only surviving bottom $\nu$ ambiguity
is the order-two class $v_0$, and its fourth square
detects it. No even multiple of $\nu$ remains hidden.

The integral seven-cycles are free on $Q,R$. By
\textup{(D.2.19g)}, the eight-cell boundary subgroup
is exactly $2\ZZ_{(2)}\{Q,R\}$. As in D.2.18,
choose preimages of $2Q,2R$ congruent to $L_8,J_3$
modulo two, and correct all remaining errors by even
cell combinations. The integral boundaries can thereby
be set to $2Q,2Q,2R$ for $L_8,M_8,J_3$, respectively,
and zero for the other four classes. This preserves
every coefficient in \textup{(D.2.19f)}.
The complete attaching table is
\[
\begin{array}{c|c}
 W_8&0\\
 C_8&v_0\\
 L_8&\kappa\\
 M_8&\kappa+b_0+c_0\\
 J_1&b_0\\
 J_2&0\\
 J_3&2\lambda+b_0.
\end{array}
 \tag{D.2.19n}
\]
Its cokernel is $\ZZ/2\{\lambda\}$, with Hurewicz
image $R$, and it kills $b_0$ and $c_0$ separately.
Cells of dimension nine or more do not change this
stable group. Equation \textup{(D.2.19c)} proves
\textup{(D.2.19a)} for the actual unstable group.

The inclusion of the Moore summand sends its top
$\eta$-lift to $b_0=0$. Any other Moore extension with
the same bottom differs by an element of $\pi_6F$
followed by the top pinch. Composing that difference
with a top $\eta$-lift gives a multiple of
$\rho\eta_6=0$. The actual extension $\mu$ of D.2.17
therefore satisfies \textup{(D.2.19b)}, independently
of its permitted top-cell choice. These statements
also hold in $T$, since $\pi_iF\cong\pi_iT$ in the
degrees used here.

\medskip\noindent
\textbf{The entire source group in the next attaching dimension.}
Write $Q_i=K_i/K_4$, with the source skeleta of D.2.18.
After the indicated cell cancellations, its seven-dimensional
quotient has the following two-local model:
\[
 Q_7\simeq D\vee M(\ZZ/2,6)_z\vee S^7_F\vee S^7_J,
 \qquad
 D=\operatorname{cofib}\left(
 \begin{pmatrix}4&\eta_5\\0&2\end{pmatrix}:
 S^5\vee S^6\longrightarrow S^5\vee S^6\right).
 \tag{D.2.19o}
\]
The six-dimensional coordinates are the Moore top $d=D$,
the free cycle $e=D+E$, and $z=A_0^2$ before their seven-cell
attachments. The free seven-classes are
$F=\sigma f$ and $J=A_0B_0+Q^4A_0$ modulo two.
We distinguish the complex $D$ in \textup{(D.2.19o)} from
the homology class $D=\sigma d$ by its argument whenever
necessary.

There is a complete and explicitly detected group
\[
 \pi_7Q_7=
 \ZZ_{(2)}\{F,J\}\oplus(\ZZ/2)\{\tau_d,\tau_e,\tau_z\}.
 \tag{D.2.19p}
\]
To see this without presuming a Moore-group extension,
collapse the bottom five-sphere of the complex $D$.
Its quotient is $S^6_d\vee M(\ZZ/2,6)_e$.
The map $\pi_7S^5\to\pi_7D$ is zero: its generator
$i\eta_5\eta_6$ is
$(i\eta_5+2e)\eta_6$, since $2\eta_6=0$.
Thus the quotient injects $\pi_7D$ into
$\pi_7(S^6\vee M(\ZZ/2,6))=(\ZZ/2)^2$.
The connecting map to $\pi_6S^5$ vanishes on both generators.
On $S^6_d$ it is the degree-four map, which kills $\eta$;
on the Moore factor it is $\eta_6$ followed after its top
pinch, which kills its bottom $\eta$ generator. Therefore
the injection is an isomorphism.

The image of any lift
$\zeta:S^7\to M(\ZZ/4,5)$ with $p\zeta=\eta_6$
is the class $\tau_d$. Its choice ambiguity is $i\eta^2$,
which has just been killed in $D$.
The other two classes are
$\tau_e=e\eta_6$ and $\tau_z=z\eta_6$.
They are detected exactly by the second squares from the
three degree-six cohomology classes $d^\vee,e^\vee,z^\vee$.
The remaining two components in \textup{(D.2.19p)} are
their integral degrees on $F,J$.
The quotient is four-connected, so these stable groups and
attaching components are the actual unstable ones. In
particular, there is no additional Whitehead component.

\medskip\noindent
\textbf{Every degree-eight source class, including the products.}
The full ordinary mod-two basis is
\[
 G_8=\sigma g,\quad H_8=\sigma h,\quad
 A_0C,\quad B_0^2,\quad Q^5A_0.
 \tag{D.2.19q}
\]
The two unit classes are dual to the suspended monomials
$u^2v,uw$. The products and the new arity-two class are
included; arities at least three start in dimension nine.
The complete operations relevant to
\textup{(D.2.19p)} are
\[
\begin{array}{c|c|c}
 \text{class}&\beta&\operatorname{Sq}_*^2\\ \hline
 G_8&F&E=d+e\\
 H_8&0&E=d+e\\
 A_0C&0&z\\
 B_0^2&0&z\\
 Q^5A_0&0&z.
\end{array}
 \tag{D.2.19r}
\]
For the unit classes use
$\operatorname{Sq}^2w=0$ and
$\operatorname{Sq}^2(uv)=u^2v+uw$.
Cartan gives
$\operatorname{Sq}_*^2(A_0C)=A_0^2$ and
$\operatorname{Sq}_*^2(B_0^2)=(\beta B_0)^2=A_0^2$.
The stable Nishida formula gives
\[
 \operatorname{Sq}_*^2Q^5A_0
   =\binom{3}{2}Q^3A_0=A_0^2,
 \tag{D.2.19s}
\]
because $\beta A_0=0$. These are the primary Serre and
Cohen--Lada--May conventions already cited above.

\medskip\noindent
\textbf{The degree-four boundary coefficient is an integral calculation.}
The Snaith splitting identifies
\[
 H_7(A;\ZZ)=\ZZ/2\{F\}\oplus\ZZ/4\{J\}.
 \tag{D.2.19t}
\]
We verify the exponent and the particular generator of its
second summand. For the first summand, Serre's mod-two
dimensions and the universal coefficient theorem make
$H_6(B;\ZZ)$ cyclic two-primary. Its first Bockstein
$\beta g=f\ne0$ fixes its exponent as two.

For weight two, use the bottom Moore chains of $X$:
$a$ has degree three, $b$ has degree four, and $db=2a$.
Use free $E\Sigma_2$ cells with
$de_1=(\tau-1)e_0$ and $de_2=(1+\tau)e_1$.
The Koszul action satisfies
$\tau(a\otimes a)=-a\otimes a$ and
$\tau(a\otimes b)=b\otimes a$.
Set
\[
\begin{gathered}
 x_0=e_0(a\otimes a),\quad r_1=e_1(a\otimes a),\quad
 t=e_0(a\otimes b),\quad t'=e_0(b\otimes a),\\
 h=e_1(a\otimes b),\qquad q_0=e_0(b\otimes b).
\end{gathered}
\]
The degree-seven boundaries are
\[
 dr_1=-2x_0,\qquad dt=-2x_0,\qquad dt'=2x_0.
\]
Hence $w_0=t-r_1$ and $w_0'=t'+r_1$ freely generate
the integral seven-cycles. The next boundaries are
\[
 dh=w_0'-w_0,\qquad dq_0=2(w_0+w_0'),
 \qquad d(e_2(a\otimes a))=0.
 \tag{D.2.19u}
\]
The other mixed $e_1$ cell has the opposite boundary.
Higher-label cells in dimension eight are cycles and do
not alter the degree-seven group. Thus its weight-two part is
\[
 \frac{\ZZ\{w_0,w_0'\}}
      {\langle w_0'-w_0,\ 2(w_0+w_0')\rangle}
 =\ZZ/4\{w_0\}.
\]
Modulo two, $w_0$ is exactly $J=A_0B_0+Q^4A_0$.
Furthermore
\[
 d(q_0-2h)=4w_0,
 \qquad \beta_2[B_0^2]=[J]\ne0.
 \tag{D.2.19v}
\]
This uses a corrected integral representative; $dq_0$
itself is not divisible by four.
The cells $A_0C$ and $Q^5A_0$ have integral cycle
representatives. The unit class $H_8$ has no
second-Bockstein component on the weight-two class $J$.
The arity separation is provided by
\cite[Theorem~1.1 and \S2]{Snaith1974}; the displayed
chain formulas use the conventions of
\cite[Chapter~III, proof of Lemma~3.10, p.~232]{I_CohenLadaMay1976}.
They prove \textup{(D.2.19t)} and the coefficient in
\textup{(D.2.19v)}, rather than inferring that coefficient
from a mod-two dimension count.

\medskip\noindent
\textbf{The actual eight-cell matrix.}
The free seven-cycles of $Q_7$ are $F,J$. Equation
\textup{(D.2.19t)} makes their boundary sublattice
$\langle2F,4J\rangle$. The first and second Bocksteins
specify a normal form preserving all the displayed
mod-two cell labels. In detail, a term $4mJ$ in the
boundary of $G_8$ is removed by replacing the free
generator $F$ by $F+2mJ$, with unchanged mod-two class.
Terms $4mF$ in the boundary of $B_0^2$ are removed by
even multiples of $G_8$. On a zero-first-Bockstein
cell the remaining $F$ coefficient is even, and a
zero second Bockstein on $J$ makes the remaining
$J$ coefficient an even multiple of the degree-four
relation. Even cell combinations therefore remove
these terms without changing a square coefficient.
Odd rescalings are invertible two-locally. This gives
boundaries $2F$ and $4J$ for $G_8$ and $B_0^2$,
and zero integral degrees for the other three cells.

The detectors in \textup{(D.2.19p)} and the complete
table \textup{(D.2.19r)} now give every attachment:
\[
\begin{array}{c|c}
 \text{eight-cell}&\text{attaching map in }\pi_7Q_7\\ \hline
 G_8&2F+\tau_d+\tau_e\\
 H_8&\tau_d+\tau_e\\
 B_0^2&4J+\tau_z\\
 A_0C&\tau_z\\
 Q^5A_0&\tau_z.
\end{array}
 \tag{D.2.19w}
\]
The calculation may be made after cancelling odd unit
pairs in a two-local finite model. Each comparison is
realizable unstably here: the quotients are four-connected,
their relevant attaching dimensions are at most seven,
and Freudenthal identifies these stable attaching maps
with ordinary ones. Pullback along localization gives
ordinary maps on the original finite source skeleta,
because the target $T$ is two-local. Retaining the
cancelled pairs gives the same mapping problem.

The correction of D.2.18 is $\mu$ on the Moore
factor, $\rho$ on the six-sphere $e$, and zero on
$z$. Choose its values on the free seven-spheres
$F,J$ to be zero. Its values on the torsion classes are
\[
 \tau_d\longmapsto\mu\zeta=0,\qquad
 \tau_e\longmapsto\rho\eta_6=0,\qquad
 \tau_z\longmapsto0.
 \tag{D.2.19x}
\]
The first two vanish by the target calculation.
They are independent of the chosen nullhomotopy on
the seven-cell $i\eta_5+2e$, because the two
representing maps already lie in its six-skeleton.
The ambiguity $i\eta^2$ in $\zeta$ is also harmless:
$\overline\alpha\eta^2=2\rho\eta_6=0$.
Thus every attachment in \textup{(D.2.19w)} is killed,
and the correction extends over $Q_8$.

\medskip\noindent
\textbf{The canonical summand also extends to the new source cells.}
The canonical map $g_2:C_2X\to T$ must not be treated
as already defined on the new eight-cell of $A$.
We verify the needed extension independently.
Choose a bottom generator $i:S^3\to X$ and put
$X'=X/S^3$. This quotient is three-connected:
$X$ is two-connected, and
$H_3(S^3;\ZZ)\to H_3X=\ZZ/2$ is surjective.
The Hurewicz and cofiber sequences give
$H_3X'=0$, hence $\pi_3X'=0$.

Freudenthal makes $B\to\Omega^2\Sigma^2B$
three-connected. Since $\pi_3B=0$, it follows that
$\pi_5Y=0$, as already checked in D.2.13.
Consequently
$h_0\Sigma^2i:S^5\to Y$ is null; equivalently
its bottom description is $u_4\eta_4$, in this
zero homotopy group. The cofiber universal property
therefore gives a pointed map
\[
 \overline h:\Sigma^2X'\longrightarrow Y,
 \qquad \overline h\Sigma^2q\simeq h_0,
 \tag{D.2.19y}
\]
where $q:X\to X'$ is the quotient.

Since $X'$ is three-connected, the stabilization
$C_2X'\to QX'$ is nine-connected: the unit on
the five-connected space $\Sigma^2X'$ is
eleven-connected, two loops subtract two, and
subsequent stages improve the bound.
Form the unsuspended homotopy pullback
\[
 P_0=(QX)\times^h_{QX'} C_2X'
          \longrightarrow QX.
 \tag{D.2.19z}
\]
It is nine-connected and receives the natural lift
of $C_2X$. Projection to $C_2X'$ followed by
$\Omega^2\overline h$ gives $P_0\to T$; its
restriction to $C_2X$ is $g_2$, by
\textup{(D.2.19y)} and naturality.
After fibration replacement the fiber over $QX$
is eight-connected. The specified lift therefore
extends across every relative cell of dimension
at most nine. This gives a canonical summand on
the required $K_8$ (indeed through source relative
dimension nine), with its original restriction
and low homology values.

Finally pull the extended correction on $Q_8$
back along $K_8\to Q_8$ and add it to this
canonical summand using multiplication in $T$.
The result is an actual ordinary map $G_8:K_8\to T$.
Its restriction through dimension six is the
chosen map of D.2.17. The argument treats the new
eight-cell explicitly and uses the unsuspended
pullback \textup{(D.2.19z)}; it does not assume
that suspension preserves a homotopy fiber sequence.
\end{proof}

\begin{claim}{D.2.20}{Proved}{The ninth-stage target tests and an ordinary extension}
Retain the notation of D.2.17--D.2.19, and let
$F=\fib(T\to K(\FF_2,2))$. The stable group and the
actual mod-two Hurewicz map satisfy
\[
 \pi_8\Sigma^\infty F\cong(\ZZ/2)^3,
 \qquad
 h:\pi_8F\longrightarrow H_8(F;\FF_2)
 \text{ is injective}.
 \tag{D.2.20h1}
\]
If $\overline\alpha:S^5\to T$ is the class of D.2.18
and $\lambda:S^7\to T$ is the generator of D.2.19, then
\[
 \overline\alpha\nu_5=0,
 \qquad \lambda\eta_7=0
 \quad\text{in }\pi_8T.
 \tag{D.2.20h2}
\]
The ordinary map of D.2.19 can be chosen to extend over
a finite nine-skeleton of a CW model of
$A=Q\Sigma K(\FF_2,2)$, preserving the homology values
of D.2.17. The statement concerns this finite stage;
it does not settle extension over the full source or
the ordinary nullification test in D.2.4.
\end{claim}
\begin{proof}
\textbf{Every degree-nine fiber class.}
Keep the fiber notation
\[
 A_4=a^2,\quad z_5=e+ab,\quad u_5=\xi a,
 \quad k=\{a,b\},\quad R=\{a,c\},
\]
and the degree-seven classes $P,Q$ of D.2.18.
The degree-eight fiber basis of D.2.19 is
\[
 W_8,C_8,L_8,M_8,J_1,J_2,J_3.
\]
Let $m,n,o,t\in H_9B$ be the labels dual to
$u^3v,u^2w,v^3,\operatorname{Sq}^4w$, respectively.
The full free double-loop basis of $H_9T$ has twenty-six
elements:
\[
\begin{gathered}
 m,n,o,t,\\
 ag,ah,a\xi b,aR,bf,bx,bk,cd,ce,c\xi a,\\
 a^2d,a^2e,a^2\xi a,abc,b^3,a^3b,\\
 \{a,f\},\{a,x\},\{b,d\},\{b,e\},\{a,k\},\xi c.
\end{gathered}
 \tag{D.2.20h3}
\]
The restricted Browder identity identifies
$\{b,\xi a\}$ with $\{a,k\}$; it is not an extra
generator. The base contributes dimension four and
the mixed terms in $H_9(B\times F)$ contribute
$2,2,3,4$. Hence $\dim H_9F=11$.
A complete fiber basis is
\[
\begin{split}
 Z_A&=A_4z_5,\qquad U_A=A_4u_5,\\
 N_9&=\xi c+aR+A_4u_5,\\
 K_1&=\{a,f\}+aR,\qquad
 K_2=\{a,x\}+bk,\\
 K_3&=\{b,d\},\qquad
 K_4=\{b,e\}+bk,\qquad K_5=\{a,k\},\\
 P_m&=m+ag+ce+a^2e+bf+a^3b,\\
 P_n&=n+ah+cd+a^2d,\qquad
 P_o=o+bx+b^3.
\end{split}
 \tag{D.2.20h4}
\]
The first two are products of primitive fiber classes;
the other nine are primitive. The top-operation diagonal
that needs particular care is
\[
 \overline\Delta(\xi c)
 =u_5\otimes A_4+A_4\otimes u_5
       +R\otimes a+a\otimes R.
 \tag{D.2.20h5}
\]
Apply top-operation additivity and Cartan to
$\Delta c=c\otimes1+1\otimes c+a\otimes a$.
The two product corrections in $N_9$ cancel exactly
the four terms in \textup{(D.2.20h5)}.
For example, the reduced diagonal of $\{a,x\}$ is
$k\otimes b+b\otimes k$, canceled by $bk$;
the other primitive assertions follow by the same
finite coproduct calculation. These are the exact
restricted-operation identities of
\cite[Chapter~III, Theorem~1.3(2)--(5), pp.~217--218]
{I_CohenLadaMay1976}.

Every displayed class has zero image under $q_*$.
For the polynomial and Browder primitives this follows
directly: the only degree-nine primitive in $H_*B$
is the label $t$, and none of those expressions has
a $t$ coefficient after applying $q_*$. The fourth-square
calculation below gives
$\operatorname{Sq}_*^4N_9=u_5$, whose image under $q_*$
is zero, whereas $\operatorname{Sq}_*^4t=d\ne0$.
Thus $N_9$ also has zero possible $t$ coefficient.
The primitive-kernel argument from D.2.18 puts these
classes in the fiber factor, and the two products
plainly belong there. The dimension count proves
completeness.

\medskip\noindent
\textbf{The complete operation table, including the top correction.}
With $\beta=\operatorname{Sq}_*^1$, the relevant operations are
\[
\begin{array}{c|c|c|c}
 &\beta&\operatorname{Sq}_*^2&\operatorname{Sq}_*^4\\ \hline
 Z_A&W_8&0&0\\
 U_A&0&0&0\\
 N_9&0&R&u_5\\
 K_1&0&0&0\\
 K_2&J_1+J_2&R&0\\
 K_3&J_1&0&0\\
 K_4&J_2&0&0\\
 K_5&0&0&0\\
 P_m&W_8&0&0\\
 P_n&0&P+Q&z_5\\
 P_o&L_8+M_8&P+Q&z_5.
\end{array}
 \tag{D.2.20h6}
\]
The $N_9$ row uses the top Nishida formula, not its
version below the top operation. Here $\xi c=Q^5c$,
$|c|=4$, and the nonlinear partition $(0,2)$ gives
\[
 \operatorname{Sq}_*^2(\xi c)
 =Q^3c+Q^4\beta c+\{c,\operatorname{Sq}_*^2c\}
 =\{c,a\}=R.
 \tag{D.2.20h7}
\]
The first term vanishes by instability and $\beta c=0$.
For the fourth square the only nonzero binomial term
is the term indexed by two:
\[
 \operatorname{Sq}_*^4(\xi c)
 =Q^3(\operatorname{Sq}_*^2c)=Q^3a=\xi a.
 \tag{D.2.20h8}
\]
Both possible nonlinear partitions have a zero factor.
The added terms $aR+A_4u_5$ have zero second and fourth
squares. These coefficients follow directly from the
primary formula in
\cite[Chapter~III, Theorem~1.3(3), pp.~217--218]
{I_CohenLadaMay1976}.

For the polynomial rows, Cartan on the Serre generators gives
\[
\begin{split}
 \beta m&=\beta n=0,\qquad \beta o=j+l,\\
 \operatorname{Sq}_*^2m&=0,\qquad
 \operatorname{Sq}_*^2n=\operatorname{Sq}_*^2o=g+h,\\
 \operatorname{Sq}_*^4m&=0,\qquad
 \operatorname{Sq}_*^4n=\operatorname{Sq}_*^4o=e.
\end{split}
\]
For example
$\operatorname{Sq}^2(u^2v)=v^3+u^2w$
and $\operatorname{Sq}^4(uv)=u^2w+v^3$.
Consequently
\[
\begin{split}
 \beta P_m&=af+af+a^4=W_8,\\
 \beta P_o&=j+l+ax+b(d+e)+ab^2=L_8+M_8,\\
 \operatorname{Sq}_*^2P_n&=\operatorname{Sq}_*^2P_o=P+Q,\\
 \operatorname{Sq}_*^4P_n&=\operatorname{Sq}_*^4P_o=e+ab=z_5.
\end{split}
\]
The bracket entries follow from its Cartan rule;
in particular $\operatorname{Sq}_*^2K_2=R$.
Thus the top operation is not the only source of
$R$ in this degree. The unit formulas use
\cite[\S2, nos.~8--9]{I_Serre1953}.

\medskip\noindent
\textbf{All second Bocksteins needed for integral normalization.}
The first-Bockstein images in degree eight span
$W_8,J_1,J_2,L_8+M_8$. Its cycle space is spanned
by those four classes and $C_8$. Thus the second
page in degree eight is one-dimensional on $[C_8]$.
There are two essential differentials:
\[
 \beta_2[P_n]=[C_8],\qquad
 \beta_2[P_m+Z_A]=[C_8].
 \tag{D.2.20h9}
\]
We check both, since omitting the second changes the
homotopy-group presentation below.
On the second page for $T$, $[a]=0$ because
$\beta b=a$, so $[P_n]=[n]+[c][d]$.
The unit $B\to T$ preserves the known differential
$\beta_2[d]=[c]$ belonging to $H_4B=\ZZ/4$.
Multiplicativity for the Pontryagin product gives
$\beta_2([c][d])=[c]^2=[C_8]$.
The term $\beta_2[n]$ lies in the base factor,
whereas $P_n$ and $C_8$ are fiber classes.
The space splitting $T\simeq B\times F$ splits
the Bockstein spectral sequence as a tensor product
over $\FF_2$. Its positive-degree base and fiber
factors have zero intersection. Therefore
$\beta_2[n]=0$ and the first equality follows.

For the second, compute
\[
 P_m+Z_A=m+ag+ce+bf,
 \qquad \beta(bg)=ag+bf.
\]
Hence $[P_m+Z_A]=[m]+[c][e]$.
The equality $\beta x=d+e$ gives $[e]=[d]$.
The same multiplicativity and base--fiber argument
forces $\beta_2[m]=0$ and proves the second equality
in \textup{(D.2.20h9)}.
These formulas and the universal coefficient theorem give
\[
 H_8F=\ZZ/4\oplus(\ZZ/2)^4.
 \tag{D.2.20h10}
\]

For completeness, the entire first-Bockstein cycle
space in degree nine has basis
\[
 U_A,N_9,K_1,K_5,K_2+K_3+K_4,P_n,P_m+Z_A.
 \tag{D.2.20h11}
\]
The second Bockstein vanishes on the first five.
For three of them this follows already from
\[
 U_A=\beta(z_5u_5),\quad
 K_1=\beta(\{b,f\}+bR),\quad
 K_5=\beta\{b,k\},
\]
followed, where necessary, by the actual projection
from $T$ to $F$. We next give integral representatives
for the other two. Thus \textup{(D.2.20h9)} accounts
for every nonzero second Bockstein relevant to the
nine-cell boundary matrix.

\medskip\noindent
\textbf{Two specified classes have genuine integral cycle lifts.}
Choose integral label chains $a_2,b_3,c_4,y_5,x_6$
with $db=2a$, $dc=dy=0$, and $dx=2y$, where the
reduction of $y$ is $d+e$. The coefficient two is
the already computed exponent of $H_5B$; even
chain changes give these compatible label choices.
For a degree-one Browder chain operation the differential is
\[
 d\lambda(s,t)=-\lambda(ds,t)-(-1)^{|s|}\lambda(s,dt).
\]
It follows that $\lambda(a,x)-\lambda(b,y)$ is an
integral cycle. Its reduction is exactly
$K_2+K_3+K_4$, since the two $bk$ terms cancel.

Since $c$ is an integral cycle of even degree,
the configuration chain $e_1(c\otimes c)$ is an
integral cycle. The terms $aR$ and $a^2\xi a$
are also products of integral cycles, so their
sum gives an integral lift of $N_9$ in $T$.
The configuration-chain conventions are those of
\cite[Chapter~III, proof of Lemma~3.10, p.~232]
{I_CohenLadaMay1976}.
Apply the actual space projection $T\simeq B\times F\to F$
to these chains. Their mod-two classes are unchanged,
because they lie in the fiber factor. Both can
therefore be chosen as zero-integral-boundary nine-cells
in a parity-preserving stable cellular model.

The class $P_m+Z_A+P_n$ has zero first and second
Bocksteins. Since \textup{(D.2.20h10)} has exponent
at most four, the same filtered integral normalization
makes it a zero-boundary nine-cell. Its operation
vector is $(P+Q,z_5)$ in the second and fourth columns.
This additional relation is retained in the following
complete component calculation.

\medskip\noindent
\textbf{All stable degree-eight components and their actual maps.}
Work two-locally, and take a finite stable approximation to
$\Sigma^\infty F$ through dimension ten before cancelling odd-unit
cell pairs. Retain its nine-skeleton, with lower skeleta $L_i$.
This preserves the required stable group and removes extraneous
nine-cells killed by unit ten-cell incidences. The lower model is
the one in D.2.18--D.2.19:
\[
 L_6=C\vee M(\ZZ/4,5)\vee S^6_{\rho_0},
 \qquad
 C=\cofib\bigl((2,\eta_4):S^4\vee S^5\to S^4\bigr).
 \tag{D.2.20t1}
\]
All groups of the $L_i$ in this calculation are stable groups.
The quotient $C\to S^5\vee S^6$ injects $\pi_8C$ into
$\ZZ/8\{\nu_5\}\oplus\ZZ/2\{\eta_6^2\}$.
Indeed the preceding term is $\pi_4^S=0$, and the following
map is $(x,y)\mapsto2x+4y$ in $\ZZ/8$.
Its kernel is cyclic of order four. Choose its generator $\omega$
so that
\[
 \pi_8C=\ZZ/4\{\omega\},\qquad
 q_*\omega=(2\nu_5,\eta_6^2),\qquad
 \delta\eta_6^2=2\omega.
 \tag{D.2.20t2}
\]
The last equality follows from $q_*\delta=(\eta_5,0)$ and
$\eta^3=4\nu$ two-locally. The stable stems and Moore
coextension formulas used here are the opened formulas in
\cite[\S2, (2.1)--(2.2), p.~149]{I_Oka1986}.

For the other summands, the Moore exact sequence and the
order-two lift $b_0$ of D.2.19 give
\[
\begin{split}
 \pi_8M(\ZZ/4,5)
   &=\ZZ/4\{n_0=\iota_5\nu_5\}
       \oplus\ZZ/2\{b_0\eta_7\},\\
 \pi_8S^6_{\rho_0}&=\ZZ/2\{c_0\eta_7\},
       \qquad c_0=\rho_0\eta_6.
\end{split}
 \tag{D.2.20t3}
\]
Here $\iota_5$ is the bottom Moore inclusion, whose
homology class is $u_5$.
The top pinch sends $b_0\eta_7$ to $\eta_6^2$;
it has order two because $b_0$ does.
Also $t\eta_6^2=i\eta^3=0$ in the mod-four Moore summand.

The seven-cell attachments
\[
 \gamma+t,\quad \gamma+2\rho_0+\delta,\quad\delta,\quad0
\]
have $\eta^2$-images $0,2\omega,2\omega,0$.
Their maps on the preceding $\eta$-groups have kernel
generated by $\eta$ on the cell combinations $P+Q$ and $R$.
Thus the cofiber sequence gives
\[
\begin{split}
 0\longrightarrow&
 \ZZ/2\{\overline\omega\}\oplus\ZZ/4\{n_0\}
 \oplus\ZZ/2\{b_0\eta_7,c_0\eta_7\}\\
 \longrightarrow&\ \pi_8L_7
 \longrightarrow (\ZZ/2)^2\longrightarrow0.
\end{split}
 \tag{D.2.20t4}
\]
The $R$ class has the actual order-two lift
$\ell_0=\lambda\eta_7$, since that seven-cell is freely wedged.
For the other lift, change the seven-cell basis from $P,Q$ to
$P+Q,Q$. The new $P+Q$ attachment is $\gamma+2\rho_0$.
Choose a lift $\vartheta$ of its top $\eta$ in the subcomplex
\[
 D=(C\vee S^6_{\rho_0})
          \mathop{\cup}_{\gamma+2\rho_0}e^7_{P+Q}.
 \tag{D.2.20t5}
\]
It exists because $\gamma\eta_6=0$ and $2\rho_0\eta_6=0$.
The element $2\vartheta$ comes from $\pi_8(C\vee S^6)$;
in $\pi_8L_7$ it therefore belongs to the subgroup generated
by $\overline\omega$ and $c_0\eta_7$. In particular it has
no $n_0$ component. This specified representative excludes
an ambiguity that an abstract group extension alone would leave.

\medskip\noindent
\textbf{The precise $\kappa\eta$ relation.}
Set $C'=C\cup_\delta e^7_Q$. Equation
\textup{(D.2.20t2)} and $\delta\eta_6\ne0$ give
\[
 \pi_8C'=\ZZ/2\{\overline\omega\}.
\]
Choose $\kappa'\in\pi_7C'$ with top-cell degree two on $Q$.
In its cone the new degree-eight class has Bockstein $Q$.
Since $\operatorname{Sq}_*^2Q=z_5$ and
$\beta z_5=A_4$, the Adem identities force its
second-square vector to be $e_6$, the only six-class in $C'$.
The quotient of $C'$ by its bottom sphere is
\[
 C'/S^4\simeq C\eta_{5}\vee S^6_{e_6},
\]
where $C\eta_5$ has cells in dimensions five and seven.
The $S^6$ component of $\kappa'$ is therefore $\eta_6$.
Consequently $\kappa'\eta_7$ has a nonzero $\eta_6^2$
component in this quotient, and hence equals
$\overline\omega$.

The normalized $\kappa$ of D.2.19 has second-square vector
$d_6$. Both $\kappa$ and $\kappa'$ project to twice the
$Q$ cell. Their difference in $\pi_7L_7$ is therefore
torsion. The complete detectors of D.2.19 show that this
difference is $b_0$ plus possibly $v_0$: the required
second-square difference is $d_6+e_6$, and no $c_0$ term
is possible. As $v_0\eta_7=i\nu\eta=0$, we obtain
\[
 \kappa\eta_7=\overline\omega+b_0\eta_7.
 \tag{D.2.20t6}
\]
This calculation uses the fixed quotient normalization of
$\gamma$ and the fixed second-square normalization of
$\kappa$ from the preceding claims.

\medskip\noindent
\textbf{The entire group before the nine-cells.}
The eight-cell attachments of D.2.19, ordered as
$W_8,C_8,L_8,M_8,J_1,J_2,J_3$, are
\[
 0,\quad v_0,\quad\kappa,\quad
 \kappa+b_0+c_0,\quad b_0,\quad0,\quad2\lambda+b_0.
\]
Their $\eta$-images generate exactly
$\langle\kappa\eta_7,b_0\eta_7,c_0\eta_7\rangle$.
By \textup{(D.2.20t6)}, these kill
$\overline\omega,b_0\eta_7,c_0\eta_7$ in
\textup{(D.2.20t4)}. They also kill $2\vartheta$, by
the representative in \textup{(D.2.20t5)}, and impose
no relation on $n_0$. The resulting torsion subgroup is exactly
\[
 T_8=\ZZ/4\{n_0\}\oplus
        \ZZ/2\{\vartheta,\ell_0\}.
 \tag{D.2.20t7}
\]
The kernel on the following free degree-seven cell groups
is freely generated by
\[
 W_8,\quad J_2,\quad2C_8,\quad2J_1,\quad2(M_8-L_8).
\]
The extension by this free group splits. More invariantly,
write $D_8=M_8-L_8$ in integral cellular coordinates.
The integral Hurewicz map of $L_8$ has kernel precisely
$T_8$ and image
\[
 \Lambda=\ZZ_{(2)}\{W_8,J_2\}
       \oplus2\ZZ_{(2)}\{C_8,J_1,D_8\}
 \subseteq H_8(L_8;\ZZ_{(2)}).
 \tag{D.2.20t8}
\]
In particular the free lift of $2C_8$ will be denoted $a_0$.

For an attaching map with zero integral boundary, the complete
torsion detectors needed here are
\[
\begin{array}{c|c|c}
 &\operatorname{Sq}_*^2&\operatorname{Sq}_*^4\\ \hline
 n_0&0&u_5\\
 \vartheta&P+Q&z_5\\
 \ell_0&R&0.
\end{array}
 \tag{D.2.20t9}
\]
The $n_0$ row is the Hopf-invariant-one cone on the
Moore bottom, and the $\ell_0$ row is the cone of $\eta$
on the $R$ seven-sphere. For $\vartheta$, its top quotient
is $\eta$ on $P+Q$, giving its second square. Its chosen
representative in \textup{(D.2.20t5)} has no $u_5$ class,
so its fourth-square $u_5$ coefficient is zero. Finally,
dualize the Adem identity
$\operatorname{Sq}^2\operatorname{Sq}^3
=\operatorname{Sq}^5+\operatorname{Sq}^4\operatorname{Sq}^1$.
The top Bockstein is zero, while
$\operatorname{Sq}_*^3(P+Q)=A_4$ and $\beta z_5=A_4$.
Its fourth-square $z_5$ coefficient must therefore be one.
For the order-four class $n_0$, these primary detectors
see only the parity of its coefficient; we retain that
limitation in the next step.

\medskip\noindent
\textbf{Three zero-boundary relations kill the entire torsion subgroup.}
The genuine integral cycles $N_9$ and $K_2+K_3+K_4$ constructed
above can be chosen as zero-boundary nine-cells. By their
operation vectors and \textup{(D.2.20t9)}, their attachments
have the respective forms
\[
 \ell_0+(2r+1)n_0,\qquad \ell_0+2s n_0.
 \tag{D.2.20t10}
\]
Neither has a $\vartheta$ component. Subtracting the two
relations gives an odd multiple of $n_0$, so they kill
the entire order-four class $n_0$, and then $\ell_0$.
No value of an undetected even coefficient has been assumed.

The additional cycle $Z_0=P_m+Z_A+P_n$ has zero first and
second Bocksteins by \textup{(D.2.20h9)}. The integral
normalization explained above makes it a zero-boundary
nine-cell. Its operation vector is $(P+Q,z_5)$, so its
attachment is $\vartheta$ plus an even multiple of $n_0$,
with no $\ell_0$ term. It kills $\vartheta$ after the
preceding two relations. Thus all of $T_8$ is killed.
Both second Bocksteins in \textup{(D.2.20h9)} are essential
here; omitting the one on $P_m+Z_A$ would leave an incorrect
order-four stable summand.

\medskip\noindent
\textbf{The remaining quotient is the integral Hurewicz image.}
Let $\mathcal B_9$ be the image of the integral ninth
cellular boundary in $H_8L_8$. Every such boundary belongs
to the lattice $\Lambda$. Since the whole Hurewicz kernel
$T_8$ has already been killed by zero-boundary cells,
the complete stable homotopy quotient is exactly
\[
 \pi_8\Sigma^\infty F=\Lambda/\mathcal B_9.
 \tag{D.2.20t11}
\]
Two attachments with the same integral boundary differ
by $T_8$, so no free-lift extension remains in this formula.
In the integral homology calculation, the four
first-Bockstein images supply order-two lifts of
$W_8,J_1,J_2,D_8$, and $C_8$ is the order-four generator.
Consequently \textup{(D.2.20t8)} identifies the quotient
with the following subgroup of $H_8F$:
\[
 \pi_8\Sigma^\infty F
 \xrightarrow[\cong]{\ h_{\ZZ}\ }
 \langle2C_8,W_8,J_2\rangle
       \subseteq\ZZ/4\{C_8\}\oplus
           \ZZ/2\{W_8,J_1,J_2,D_8\}.
 \tag{D.2.20t12}
\]
Changing an order-two lift by $2C_8$ does not change
this subgroup. It is $(\ZZ/2)^3$. In particular, if
$a_0$ denotes the class with integral Hurewicz image
$2C_8$, then
\[
 \ker\bigl(h_{\FF_2}:\pi_8\Sigma^\infty F
                     \to H_8(F;\FF_2)\bigr)
       =\ZZ/2\{a_0\}.
 \tag{D.2.20t13}
\]
All nine-cell attachments have been accounted for:
the zero-boundary ones killed $T_8$, and their integral
boundary subgroup gives exactly \textup{(D.2.20t11)}.
Higher cells cannot change this stable group.

\medskip\noindent
\textbf{The ordinary Hurewicz map is injective.}
Write $F=\Omega^2 Z$. The double suspension--loop unit
\[
 F\longrightarrow\Omega^2\Sigma^2F
\]
has a homotopy retraction by the adjunction triangle
identity. Since $\Sigma^2F$ is five-connected, its
tenth homotopy group is stable. Therefore stabilization
gives a split injection
\[
 \pi_8F\lhook\joinrel\longrightarrow\pi_8\Sigma^\infty F.
 \tag{D.2.20t14}
\]
In particular the stable zeros $n_0=0$ and $\ell_0=0$
already prove the actual zeros in \textup{(D.2.20h2)}.

The remaining stable mod-two-Hurewicz-kernel element
$a_0$ is not an actual sphere map. Its integral
Hurewicz image is $2C_8$, with $C_8$ of order four.
In its stable cone a new degree-nine mod-two class $x_9$
satisfies
\[
 \beta x_9=C_8\pmod{\operatorname{im}\beta|_{H_9F}}.
\]
The old first-Bockstein images are
$W_8,J_1,J_2,L_8+M_8$; all have zero fourth square.
Since $\operatorname{Sq}_*^4C_8=A_4$, the Adem identity
$\operatorname{Sq}^1\operatorname{Sq}^4
=\operatorname{Sq}^5$ gives
\[
 \operatorname{Sq}_*^5x_9
   =\operatorname{Sq}_*^4\beta x_9=A_4\ne0.
 \tag{D.2.20t15}
\]
If $a_0$ came from an actual map $S^8\to F$, its
ordinary space cone would have this operation.
But $\operatorname{Sq}^5$ on degree-four cohomology
is zero by instability. This contradiction excludes
$a_0$ from the image of \textup{(D.2.20t14)}.
By \textup{(D.2.20t13)}, that image meets the stable
mod-two Hurewicz kernel only in zero. Hence
$\pi_8F\to H_8(F;\FF_2)$ is injective, as claimed.
No choice of a stable quadratic complement and no
identification of stable and ordinary eighth homotopy
groups has been used.

\medskip\noindent
\textbf{A source decomposition that treats every nine-cell at once.}
We now use the target conclusions to extend the ordinary map,
keeping the choices made through dimension eight explicit.
Put $Q_i=K_i/K_4$ as in D.2.18--D.2.19, and work with the
two-local quotient cell models there. The five eight-cell
attachments in \textup{(D.2.19w)} give, by cell slides,
\[
 Q_8\simeq D'\vee M(\ZZ/2,7)_F\vee M(\ZZ/4,7)_J
                  \vee C_1^8\vee S^8_s,
 \tag{D.2.20s1}
\]
where
\[
\begin{split}
 D&=\operatorname{cofib}\left(
  \begin{pmatrix}4&\eta_5\\0&2\end{pmatrix}:
  S^5\vee S^6\longrightarrow S^5\vee S^6\right),\\
 D'&=D\cup_{\tau_d+\tau_e}e^8_H,\qquad
 C_1^8=\operatorname{cofib}\bigl((2,\eta_6):S^6\vee S^7\to S^6\bigr).
\end{split}
 \tag{D.2.20s2}
\]
Indeed slide the cell $G_8$ over $H_8$, leaving the
degree-two attachment on $F$; slide $B_0^2$ over $A_0C$,
leaving the degree-four attachment on $J$; and slide
$Q^5A_0$ over $A_0C$, leaving the free eight-sphere.
The integral orientation of the last sphere is
$s=Q^5A_0-A_0C$. The cell $A_0C$ supplies the top
cell of $C_1^8$, while $H_8$ supplies the top cell of $D'$.
These are actual comparisons for the mapping problem:
the quotients are four-connected, their dimensions are
at most eight, and Freudenthal realizes the stable cell
slides unstably. Maps into the two-local space $T$ may
then be pulled back to the original finite skeleta.

In D.2.19 the correction is zero on the source six-sphere
$z$ and on the free seven-spheres $F,J$. Its nullhomotopies
on the eight-cells may therefore be chosen so that the
correction is constant on all the factors of
\textup{(D.2.20s1)} except $D'$. This makes its restrictions
to the two Moore factors, $C_1^8$, and the last sphere
literally zero. On $D'$ it retains the Moore extension
$\mu$, the bottom class $\overline\alpha$, and the
six-sphere value $\rho$. Its extension across the
$H$-cell can be any of the allowed nullhomotopies;
we will not assume that its effect on the next
homotopy group is an arbitrary free choice.

\medskip\noindent
\textbf{Every class in the remaining attaching group has zero
mod-two Hurewicz image.}
In $D$, the two classes $\tau_d,\tau_e$ are independent
order-two elements, by the complete calculation
\textup{(D.2.19p)}. Thus $\tau_d+\tau_e$ is nonzero
and has order two. The cofiber sequence
\[
 S^7\xrightarrow{\tau_d+\tau_e}D\longrightarrow D'
             \xrightarrow{p_H}S^8
 \tag{D.2.20s3}
\]
gives
\[
 \operatorname{im}\bigl((p_H)_*:\pi_8D'\to\pi_8S^8\bigr)
                  =2\ZZ_{(2)}.
 \tag{D.2.20s4}
\]
This is first computed in the stable exact sequence;
four-connectivity identifies all the relevant groups
with their ordinary counterparts. The complex $D$
has dimension seven and $H_7(D;\ZZ_{(2)})=0$:
its top cellular differential has nonzero degree two
on the free six-cell. The homology sequence of
\textup{(D.2.20s3)} consequently identifies
$H_8(D';\ZZ_{(2)})$ with $\ZZ_{(2)}$, by the
pinch onto the $H$-cell. It follows from
\textup{(D.2.20s4)} that every actual sphere map
$S^8\to D'$ has even degree on this sole integral
eight-dimensional homology generator. Therefore
\[
 h_{\FF_2}:\pi_8D'\longrightarrow H_8(D';\FF_2)
                  \quad\text{is zero}.
 \tag{D.2.20s5}
\]
There is no additional Tor summand in this last
homology group, since $H_7D'=0$. This argument
includes the entire torsion subgroup of $\pi_8D'$;
it does not infer its vanishing from incomplete
primary-square detectors.

\medskip\noindent
\textbf{The correction kills all actual nine-cell attachments.}
Let $f:D'\to T$ be the correction, with any of its
allowed fillings of the $H$-cell. It lifts to
$F=\fib(T\to B)$, since
$\Map_*(D',K(\FF_2,2))$ is weakly contractible.
Naturality and \textup{(D.2.20s5)} make the image
of every element of $\pi_8D'$ have zero mod-two
Hurewicz in $F$. The target result proved above
makes the actual Hurewicz homomorphism on $\pi_8F$
injective. Thus
\[
 f_*:\pi_8D'\longrightarrow\pi_8T
                       \quad\text{is zero}.
 \tag{D.2.20s6}
\]
The correction is constant on the other wedge
factors in \textup{(D.2.20s1)}. There are no cross
Whitehead terms in dimension eight by their
connectivities. Consequently its homomorphism
on all of $\pi_8Q_8$ is zero. Every nine-cell
attaching map has zero image, regardless of any
component not measured by primary operations.
Choose nullhomotopies on these cells to extend
the correction to $Q_9$.

Finally, the canonical summand already constructed
in \textup{(D.2.19y)--(D.2.19z)} extends through
source relative dimension nine. We keep that same
map and add the pullback of the correction along
$K_9\to Q_9$. This gives an actual ordinary map
$G_9:K_9\to T$ extending the allowed $G_8$ and
retaining every six-dimensional value in D.2.17.
The canonical summand has only been supplied
through dimension nine, and no extension over
the tenth or later cells follows from this proof.

\end{proof}

\begin{claim}{D.2.21}{Proved}{The exact tenth-stage criterion for the canonical summand}
Keep $X=\Sigma K(\FF_2,2)$, $Y=\Sigma^2K(\FF_2,2)$,
$T=\Omega^2Y$, and the factorization
$\overline h:\Sigma^2X'\to Y$, $X'=X/S^3$, of
\textup{(D.2.19y)}. Let $\chi:S^4\to T$ be its double
adjunct on the bottom four-sphere of $X'$; thus
$h_{\FF_2}(\chi)=a^2$. Write
$\widehat\chi:S^6\to Y$ for its double adjunct.

Fix the canonical map $g_2:C_2X\to T$ and the relative
eight-stage extension supplied by \textup{(D.2.19z)}.
There is an extension through relative dimension ten,
allowing the nine-cell fillings to vary, if and only if
\[
 [\widehat\chi,\widehat\chi]\in4\pi_{11}Y.
 \tag{D.2.21a}
\]
Equivalently, the homotopy Browder square of $\chi$ is zero
in $\pi_9T/4\pi_9T$. This criterion concerns the specified
canonical summand and fixed eight-stage choice. It does not
allow arbitrary changes on the eight-cells, and it does not
decide the separate correction's tenth-stage extension.
\end{claim}
\begin{proof}
Use the labelled configuration models for $C_2X$ and $QX$;
the approximation equivalences apply because $X$ is connected
\cite[Construction~2.4; Theorem~6.1]{May1972}.
The spaces of ordered configurations are connected. Their
maps to the corresponding contractible free symmetric-group
spaces therefore admit relative free CW models with cells
of dimension at least two: attach equivariant two-cells to
kill the fundamental group and then cells of higher dimension.
The label filtration consequently has no relative cells of
weight at least three below dimension eleven, since every
positive label of $X$ has dimension at least three.
Weight one is unchanged. Thus the binary relative complex
is the entire ordinary relative cell model through dimension
ten; a stable splitting alone is not being used to identify
these attaching maps.

Write $a_3,b_4,c_5$ for the three bottom label cells of $X$.
Their integral cellular differentials are
$\partial b_4=2a_3$, $\partial a_3=\partial c_5=0$.
For two labels the configuration spaces are the antipodal
circle and contractible antipodal infinite sphere. Choose
one free $\Sigma_2$-cell $e_i$ in each dimension of the
latter. The relative cells have $i\ge2$, and their
configuration boundaries alternate $1+\tau$ in even
degree and $\tau-1$ in odd degree. The tensor permutation
includes the Koszul sign. The complete list through degree
ten is
\[
\begin{array}{c|l}
8&q=e_2(a_3\otimes a_3)\\
9&r=e_3(a_3\otimes a_3),\quad
  t=e_2(a_3\otimes b_4),\quad
  t'=e_2(b_4\otimes a_3)\\
10&u=e_4(a_3\otimes a_3),\quad
  s=e_3(a_3\otimes b_4),\quad
  s'=e_3(b_4\otimes a_3),\\
 &v=e_2(b_4\otimes b_4),\quad
  e_2(a_3\otimes c_5),\quad e_2(c_5\otimes a_3).
\end{array}
 \tag{D.2.21b}
\]
The relative integral differentials are exactly
\[
\begin{gathered}
 \partial r=-2q,\qquad \partial t=-2q,
 \qquad \partial t'=2q,\\
 \partial s=2r-t+t',\qquad
 \partial s'=-2r+t-t',\qquad
 \partial v=2t+2t',
\end{gathered}
 \tag{D.2.21c}
\]
and the other three degree-ten differentials vanish.
These formulas include the internal differential
$\partial b_4=2a_3$; omitting it would give the wrong
divisibility condition. The full saturated kernel in
degree nine has basis $t-r,t'+r$. The nonzero boundary
columns in that basis are $(-1,1),(1,-1),(2,2)$, with
Smith form $\operatorname{diag}(1,4)$. In particular its
relative ninth homology is $\ZZ/4$.

The quotient $X\to X'$ sends every relative cell in
\textup{(D.2.21b)} except $v$ into $C_2X'$: a label in
the collapsed bottom sphere is removed. It supplies the
specified eight-stage map and a choice of the nine-cell
fillings for which every tenth-stage obstruction except
the coefficient on $v$ is zero.

We identify that coefficient homotopically, including its
normalization. On the bottom four-sphere of $X'$, the first
stabilization is $\Omega^2S^6\to\Omega^3S^7$.
The pair $S^6\to\Omega S^7$ is eleven-connected. Its first
relative generator is the twelve-cell of $J_2S^6$.
The latter is obtained from $S^6\times S^6$ by identifying
the two wedge summands with $S^6$, so its attaching map is
exactly $[\iota_6,\iota_6]$. Double looping shifts the
relative dimension to ten. This gives the Browder square
as the boundary of $v$; the integral boundary $2e_1(b_4
\otimes b_4)$ is its usual Hopf-invariant-two normalization,
not an additional factor of two on the Whitehead square.
The next configuration boundary is $\tau-1=0$ on these
even-dimensional labels and does not change this generator.
Thus the obstruction coefficient is
\[
 \theta=\lambda_1(\chi,\chi)\in\pi_9T,
 \qquad \theta\longleftrightarrow
 [\widehat\chi,\widehat\chi]\in\pi_{11}Y.
 \tag{D.2.21d}
\]

Finally keep the eight-stage extension fixed. The possible
changes of the three nine-cell fillings are arbitrary
$R,T_0,T_1\in M=\pi_9T$. Cellular obstruction theory and
\textup{(D.2.21c)} change the coefficients on $s$ and $v$ by
\[
 d_s=2R-T_0+T_1,\qquad d_v=2T_0+2T_1;
 \qquad d_v-2d_s=4(T_0-R).
 \tag{D.2.21e}
\]
The change on $s'$ is $-d_s$ and the other three coefficients
are unchanged. Consequently killing every coefficient
forces $\theta\in4M$. Conversely, if $\theta=4m$, choose
$R=0$ and $T_0=T_1=-m$; this preserves the zero
coefficients and cancels $\theta$. The identity holds for
every abelian group $M$, including its two-torsion, and
proves both directions. The integral chain matrix and
cochain identity are independently checked in
\texttt{code/D2\_canonical\_relative\_chain.py}.
No secondary contribution from changing an eight-cell
filling has been discarded: those fillings were fixed
throughout the assertion.
\end{proof}

\begin{claim}{D.2.22}{Proved}{The tenth-stage target: exponent, integral Hurewicz kernel, and eta composites}
Keep $B=K(\FF_2,2)$, $Y=\Sigma^2B$, $T=\Omega^2Y$, and the
three-connected fiber $F$ of \textup{D.2.18--D.2.20}.
Let $\rho\in\pi_6T$ be the order-four element of
\textup{D.2.18}. Then
\[
 4\pi_9T=0,\qquad
 \ker\bigl(h_{\ZZ}:\pi_9T\longrightarrow H_9(T;\ZZ)\bigr)
   =\ZZ/2\{\rho\nu_6\},\qquad \rho\nu_6\ne0,
 \tag{D.2.22a1}
\]
and
\[
 \pi_8T\circ\eta_8=0.
 \tag{D.2.22a2}
\]
The group $\pi_9T$ is either $(\ZZ/2)^4$ or
$\ZZ/4\oplus(\ZZ/2)^2$. The calculation retains the unresolved
common even-nu attachment coefficient that distinguishes these
two alternatives. For the bottom class $\chi\in\pi_4F$ used in
\textup{D.2.21}, one has the precise restriction
\[
 \lambda_1(\chi,\chi)\in\{0,\rho\nu_6\}.
 \tag{D.2.22a18}
\]
Its value is not selected. The criterion of \textup{D.2.21}
therefore reduces to vanishing of this actual square. No extension
of the combined map through the tenth skeleton is asserted.
\end{claim}
\begin{proof}
All displayed homology operations have coefficients $\FF_2$ unless
integral coefficients are indicated. The spaces have finite
two-primary homotopy groups in the degrees used, by the Serre-class
argument already applied to $Y$. Thus the two-local component
calculation below computes the actual groups in the statement.

\medskip\noindent
\textbf{The complete next fiber homology table.}
Continue the label notation of \textup{D.2.20}. Put
\[
 \begin{gathered}
 A=a^2,\quad z=e+ab,\quad u_F=\xi a,\quad
 k=\{a,b\},\quad R=\{a,c\},\\
 d_6=b^2,\quad e_6=b^2+k,\quad f_6=f+ac+a^3.
 \end{gathered}
 \tag{D.2.22a3}
\]
Here the unadorned $d,e$ remain the degree-five label classes,
dual to $w_B$ and $u_Bv_B$; they are not $d_6,e_6$.
The previously computed classes are
\[
\begin{gathered}
 P=g+ae+bc+a^2b,\qquad Q=h+ad,\\
 W=a^4,\quad C=c^2,\quad L=l+bd,\quad
 M=j+ax+be+ab^2,\\
 J_1=\{a,d\},\quad J_2=\{a,e\}+ak,\quad
 J_3=\{b,c\}+ak.
\end{gathered}
\]
Write $p,q,r,s$ for the new degree-ten labels dual to
$u_B^5,u_B^2v_B^2,u_Bv_Bw_B,w_B^2$.
The free double-loop basis has twelve indecomposables in degree ten:
\[
\begin{gathered}
 p,q,r,s,\\
 \{a,g\},\{a,h\},\{b,f\},\{b,x\},
 \{c,d\},\{c,e\},\{b,k\},\{a,R\}.
\end{gathered}
\]
The restricted identities identify $\{a,\xi b\}$ with $\{b,k\}$
and $\{c,\xi a\}$ with $\{a,R\}$.
The thirty products are
\[
\begin{gathered}
 ai,aj,al,a\{a,d\},a\{a,e\},a\{b,c\},\\
 bg,bh,b\xi b,bR,cf,cx,ck,\\
 d^2,de,du_F,e^2,eu_F,u_F^2,\\
 a^2f,a^2x,a^2k,abd,abe,abu_F,\\
 ac^2,b^2c,a^3c,a^2b^2,a^5.
\end{gathered}
\]
Thus $\dim H_{10}T=42$. The base and mixed summands of
$H_{10}(B\times F)$ have dimensions $4+2+4+3+4+7=24$,
so $\dim H_{10}F=18$.

A complete fiber basis consists of
\[
\begin{gathered}
 D_A=Ad_6,\quad E_A=Ae_6,\quad F_A=Af_6,
 \quad Z_2=z^2,\quad ZU=zu_F,\quad U_2=u_F^2,
 \quad D_2=d^2,\\
 Q_{10}=q+cx+aM+bP+abz,\\
 P_p=p+ai+cf+a^2f+ac^2+a^3c+a^5,
 \qquad P_r=r+al+bh+ed,
\end{gathered}
\tag{D.2.22a4}
\]
and the eight classes
\[
\begin{aligned}
 B_1&=\{a,g\}+a\{a,e\}+ck+bR+Ak,&
 B_2&=\{a,h\}+a\{a,d\},\\
 B_3&=\{b,f\}+ck+a\{b,c\}+Ak,&
 B_4&=\{b,x\},\\
 B_5&=\{c,d\}+a\{a,d\},&
 B_6&=\{c,e\}+bR+a\{b,c\}+a\{a,e\}+Ak,\\
 B_7&=\{b,k\},&B_8&=\{a,R\}.
\end{aligned}
\tag{D.2.22a5}
\]
The classes $D_2,P_p,P_r,B_1,\ldots,B_8$ are primitive,
and $\overline\Delta Q_{10}=z\otimes z$.
For example
\[
\begin{aligned}
 \overline\Delta\{a,g\}
 &=a\otimes\{a,e\}+\{a,e\}\otimes a
   +k\otimes c+c\otimes k+b\otimes R+R\otimes b,\\
 \overline\Delta\{b,f\}
 &=k\otimes c+c\otimes k
   +a\otimes\{b,c\}+\{b,c\}\otimes a,\\
 \overline\Delta\{c,e\}
 &=b\otimes R+R\otimes b
   +a\otimes(\{b,c\}+\{a,e\})\\
 &\quad +(\{b,c\}+\{a,e\})\otimes a
   +A\otimes k+k\otimes A.
\end{aligned}
\]
The other bracket diagonals have respectively the two terms canceled
by the displayed product correction; $\{b,x\}$ is primitive because
$\overline\Delta x=b\otimes b$ and $\{b,b\}=0$.
The polynomial diagonal of $Q_{10}$ follows by expanding it as
$q+cx+aj+a^2x+abe+a^2b^2+bg+b^2c$.
All these classes have zero base image; there is no primitive in
$H_{10}B$. Their distinct leading monomials prove independence,
and the dimension count proves completeness.

Write $K_\Sigma=K_2+K_3+K_4$ and
$Z_0=P_m+Z_A+P_n$ in the degree-nine notation of
\textup{D.2.20}. The entire operation table is
\[
\begin{array}{c|ccc}
 &\beta&\operatorname{Sq}_*^2&\operatorname{Sq}_*^4\\\hline
 D_A,E_A,F_A,Z_2&0&W&0\\
 ZU&U_A&0&0\\
 U_2&0&0&0\\
 D_2&0&0&d_6\\
 Q_{10}&Z_0&M+C+W&f_6\\
 P_p&0&W&f_6\\
 P_r&0&0&0\\
 B_1&K_1&J_2&0\\
 B_2&0&J_2&0\\
 B_3&K_1&0&0\\
 B_4&K_\Sigma&J_1+J_2+J_3&d_6+e_6\\
 B_5&0&J_1+J_3&d_6+e_6\\
 B_6&0&J_2&0\\
 B_7&K_5&0&0\\
 B_8&0&0&0
\end{array}
\tag{D.2.22a6}
\]
It follows directly from Cartan and the CLM Browder identities
\cite[Chapter III, Theorems 1.2--1.4]{I_CohenLadaMay1976}.
For the new labels, the triples $(\beta,\operatorname{Sq}_*^2,
\operatorname{Sq}_*^4)$ are $(0,0,f)$ for $p$,
$(m+n,j,f)$ for $q$, $(0,0,0)$ for $r$, and $(t,l,x)$ for $s$.
In particular $\beta Q_{10}=Z_0$, an independent check on the
essential second-Bockstein correction in \textup{D.2.20}.
The coproducts and twenty-one operation values on the seven
polynomial classes are also independently checked by
\texttt{code/FGH\_H10\_poly\_check.py} and the root's separate
polynomial certificate; the bracket calculations above do not
depend on those certificates.

\medskip\noindent
\textbf{The order-eight class and both third Bocksteins.}
The universal-coefficient theorem, the known $H_8F$, and the
rank-five first Bockstein in \textup{(D.2.22a6)} give six cyclic
summands in $H_9(F;\ZZ)$, five of order two. The remaining
higher-exponent class is represented by $N$.
We determine its exact order, including the second order-eight
boundary that will be essential below.

Use the integral chain summand of $B$ with $c$ in degree four,
$d$ in degree five, and $\partial d=4c$. In the arity-two
configuration-circle complex put
\[
 z_c=e_1(cc),\quad v=e_0(cd-dc),\quad
 q_0=e_0(dd),\quad h=e_1(cd),\quad h'=e_1(dc).
\]
The Koszul transposition signs give the exact formulas
\[
\begin{gathered}
 \partial q_0=4v,\qquad
 \partial h=-v-4z_c,\qquad
 \partial h'=v-4z_c,\\
 \partial e_1(dd)=-2q_0-4h+4h'.
\end{gathered}
\tag{D.2.22a7}
\]
The degree-nine cycles have basis $v,z_c$, and the boundary
columns are $(4,0),(-1,-4),(1,-4)$, with Smith invariants
$1,8$. Thus the class of $z_c$ has exact order eight.
Both mixed-label circle chains are needed: using only one
would miss the relation $\partial(h+h')=-8z_c$.

This chain calculation is a retract of the actual arity-two
summand. Indeed a bounded-below free integral chain complex
has a Smith decomposition, so $\widetilde C_*(B;\ZZ)$ retracts
as a chain complex onto the displayed Moore-four summand.
Its tensor square is equivariant for the signed transposition.
The configuration chains are free over $\ZZ[\Sigma_2]$;
the equivariant shuffle equivalence therefore identifies the
derived tensor with the chains of the actual filtration quotient.
The integral stable splitting of Snaith
\cite[Theorem 1.1 and Section 2]{Snaith1974} makes this a
homology retract of $T$. No topological Moore-space retract
of $B$ is assumed.

The integral lift of $N$ has weight-two component $\xi c$;
the other terms $aR$ and $a^2\xi a$ have order dividing two.
Thus it has order eight in $H_9T$. In the actual product splitting
$T\simeq B\times F$, all summands outside $H_9F$ have exponent
at most four. This includes the mixed Kunneth summands, notably
$H_4B\otimes H_5F\cong\ZZ/4$; the base summand itself has
exponent two by the integral Eilenberg--Mac Lane calculation
\cite[\S3.2]{I_Clement2002}. Projection therefore retains order eight in the
fiber. Consequently
\[
 H_9(F;\ZZ)\cong\ZZ/8\oplus(\ZZ/2)^5.
 \tag{D.2.22a8}
\]

Choose the integral labels so that $e=d+y$, $\partial y=0$,
and $\partial x=2y$, allowing odd-unit changes. The same
Moore-four retraction sends both $d$ and $e$ to $d$.
Hence the weight-two parts $\{c,d\}$ of $B_5$ and
$\{c,e\}$ of $B_6$ both project to $h+h'$.
Their higher-weight corrections disappear. This proves
\[
 \beta_3[B_5]=[N],\qquad \beta_3[B_6]=[N].
 \tag{D.2.22a9}
\]
For an exhaustive check, the first-Bockstein cycle basis in
degree ten is
\[
 D_A,E_A,F_A,Z_2,U_2,D_2,P_p,P_r,
 B_2,B_5,B_6,B_8,B_1+B_3.
\]
Only $B_5,B_6$ project to $h+h'$; the others project either
to zero or to $q_0$. The latter is already a first-Bockstein
boundary by \textup{(D.2.22a7)}. Thus these are exactly the
two nonzero third differentials into $N$. Projection also
shows that $N$ cannot be hit earlier and that these two
classes survive to the third page.

\medskip\noindent
\textbf{The four-connected cover and the complete low Koszul check.}
Let $E=F\langle5\rangle$ be the fiber of the first Postnikov
map $p:F\to K(\FF_2,4)$, representing $A^*$.
Serre's polynomial description \cite[Section 2]{I_Serre1953}
gives the following complete base-generator list through degree eleven:
\[
\begin{gathered}
 x_4=x,\ x_5=\operatorname{Sq}^1x,\ x_6=\operatorname{Sq}^2x,
 \ x_7=\operatorname{Sq}^3x,\ y_7=\operatorname{Sq}^2\operatorname{Sq}^1x,\\
 x_8=\operatorname{Sq}^3\operatorname{Sq}^1x,
 \ x_9=\operatorname{Sq}^4\operatorname{Sq}^1x,
 \ x_{10}=\operatorname{Sq}^4\operatorname{Sq}^2x,\\
 x_{11}=\operatorname{Sq}^5\operatorname{Sq}^2x,
 \ y_{11}=\operatorname{Sq}^4\operatorname{Sq}^2\operatorname{Sq}^1x.
\end{gathered}
\]
Their images through degree ten are
\[
\begin{array}{c|l}
 x_4&A^*\\
 x_5&z^*\\
 x_6&d_6^*+e_6^*+f_6^*\\
 x_7&V^*+P^*\\
 y_7&P^*+Q^*\\
 x_8&L^*+M^*\\
 x_9&P_n^*+P_o^*\\
 x_{10}&D_2^*+Q_{10}^*+P_p^*.
\end{array}
\tag{D.2.22a10}
\]
For degree eleven let $v_{11}$ be the label dual to
$u_B\operatorname{Sq}^4w_B$, and retain the degree-nine label
$t$ dual to $\operatorname{Sq}^4w_B$. The primitive fiber
classes $\xi d$ and $v_{11}+at$ satisfy
\[
\begin{array}{c|cc}
 &\beta&\operatorname{Sq}_*^4\\\hline
 \xi d&D_2&V\\
 v_{11}+at&0&Q.
\end{array}
\tag{D.2.22a11}
\]
The top CLM formulas give the first row: the only fourth-square
term is $Q^4b=\xi b$. Label Cartan gives the second row.
There is no primitive in $H_{11}B$, so both lie in the fiber.
They detect $x_{11},y_{11}$ as the identity matrix.
All base products through degree eleven are independently detected
by $C,Z_A,D_A,E_A,F_A,Q_{10}$ and the products
$AH_7F,zH_6F$. Thus $p^*$ is injective through degree eleven.

Injectivity alone would not exclude mixed Serre differentials.
Instead use the convergent cohomological Eilenberg--Moore sequence
\cite[p.334]{FGHMay1968} and its polynomial Koszul complex.
A Koszul generator for a base generator of degree $j$ has total
degree $j-1$. Length three starts in degree $3+4+5=12$;
length two with a positive $H^*F$ coefficient starts in degree
eleven. Hence the entire positive-length calculation through ten
has scalar length-two coefficients and length-one coefficients
of degree at most seven. The exact ranks are
\[
\begin{array}{c|ccc}
 \text{total degree}&\dim K_1&\operatorname{rank}(K_1\to K_0)
 &\operatorname{rank}(K_2\to K_1)\\\hline
 7&2&2&0\\
 8&4&3&1\\
 9&7&6&1\\
 10&13&10&3.
\end{array}
\tag{D.2.22a12}
\]
In degrees three through six the scalar images are independent.
For degree seven the images are $C^*,L^*+M^*$.
For degree eight, $A^*H^5F$ has rank two and $p^*x_9$
adds one; the single relation is the Koszul relation for
$x_4,x_5$. For degree nine, $A^*H^6F$ and $z^*H^5F$
have disjoint ranks three and two, and $p^*x_{10}$ adds one
rank detected on primitive $D_2$; the single relation is
the one for $x_4,x_6$.

For degree ten, $A^*H^7F$ and $z^*H^6F$ have disjoint ranks
four and three. The product $u_F^*p^*x_6$ adds one, while
$z^*p^*x_6$ is already in the latter rank-three space.
The two primitive tests in \textup{(D.2.22a11)} add two more
ranks. The three relations are those for
$(x_4,x_7),(x_4,y_7),(x_5,x_6)$.
All product ranks follow by evaluation on products of the
primitive homology factors; the equal-factor product $z^{*2}$
is detected on $Q_{10}$, not on $z^2$.
The scalar length-two wedges in degrees seven through ten are
\[
 t_3t_4;\quad t_3t_5;\quad
 t_3t_6,t_3t'_6,t_4t_5;\quad
 t_3t_7,t_4t_6,t_4t'_6.
\]
Their differentials are independent. This proves that positive
Tor vanishes through total degree ten. Higher base generators
have Koszul degree at least eleven and cannot contribute.
There is consequently no incoming higher Eilenberg--Moore
differential in this range, and restriction gives
\[
 H^{\le10}E\cong
 \bigl(H^*F/(p^*\widetilde H^*K(\FF_2,4))H^*F\bigr)^{\le10}.
 \tag{D.2.22a13}
\]

\medskip\noindent
\textbf{The cover table used by the component calculation.}
Dually $H_iE\to H_iF$ is injective through degree ten, with
the following complete image basis:
\[
\begin{array}{c|l}
5&u_F\\
6&D=d_6+e_6,\quad E_6=d_6+f_6\\
7&S=V+P+Q,\quad R\\
8&W,\quad D_8=L+M,\quad J_1,J_2,J_3\\
9&N,K_1,K_2,K_3,K_4,K_5,P_m,S_9=P_n+P_o\\
10&Z_2,U_2,T_{10}=D_2+P_p,P_r,B_1,\ldots,B_8.
\end{array}
\tag{D.2.22a14}
\]
Their dimensions are $1,2,2,5,8,12$, and every displayed
class is primitive. Below degree nine the nonzero first,
second, and fourth squares are precisely
\[
 \beta S=D+E_6,\quad \operatorname{Sq}_*^2S=u_F,
 \quad \operatorname{Sq}_*^2D_8=E_6,
 \quad \operatorname{Sq}_*^2J_1=\operatorname{Sq}_*^2J_3=D,
 \quad\beta J_3=R.
\]
In degree nine the table is
\[
\begin{array}{c|ccc}
 &\beta&\operatorname{Sq}_*^2&\operatorname{Sq}_*^4\\\hline
 N&0&R&u_F\\
 K_1&0&0&0\\
 K_2&J_1+J_2&R&0\\
 K_3&J_1&0&0\\
 K_4&J_2&0&0\\
 K_5&0&0&0\\
 P_m&W&0&0\\
 S_9&D_8&0&0
\end{array}
\tag{D.2.22a15}
\]
and in degree ten it is
\[
\begin{array}{c|ccc}
 &\beta&\operatorname{Sq}_*^2&\operatorname{Sq}_*^4\\\hline
 Z_2&0&W&0\\
 U_2&0&0&0\\
 T_{10}&0&W&E_6\\
 P_r&0&0&0\\
 B_1&K_1&J_2&0\\
 B_2&0&J_2&0\\
 B_3&K_1&0&0\\
 B_4&K_\Sigma&J_1+J_2+J_3&D\\
 B_5&0&J_1+J_3&D\\
 B_6&0&J_2&0\\
 B_7&K_5&0&0\\
 B_8&0&0&0.
\end{array}
\tag{D.2.22a16}
\]
All these operations are restrictions of the fiber tables,
so there is no unresolved Steenrod extension from the spectral
sequence.

The integral groups are
\[
\begin{gathered}
 H_5E=\ZZ/4,\quad H_6E=\ZZ/2,\quad H_7E=\ZZ/2,
 \quad H_8E=(\ZZ/2)^4,\\
 H_9E=\ZZ/8\oplus(\ZZ/2)^3.
\end{gathered}
\tag{D.2.22a17}
\]
The lower groups follow from the displayed first Bocksteins and
the known second Bockstein into $u_F$.
For the last group, the second Bockstein page in degree nine
is the single span of $N$, mapping injectively to that for $F$.
An incoming second differential would contradict its survival
in $F$. The classes $B_5,B_6$ cannot be earlier boundaries,
since their images have nonzero third differential in $F$.
Naturality therefore gives both
$\beta_3B_5=\beta_3B_6=N$ in $E$.
The complete degree-ten first-Bockstein cycle basis is
\[
 Z_2,U_2,T_{10},P_r,B_2,B_5,B_6,B_8,B_1+B_3.
\]
Under the Moore-four quadratic projection these map respectively
to $q_0,0,q_0,q_0,0,h+h',h+h',0,0$.
Thus only $B_5,B_6$ have nonzero third Bockstein, exactly as
required for \textup{(D.2.22a17)} and the ten-cell relations below.

In this particular degree the integral inclusion $H_9E\to H_9F$
is also injective. An element in its kernel reduces to zero modulo
two, by \textup{(D.2.22a13)}, so it lies in $2\ZZ/8\{N\}$.
The image of $N$ has odd coefficient on the unique order-eight
summand in \textup{(D.2.22a8)}, and hence has exact order eight.
No nonzero element of that even subgroup can vanish. Since
$F\to T$ is split, the integral Hurewicz kernel computed using
$E$ is the kernel in the statement.

\medskip\noindent
\textbf{The actual homotopy group and a stable cell model.}
The cover used above has an actual double-loop realization.
Put $Z=\operatorname{fib}(Y\to K(\mathbb F _2,4))$.
Then $F=\Omega^2Z$, the space $Z$ is five-connected, and
$\pi_6Z=\ZZ/2$.  Double-looping its first Postnikov map
$Z\to K(\mathbb F _2,6)$ realizes the map
$F\to K(\mathbb F _2,4)$, whose fiber is $E$.
Thus $E$ is a four-connected double-loop space.
The unit $E\to\Omega\Sigma E$ has a homotopy retraction by
the loop-space counit; Freudenthal makes it nine-connected.
Since $\Sigma E$ is five-connected, this proves
\[
 \pi_9T\cong\pi_9F\cong\pi_9E
       \cong\pi_9\Sigma^\infty E.
 \tag{D.2.22t1}
\]
This argument does not assert that stabilization of $F$ preserves
its specified double-loop operations.

Work two-locally.  Choose a finite stable approximation to
$\Sigma^\infty E$ through dimension eleven before cancelling
odd-unit cell pairs, and retain its ten-skeleton.
Write $L_i$ for its stable skeleta.  This convention removes
extraneous ten-cells killed by unit eleven-cell incidences and
preserves the required $\pi_9$.
Let $M=M(\ZZ/4,5)$, with bottom inclusion $i_5$ and top
six-class $D$, and put $r=D+E_6$ and $t=i_5\eta_5$.
The lower integral boundaries and operations in the preceding
tables give
\[
 \begin{split}
 D_0&=(M\vee S^6_r)\mathop{\cup}_{t+2r}e^7_S,
       \qquad L_7=D_0\vee S^7_R,\\
 D_0&=\cofib\left(
 \begin{pmatrix}4&\eta\\0&2\end{pmatrix}:
 S^5\vee S^6\longrightarrow S^5\vee S^6\right).
 \end{split}
 \tag{D.2.22t2}
\]
In particular $\beta S=r$ and $\operatorname{Sq}_*^2S=u_5$.
All cell assertions here concern stable cofibers; no minimal
unstable CW decomposition is being assumed.

Let $b\in\pi_7M$ be the order-two lift of $\eta$ on the
Moore top, and let $c=r\eta_6$.
The cofiber exact sequence gives
\[
 \pi_7L_7=\ZZ_{(2)}\{R\}\oplus
                    \ZZ/2\{b,c\}.
\]
The complete eight-cell attachment list, ordered as
$W,D_8,J_1,J_2,J_3$, is
\[
 0,\qquad b+c,\qquad b,\qquad0,\qquad2R+b.
 \tag{D.2.22t3}
\]
The $b,c$ coordinates are detected independently by the
second-square vectors $D,r$; the integral coordinate is fixed
by $\beta J_3=R$.
Subtract the $J_1$ cell from each of $D_8,J_3$.
These stable cofiber row changes give an actual splitting
\[
 \begin{split}
 L_8&=K_8\vee M(\ZZ/2,7)_R\vee S^8_W\vee S^8_{J_2},\\
 K_8&=D_0\cup_b e^8_{B_8}\cup_c e^8_{C_8},
       \qquad B_8=J_1,\quad C_8=D_8-J_1.
 \end{split}
 \tag{D.2.22t4}
\]
The subscripts $B_8,C_8$ here distinguish these cells from
the original space $B$ and from the cover's ten-classes $B_i$.

\medskip\noindent
\textbf{All ninth homotopy components of the eight-skeleton.}
First, $\pi_9M=\ZZ/4\{m\}$, with top quotient
$q_*m=2\nu_6$.  Write $z=r\nu_6$.
The full matrix cofiber sequence identifies the quotient of
$\pi_9D_0$ by $\ZZ/2\{z\}$ with the kernel of
\[
 \ZZ/8\{\nu_5\}\oplus\ZZ/2\{\eta_6^2\}
       \longrightarrow\ZZ/8,
       \qquad (x,y)\longmapsto4x+4y.
 \tag{D.2.22t5}
\]
This kernel is cyclic of order eight generated by
$(\nu_5,\eta_6^2)$; choose a lift $s_0$.
The possible discrepancy in $2s_0=m$ is $z$.
It is zero: the fourth-square detector of both $2s_0$ and
$m$ is zero, whereas the detector of $z$ is $r$.
The same argument excludes $8s_0=z$.
Changing $s_0$ by $z$ if necessary normalizes its
fourth-square vector to $D$.  Consequently
\[
 \pi_9D_0=\ZZ/8\{s_0\}\oplus\ZZ/2\{z\},
       \qquad2s_0=m.
 \tag{D.2.22t6}
\]
Here a square detector means the homology square on the top
class of the stable mapping cone.  In this range it is additive
and vanishes on twice an attaching map: the map from the cone
of $2f$ to the cone of $f$ is zero on the top mod-two class
and is the identity on the relevant lower homology.
The nonzero detector for $z$ is the Hopf-invariant-one
detector of $\nu$.

The two eight-attachments have composites
\[
 b\eta_7^2=2m=4s_0,
       \qquad c\eta_7^2=4z=0.
\]
There is no additional $\eta$-lift in their cofiber sequence:
the classes $b\eta_7,c\eta_7$ are independent before those
eight-cells are attached.  Therefore
\[
 \begin{split}
 \pi_9K_8&=\ZZ/4\{s\}\oplus\ZZ/2\{z\},\\
 \pi_9L_8&=\ZZ/4\{s\}\oplus\ZZ/2\{z\}
     \oplus\ZZ/4\{\tau\}
     \oplus\ZZ/2\{w\eta_8,j\eta_8\}.
 \end{split}
 \tag{D.2.22t7}
\]
The class $\tau$ is the Moore $\eta$ coextension in
$M(\ZZ/2,7)_R$, with $2\tau=R\eta_7^2$;
$w,j$ are the two free eight-sphere inclusions.
The stable stem and Moore identities used in these exact
sequences are the opened formulas
$\pi_1^S=\ZZ/2\{\eta\}$,
$\pi_2^S=\ZZ/2\{\eta^2\}$,
$\pi_3^S{}_{(2)}=\ZZ/8\{\nu\}$,
$\eta^3=4\nu$ and $2\widetilde\eta=i\eta^2$ in
\cite[\S2, (2.1)--(2.2), p.~149]{I_Oka1986}.

\medskip\noindent
\textbf{The exact $\eta$ composites of the free lifts.}
Set $J=M\cup_b e^8_{B_8}$.
Its cofiber sequence gives $\pi_9J=\ZZ/2\{\overline m\}$,
and $J/S^5$ is $C\eta$ with bottom cell in dimension six.
Let $e\in\pi_8J$ have top-cell degree two on $B_8$.
Its quotient is the top-degree-two class of $C\eta$.
The actual identity
$\eta_{C\eta}=i\nu q$, opened in
\cite[(8.4), p.~1279]{EH5AngeltveitRognes2005}, gives
\[
 e\eta_8=\overline m\quad\hbox{in }\pi_9J,
       \qquad e\eta_8=2s\quad\hbox{in }\pi_9L_8.
 \tag{D.2.22t8}
\]
Indeed its quotient is $2i\nu$, the nonzero image of
$\overline m$.  The free lift of $2C_8$ lies in the actual
$C\eta$ subcomplex on $r,C_8$ and has $\eta$ composite
$2r\nu=2z=0$.  The other two free composites are
$w\eta_8,j\eta_8$.

The degree-eight torsion of $L_8$ is
$\ZZ/4\{n=i_5\nu_5\}\oplus\ZZ/2\{\ell=R\eta_7\}$;
its free Hurewicz lattice has basis
$W,J_2,2J_1,2(D_8-J_1)$.
The nine-cells supply all four integral boundary directions.
The zero-boundary combinations $N$ and
$K_\Sigma=K_2+K_3+K_4$ kill $n,\ell$.
On taking $\eta$ composites, $n\eta=0$,
$\ell\eta=2\tau$, and \textup{(D.2.22t8)} gives $2s$.
Torsion additions to the integral attachments contribute only
$2\tau$, already killed.  Thus the image of
\textup{(D.2.22t7)} in $\pi_9L_9$ is exactly
\[
 V=(\ZZ/2)^5\{s,z,\tau,w\eta_8,j\eta_8\}.
 \tag{D.2.22t9}
\]
The quotient $\pi_9L_9/V$ is free of rank four, the kernel
of the full integral nine-cell attaching matrix.  In particular
$V$ is the entire torsion subgroup.  Its complete detectors are
\[
\begin{array}{c|c|c}
 &\operatorname{Sq}_*^2&\operatorname{Sq}_*^4\\ \hline
 s&0&D\\
 z&0&r\\
 \tau&J_1+J_3&0\\
 w\eta_8&W&0\\
 j\eta_8&J_2&0.
\end{array}
 \tag{D.2.22t10}
\]
The five vectors are independent.  Thus specifying these
coordinates leaves no unmeasured torsion component.

\medskip\noindent
\textbf{The free lattice, retaining the even-$\nu$ coefficients.}
Replace $N$ by $N'=N-K_\Sigma$, rescale $n$ by an odd
unit, and change $\ell$ by an even multiple of $n$.
We may then take the attachments of $N'$ and
$Q=K_\Sigma$ to be $n$ and $\ell$, respectively.
The order-eight homology generator changes only by an
order-two class.  The other two zero-boundary nine-cycles
have attachments
\[
 K_1:\ 2a n,
       \qquad K_5:\ 2b n,
       \qquad a,b\in\{0,1\}.
 \tag{D.2.22t11}
\]
These two bits are retained.  A mod-two spherical representative
alone would not set either bit to zero: its integral Hurewicz
class could differ from $K_i$ by $2N'$.

Choose integral order-two representatives $K_1,K_5,Q$ by
the half-boundaries of $B_3,B_7,B_4$, respectively, and use
the order-eight generator $N'$ supplied by the integral
configuration-chain calculation above.
The free Hurewicz lattice of $\pi_9L_9$ then has a basis
$p,q,x,y$ satisfying
\[
 \begin{aligned}
 h(p)&=4N',& h(q)&=2Q,\\
 h(x)&=K_1-2aN',&h(y)&=K_5-2bN'.
 \end{aligned}
 \tag{D.2.22t12}
\]
Normalize $p,q$ so that their second- and fourth-square
detectors are zero.  To justify this normalization, the Adem
relation $\operatorname{Sq}^2\operatorname{Sq}^2
=\operatorname{Sq}^3\operatorname{Sq}^1$ forces either
second-square vector into
$\langle W,J_2,J_1+J_3\rangle$.
For $p$ the top third square is zero; for $q$ it is $R$,
whose Bockstein is zero.  Either fourth-square vector lies
in $\langle D,r\rangle$.
The five entries of \textup{(D.2.22t10)} remove precisely
these five coordinates.  Even multiples of all attaching maps,
including $2x,2y$, have zero primary detectors by the
degree-two mapping-cone argument above.

\medskip\noindent
\textbf{All twelve ten-cell relations.}
The exact integral normalization is controlled by
$H_9E=\ZZ/8\{N'\}\oplus\ZZ/2\{K_1,K_5,Q\}$.
A first-Bockstein cycle with zero third Bockstein can be
made an exact zero-boundary ten-cell: its remaining $N'$
coefficient is divisible by sixteen and is removed by twice
an eight-$N'$ cell, while its order-two coefficients are
removed by even multiples of the corresponding degree-two
cells.  These changes preserve all primary detectors.
Since \emph{both} $B_5$ and $B_6$ have third Bockstein $N$,
they may each be normalized to boundary $8N'$.
Odd units do not change the ensuing two-primary relations.

The zero-boundary cells $Z_2,B_2,T_{10}$ give
\[
 w\eta_8=0,
       \qquad j\eta_8=0,
       \qquad w\eta_8+s+z=0.
 \tag{D.2.22t13}
\]
The two eight-$N'$ cells, using the normalized $p$, give
\[
 B_5:\ 2p+\tau+s=0,
       \qquad B_6:\ 2p+j\eta_8=0.
 \tag{D.2.22t14}
\]
Consequently $2p=0$ and $\tau=z=s$.
The zero-boundary cycle $B_1+B_3$ repeats $j\eta_8$.
The classes $U_2,P_r,B_8$ have zero detectors and impose
no new relation, by completeness of \textup{(D.2.22t10)}.
The third-Bockstein-zero combination $B_5+B_6$ repeats
$\tau+s+j\eta_8$.  Exactly one nonzero element of $V$
survives, represented by $s$.

The remaining integral cells have boundaries
$2Q,2K_1,2K_5$, which in \textup{(D.2.22t12)} are
$q,2x+ap,2y+bp$.  Their complete relations are
\[
 \begin{array}{ll}
 B_4:&q+\tau+s+j\eta_8=0,\\
 B_3:&2x+ap=0,\\
 B_7:&2y+bp=0.
 \end{array}
 \tag{D.2.22t15}
\]
Here $q$ was normalized above, and $2x,2y$ have zero
primary detectors; $B_1$ differs from $B_3$ by the already
killed $j\eta_8$.  All twelve ten-cells and all independent
integral boundary directions have now been used.
The resulting full presentation is
\[
 \pi_9E=
 \langle p,x,y,s\mid
       2p=2s=0,\quad2x=ap,\quad2y=bp\rangle.
 \tag{D.2.22t16}
\]
In particular, independently of the two undetermined bits,
\[
 4\pi_9T=0,
 \qquad
 \ker\bigl(h_{\ZZ}:\pi_9E\to H_9E\bigr)
       =\ZZ/2\{s\}\ne0.
 \tag{D.2.22t17}
\]
Its Hurewicz image is given exactly by
\textup{(D.2.22t12)} and has order eight.
The group is $(\ZZ/2)^4$ if $a=b=0$, and
$\ZZ/4\oplus(\ZZ/2)^2$ otherwise.
These alternatives are retained rather than selected by
an abstract group-order argument.
The previously proved injectivity of $H_9E\to H_9F$
and the splitting $F\to T$ identify the same kernel for
the integral Hurewicz map on $\pi_9T$.

Finally the $r$ sphere maps to an odd multiple of the
class $\rho\in\pi_6T=\ZZ/4$ of D.2.18.
Since $z=r\nu_6$ and $z=s$ in the final group, the
surviving element has the specified actual representative
\[
 s=\rho\nu_6\ne0.
 \tag{D.2.22t18}
\]
The same eight- and nine-cell exact sequences kill $n,\ell$
and the free directions $2J_1,2(D_8-J_1)$, while leaving
the $W,J_2$ directions modulo two.
Thus the actual group $\pi_8E=\pi_8T\cong(\ZZ/2)^2$ is
represented by those two eight-spheres.  Their $\eta$
composites are precisely the two classes killed in
\textup{(D.2.22t13)}.  Hence
\[
 \pi_8T\circ\eta_8=0.
 \tag{D.2.22t19}
\]
This excludes that particular proposed change of a relative
eight-cell filling.  It does not exhaust every possible change
of the preceding map and does not evaluate
$\lambda(\chi,\chi)$.


\medskip\noindent
\textbf{The common attachment coefficient and the remaining scope.}
The component calculation retains bits $a,b$ for the even-nu
parts of the $K_1,K_5$ attachments. They are equal, although
their common value is not determined here.
Indeed the actual homotopy Browder bracket of the bottom
$a\in\pi_2T$, of order two, with $\rho\in\pi_6T$ has order
dividing two. Its mod-two Hurewicz image is
\[
 \{a,b^2+k+p_3\}=K_1+K_5.
\]
The derivation law gives $\{a,b^2\}=0$ and
$\{a,p_3\}=\{a,f\}+aR=K_1$.
In the component presentation an integral Hurewicz class reducing
to $K_1+K_5$ has order dividing two only when $a=b$;
otherwise its coefficient on the order-eight class is $2$ modulo
$4$. Thus only the common bit remains. Neither choice changes
\textup{(D.2.22a1)--(D.2.22a2)}.
No identification of the actual square in \textup{D.2.21} with
the nonzero Hurewicz-kernel generator has been made. Its integral
Hurewicz image can nevertheless be determined exactly.

In integral homology the fiber generator is
\[
 h_{\ZZ}(\chi)=a^2-2c\in H_4T,
 \tag{D.2.22a19}
\]
not the uncorrected Pontryagin square $a^2$.
To verify the normalization, use the homomorphism
$K(\ZZ,2)=\mathbb CP^\infty\to K(\ZZ/2,2)=B$.
Its degree-four image is an odd generator $c$ because the
mod-two square of the fundamental class pulls back nontrivially.
The degree-two generator of $H_*\mathbb CP^\infty$ has
Pontryagin square twice the degree-four generator, as follows
from its primitive integral cohomology generator.
Consequently $q_*(a^2)=2c$. The kernel in degree four is
the order-two class in \textup{(D.2.22a19)}, which reduces to
$A$ and is the Hurewicz image of the bottom sphere of $F$.

The integral Poisson derivation law now gives
\[
 \{a^2-2c,a^2-2c\}=4\{c,c\}=8\xi c=0.
 \tag{D.2.22a20}
\]
Every omitted cross term has a factor $2a=0$; the even-label
self-bracket is $\{c,c\}=2\xi c$, and the exact order-eight
calculation above proves the last equality.
Naturality of Hurewicz for the homotopy Browder operation
therefore places the actual square in the kernel already
computed, proving \textup{(D.2.22a18)}. Since $4\pi_9T=0$,
four-divisibility in \textup{D.2.21} is exactly vanishing.
This does not distinguish the two values in
\textup{(D.2.22a18)}. The later geometric arguments settle
both retained questions: \textup{D.2.23} proves that the actual
square is zero, and \textup{D.2.31} proves that the common
attachment bit is zero, so $\pi_9T\cong(\ZZ/2)^4$.
The two-case calculation above is retained to identify exactly
which information those later constructions supply.
\end{proof}

\begin{claim}{D.2.23}{Proved}{The actual canonical Browder square vanishes}
For the class $\chi\in\pi_4F$ of \textup{D.2.21},
\[
 \lambda_1(\chi,\chi)=0\in\pi_9F,
 \qquad [\widehat\chi,\widehat\chi]=0\in\pi_{11}Y.
 \tag{D.2.23a}
\]
Consequently the fixed canonical summand of \textup{D.2.21}
extends through relative dimension ten. This assertion does not
extend the separate correction or the combined map $G_9$.
\end{claim}
\begin{proof}
The integral Hurewicz calculation and complete kernel computation
of \textup{D.2.22} put the actual square in
$\{0,\rho\nu_6\}$, with $\rho\nu_6\ne0$.
We distinguish these two actual maps by a cohomology relation
and their ordinary mapping cones.

Use the homology bases of \textup{D.2.22}, and write
$e_6^*,J_1^*$ for their dual cohomology vectors. The entire
eighteen-row table \textup{(D.2.22a6)} gives
\[
 \Psi=\operatorname{Sq}^4e_6^*
              +\operatorname{Sq}^2J_1^*=0
                 \quad\text{in }H^{10}(F;\FF_2).
 \tag{D.2.23b}
\]
Indeed each summand evaluates to one precisely on $B_4,B_5$
and to zero on the other sixteen basis vectors. In particular
the fourth squares of $D_2,Q_{10},P_p$ are respectively
$d_6,f_6,f_6$, with zero $e_6$ coefficient.

For any $f:S^9\to F$ with zero mod-two Hurewicz image,
the two lower cohomology classes extend uniquely to
$C_f=F\cup_f e^{10}$. The expression in
\textup{(D.2.23b)} is zero on the old $H_{10}F$ and hence
has a well-defined value on a new top homology class.
This value is independent of adding an old homology class.
For $f=\rho\nu_6$ the ordinary map of cones from
$S^6\cup_{\nu_6}e^{10}$ gives
\[
 \operatorname{Sq}_*^4(v)=h(\rho)=e_6+f_6,
 \qquad \operatorname{Sq}_*^2(v)=0.
 \tag{D.2.23c}
\]
The first is the Hopf-invariant-one detector of $\nu$,
and the second follows from the zero degree-eight homology
of this two-cell source. Thus the value of $\Psi$ is one
on this cone. Its value on the null map is zero.

Now form the cone of the actual Browder square. Since $F$ is
an actual double-loop space, $\chi$ has a free $E_2$ extension
$C_2S^4\to F$. In $C_2S^4\to C_3S^4$ the only relative
cell through dimension ten is the binary ten-cell. Its
attaching map is the Browder square, with the normalization
proved in \textup{D.2.21}. Push out this relative ten-stage
along the free extension. The result is the ordinary
homotopy cofiber
\[
 C_{\lambda_1(\chi,\chi)}=F\cup_{\lambda_1(\chi,\chi)}e^{10}.
 \tag{D.2.23d}
\]
No $E_3$ structure on $F$ is assumed in this pushout.

Let $x\in H_4S^4$ be the bottom class. In the source the
absolute class $v=Q^6x$ has coefficient one on the relative
ten-cell. This follows either from its binary configuration
representative or the free three-fold loop homology basis
\cite[Chapter III, Theorems 3.1--3.2]{I_CohenLadaMay1976}.
There is no degree-ten class in $H_*C_2S^4$, and the mod-two
Browder attaching boundary is zero, so this top coefficient
is unambiguous. The source comparison $C_3S^4\to QS^4$ is
ten-connected. It is therefore injective on $H_8$, and the
ordinary stable Nishida formula gives, already in $C_3S^4$,
\[
 \operatorname{Sq}_*^2(Q^6x)
 =\binom42 Q^4x+Q^5\operatorname{Sq}_*^1x=0,
 \qquad \operatorname{Sq}_*^4(Q^6x)=0.
 \tag{D.2.23e}
\]
Here $\binom42=0$ modulo two and every positive square of
$x$ is zero. The second equality also follows from $H_6C_3S^4=0$.
These operations are unchanged in the relative ten-stage,
whose comparison with the full source is an isomorphism
in degrees six and eight. The Nishida input is
\cite[Chapter III, Theorem 1.3]{I_CohenLadaMay1976}; the
comparison with the infinite-loop source avoids any
unstated top-operation correction.

The pushout map takes $v$ to a class with coefficient one
on the new cone cell and with both displayed squares zero,
by ordinary Steenrod naturality. Therefore the value of
$\Psi$ on this cone is zero. Since the only two possibilities
were $0$ and $\rho\nu_6$, and the latter has value one by
\textup{(D.2.23c)}, the actual square is zero. The extension
conclusion follows from \textup{D.2.21}.
\end{proof}

\begin{claim}{D.2.24}{Proved}{A primary operation excludes the remaining dangerous unit pattern}
Put $B=K(\FF_2,2)$, $X=\Sigma B$, $A=QX$, $Y=\Sigma^2B$,
and $T=\Omega^2Y$. Retain the labels and classes of
\textup{D.2.17--D.2.22}, in particular
\[
 k=\{a,b\},\qquad p_3=f+ac+a^3,\qquad
 d_6=b^2,\quad e_6=b^2+k,\quad f_6=p_3.
 \tag{D.2.24a}
\]
Let $d\in H_5(B;\FF_2)$ be dual to
$w=\operatorname{Sq}^2\operatorname{Sq}^1u$ and zero on
$u\operatorname{Sq}^1u$, and write $D=\sigma d\in H_6X$.
For every ordinary pointed map $g:X\to T$, the coefficient
of $e_6$ in $g_*D$ is zero. Equivalently, its coefficient
on $k$ in the primitive basis $b^2,k,p_3$ is zero.

Consequently the pattern $G_*D=b^2+k$ constructed through
dimension nine in \textup{D.2.17--D.2.20} cannot extend over
the unit ten-skeleton. This conclusion applies to every map
with that pattern, independently of its earlier cell fillings.

For an arbitrary ordinary map $G:A\to T$, there is
$s\in\FF_2$ such that, on the unit classes and the indicated
extra product,
\[
 G_*D=s p_3,\qquad G_*E=0,\qquad G_*(A_0^2)=0,
 \quad E=\sigma e,\quad A_0=\sigma a.
 \tag{D.2.24b}
\]
Its once-suspended adjunct $\Sigma A\to\Omega Y$ therefore
has zero coefficient on the old commutator
$\bar a\bar b+\bar b\bar a$, where $\bar a=\sigma a$ and
$\bar b=\sigma b$ are the loop-homology letters, in its
entire degree-seven homology image. This is a statement about
ordinary maps and the first attaching-map test; it does not
assert persistence through later nullification stages.
\end{claim}
\begin{proof}
All homology and cohomology in this proof has coefficients
$\FF_2$. Let $F=\fib(T\to B)$ be the three-connected fiber
used above. Its homology injects into that of $T$.
Write $e_6^*$ and $J_1^*$ for the dual basis vectors to
$e_6$ and $J_1=\{a,d\}$ in $H_6F$ and $H_8F$.
The complete eighteen-dimensional table
\textup{(D.2.22a6)} gives the exact cohomology identity
\[
 \Psi_F:=\operatorname{Sq}^4(e_6^*)+
                   \operatorname{Sq}^2(J_1^*)=0
                   \quad\hbox{in }H^{10}F.
 \tag{D.2.24c}
\]
Indeed the only degree-ten basis vectors whose fourth square
has $e_6$ coefficient one are $B_4,B_5$. These are exactly
the two basis vectors whose second square has $J_1$
coefficient one. Every other basis vector contributes zero
to both functionals. Thus this is an identity in the full
cohomology group, with no untested complementary summand.
The table follows from the opened CLM formulas
\cite[Chapter III, Theorems 1.2--1.4]{I_CohenLadaMay1976};
no homotopy-group presentation or secondary attachment
coefficient is needed for \textup{(D.2.24c)}.

Write $v=\operatorname{Sq}^1u$, $w=\operatorname{Sq}^2v$,
and $t=\operatorname{Sq}^4w$ in the polynomial presentation
of $H^*B$. Serre's generator theorem and Cartan give
\[
\begin{aligned}
 \operatorname{Sq}^4w&=t,&
 \operatorname{Sq}^4(uv)&=u^2w+v^3,\\
 \operatorname{Sq}^2(u^2v)&=u^2w+v^3,&
 \operatorname{Sq}^2(uw)&=u^2w+v^3.
\end{aligned}
 \tag{D.2.24d}
\]
The degree-five cohomology basis is $w,uv$, and the
degree-seven basis is $u^2v,uw$. The complete degree-nine
basis is $u^3v,u^2w,v^3,t$; see
\cite[\S2, nos.8--9]{I_Serre1953}.
Let $t_9\in H_9B$ be dual to $t$ and zero on the other
three basis monomials. By suspension compatibility,
the class $T_{10}=\sigma t_9\in H_{10}X$ satisfies
\[
 \operatorname{Sq}_*^4T_{10}=D,
 \qquad \operatorname{Sq}_*^2T_{10}=0.
 \tag{D.2.24e}
\]

Every pointed $g:X\to T$ lifts to $F$: its composite
with $T\to K(\FF_2,2)$ is null because $H^2X=0$.
Use the same letter for this lift. Naturality and the
ordinary homology--cohomology pairing now give
\[
\begin{split}
 0
 &=\langle\Psi_F,g_*T_{10}\rangle\\
 &=\langle e_6^*,g_*\operatorname{Sq}_*^4T_{10}\rangle
   +\langle J_1^*,g_*\operatorname{Sq}_*^2T_{10}\rangle\\
 &=\langle e_6^*,g_*D\rangle.
\end{split}
 \tag{D.2.24f}
\]
This proves the universal assertion. It uses neither an
$H$-map hypothesis nor naturality of Dyer--Lashof operations
under an arbitrary ordinary map. The same argument already
applies on a unit ten-skeleton: choose a lift of $T_{10}$
to its homology; inclusion is an isomorphism in the lower
degrees occurring in \textup{(D.2.24e)}. Since
$b^2+k=e_6$, the previously constructed dangerous pattern
cannot extend over this skeleton.

For completeness, we recall the low-degree constraints needed
for the last assertion. For $G:A\to T$, \textup{D.2.16}
and the primitive basis in degree five give $G_*(\sigma c)=0$.
Also $G_*A_0=0$ and $G_*(\sigma b)=s a^2$ for some $s$.
The complete primitive bases in degrees six and seven are
\[
 \langle b^2,k,p_3\rangle,
 \qquad
 \langle\xi b,\{a,c\},P,Q\rangle,
 \tag{D.2.24g}
\]
where $P,Q$ are as in \textup{D.2.22}. Naturality of
$\beta$, $\operatorname{Sq}_*^2$, and
$\operatorname{Sq}_*^3$ on the unit classes gives,
on writing $G_*D=d_1b^2+d_2k+d_3p_3$ and
$G_*E=e_1b^2+e_2k+e_3p_3$,
\[
 d_1=e_1,\qquad d_2=e_2=\epsilon,\qquad
 d_1+d_3=s,\qquad e_1=e_3.
 \tag{D.2.24h}
\]
The exact second-Bockstein calculation of
\textup{(D.2.17e)} has
$\beta_2[b^2]=\beta_2[k]=[\xi a]$ and $[p_3]=0$.
On the source $\beta_2[D]=[\sigma c]$.
Since $G_*(\sigma c)=0$, naturality gives
$d_1=d_2=\epsilon$. Hence
\[
 G_*D=\epsilon(b^2+k)+(s+\epsilon)p_3,
 \qquad
 G_*E=\epsilon(b^2+k+p_3).
 \tag{D.2.24i}
\]
Apply \textup{(D.2.24f)} to the unit restriction of $G$.
It forces $\epsilon=0$.

There is no additional image on $A_0^2$. The source class
$R_0=Q^4A_0$ satisfies $\beta R_0=A_0^2$ and
$\operatorname{Sq}_*^2R_0=0$. The second identity excludes
a $\xi b$ coefficient in $G_*R_0$, so its Bockstein is
a multiple of $p_3$. Applying the second square, and using
$\operatorname{Sq}_*^2p_3=a^2$, excludes that multiple.
Thus $G_*(A_0^2)=0$, completing \textup{(D.2.24b)}.

Finally $H_6A$ has basis $D,E,A_0^2$. Homology suspension
from $T$ to $\Omega Y$ sends $k$ to the old tensor
commutator, sends $b^2$ to zero, and sends $p_3$ to
the linear $f$ label. Formula \textup{(D.2.24b)} therefore
has no commutator coefficient on any degree-seven input
of the once-suspended adjunct. This proves the stated
scope without making an assertion about subsequent cones.
\end{proof}

\begin{claim}{D.2.25}{Proved}{Two further first-cone commutator functionals vanish}
Let $B=K(\FF_2,2)$, $A=Q\Sigma B$, $Y=\Sigma^2B$,
and $T=\Omega^2Y$.  For a pointed map
$h:\Sigma^kA\to Y$, with $k\ge2$, let
\[
 G:\Sigma^{k-2}A\longrightarrow T
\]
be its twice-adjoint.  On $H_6T$ use the labels
$k_6=\{a,b\}$ and $p_3=f+ac+a^3$ of D.2.17.
Define the commutator functional of $h$ to be
\[
 H_8(\Sigma^kA;\FF_2)
   \cong H_6(\Sigma^{k-2}A;\FF_2)
   \xrightarrow{G_*}H_6T
   \xrightarrow{\text{coefficient of }k_6}\FF_2.
 \tag{D.2.25a}
\]
This functional is identically zero for $k=3$, for $k=5$,
and for every $k\ge6$.
Under homology suspension, $k_6$ maps to the old tensor
commutator in $H_7\Omega Y$, so the conclusion concerns the
entire relevant input homology, rather than only the kernel
of an individual attaching homology map.

Together with D.2.24, this treats the exponents $2,3$ and
all exponents at least $5$.  No assertion for exponents
$0,1,4$, for arbitrary simultaneous attachments involving
those exponents, or for later nullification stages is made.
\end{claim}
\begin{proof}
All homology and cohomology below have coefficients $\FF_2$.
Write $F=\operatorname{fib}(T\to B)$, and use its classes
\[
 d_6=b^2,\qquad e_6=b^2+k_6,\qquad f_6=p_3.
\]
The coefficient of $k_6$ is exactly the coefficient of $e_6$
in this basis.  Since every source $\Sigma^{k-2}A$ here has
zero $H^2$, its composite with $T\to K(\FF_2,2)$ is null;
choose a lift of $G$ to $F$.
The splitting $T\simeq B\times F$ makes the fiber
coordinates unambiguous.  We use the exact identity proved
from the complete degree-ten table in D.2.22--D.2.24:
\[
 \operatorname{Sq}^4(e_6^*)+
           \operatorname{Sq}^2(J_1^*)=0,
       \qquad J_1=\{a,d\}\in H_8F.
 \tag{D.2.25b}
\]

For $k=3$, the source is $\Sigma A$.
The only class of $H_6\Sigma A$ is the suspended unit
$C_6=\sigma^2c$, where $c\in H_4B$ is dual to $u^2$.
Indeed the unit $\Sigma B\to Q\Sigma B$ is an isomorphism
on degree-five homology; the first additional stable weight
starts in degree six.
Restrict $G$ to the actual unit $\Sigma^2B$.
Let $f\in H_6B$ and $i\in H_8B$ be dual to $u^3$ and
$u^4$, respectively, and set
$F_8=\sigma^2f$, $I_{10}=\sigma^2i$.
Serre's polynomial description and Cartan give
\[
 \begin{aligned}
 \operatorname{Sq}_*^4 I_{10}&=C_6,
       &\operatorname{Sq}_*^2 I_{10}&=F_8,\\
 \beta F_8&=0,
       &\operatorname{Sq}_*^2 F_8&=0.
 \end{aligned}
 \tag{D.2.25c}
\]
For the first row, use
$\operatorname{Sq}^4(u^2)=u^4$,
$\operatorname{Sq}^2(u^3)=u^4+uv^2$, and
$\operatorname{Sq}^2(v^2)=0$.
For the second, use
$\operatorname{Sq}^1w=\operatorname{Sq}^1(uv)=v^2$
and $\operatorname{Sq}^2(u^2)=v^2$.
These are the same opened Serre generators and stable
suspension conventions as in
\cite[\S1, nos.~2.2--2.3; \S2, nos.~8--9]{I_Serre1953}.

The source $F_8$ is primitive.  In the complete basis
$W_8,C_8,L_8,M_8,J_1,J_2,J_3$ of $H_8F$, only $C_8$
is nonprimitive, with reduced diagonal $A_4\otimes A_4$.
The table in D.2.19 therefore shows that every primitive
Bockstein cycle has the form
\[
 wW_8+j_2J_2+t(L_8+M_8)+j_1J_1.
\]
Its second square is
\[
 t(d_6+f_6)+j_1(d_6+e_6).
\]
Independence of $d_6,e_6,f_6$ shows that this square
vanishes only if $t=j_1=0$.
By \textup{(D.2.25c)}, $G_*F_8$ has zero $J_1$
coefficient.  Evaluate \textup{(D.2.25b)} on $G_*I_{10}$:
\[
 0=\langle e_6^*,G_*C_6\rangle+
           \langle J_1^*,G_*F_8\rangle
   =\langle e_6^*,G_*C_6\rangle.
 \tag{D.2.25d}
\]
This proves the claimed identical vanishing for $k=3$.

For $k=5$, the source $\Sigma^3A$ is five-connected and
its bottom integral homology group is $\ZZ/2$ in degree six.
Hurewicz thus identifies its bottom homotopy group with
$\ZZ/2$.  By D.2.18 the actual target group is
$\pi_6T=\ZZ/4\{\rho\}$, with mod-two Hurewicz image
$b^2+k_6+p_3$.
The image of any homomorphism $\ZZ/2\to\ZZ/4$ lies in
$\{0,2\rho\}$, whose mod-two Hurewicz image is zero.
Thus $G_*$ is zero on all of $H_6\Sigma^3A$.
For $k\ge6$, the source $\Sigma^{k-2}A$ has no
degree-six homology.  This finishes each asserted case.
\end{proof}

\begin{claim}{D.2.26}{Proved}{The exact quadratic test for a simultaneous first cone}
All groups below have coefficients in $\FF_2$.
For every three-connected CW space $C$, Pontryagin multiplication
induces a natural injection
\[
 \frac{(H_4C\otimes H_5C)\oplus(H_5C\otimes H_4C)}
 {\im(\overline\Delta:H_9C\longrightarrow
                (H_4C\otimes H_5C)\oplus(H_5C\otimes H_4C))}
 \lhook\joinrel\longrightarrow H_7(\Omega C).
 \tag{D.2.26a}
\]
Here the factors are identified with $H_3(\Omega C)$ and
$H_4(\Omega C)$ by the canonical homology desuspensions.

Put $Y=\Sigma^2K(\FF_2,2)$, and let $u\in H^4Y$ be the
bottom class and $v=\operatorname{Sq}^1u\in H^5Y$.
Write $a,b$ for the degree-three and degree-four letters in
$H_*(\Omega Y)$, and $\kappa=ab+ba$.
Let $h=\bigvee_i h_i:S=\bigvee_i S_i\to Y$, where every $S_i$
is two-connected and $h_i$ is zero on $H_4,H_5$, and set
$C=Y\cup_h CS$. The old commutator survives in $H_7(\Omega C)$
exactly when its tensor $a_4\otimes b_5+b_5\otimes a_4$ is
outside the image of the diagonal in \textup{(D.2.26a)}.
If $\alpha\in H^4C$ lifts $u$, the condition
\[
 \alpha\operatorname{Sq}^1\alpha=0
 \tag{D.2.26b}
\]
is independent of the lift and implies this survival. It is
equivalent to survival when $H_4C,H_5C$ are precisely the old
one-dimensional groups.

There is an exact simultaneous test for \textup{(D.2.26b)}.
Fix singular cocycles $a^4,b^5$ representing $u,v$ and a cochain
$c^8$ on $Y$ with $dc^8=a^4\smile b^5$. For each $i$ choose
$p_i^3$ with $dp_i^3=h_i^*a^4$, and put
\[
 \theta_i=[h_i^*c^8-p_i^3\smile h_i^*b^5]\in H^8S_i.
 \tag{D.2.26c}
\]
Then \textup{(D.2.26b)} holds if and only if there is one
$\ell\in H^8Y$ such that $\theta_i=h_i^*\ell$ for every $i$.
Distinct choices of $\ell$ for the separate cones do not suffice.

If, moreover, $S_i=\Sigma R_i$ with $R_i$ two-connected and
$h_i$ also zero on $H_6$, define $L_i:H_8S_i\to\FF_2$ by the
ordinary adjunct $\bar h_i:R_i\to\Omega Y$ and the fixed
James decomposition:
\[
 (\bar h_i)_*z=(h_i)_*(\sigma z)+L_i(\sigma z)\kappa,
 \quad
 H_7(\Omega Y)=H_8Y\oplus\FF_2\{ab,ba\}.
 \tag{D.2.26d}
\]
The first term lies in the word-length-one summand. In this
suspended-source case the exact necessary and sufficient condition
for survival is
\[
 L=\bigoplus_iL_i\text{ vanishes on }
 \ker\!\left(\bigoplus_iH_8S_i\xrightarrow{h_*}H_8Y\right),
 \quad\text{equivalently }L_i=\ell\circ(h_i)_*\text{ for one }\ell.
 \tag{D.2.26e}
\]
The assertions apply to arbitrary small wedges and CW sources.
They are criteria for the first cone, and assert neither that
all the attaching maps satisfy them nor persistence at later stages.
\end{claim}
\begin{proof}
For a finite-type three-connected $C$, take a minimal
Adams--Hilton model $(T(V),d)$ over $\FF_2$.
Here $V=s^{-1}\widetilde H_*C$ starts in degree three, $d$ has
no linear part, and its quadratic part is the desuspended reduced
diagonal; these are the exact assertions of
\cite[\S3.1, p.300]{FGHHalperinLemaire1995}.
The degree-seven length-two words are $V_3V_4+V_4V_3$ and are
cycles. Words of length at least three start in degree nine.
A degree-eight length-two word has letters of degree at most five,
whose differentials vanish. Hence all boundaries in this
degree-seven subspace are exactly $d(V_8)$, which is quadratic
and equals the diagonal of $H_9C$. This proves
\textup{(D.2.26a)}. Its definition by the actual Pontryagin product
proves naturality and rules out an unresolved filtration extension.

For general $C$, use a CW model with cells only in dimensions zero
and at least four and take its filtered finite subcomplexes.
All are three-connected. Homology, field tensor products, the
diagonal and its cokernel commute with this union. So does loop
homology: a finite singular cycle or bounding chain in the loop
space has compact adjoint and lies in a finite subcomplex.
The natural injections therefore pass to the filtered colimit.

The cone $C$ is three-connected, and the old $H_4,H_5$ inject
by the homology cofiber sequence. Naturality gives the exact tensor
criterion. If \textup{(D.2.26b)} holds, the functional
$\alpha\otimes\operatorname{Sq}^1\alpha$ annihilates the
diagonal image and evaluates to one on the old tensor, proving
survival. If the two groups are one-dimensional, cocommutativity
makes that tensor span every possible diagonal image, proving
the converse under the stated extra hypothesis.

The ordinary coaction $C\to C\vee\Sigma S$ proves that
connecting-image classes annihilate positive cohomology of $C$:
projection to $C$ is homotopic to the identity, and pullback from
$\Sigma S$ is the connecting homomorphism, while products of
positive classes on the separate wedge summands are zero.
Any two lifts differ by a connecting-image class. Thus products
of lifts of fixed positive classes are independent of the lifts.

Let $\delta:H^8S\to H^9C$ be the connecting homomorphism.
We prove using genuine singular cup products that, for lifts
$\alpha,\beta$ of $u,v$,
\[
 \alpha\beta=\delta((\theta_i)_i).
 \tag{D.2.26f}
\]
First \textup{(D.2.26c)} is a cocycle by Leibniz. Changing $p_i$
by a closed cochain changes it by a product with $h_i^*b^5$,
which is exact, so its class is unchanged.

Use a mapping-cylinder cone cover $U,V$, with $U\simeq Y$,
$V$ contractible and $U\cap V$ a collar of $S$.
The small singular chains $C_*(U)+C_*(V)$ include
quasi-isomorphically in $C_*(C)$ by subdivision and retain the
actual Alexander--Whitney cup: faces of each small simplex stay
in the same cover member. Their dual has the exact sequence
\[
 0\to C^*_{\rm small}(C)\to C^*(U)\oplus C^*(V)
       \to C^*(U\cap V)\to0.
 \tag{D.2.26g}
\]
The last map is the difference of restrictions and is onto by
extension of cochains on singular simplices.
Choose global cocycles for $\alpha,\beta$ restricting exactly
to the pullbacks of $a^4,b^5$ on $U$; extend correcting cochains
from $U$ to achieve these representatives. On $V$ write
$\alpha=dP^3$. Their cup product has primitive $c^8$ on $U$
and primitive $P^3\smile\beta$ on $V$. The difference on the
collar, restricted to $S$, is precisely
$h^*c^8-p^3\smile h^*b^5$. The connecting construction in
\textup{(D.2.26g)} gives \textup{(D.2.26f)}.
No multiplication on an abstract mapping-cone DGA is assumed.

Take $\beta=\operatorname{Sq}^1\alpha$. Since
$H^8S=\prod_iH^8S_i$, the cofiber exact sequence now gives
\[
 \alpha\operatorname{Sq}^1\alpha=0
 \quad\Longleftrightarrow\quad
 (\theta_i)_i\in\im(H^8Y\xrightarrow{h^*}\prod_iH^8S_i).
 \tag{D.2.26h}
\]
This is exactly the common-functional assertion, by universal
coefficients over a field. It includes cancellations between
different summands of $H_8S$. For example, two input vectors
with the same nonzero target image but different functional
values are harmless to the separate kernel tests and violate
the simultaneous one. This is a linear-algebra observation,
not a claimed realization by attaching maps in this problem.

For the last assertion put $V=\widetilde H_*(\Sigma B)$,
$W=\widetilde H_*(\bigvee_iR_i)$. The James identification
gives $H_*(\Omega Y)=T(V)$ with primitive letters, since
$\Sigma B$ is itself a suspension. In degrees three, four and
five, homology suspension is an isomorphism; hence the low
vanishing hypotheses give $(\bar h)_*W_3=(\bar h)_*W_4
=(\bar h)_*W_5=0$. Coalgebra naturality in degree seven
then forces the $ab$ and $ba$ coefficients to agree, proving
that \textup{(D.2.26d)} is well-defined.

Both source and target are AH-suspensions. The opened relative
cone construction in the proof of
\cite[Proposition 7.4, p.309]{FGHHalperinLemaire1995} gives
\[
 AH(C)=(T(V\oplus sW),D),\qquad
 D|_V=0,\qquad D(sw)=(\bar h)_*w.
 \tag{D.2.26i}
\]
Old generators start in degree three, new ones in degree four.
The differentials of new generators in degrees four, five and
six vanish by the preceding low vanishing. Every decomposable
degree-eight chain therefore has zero differential: its factors
have these degrees or are old generators; a degree-seven new
generator would require a degree-one companion, and words of
length three start in degree nine. Thus all degree-seven
boundaries are exactly $D(sW_7)$. The old $\kappa$ is a boundary
if and only if a finite sum of the vectors in
\textup{(D.2.26d)} has zero linear component and nonzero
$\kappa$ component. This is precisely the negation of
\textup{(D.2.26e)}. In particular new $H_5C$ classes permitted
by a once-suspended source create no omitted boundary here.

In fact the two normalizations agree modulo one common pullback:
\[
 (\theta_i)_i-L\in\im(h^*:H^8Y\to H^8S).
 \tag{D.2.26j}
\]
For $z\in\ker h_*$ choose a lift $\widetilde z\in H_9C$.
After cancelling the linear differential in
\textup{(D.2.26i)}, its quadratic differential pairs with
$\alpha\otimes\beta$ as $L(z)$, by the reduced-diagonal
identification of \cite[\S3.1]{FGHHalperinLemaire1995}.
But \textup{(D.2.26f)} gives the same pairing as
$\langle\delta\theta,\widetilde z\rangle=\theta(z)$.
Thus $\theta-L$ vanishes on the whole kernel and factors
through $h_*$, proving \textup{(D.2.26j)}.

The cited paper permits characteristic two, with a finite-type
convention. For an arbitrary wedge, every input vector or
boundary witness uses finitely many source cells and summands.
Choose finite subpairs of the mapping-cylinder cone, enlarging
the finite target to contain chains witnessing the finitely many
required low homology vanishings. The finite-type relative
construction proves the same criterion for each finite witness,
and passage to their filtered union proves it as stated.
Finally a functional on the source sum vanishes on $\ker h_*$
exactly when it factors through $\im h_*$; extend it linearly to
$H_8Y$ to obtain the single $\ell$. No finite-dimensional duality
or countability assumption enters this step.
\end{proof}

\begin{claim}{D.2.27}{Proved}{The unsuspended first-cone functional vanishes with one common normalization}
Let $B=K(\FF_2,2)$, $A=Q\Sigma B$, and $Y=\Sigma^2B$.
Write $\alpha\in H^4(Y;\FF_2)$ for the suspended fundamental
class and $\beta=\operatorname{Sq}^1\alpha$.
There is one choice of the common primitive $c_8$ in
\textup{D.2.26}, with $dc_8=a_4\smile b_5$, for which
\[
 \theta_h=0\in H^8(A;\FF_2)
 \qquad\text{for every pointed map }h:A\longrightarrow Y.
 \tag{D.2.27a}
\]
Consequently, after any small simultaneous family of such
unsuspended attachments, a lift $\widehat\alpha$ to the common
cone satisfies
$\widehat\alpha\operatorname{Sq}^1\widehat\alpha=0$.
In particular the old degree-seven loop commutator survives.
This statement does not address exponents one or four, or
propagation through subsequent cones.

More generally, the same conclusion holds for one small
simultaneous family of maps $\Sigma^{k_i}A\to Y$ with
\[
 k_i\in\{0,2,3\}\cup\{5,6,7,\ldots\}.
 \tag{D.2.27a'}
\]
All these exponents use the same common primitive $c_8$;
the assertion permits mixtures of the indicated exponents.

The proof supplies an actual delooped relation. Put
$Z=\operatorname{fib}(Y\to K(\FF_2,4))$, using $\alpha$,
and $F=\Omega^2Z$. There are unique classes
$\vartheta_8\in H^8Z$, $\lambda_{10}\in H^{10}Z$ whose
twofold cohomology suspensions are $e_6^*,J_1^*$ of
\textup{D.2.24}. For the bottom $v_6\in H^6Z$,
\[
 \operatorname{Sq}^4\vartheta_8+
 \operatorname{Sq}^2\lambda_{10}
       =\epsilon v_6^2\quad(\epsilon\in\FF_2).
 \tag{D.2.27b}
\]
The common choice in \textup{(D.2.27a)} makes the universal
relative cup class on $Z$ exactly $\vartheta_8$.
\end{claim}
\begin{proof}
All cohomology and homology in this proof has coefficients
$\FF_2$. We first justify every delooping interface; no
cohomology-suspension isomorphism outside its indicated range
is assumed.

By the preceding low homotopy calculations,
$\pi_4Y=\ZZ/2$ and $\pi_5Y=0$. Thus $Z$ is five-connected,
$W=\Omega Z$ is four-connected, and $F=\Omega W$ is the
three-connected fiber already used in \textup{D.2.18--D.2.25}.
Moreover $H^6Z$ is one-dimensional because
$\pi_6Z=\pi_6Y=\ZZ/2$.

Consider the path fibrations for $Z$ and $W$.
Their Serre spectral sequences give the following low-degree
statements directly. Cohomology suspension identifies $H^8Z$
with $H^6F$. It also identifies $H^{10}Z$ with $H^9W$.
For $H^9W\to H^8F$, the only differential before transgression
out of the fiber group is
\[
 d_5:H^8F\longrightarrow H^5W\otimes H^4F.
 \tag{D.2.27c}
\]
The loop action on the path fibration identifies its dual
with multiplication of the transgressed bottom class by the
bottom fiber class. Explicitly the bottom tensor maps to
the Pontryagin product $A_4^2=W_8\in H_8F$.
This is the usual module property of the path-fibration
differential: the differential of a transgressive base
class tensored with a fiber cycle is its transgression
multiplied by that cycle.
Thus the kernel of \textup{(D.2.27c)} is exactly the
$W_8$-annihilator. The targets of $d_6,d_7,d_8$ vanish by
connectivity, and $d_9$ identifies this kernel with $H^9W$;
there is no earlier mixed source entering that base group.
Since $J_1^*$ vanishes on $W_8$, it has a unique lift to
$H^9W$ and then to $H^{10}Z$. This gives $\lambda_{10}$.
The lower isomorphism gives $\vartheta_8$.

For clarity, the kernel needed in degree twelve is equally
controlled. Before transgression into $H^{12}Z$, the only
possible mixed source in its path spectral sequence has
bidegree $(6,5)$; its image is $H^6Z\cdot H^6Z$.
Before transgression into $H^{11}W$, the only mixed sources
have bidegrees $(6,4),(5,5)$; their images lie in
$H^5W\cdot H^6W$. Cohomology suspension kills products,
so these are the exact kernels of the two respective
suspensions in these degrees. The image of
$H^{12}Z\to H^{11}W$ is primitive for loop multiplication.
Every degree-five and degree-six class of $W$ is primitive
by connectivity, and the reduced coproduct on their mixed
products is
\[
 x_5y_6\longmapsto x_5\otimes y_6+y_6\otimes x_5.
 \tag{D.2.27d}
\]
It is injective on the mixed-product subspace: the two tensor
bidegrees differ, so a zero coproduct forces the corresponding
tensor in $H^5W\otimes H^6W$ to vanish and hence its product
to vanish. Thus the second kernel has zero intersection with
the first suspension image. We have proved
\[
 \ker\bigl(H^{12}Z\longrightarrow H^{10}F\bigr)
       =\langle v_6^2\rangle.
 \tag{D.2.27e}
\]
These arguments are just the low path-spectral-sequence edge
calculations; they do not require a collapse in higher degrees.
The exact full-fiber relation
$\operatorname{Sq}^4e_6^*+\operatorname{Sq}^2J_1^*=0$
of \textup{D.2.24}, and suspension compatibility of squares,
now give \textup{(D.2.27b)}.

We next identify its class with the actual relative cup
functional, keeping one normalization for all maps.
Write $q:Z\to Y$. Choose actual cocycles $a_4,b_5$ and a
cochain $c_8$ on $Y$ with $dc_8=a_4\smile b_5$.
Since $q^*a_4$ is exact, choose $P_3$ with
$dP_3=q^*a_4$ and form
\[
 \Theta_8=[q^*c_8-P_3\smile q^*b_5]\in H^8Z.
 \tag{D.2.27f}
\]
Changing $P_3$ does not change this class: its difference
is a closed degree-three class multiplied by the exact
$q^*b_5$. Every $h:A\to Y$ lifts to $Z$, since
$h^*\alpha=0$ by \textup{D.2.12}. The actual cochain formula
of \textup{D.2.26} gives
\[
 \theta_h=\widetilde h^{,*}\Theta_8
 \quad\text{for every fiber-model lift }\widetilde h:A\to Z.
 \tag{D.2.27g}
\]
The independence just checked also handles the choice of lift.

Evaluate \textup{(D.2.27f)} on the counit
$\varepsilon:\Sigma^2F\to Z$. Its composite into $Y$ has
zero homology in degrees four and five by connectivity,
and in degree six because the bottom class $A_4=a^2$
has zero homology suspension. Thus \textup{(D.2.26j)} applies.
Its once-adjunct $\Sigma F\to\Omega Y$ has images on the
three degree-six fiber classes given by
\[
 d_6\longmapsto0,\qquad e_6\longmapsto ab+ba,
       \qquad f_6\longmapsto \text{the linear }f\text{ label}.
 \tag{D.2.27h}
\]
Here $a,b$ are the degree-three and degree-four letters in
$H_*\Omega Y$. The exact per-vector comparison
\textup{(D.2.26j)} therefore gives
$\sigma^2\Theta_8=e_6^*+t f_6^*$ for some $t\in\FF_2$.
The image of $\sigma^2q^*:H^8Y\to H^6F$ is precisely
$\langle f_6^*\rangle$. Changing $c_8$ by that one fixed
old cohomology class removes $t$. Since
$\sigma^2:H^8Z\to H^6F$ is an isomorphism, the resulting
universal class is exactly $\vartheta_8$.
This is a single change on $Y$, shared by every attaching map.

The homology of $A=Q\Sigma B$ is polynomial on primitive
generators: the unit labels are primitive because $\Sigma B$
is a suspension, and their Dyer--Lashof generators remain
primitive. Its dual Hopf algebra consequently has zero square
on every positive-dimensional cohomology class. This also
follows by pairing a square with the primitive-polynomial
coproduct: every nontrivial equal-factor term has even
coefficient. We use the opened free infinite-loop-space
description and Nishida formulas
\cite[Chapters I and III]{I_CohenLadaMay1976}.
Pulling \textup{(D.2.27b)} back along $\widetilde h$ gives
\[
 \operatorname{Sq}^4\widetilde h^{,*}\vartheta_8+
 \operatorname{Sq}^2\widetilde h^{,*}\lambda_{10}=0
                         \quad\text{in }H^{12}A.
 \tag{D.2.27i}
\]

We finish with five explicit witnesses, so no rank assertion
rests only on a computer calculation. Put
$A_0=\sigma a_2$, $B_0=\sigma b_3$, $C_5=\sigma c_4$,
$D_6=\sigma d_5$. Let $I_9,J_9$ be the unit classes dual
to $u^4,uv^2$, and let $G_8,H_8$ be the unit classes dual
to $u^2v,uw$. The complete basis of $H_8A$ is
\[
 Q^5A_0,\quad A_0C_5,\quad H_8,\quad G_8,\quad B_0^2.
 \tag{D.2.27j}
\]
For a base monomial $m$, write $\sigma(m^\vee)$ for its
unit homology dual, and put $t=\operatorname{Sq}^4w$.
The following degree-twelve classes all have zero second
square; their fourth squares are the displayed basis:
\[
\begin{array}{c|c|c}
 z\in H_{12}A&\operatorname{Sq}_*^2z&\operatorname{Sq}_*^4z\\ \hline
 Q^7C_5&0&Q^5A_0\\
 A_0(I_9+J_9)&0&A_0C_5\\
 \sigma((ut)^\vee)&0&H_8\\
 \sigma((u^4v)^\vee)+\sigma((u^3w)^\vee)&0&G_8\\
 D_6^2&0&B_0^2
\end{array}
 \tag{D.2.27k}
\]
Here the first row is the ordinary source Nishida identity:
the only nonzero term of $\operatorname{Sq}_*^4Q^7C_5$
is $Q^5\operatorname{Sq}_*^2C_5=Q^5A_0$, while the
leading coefficient for its second square is
$\binom52=0$ modulo two and $\beta C_5=0$.
The second row uses
$\operatorname{Sq}_*^2I_9=\operatorname{Sq}_*^2J_9=F_7$,
$\operatorname{Sq}_*^4I_9=C_5$, and
$\operatorname{Sq}_*^4J_9=0$.
The last row uses $\beta D_6=0$ and
$\operatorname{Sq}_*^2D_6=B_0$.
For both remaining rows, the complete relevant Cartan lists
in Serre's polynomial algebra are
\[
\begin{aligned}
 \operatorname{Sq}^2(u^3v)&=u^4v+uv^3+u^3w,&
 \operatorname{Sq}^2(u^2w)&=v^2w,\\
 \operatorname{Sq}^2(v^3)&=v^2w,&
 \operatorname{Sq}^2t&=0,\\
 \operatorname{Sq}^4(u^2v)&=u^4v+v^2w,&
 \operatorname{Sq}^4(uw)&=ut.
\end{aligned}
 \tag{D.2.27l}
\]
The degree-nine basis is $u^3v,u^2w,v^3,t$ and the
degree-seven basis is $u^2v,uw$, so these lists have no
omitted unit coefficient; compare
\cite[\S2, nos.~8--9]{I_Serre1953}.
All Dyer--Lashof calculations here occur inside the source;
no ordinary map is assumed to preserve those operations.

Pair \textup{(D.2.27i)} with each row of
\textup{(D.2.27k)}. It forces
$\widetilde h^{,*}\vartheta_8$ to vanish on the entire
basis \textup{(D.2.27j)}. Equations
\textup{(D.2.27g)--(D.2.27h)} prove
\textup{(D.2.27a)} with the one fixed normalization.
The tuple for any small simultaneous unsuspended family
is therefore zero. The exact common-cone criterion of
\textup{D.2.26} proves the asserted cup-product vanishing
and old-commutator survival.

For the mixed-family assertion, let $h:\Sigma^kA\to Y$
with $k\ge2$, lift it to $Z$, and let
$G:\Sigma^{k-2}A\to F$ be the double adjunct of that lift.
With exactly the universal normalization already fixed,
naturality of cohomology suspension gives
\[
 \theta_h=\widetilde h^{,*}\vartheta_8
       =\sigma^2(G^*e_6^*).
 \tag{D.2.27m}
\]
The entire functional $G^*e_6^*$ is zero for $k=2$ by
\textup{D.2.24}, and for $k=3$ and every $k\ge5$ by
\textup{D.2.25}. Thus every component in
\textup{(D.2.27a')} has zero Hopf functional with this one
choice of $c_8$. The simultaneous criterion applies to the
zero tuple, including all cross-attachment cancellations.
Nothing in this argument
identifies the new target after another cone with the
original $Y$, so it supplies no iteration claim.
\end{proof}

\begin{claim}{D.2.28}{Proved}{The once-suspended functional has only one possible quadratic coefficient}
Retain $B=K(\FF_2,2)$, $A=Q\Sigma B$, $Y=\Sigma^2B$,
$Z=\operatorname{fib}(Y\to K(\FF_2,4))$, and the single
normalization of the relative cup class in \textup{D.2.27}.
Put $W=\Omega Z$ and
$\vartheta_7=\sigma\vartheta_8\in H^7W$.
For every ordinary pointed map
\[
 g:\Sigma B\longrightarrow W
\]
one has $g^*\vartheta_7=0$.

For an arbitrary map $h:\Sigma A\to Y$, choose a lift to
$Z$ and let $G:A\to W$ be its adjunct. Let
$\alpha_3\in H^3A$ be the fundamental class and
$\beta_4=\operatorname{Sq}^1\alpha_3$.
Then, with the same fixed normalization for all maps,
\[
 G^*\vartheta_7=\epsilon\alpha_3\beta_4,
 \qquad
 \theta_h=\sigma(\epsilon\alpha_3\beta_4)
       \quad\text{for some }\epsilon\in\FF_2.
 \tag{D.2.28a}
\]
The possible value $\epsilon=1$ is not excluded here.
Thus this is a complete reduction to one coefficient,
not a vanishing theorem for every once-suspended attachment
or an assertion about subsequent nullification stages.
\end{claim}
\begin{proof}
All coefficients are $\FF_2$. Let
$\lambda_9=\sigma\lambda_{10}\in H^9W$.
The square term in \textup{(D.2.27b)} vanishes under
cohomology suspension, giving the exact identity
\[
 \operatorname{Sq}^4\vartheta_7+
       \operatorname{Sq}^2\lambda_9=0.
 \tag{D.2.28b}
\]
There is a second relation in lower degrees.
Let $\mathcal R_8\in H^8W$ be the unique class whose
cohomology suspension is dual to
$R=\{a,c\}\in H_7F$, where $F=\Omega W$.
The low path-fibration ranges in \textup{D.2.27} make
$H^8W\to H^7F$ an isomorphism and
$H^9W\to H^8F$ injective.
The full degree-eight fiber table gives
\[
 \operatorname{Sq}^2e_6^*=J_1^*+J_3^*,
 \qquad \operatorname{Sq}^1R^*=J_3^*.
 \tag{D.2.28c}
\]
For the second identity, $J_3$ is the only basis class
whose Bockstein has $R$ coefficient one.
For the first, $J_1,J_3$ are exactly the basis classes
whose second square has $e_6$ coefficient one.
Suspension compatibility and the displayed injectivity
therefore prove the actual relation
\[
 \lambda_9=\operatorname{Sq}^2\vartheta_7+
               \operatorname{Sq}^1\mathcal R_8.
 \tag{D.2.28d}
\]
Substitute this into \textup{(D.2.28b)} and use the Adem
identity $\operatorname{Sq}^2\operatorname{Sq}^2
=\operatorname{Sq}^3\operatorname{Sq}^1$:
\[
 \bigl(\operatorname{Sq}^4+
       \operatorname{Sq}^3\operatorname{Sq}^1\bigr)
       \vartheta_7
       +\operatorname{Sq}^2\operatorname{Sq}^1
          \mathcal R_8=0\quad\text{in }H^{11}W.
 \tag{D.2.28e}
\]
This proof uses only the established low suspension
interfaces and the already delooped relation. It does not
replace $H^{11}W$ by an indecomposable quotient or omit
possible higher bar-filtration classes.

Let $u,v,w$ be the Serre generators of degrees $2,3,5$,
with $v=\operatorname{Sq}^1u$ and
$w=\operatorname{Sq}^2v$.
For a map $g:\Sigma B\to W$ write, in the complete
cohomology bases of the source,
\[
 \begin{split}
 g^*\vartheta_7&=a\sigma(u^3)+b\sigma(v^2),\\
 g^*\mathcal R_8&=c\sigma(u^2v)+d\sigma(uw).
 \end{split}
 \tag{D.2.28f}
\]
Cartan and the polynomial description
\cite[\S2, nos.~8--9]{I_Serre1953} give
\[
\begin{aligned}
 (\operatorname{Sq}^4+\operatorname{Sq}^3
                  \operatorname{Sq}^1)(u^3)&=u^5,\\
 (\operatorname{Sq}^4+\operatorname{Sq}^3
                  \operatorname{Sq}^1)(v^2)&=w^2,\\
 \operatorname{Sq}^2\operatorname{Sq}^1(u^2v)&=0,\\
 \operatorname{Sq}^2\operatorname{Sq}^1(uw)&=w^2+u^2v^2.
\end{aligned}
 \tag{D.2.28g}
\]
In the first row,
$\operatorname{Sq}^4u^3=u^5+u^2v^2$ and
$\operatorname{Sq}^3\operatorname{Sq}^1u^3
=\operatorname{Sq}^3(u^2v)=u^2v^2$.
The last row follows from
$\operatorname{Sq}^1(uw)=vw+uv^2$.
Pulling \textup{(D.2.28e)} back gives
\[
 a\sigma(u^5)+(b+d)\sigma(w^2)
       +d\sigma(u^2v^2)=0.
 \tag{D.2.28h}
\]
These three monomials are independent, so $a=b=d=0$.
This proves the universal unit assertion without first
requiring the map to extend to $A$.

For the full source, the complete basis of $H_7A$ is
\[
 Q^4A_0,\quad A_0B_0,\quad F_7,\quad X_7,
 \tag{D.2.28i}
\]
where $A_0=\sigma a_2$, $B_0=\sigma b_3$,
$F_7=\sigma((u^3)^\vee)$, and
$X_7=\sigma((v^2)^\vee)$.
The unit assertion applied to $G|_{\Sigma B}$ kills
the last two coefficients of $G^*\vartheta_7$.
Set $C_5=\sigma((u^2)^\vee)$.
The source Nishida identities give
\[
 \operatorname{Sq}_*^2Q^6C_5=0,
 \qquad
 \operatorname{Sq}_*^4Q^6C_5=Q^4A_0.
 \tag{D.2.28j}
\]
For the second square the leading coefficient is
$\binom42=0$ modulo two and $\beta C_5=0$.
For the fourth square only the term
$Q^4\operatorname{Sq}_*^2C_5=Q^4A_0$ survives.
These are operations inside the free infinite-loop source,
using \cite[Chapters I and III]{I_CohenLadaMay1976};
no ordinary map is assumed to preserve Dyer--Lashof operations.
Pairing \textup{(D.2.28b)} after pullback with $Q^6C_5$
kills the first coefficient in \textup{(D.2.28i)}.

The cohomology class dual to the remaining monomial $A_0B_0$
is exactly $\alpha_3\beta_4$. The unit classes and the
Dyer--Lashof generator in \textup{(D.2.28i)} are primitive,
whereas the reduced diagonal of $A_0B_0$ is
$A_0\otimes B_0+B_0\otimes A_0$.
This proves the first equality in \textup{(D.2.28a)}.
Its second equality is naturality of cohomology suspension
for the actual lift to $Z$ and the universal class of
\textup{D.2.27}.

The retained coefficient cannot be discarded by these
relations alone. In the source write
$\gamma_5=\operatorname{Sq}^2\alpha_3$.
The pair
\[
 q_7=\alpha_3\beta_4,
       \qquad r_8=\alpha_3\gamma_5
 \tag{D.2.28k}
\]
satisfies \textup{(D.2.28e)} algebraically: positive cup
squares vanish in $H^*A$, so
$\operatorname{Sq}^1q_7=0$, and Cartan gives
$\operatorname{Sq}^2\operatorname{Sq}^1r_8
=\operatorname{Sq}^4q_7$.
This is only compatibility of cohomology classes, not
realization by a map. It explains precisely why the
proof stops at the single unresolved coefficient.
\end{proof}

\begin{claim}{D.2.29}{Proved}{A fourth-suspension unit map and its canonical obstruction}
Let $B=K(\FF_2,2)$, $X=\Sigma B$, $Y=\Sigma^2B$, and
$C_mX=\Omega^m\Sigma^mX$.  Let $u_4:S^4\to Y$ be the
bottom map.  The generalized Whitehead product defines an actual map
\[
 h_0=[u_4,\id_Y]:\Sigma^5B=\Sigma^4X\longrightarrow Y.
 \tag{D.2.29a}
\]
Its double adjunct sends $\sigma^3b$, where $b\in H_3B$,
to $k=\{a,b\}\in H_6(\Omega^2Y;\FF_2)$, with no linear
$p_3$ component.  In the cofiber of $h_0$, the unique extensions
of the old degree-four class and its first square have nonzero product.
Thus the potentially dangerous fourth-suspension pattern is realized
on the \emph{entire unit}, not only on a finite skeleton.

The iterated counits give compatible canonical extensions
\[
 h_m:\Sigma^4C_mX\longrightarrow Y\qquad(0\le m\le4).
 \tag{D.2.29b}
\]
The specified canonical $h_4$ has no ordinary pointed extension
to $\Sigma^4C_6X$, and therefore none to $\Sigma^4QX$.
This conclusion fixes the restriction on $\Sigma^4C_4X$.
It does \emph{not} exclude an arbitrary ordinary extension to
$\Sigma^4QX$ that agrees only with $h_0$ on $\Sigma^4X$.
\end{claim}
\begin{proof}
All homology in the proof has coefficients $\FF_2$, and
$\beta=\operatorname{Sq}_*^1$.  Set $T=\Omega^2Y$ and
use the labels $a,b$ of degrees two and three, respectively,
with $\beta b=a$.  Write $B_0=\sigma b\in H_4X$.

\medskip\noindent
\textbf{The actual unit map and its cone.}
The source of the generalized Whitehead product of
$S^4=\Sigma S^3\to Y$ and $Y=\Sigma(\Sigma B)\to Y$ is
\[
 \Sigma(S^3\wedge\Sigma B)=\Sigma^5B.
\]
This gives \textup{(D.2.29a)} with the stated suspension indices.
The construction and its identification with the adjoint commutator
are \cite[Definitions~2.2--2.3 and Theorem~2.4,
pp.~10--13]{EH5Arkowitz1962}.  Its double adjunct is the
based generalized Browder construction on the bottom map
$S^2\to B\to T$ and the unit $B\to T$.
The mapping-cylinder construction in the remark following
Definition~2.3 applies to arbitrary based CW models, so the
infinite source creates no finite-polyhedron restriction.
For the primitive bottom classes $a,b$, its homology values are
\[
 (h_0^{\flat\flat})_*(\sigma^3a)=\{a,a\}=0,
 \qquad
 (h_0^{\flat\flat})_*(\sigma^3b)=\{a,b\}=k.
 \tag{D.2.29c}
\]
These are the geometric Browder operation and its suspension
compatibility in \cite[Chapter~III, Theorems~1.2--1.4,
pp.~216--219]{I_CohenLadaMay1976}.
The formula here is intentionally restricted to primitive inputs.
For nonprimitive labels the based family has convolution corrections;
no formula $x\mapsto\{a,x\}$ on every label is used.
On the bottom sphere the map is the ordinary Whitehead square
$[u_4,u_4]=-2\alpha$, with the orientation convention of D.2.17.
Its zero mod-two Hurewicz value in \textup{(D.2.29c)} does
not assert that this bottom homotopy element is zero.

Under one homology suspension, $k$ maps to the tensor
commutator $[\bar a,\bar b]\in H_7\Omega Y$; under the next,
that commutator maps to zero.  The group
$H_8(\Sigma^5B)$ is one-dimensional on $\sigma^5b$.
Thus the degree-eight attaching homology map is zero, whereas
the exact suspended-source cone criterion of D.2.26 evaluates
the old cup product to one on the connecting class over
$\sigma^5b$.  The source is six-connected, so the old
degree-four and degree-five cohomology extensions are unique.
Their product is therefore nonzero, with no choice-of-lift
qualification.  The absence of a $p_3$ component follows
directly from the second equality in \textup{(D.2.29c)}.

\medskip\noindent
\textbf{The compatible canonical stages.}
For $0\le m\le4$, let $\varepsilon^{(m)}$ be the counit
of $\Sigma^m\dashv\Omega^m$, evaluated at $\Sigma^mX$, and set
\[
 r_m=\Sigma^{4-m}\varepsilon^{(m)}:
       \Sigma^4C_mX\longrightarrow\Sigma^4X,
 \qquad h_m=h_0r_m.
 \tag{D.2.29d}
\]
The triangle identities give extension over the unit, and
naturality of the iterated counits gives compatibility as $m$
increases.  The fourfold adjunct of $h_4$ is exactly the
fourfold loop map
\[
 g_4=\Omega^4h_0:C_4X\longrightarrow\Omega^4Y.
 \tag{D.2.29e}
\]
In particular it preserves the $E_4$ homology operations.
If $q=(g_4)_*(B_0)\in H_4\Omega^4Y$, the twofold
homology suspension satisfies
\[
 \sigma^2q=k\in H_6T,
 \tag{D.2.29f}
\]
by adjunction and \textup{(D.2.29c)}.

\medskip\noindent
\textbf{A Bockstein class in the actual target.}
The element
\[
 \xi_1k=Q^7k\in H_{13}T
 \quad\hbox{satisfies}\quad
 \xi_1k\notin\beta H_{14}T.
 \tag{D.2.29g}
\]
Here is a complete basis argument for the nonboundary assertion.
The free $E_2$ homology theorem gives a polynomial basis whose
indecomposables are the ordinary Hall Lie words in the labels,
and their iterated outer restricted squares $\xi_1$;
see \cite[Chapter~III, Theorems~3.1--3.2 and Lemma~3.8]{I_CohenLadaMay1976}.
The Hall word $k=\{a,b\}$ is nonzero, and its outer square
$\xi_1k$ is an independent restricted-square generator.
Every outer $\xi_1$ output has odd degree $2|z|+1$.
There is consequently no such generator in degree fourteen.

A degree-fourteen polynomial monomial with at least two positive
factors has decomposable Bockstein, by Cartan.  A label has a
linear combination of labels as Bockstein.  Finally, expanding
the Bockstein of an ordinary Hall bracket retains an outer
bracket.  If a restricted square occurs inside such a bracket,
the restricted-Lie identity
$\{v,\xi_1w\}=\{\{v,w\},w\}$ reduces it to ordinary
Hall words.  The top Bockstein formula
\[
 \beta\xi_1z=(|z|\bmod2)z^2+\{\beta z,z\}
 \tag{D.2.29h}
\]
and the bracket Cartan formula introduce only ordinary brackets
and products in this expansion.  No term can become the
independent \emph{outer} restricted square $\xi_1k$.
These exhaust the polynomial basis in degree fourteen and prove
\textup{(D.2.29g)}.  The identities used here are
\cite[Chapter~III, Theorems~1.2--1.4,
pp.~216--219]{I_CohenLadaMay1976}.

\medskip\noindent
\textbf{The sixth-stage contradiction.}
Suppose an ordinary pointed map
$G:C_6X\to\Omega^4Y$ restricts to $g_4$ on $C_4X$.
The class $Q^7B_0\in H_{11}C_4X$ is the top $E_4$
operation on $B_0$.  In $C_6X$, $Q^8B_0$ is below the
top operation, so the ordinary Nishida Bockstein formula gives
\[
 \beta Q^8B_0=Q^7B_0.
 \tag{D.2.29i}
\]
By preservation of operations for the specified $E_4$ map
$g_4$, and homology-suspension compatibility,
\[
 \begin{split}
 \sigma^2G_*(Q^7B_0)
   &=\sigma^2(g_4)_*(Q^7B_0)\\
   &=\sigma^2Q^7q
     =Q^7(\sigma^2q)=Q^7k=\xi_1k.
 \end{split}
 \tag{D.2.29j}
\]
Naturality of the Bockstein for the \emph{ordinary} map $G$
and for homology suspension, together with
\textup{(D.2.29i)}, makes the same class equal to
\[
 \beta\bigl(\sigma^2G_*(Q^8B_0)\bigr),
\]
contrary to \textup{(D.2.29g)}.
No preservation of multiplication or Dyer--Lashof operations
has been assumed for $G$.

The use of $C_6X$ is deliberate.  At $C_5X$, $Q^8B_0$
is the top operation; its Bockstein formula has a Browder
correction involving $\beta B_0$ and $B_0$.
The argument does not discard that correction and makes no
nonextension assertion at the fifth stage.
Finally, if a map agrees only on the unit, its restriction to
$C_4X$ need not be the $E_4$ map $g_4$.
Equation \textup{(D.2.29j)} is then unavailable.
The asserted obstruction therefore has exactly the fixed
canonical scope stated in the claim.
\end{proof}

\begin{claim}{D.2.30}{Proved}{The remaining once-suspended coefficient is realized through dimension ten}
Put $B=K(\FF_2,2)$, $A=Q\Sigma B$, $Y=\Sigma^2B$,
$Z=\fib(Y\to K(\FF_2,4))$, and $W=\Omega Z$.
Use the common class $\vartheta_7\in H^7W$ of D.2.28.
There is an actual ordinary map on a two-local ten-skeleton
of $A$,
\[
 g_{10}:K_{10}\longrightarrow W,
 \qquad
 g_{10}^*\vartheta_7=
       (\alpha_3\operatorname{Sq}^1\alpha_3)|_{K_{10}}\ne0.
 \tag{D.2.30a}
\]
Here the finite cellular model is chosen through dimension eleven
before cancelling odd-unit cell pairs, as in D.2.19--D.2.20.
The map can be pulled back to a corresponding ordinary finite
skeleton by cellular approximation and localization.

This realizes the value $\epsilon=1$ left in D.2.28 by one
coordinated choice of extensions. It does not extend the map
over the eleventh skeleton or over all of $A$, and it gives
no assertion about subsequent nullification stages.
\end{claim}
\begin{proof}
All homology and cohomology below have coefficients $\FF_2$,
except for the displayed homotopy groups and Moore spaces.
Write $h$ for the ordinary mod-two Hurewicz homomorphism.
Write $F=\Omega W$. Its classes $d,e,f$ in degree six are
\[
 d=b^2,\qquad e=b^2+\{a,b\},\qquad f=p_3,
 \qquad r=e+f.
\]
We use the same letters for their homology suspensions in
degree seven of $W$. Let $R_8\in H_8W$ be the suspension
of $R=\{a,c\}\in H_7F$.

\medskip\noindent
\textbf{The actual target groups and the missing Hurewicz direction.}
The target calculations D.2.18--D.2.19 and D.2.22 give
\[
 \begin{aligned}
 \pi_7W&=\ZZ/4\{\rho\},& h(\rho)&=r,\\
 \pi_8W&=\ZZ/2\{\lambda\},& h(\lambda)&=R_8,\\
 \pi_9W&\cong(\ZZ/2)^2,&\ker h
       &=\ZZ/2\{w_9\}.
 \end{aligned}
 \tag{D.2.30b}
\]
The last assertion uses the actual image calculation in
\textup{(D.2.22t19)}. The two actual eight-spheres in $F$
have Hurewicz images $A_4^2$ and $J_2$.
Homology suspension kills $A_4^2$ and retains $J_2$;
the path-fibration calculation of D.2.27 identifies
$H_9W$ with $H_8F/\langle A_4^2\rangle$.
Thus $w_9$ is the adjoint of the actual sphere with
Hurewicz image $A_4^2$. In particular, Hurewicz injectivity
for $F$ has not been transported incorrectly to $W$.

Let $a_5\in H^5W$ be the bottom class. In the actual
mapping cone $C_{w_9}$ it has a unique extension
$\widehat a_5$, and
\[
 a_5^2=0\text{ in }H^{10}W,
 \qquad
 \widehat a_5^{\,2}=\text{the nonzero relative top class}
       \text{ in }H^{10}C_{w_9}.
 \tag{D.2.30c}
\]
Here is the precise detector. In a minimal Adams--Hilton
model for the four-connected space $W$, the bottom
generator has degree four and represents $A_4$.
The adjoint of $w_9$ represents $A_4^2$ in
$H_8\Omega W=H_8F$. Attaching the ten-cell adds a
degree-nine generator whose differential represents this
class. Subtracting an old chain makes that differential
exactly $A_4^2$. Its quadratic part is the reduced diagonal
of the new top class, by the actual cell-attachment and
cup-product comparison in
\cite[\S3.1, p.~300]{FGHHalperinLemaire1995}.
This gives its square coefficient one.
The old square $a_5^2$ is zero: a nonzero old coefficient
would already make $A_4^2$ a boundary. In this degree no
other word can cancel it, since positive generators begin
in degree four and the bottom space is one-dimensional.
The nonzero established class $A_4^2\in H_8F$ excludes
that possibility. Thus the new square is precisely the
relative top class, proving \textup{(D.2.30c)}.

\medskip\noindent
\textbf{A coordinated ordinary map on the eight-skeleton.}
Use the finite two-local cellular models of
D.2.19--D.2.20 and put $Q_i=K_i/K_4$.
Choose the model through dimension eleven before
cancelling odd-unit pairs. Its mod-two cellular
differential is then zero in the dimensions used below.
Cancellation is performed on the actual two-local
cellular approximation, with higher attachments
transported along the equivalence; no unstable map is
replaced merely by its stable homology matrix.
The target $W$ is two-local, so the original mapping
problem is recovered by localization and cellular
approximation.

The actual eight-stage decomposition
\textup{(D.2.20s1)} is
\[
 Q_8\simeq D'\vee M(\ZZ/2,7)_F\vee M(\ZZ/4,7)_J
                \vee C_1^8\vee S^8_s.
 \tag{D.2.30d}
\]
Here $D'$ is the unit block, $J$ is the seven-sphere
with mod-two label $Q^4A_0+A_0B_0$,
\[
 C_1^8=\cofib\bigl((2,\eta_6):S^6\vee S^7\to S^6\bigr),
 \qquad s=Q^5A_0-A_0C_5.
\]
The top eight-cell of $C_1^8$ is $A_0C_5$.
The top cell of $M(\ZZ/4,7)_J$ is the slide
$B_0^2-A_0C_5$. All these comparisons are actual:
the quotients are four-connected and have dimension
at most eight, within the ordinary realization range
of the stable cell slides.

Set the map to zero on $D'$ and $M(\ZZ/2,7)_F$.
Send the bottom sphere of $M(\ZZ/4,7)_J$ to $\rho$.
It extends across the top cell because $4\rho=0$.
Any extension sends its top mod-two class to a multiple
of $R_8$: that source class has zero Bockstein and
zero second square, and the complete target table
in degree eight has simultaneous kernel exactly
$\langle R_8\rangle$.
Changing the extension by the actual sphere $\lambda$
toggles this coefficient. Choose the top value zero.

On $C_1^8$, send the six- and seven-skeleton to zero
and send the top eight-cell, through its pinch, to
$\lambda$. Send $s$ to $-\lambda$.
These ordinary maps define an actual $g_8:Q_8\to W$.
In the original source labels their values are
\[
 \begin{aligned}
 (g_8)_*(Q^4A_0)&=0,&(g_8)_*(A_0B_0)&=r,\\
 (g_8)_*(A_0C_5)&=R_8,&(g_8)_*(Q^5A_0)&=0,\\
 (g_8)_*(B_0^2)=R_8.
 \end{aligned}
 \tag{D.2.30e}
\]
The last value uses the stated slide of the Moore top.
The unit degree-seven classes have value zero.
Since $\vartheta_7$ is dual to the $e$ coordinate,
the complete source degree-seven basis of D.2.28 gives
\[
 g_8^*\vartheta_7=
          (\alpha_3\operatorname{Sq}^1\alpha_3)|_{Q_8},
 \qquad g_8^*a_5=0.
 \tag{D.2.30f}
\]
For the first formula, the notation on the right means
the unique quotient class pulling back to that class
on $K_8$. Collapsing $K_4$ is an isomorphism on
cohomology in degree seven.

\medskip\noindent
\textbf{Every ninth and tenth attachment extends.}
Each nine-cell has zero mod-two cellular boundary.
Its attaching map has zero mod-two Hurewicz in $Q_8$,
so its composite with $g_8$ has zero Hurewicz in $W$.
The middle line of \textup{(D.2.30b)} is injective.
All these composites are therefore null, and choices
of their nullhomotopies give an actual $g_9:Q_9\to W$.

The same cellular argument makes every ten-cell
composite belong to $\{0,w_9\}$.
Suppose one were $w_9$. Its actual attachment supplies
a map of cones
\[
 Q_9\cup e^{10}\longrightarrow C_{w_9}
 \tag{D.2.30g}
\]
whose relative top-cell degree is one. The source
ten-cell is a nonzero mod-two relative cohomology
coordinate, since its mod-two cellular boundary is
zero. Hence the relative top class in
\textup{(D.2.30c)} pulls back nontrivially.
On the other hand, $g_9^*a_5=0$ by
\textup{(D.2.30f)}, and attaching a ten-cell does not
change degree-five cohomology. The pullback of
$\widehat a_5$ in \textup{(D.2.30g)} is therefore
zero, so its square is zero. This contradicts
\textup{(D.2.30c)}.

Every tenth-cell obstruction is consequently zero.
The cell extensions attach to the same fixed $Q_9$,
so they assemble to an ordinary map $Q_{10}\to W$.
Its composite with $K_{10}\to Q_{10}$ is the map
in \textup{(D.2.30a)} and retains all the values
in \textup{(D.2.30e)}.
The next attaching group is $\pi_{10}W=\pi_9F$,
where D.2.22 has an integral-Hurewicz kernel and
an additional direction detected by the order-eight
homology class. This proof does not remove those
eleventh-cell obstructions.
\end{proof}

\begin{claim}{D.2.31}{Proved}{The remaining extension bit vanishes}
Keep $B=K(\FF_2,2)$, $Y=\Sigma^2B$, and $T=\Omega^2Y$.
The actual homotopy group left partly undetermined in
\textup{D.2.22} is
\[
       \pi_9T=\pi_{11}Y\cong(\ZZ/2)^4.
       \tag{D.2.31a}
\]
In particular, both even-nu attachment coefficients of that
calculation are zero. There is an actual element
$\gamma\in\pi_9T$ of order two with
\[
       h_{\FF_2}(\gamma)=K_5=\{a,\{a,b\}\}.
       \tag{D.2.31b}
\]
The integral Hurewicz kernel remains the nonzero group
$\ZZ/2\{\rho\nu_6\}$ of \textup{D.2.22}.
\end{claim}
\begin{proof}
Write $X=\Sigma B$ and let
\[
       M=M(\ZZ/2,3)=P^4(2).
\]
Suspending the bottom Moore map $M(\ZZ/2,2)\to B$
gives a map $j:M\to X$. Denote the two mod-two homology
generators of $M$ by $u_3,v_4$; their images are respectively
$\sigma a$ and $\sigma b$. We work two-locally in the
construction. The homotopy groups of $Y$ in question are
two-primary, as already established in \textup{D.2.22}, so
this computes the actual group in the claim.

The spherical commutator construction of
\cite[Proposition~2.1, pp.~3249--3250]{I_ChenWu2013}
gives an actual map
\[
 \alpha:S^7\longrightarrow M\wedge M,
 \qquad \alpha_*(\iota_7)=uv+vu.
 \tag{D.2.31c}
\]
Let $E_M:M\to\Omega\Sigma M$ be the James unit. Its
\emph{left-nested} triple Samelson product is the actual map
\[
 S_3=[[E_M,E_M],E_M]:M^{\wedge3}\longrightarrow\Omega\Sigma M.
 \tag{D.2.31d}
\]
This is also the map used in
\cite[Section~1 and the proof of Theorem~1.6]{I_Wu1998}.
Since $M$ is a suspension, both homology generators are
primitive, and the induced map on tensor generators is
$\beta_3(xyz)=[[x,y],z]$. Consider the composite
\[
\begin{split}
 M\wedge S^7
 &\xrightarrow{\ \id\wedge\alpha\ } M^{\wedge3}
 \xrightarrow{\ S_3\ }\Omega\Sigma M
 \xrightarrow{\ \Omega\Sigma j\ }\Omega Y.
\end{split}
 \tag{D.2.31e}
\]
Its source is $M(\ZZ/2,10)$. The bottom sphere $S^{10}$
represents the order-two generator of its first nonzero
homotopy group, by the ordinary Hurewicz theorem. Restriction
of \textup{(D.2.31e)} to this sphere therefore represents an
element of order at most two in $\pi_{10}\Omega Y$.
Equivalently, its adjoint is the restriction to the bottom
$S^{11}$ of an actual map
$\Sigma(M\wedge S^7)=P^{12}(2)\to Y$.
Let $\gamma\in\pi_9T$ be the corresponding double adjunct.

The once-adjoint Hurewicz image is computed directly:
\[
 \beta_3(uuv+uvu)=uuv+vuu.
 \tag{D.2.31f}
\]
Indeed $\beta_3(uuv)=0$ over $\FF_2$, whereas
$\beta_3(uvu)=uuv+vuu$. Thus in $H_{10}(\Omega Y;\FF_2)$
the image of this sphere is the nonzero class
\[
      (\sigma a)^2(\sigma b)+(\sigma b)(\sigma a)^2.
      \tag{D.2.31g}
\]
No order for a sphere chosen by an unrelated Moore coextension
has been assumed: the order-two assertion here follows from
the actual source of \textup{(D.2.31e)}.

The complete presentation and integral Hurewicz lattice in
\textup{D.2.22}, for either value of its retained common bit,
show that
\[
 \operatorname{im}\bigl(h_{\FF_2}:\pi_9T\to H_9T\bigr)
       =\langle K_1,K_5\rangle.
 \tag{D.2.31h}
\]
In fact its generators $p,s$ have zero mod-two Hurewicz
image, while $x,y$ have images $K_1,K_5$. The homology
suspension from $T$ to $\Omega Y$ satisfies
\[
\begin{aligned}
 \sigma K_1&=(\sigma a)(\sigma f)+(\sigma f)(\sigma a),\\
 \sigma K_5&=(\sigma a)^2(\sigma b)+(\sigma b)(\sigma a)^2.
\end{aligned}
 \tag{D.2.31i}
\]
Here the product correction in $K_1=\{a,f\}+aR$ suspends
to zero, and a double-loop Browder bracket suspends to the
ordinary loop commutator. These are the homology-suspension
identities used throughout the first-cone calculation; see
\cite[Chapter~III, Theorem~1.4]{I_CohenLadaMay1976}.
The two expressions are independent in the James tensor
algebra $H_*(\Omega\Sigma X;\FF_2)$: one has word length
two and the other has word length three, and both are nonzero.
Naturality under adjunction, together with
\textup{(D.2.31g)--(D.2.31i)}, forces
$h_{\FF_2}(\gamma)=K_5$. Hence $\gamma$ has exact order two.

For completeness, write the two coefficients in the
\textup{D.2.22} presentation as $\epsilon_1,\epsilon_5$,
to distinguish them from the homology classes $a,b$:
\[
 \pi_9T=\langle p,x,y,s\mid
       2p=2s=0,\quad2x=\epsilon_1p,\quad
       2y=\epsilon_5p\rangle.
 \tag{D.2.31j}
\]
Every element with mod-two Hurewicz image $K_5$ has double
$\epsilon_5p$: the difference from $y$ is a combination
of $p,s,2x,2y$, all of whose doubles vanish.
Since $2\gamma=0$ and $p\ne0$, we obtain
$\epsilon_5=0$. The actual order-two Browder representative
of $K_1+K_5$ already proved in \textup{D.2.22} gives
$\epsilon_1=\epsilon_5$, so $\epsilon_1=0$ as well.
The complete presentation \textup{(D.2.31j)} now gives
\textup{(D.2.31a)}. Its previously computed integral
Hurewicz kernel is unchanged.
\end{proof}

\begin{claim}{D.2.32}{Proved}{The coordinated once-suspended map stops at the eleventh skeleton}
Use the spaces and the two-local cellular models of D.2.30, and put
$R_i=K_i/K_6$. There is a specified ordinary map
\[
 g_9:R_9\longrightarrow W=\Omega\fib(\Sigma^2K(\FF_2,2)
                         \to K(\FF_2,4))
\]
with the degree-eight values of D.2.30, whose composite from $K_9$
pulls $\vartheta_7$ back to $\alpha_3\operatorname{Sq}^1\alpha_3$.
It extends over $R_{10}$, but no such extension extends over $R_{11}$.
For the normalized eleven-cell relation described below, every
ten-cell filling has the same nonzero obstruction
\[
 p\in\pi_{10}W\cong(\ZZ/2)^4,
 \qquad h_{\ZZ}(p)=4N',\qquad \Psi(p)=0.
 \tag{D.2.32a}
\]
Here $h_{\ZZ}$ is computed in the equivalent group $\pi_9E$ of
D.2.22, $N'$ generates its order-eight homology summand, and
$\Psi$ is the suspended cone detector of D.2.23.

The lower map is fixed explicitly in the proof. This does not
classify other lower maps even when they vanish on $K_6$
(for example, a different choice on the $J$-Moore top), nor
maps nonzero on $K_6$. It does not prove that the coefficient
in D.2.28 always vanishes or assert survival through later
nullification stages.
\end{claim}
\begin{proof}
All undisplayed homology coefficients are $\FF_2$; integral
calculations are two-local. Retain the source labels of
D.2.19--D.2.20 and abbreviate $A_0,B_0,C_5,D_6,E_6,F_7,X_7$
by $A,B,C,D,E,F,X$ within this proof. Set $R=Q^4A$,
$P=AC$, $N=Q^5A$, and $J=AB+R$ in mod-two degree seven.

\medskip\noindent
\textbf{The ordinary range and the fixed lower map.}
The source $R_i$ is six-connected. Its maps to $W$ lift to the
five-connected cover $W'$; write $E=\Omega W'$ as in D.2.22.
The evaluation $\Sigma E\to W'$ is eleven-connected: its
homotopy fiber is $E*E$, which is ten-connected. Thus relative
lifts through dimension eleven are available. Moreover $\Sigma E$
is five-connected, so its tenth homotopy group is stable. All
source attachments used below are in the ordinary stable range
of a six-connected source. These facts justify the finite stable
cell calculations with their specified ordinary maps; no
unrestricted stable replacement of the mapping problem is used.

Collapse $K_6$ in the established source eight- and nine-cell
tables. Before the nine-cells, the factors are
\[
 M(\ZZ/2,7)_F\vee M(\ZZ/4,7)_J\vee S^7_X\vee S^7_R
           \vee S^8_H\vee S^8_P\vee S^8_N.
\]
The nine-attachments, in the previous unit notation $I_9,J_9,L_9$,
are
\[
\begin{array}{c|c@{\qquad}c|c}
 I_9&\eta F&J_9&2H+\eta F\\
 L_9&2H&AD&\eta J+\eta R\\
 AE&0&BC&2P+\eta J+\eta R\\
 Q^6A&2N&Q^5B&\eta R\\
 A^3&0&&
\end{array}
 \tag{D.2.32b}
\]
The former even-$\nu$ and lower torsion ambiguities disappear
on collapsing $K_6$. The remaining integral degrees and second
squares detect every component in this range. Cell slides give
two factors on which the chosen map will be nonzero:
\[
 \begin{aligned}
 J_*&=M(\ZZ/4,7)_J\cup_{i\eta}e^9_U,
                 &&U=AD-Q^5B,\\
 P_*&=M(\ZZ/2,8)_P,
                 &&\text{top cell }BC-AD.
 \end{aligned}
 \tag{D.2.32c}
\]
The other factors are
$S^7_R\cup_\eta e^9_{Q^5B}$,
$M(\ZZ/2,7)_F\cup_{i\eta}e^9_{I_9}$,
$M(\ZZ/2,8)_H$, $M(\ZZ/2,8)_N$, $S^7_X$, and the
three free nine-spheres with labels $J_9-I_9-L_9,AE,A^3$.

Use the target $L_8$ of D.2.22. Send $J_*$ to $\Sigma L_8$
by the Moore extension of $\rho$ and the nullhomotopy of
$\rho\eta$ supplied by its $c$-cell. Send $P_*$ to the
suspended target $R$-Moore factor. Send every other factor
to zero. Composing with evaluation defines the specified
$g_9$. Its Moore top has zero Hurewicz, its $P$ value is
$R_8$, and its $N$ value is zero. These are exactly the
coordinated choices of D.2.30, including the original
$B^2$ value $R_8$. In particular it extends over dimension
ten by that theorem. The specified nine-map already lands
in $\Sigma L_8$, a fact needed to control integral fillings.

\medskip\noindent
\textbf{All integral degree-ten cycles.}
Let $m,n,o,t$ denote the four unit ten-classes dual to
$u^3v,u^2w,v^3,t_9$. The complete degree-ten basis is
\[
 AR,Q^6B,Q^7A,A^2B,AF,AX,BD,BE,C^2,m,n,o,t.
\]
Its nonzero first Bocksteins are
\[
\begin{gathered}
 \beta(AR)=\beta(A^2B)=A^3,\qquad
 \beta(Q^6B)=Q^5B,\qquad \beta(AX)=AD+AE,\\
 \beta(BD)=AD,\qquad\beta(BE)=AE,\qquad
 \beta(o)=J_9+L_9.
\end{gathered}
 \tag{D.2.32d}
\]
They have rank five. Define seven cycles
\[
 \begin{array}{lll}
 K_1=AR+A^2B,&K_2=C^2,&K_3=Q^7A,\\
 K_4=AF,&K_5=m+n,&K_6=t,\\
 K_7=AX+BD+BE.&&
 \end{array}
 \tag{D.2.32e}
\]
The full degree-eleven Bockstein table is equivalently the
following list of its eight nonzero entries, together with
the nine listed zero entries:
\[
\begin{array}{c|cccccccc}
 z&BR&Q^6C&Q^8A&AG&BF&\sigma(u^2v^2)^\vee
                         &\sigma(w^2)^\vee&BX\\ \hline
 \beta z&K_1&K_2&K_3&K_4&K_4&K_5&K_6&K_7
\end{array}
 \tag{D.2.32f}
\]
The zero entries are
$AQ^5A,Q^7B,A^2C,AH,AB^2,CD,CE,
\sigma(u^5)^\vee,\sigma(uvw)^\vee$.
Thus no nonprimitive degree-eleven source class is omitted.
Cartan proves the product entries, and
$\beta Q^s=(s-1)Q^{s-1}$ proves the Dyer entries, including
$Q^5C=C^2$. The unit entries follow from Serre's polynomial
algebra and $\operatorname{Sq}^1t_9=w^2$
\cite[\S2, nos.~8--9]{I_Serre1953}
\cite[Chapter~I, Theorem~1.1]{I_CohenLadaMay1976}.
These are operations inside the source, not operations
preserved by arbitrary ordinary maps.

The known $H_8A=(\ZZ/2)^3$ and the mod-two dimensions
$\dim H_9A=9$, $\dim H_{10}A=13$ show by UCT that
$H_9A$ has six cyclic summands and $H_{10}A$ has seven.
All these integral groups are finite two-primary groups.
Rank seven in (D.2.32f) forces every summand of $H_{10}A$
to have exponent two. Consequently
\[
 H_{10}(A;\ZZ_{(2)})\cong(\ZZ/2)^7.
 \tag{D.2.32g}
\]
The integral cycle lattice of $R_{10}$ is free of rank
seven, with a basis reducing to (D.2.32e). Its eleven-boundary
image is exactly twice the lattice. Odd-unit changes and
even cell additions therefore normalize seven eleven-cell
boundaries to $2K_1,\ldots,2K_7$ and every other boundary
to zero, while preserving all mod-two labels.

The last relation can be retained exactly. Choose the
elementary integral chains with
$dB=2A$, $dD=4C$, $dE=-4C$, $dX=2(D+E)$.
The integral nine-cycle labels used in (D.2.32b) are
$AD'=AD+2BC$ and $AE'=AE-2BC$: thus
$d(BD)=2AD'$, $d(BE)=2AE'$, and
$d(AX)=-2(AD'+AE')$. Their reductions are the displayed
mod-two labels $AD,AE$; the following product-chain
identity is unchanged by this normalization.
The actual Pontryagin chain product gives
\[
 d(AX+BD+BE)=0,\qquad d(BX)=2(AX+BD+BE)=2K_7.
 \tag{D.2.32h}
\]
A cellular chain comparison transports these identities
to the finite model. The nonzero mod-two class makes
$K_7$ a primitive two-local cycle coordinate. This
normalization does not assume that $K_7$ is spherical
before its own ten-cell is attached.

\medskip\noindent
\textbf{The Moore coextension detected by the seventh cycle.}
The complete second-square table on the seven cycles is
\[
 \operatorname{Sq}_*^2K_i=0\ (i=1,2,3,4,6),\qquad
 \operatorname{Sq}_*^2K_5=G_8+H_8,\qquad
 \operatorname{Sq}_*^2K_7=P+B^2.
 \tag{D.2.32i}
\]
In (D.2.32c), $P+B^2$ is precisely the $J$-Moore top:
the two $P$ coordinates cancel. All seven cycles have
zero integral ninth-cell degrees.
The Chang cofiber sequence gives
\[
 \pi_9J_*\cong\ZZ_{(2)}\{\gamma\}\oplus\ZZ/2\{\delta\},
 \qquad h(\gamma)=2U,
\]
where $\delta$ is the image of the Moore eta-coextension
and is detected by the second square from the Moore top.
This is also the explicit generator of
\cite[Proposition~4.1 and Table~3]{I_Li2024}.
Thus $K_7$ projects to exactly $\delta$, while the other
six cycles project to zero. Their projections to $P_*$
are all zero: $\pi_9P_*$ is the bottom eta line, and
none has the corresponding second-square coefficient.
All other source factors were assigned the zero map.
There is no cross Whitehead term; such terms begin in
dimension thirteen for these source factors.

Here is the exact target image of $\delta$, with all
indeterminacies. Desuspend once and use
\[
 D_0=\cofib\left(
 \begin{pmatrix}4&\eta\\0&2\end{pmatrix}:
 S^5\vee S^6\longrightarrow S^5\vee S^6\right)
\]
from D.2.22. Its bottom classes satisfy
$2\rho=-\alpha\eta$; write $n=\alpha\nu$, of order four.
Let $\mu:M(\ZZ/4,6)\to D_0$ extend $\rho$.
The Moore comparison $f:M(\ZZ/2,6)\to M(\ZZ/4,6)$ has
bottom degree two and top degree one
\cite[(3.1)--(3.2)]{I_Li2024}.
The maps $\mu f$ and $-\alpha\bar\eta$ agree on the bottom,
so their difference factors through the top seven-sphere.
The group $\pi_7D_0$ is generated by $b,c$, both of order two.
Compose with $\widetilde\eta_2:S^8\to M(\ZZ/2,6)$.
The actual extension/coextension computation of
\cite[\S2, p.~191]{I_Gray2014} gives the stable bracket
$\bar\eta\widetilde\eta_2\in\{2\nu,6\nu\}$.
Both choices have image $2n$. Therefore
\[
 \mu\widetilde\eta_4=2n
       \pmod{\langle b\eta,c\eta\rangle},
 \qquad \mu\widetilde\eta_4=2n\text{ in }\pi_8L_8.
 \tag{D.2.32j}
\]
Changing the Moore coextension by bottom $\eta^2$ adds
$\rho\eta^2=c\eta$, also killed in $L_8$.
Thus the displayed equality includes every relevant
extension and coextension ambiguity.

\medskip\noindent
\textbf{Every eleven-cell obstruction, including the cone detector.}
The first six cycle attachments are null already in
$\Sigma L_8$, so choose their fillings there. They contribute
no target degree-ten homology. The seventh has composite
$2n$ by (D.2.32j); fill it with two copies of the normalized
target $N'$-cell, which attaches by $n$ in $L_9$.
This reference filling has integral coordinate $2N'$.
The map and the homotopy identifying its attachment with
$2n$ both lie in $\Sigma L_8$, so no additional top
homology coordinate is introduced.

It follows from the complete boundary normalization that
every eleven-cell obstruction has zero integral Hurewicz,
except the $BX$ coordinate, whose Hurewicz is $4N'$.
The only integral-Hurewicz kernel is $\langle s\rangle$,
$s=\rho\nu$, by D.2.22. The actual suspended cone
detector removes this last ambiguity on every cell.

For the chosen lower map the pullbacks, on the original
source, are
\[
 q_7=\alpha\beta,\qquad \ell_9=\alpha\delta,
 \qquad \beta=\operatorname{Sq}^1\alpha,
        \delta=\operatorname{Sq}^2\beta.
\]
The extra $B^2$-dual in the $\mathcal R_8$ pullback has
zero first Bockstein and does not alter $\ell_9$.
With $\gamma=\operatorname{Sq}^2\alpha$ and
$X^*=\operatorname{Sq}^3\beta$, Cartan gives
\[
 \operatorname{Sq}^4q_7=
       \gamma\delta+\beta X^*
       =\operatorname{Sq}^2\ell_9.
 \tag{D.2.32k}
\]
In particular the cross term
$\operatorname{Sq}^1\alpha\operatorname{Sq}^1\delta$
in the second expression has been retained.
Although a degree-seven quotient class can have a relative
ambiguity, $H^{11}(R_{11})\to H^{11}(K_{11})$ is an
isomorphism. Thus (D.2.32k) also determines the top value
on each actual quotient-cell attachment.

For an attaching composite $\omega:S^{10}\to W$, the
classes $\vartheta_7,\lambda_9$ extend uniquely over its
actual cone. The top value of
$\operatorname{Sq}^4\vartheta_7+\operatorname{Sq}^2\lambda_9$
is zero on the normalized $p$ and one on $s$.
There is an actual map from the corresponding source
one-cell attachment to this cone with relative top degree
one. The top source coordinate is nonzero modulo two,
since its integral cellular boundary is even. Equation
(D.2.32k) makes the pulled-back detector zero. Hence
every obstruction has zero detector: all zero-Hurewicz
obstructions vanish, and the $BX$ obstruction is exactly
$p$, rather than $p+s$.
The normalization here is explicitly
$h_{\ZZ}(p)=4N'$ and $\Psi(p)=0$; these conditions remove
its possible $s$ ambiguity.

Finally, changing the ten-cell filling of $K_i$ by
$z_i\in\pi_{10}W$ changes its normalized relation by
$2z_i$, and changes every zero-boundary relation by zero.
This is the actual cell-coaction law; in this range no
Whitehead term involving the added ten-sphere can occur.
D.2.31 proves
\[
                  2\pi_{10}W=0.
 \tag{D.2.32l}
\]
Thus no choice of any ten-cell filling removes the
nonzero $BX$ obstruction. All other eleven-cell
obstructions have already been accounted for. This
proves the claimed nonextension for the specified
coordinated nine-map, with the stated limitations.
\end{proof}

\begin{claim}{D.2.33}{Proved}{Every once-suspended map has the same remaining Toda test}
Keep the spaces and the single normalization of
\textup{D.2.27--D.2.28}, and let
$\lambda\in\pi_8W\cong\ZZ/2$ be the class whose mod-two
Hurewicz image is $R_8$. The Toda bracket
\[
 \Delta=\langle\lambda,\eta_8,2\rangle
       \subseteq\pi_{10}W
 \tag{D.2.33a}
\]
is a singleton, and its value is either $0$ or the element
$p$ of \textup{D.2.32}. If an ordinary map
$h:\Sigma A\to Y$ has the nonzero coefficient
$\epsilon=1$ of \textup{D.2.28}, then
\[
                         \Delta=p.
 \tag{D.2.33b}
\]
More precisely, its quadratic stable summand has a definite
Moore-top choice $\tau\in\FF_2$, relative to the reference
Moore map fixed in \textup{D.2.32}; its eleven-cell
obstruction, modulo the integral-Hurewicz kernel
$\langle s\rangle=\langle\rho\nu\rangle$, is
\[
                    p+\tau\Delta.
 \tag{D.2.33c}
\]
Thus $\epsilon=1$ requires both $\Delta=p$ and $\tau=1$.
In particular, $\Delta=0$ would imply that the normalized
functional is zero for every once-suspended attachment.
The value of this bracket is determined in \textup{D.2.34} below;
the present necessary condition does not construct a map on the full
source or prove a statement about subsequent nullification stages.
\end{claim}
\begin{proof}
Write $X=\Sigma B$, $A=QX$, and
$Z=\fib(Y\to K(\FF_2,4))$, so $W=\Omega Z$.
All cell calculations are two-local. The homotopy groups used
are the actual finite two-primary groups already computed.

\medskip\noindent
\textbf{Exactly the stabilization needed for an obstruction.}
The space $Z$ is five-connected, with
$\pi_6Z=\ZZ/2\{\chi_Z\}$. The double adjunct of its
Whitehead square is the Browder square in \textup{D.2.23},
so $[\chi_Z,\chi_Z]=0$ in $\pi_{11}Z$.
Use the reduced-product model $JZ\simeq\Omega\Sigma Z$
\cite[Section~1 and Theorem~5.6]{I_James1955}.
The pair $(JZ,Z)$ is eleven-connected. Its first relative
group is
\[
 \pi_{12}(JZ,Z)\cong H_{12}(JZ,Z)
       \cong H_6Z\otimes H_6Z\cong\ZZ/2.
 \tag{D.2.33d}
\]
This follows directly from relative Hurewicz and the
word-length filtration: two letters first occur in degree
twelve, and three letters cannot affect that group.
Naturality from $S^6\to Z$ identifies its boundary with
$[\chi_Z,\chi_Z]$, since the twelve-cell of $J_2S^6$
attaches by the ordinary Whitehead square. The exact sequence
therefore makes $\pi_{11}Z\to\pi_{12}\Sigma Z$ an
isomorphism. The latter group is already stable because
$\Sigma Z$ is six-connected. Hence
\[
 \pi_{11}Z\xrightarrow{\ \cong\ }
       \pi_{11}\Sigma^\infty Z.
 \tag{D.2.33e}
\]
In particular $W\to\Omega QZ$ is ten-connected and
injective on $\pi_{10}$. This is a statement about the
specified obstruction group, not an equivalence of
unrestricted mapping sets.

Suppose now that a full map $g:A\to W$ has
$g^*\vartheta_7=\alpha_3\beta_4$. Stabilize its adjunct
$\Sigma A\to Z$ and use the natural Snaith splitting
\cite[Theorem~1.1 and Section~2]{Snaith1974}. Its quadratic
summand gives a map
\[
       D_2X=(E\Sigma_2)_+\wedge_{\Sigma_2}X^{\wedge2}
                       \longrightarrow\Omega QZ.
 \tag{D.2.33f}
\]
The actual space $D_2X$ is five-connected. The ten-connectivity
just proved lifts its ten-skeleton to $W$. Each eleven-cell
obstruction has zero image in $\pi_{10}\Omega QZ$, since
\textup{(D.2.33f)} is defined on the full summand.
Injectivity in \textup{(D.2.33e)} makes every actual obstruction
zero. We have therefore obtained an actual map on the eleven-stage
of $D_2X$ preserving the quadratic $\vartheta_7$ functional.
Here $\vartheta_7$ is the suspension of the ordinary class
$\vartheta_8$ on $Z$, so its pairing is preserved by this
stabilization and adjunction. No lifting of an arbitrary
twelve-dimensional stable map is being assumed.

This also separates all the unit lower data from the present
test. In the stable Snaith grading the classes
$AX,BD,BE,BX$ have weight two. The classes on the unit
$X$, including possible nonzero data on its five- and
six-cells, have weight one. They cannot contribute to the
restriction \textup{(D.2.33f)}.

\medskip\noindent
\textbf{The complete quadratic nine-stage.}
Use a finite two-local cellular model of this summand through
dimension eleven, with odd-unit pairs cancelled. Its attaching
groups through dimension ten are in the ordinary stable range,
since its bottom dimension is six. Abbreviate the unit homology
labels as in \textup{D.2.32}, and put
$R=Q^4A$, $P=AC$, $N=Q^5A$, $J=AB+R$.
The weight-two part of the previously computed eight-stage is
\[
 C\vee M(\ZZ/4,7)_J\vee S^8_s,
 \qquad
 C=C_1^8=\cofib((2,\eta_6):S^6\vee S^7\to S^6),
 \qquad s=N-P.
 \tag{D.2.33g}
\]
Its top Moore cell is $B^2-P$, while $P$ is the top cell
of $C$. Write
$\pi_8C=\ZZ_{(2)}\{\gamma\}\oplus\ZZ/2\{\delta\}$,
with quotient coordinates $\gamma\mapsto(0,2)$ and
$\delta\mapsto(\eta_7,0)$ as in \textup{D.2.18}.
The five nine-attachments are exactly
\[
\begin{array}{c|ccccc}
 &AD&AE&BC&Q^6A&Q^5B\\\hline
 &\tau_J+\delta&0&\gamma+\tau_J+\delta&
                    \gamma+2s&\delta.
\end{array}
 \tag{D.2.33h}
\]
Integral degrees and the second squares detect every component.
The even-$\nu$ components retained in the full-source table
belong to its unit summand and do not occur here. Cell slides
give the four factors
\[
\begin{split}
 J_*&=M(\ZZ/4,7)_J\cup_{i\eta}e^9_{AD-Q^5B},\\
 C_*&=C\cup_\delta e^9_{Q^5B}
                    \cup_\gamma e^9_{BC-AD},\\
 M(\ZZ/2,8)_s,&\qquad S^9_{AE}.
\end{split}
 \tag{D.2.33i}
\]
All these comparisons are ordinary in this connectivity and
dimension range.

The four integral degree-ten cycle coordinates in this
summand reduce to
\[
             C^2,\quad Q^7A,\quad AF,\quad
                    K=AX+BD+BE.
 \tag{D.2.33j}
\]
The complete Bockstein table in \textup{D.2.32}, restricted
to weight two, gives
$H_{10}(D_2X;\ZZ_{(2)})\cong(\ZZ/2)^4$.
Indeed the mod-two dimensions in degrees nine and ten are
five and seven, the integral eighth group has two cyclic
summands, and the degree-eleven first Bockstein has the four
independent images in \textup{(D.2.33j)}.
For the eighth group, the established weight-two seventh
group is $\ZZ/4$, the mod-two eighth dimension is three,
and the two independent first-Bockstein images from degree
nine fix both exponents as two.
The integral eleven-boundary lattice is twice the cycle lattice.
The exact Pontryagin chain calculation retains
\[
                 d(BX)=2K.
 \tag{D.2.33k}
\]
The cycle $K$ has zero integral ninth-cell degrees and
\[
        \operatorname{Sq}_*^2K=P+B^2,
                  \qquad \operatorname{Sq}_*^4K=0.
 \tag{D.2.33l}
\]

Its $J_*$ projection is precisely the Moore eta-coextension,
by the complete Chang calculation of \textup{D.2.32}.
Its projections to the other three factors vanish.
For $M(\ZZ/2,8)_s$ this follows from the zero corresponding
second-square coefficient, and for the nine-sphere it follows
from the zero integral degree. Here is the full check for $C_*$,
so the nonzero bottom-six branch is not omitted.
The first three stable stems and the Chang cofiber sequence give
\[
 \pi_9C=(\ZZ/2)\{i\nu_6,\delta\eta_8\},
 \qquad \gamma\eta_8=0.
\]
The last equality is the cone-of-$\eta$ fourth-square argument
already used in \textup{D.2.19}; the only possible residual
bottom $\nu$ has nonzero fourth-square detector.
After the two nine-cells the exact sequence yields
\[
 \pi_9C_*\cong\ZZ_{(2)}\{\kappa\}\oplus\ZZ/2\{i\nu_6\},
               \qquad h_{\ZZ}(\kappa)=2[Q^5B].
 \tag{D.2.33m}
\]
The integral degree of $K$ excludes $\kappa$, and its fourth
square excludes $i\nu_6$. Thus there is no undetected component.

\medskip\noindent
\textbf{The nonzero bottom-six branch is excluded.}
The class $A^2$ is the first Bockstein of $Q^4A$.
Every first Bockstein from $H_7W$ is zero, by the complete
low target table. Thus the bottom six-sphere of the quadratic
summand has zero mod-two Hurewicz after mapping to $W$.
Its value is either $0$ or $2\alpha$ in
$\pi_6W=\ZZ/4\{\alpha\}$.

The other ten-cycle $C^2$ has zero integral ninth-cell degrees and
\[
 \operatorname{Sq}_*^2(C^2)=0,
       \qquad \operatorname{Sq}_*^4(C^2)=A^2.
 \tag{D.2.33p1}
\]
Cartan gives the first equality from $\beta C=0$, and the
second from $\operatorname{Sq}_*^2C=A$.
The complete groups of the four factors in
\textup{(D.2.33i)} now identify its attachment as precisely
the bottom $i\nu_6$ in $C_*$.
There is no hidden eta-squared component on a seven-sphere:
the $J_*$ Chang group has only its free generator and the
Moore coextension detected by the second square, and the
other two factors have only the stated eta and degree
coordinates. Equation \textup{(D.2.33m)} accounts for the
entire $C_*$ component.

Suppose the bottom maps to $2\alpha$. Represent the
nine-stage map in $\Sigma L_8$ as justified in the next
paragraph. The $C^2$ attachment then maps to exactly
$2n=2\alpha\nu$ there. There is also an exact integral
relation, not merely its first Bockstein:
\[
                         d(Q^6C)=\pm2C^2.
 \tag{D.2.33p2}
\]
Choose an integral cycle representing the odd-dimensional
class $C_5$. In the arity-two chain model the ordered
configuration one-cell satisfies
$d(e_1(C\otimes C))=-2e_0(C\otimes C)$, since the
transposition of two degree-five labels has Koszul sign
$-1$. Its mod-two classes are exactly $Q^6C$ and $C^2$.
A parity-preserving cellular normalization retains this
relation. Filling the ten-cell by $2N'$ makes the
corresponding eleven-obstruction have integral Hurewicz
$4N'\ne0$. Any other filling changes the obstruction by
twice an element of $\pi_{10}W$, hence by zero.
The $J$-top toggle has zero coefficient on this cycle,
as its second square is zero. This contradicts the
actual eleven-stage extension obtained from
\textup{(D.2.33f)}. The bottom must therefore map to zero.
It may be made constant by a homotopy and collapsed.
This rules out the additional lower-six branch rather
than identifying it with the reference map.

\medskip\noindent
\textbf{Non-J data do not change the BX p contribution.}
Any map from the five-connected nine-stage in
\textup{(D.2.33i)} to $W$ lifts to its five-connected cover
$W'$, then through the eleven-connected evaluation
$\Sigma E\to W'$. The actual comparison
$\Sigma L_8\to\Sigma E$ is nine-connected in the low
cellular realization range. Thus the nine-stage map can be
represented in $\Sigma L_8$. The image of
$\pi_{10}\Sigma L_8$ in $\pi_{10}W$ is contained in
$\langle s\rangle$, by the complete torsion calculation
\textup{(D.2.22t7)--(D.2.22t16)}. In particular it has no
component on $p$, whose integral Hurewicz is $4N'$.
This factorization also justifies the reference filling
used for \textup{(D.2.33p2)}.

After the bottom six-sphere is collapsed, $C_*$ becomes
the wedge of $S^7_R\cup_\eta e^9$ and $M(\ZZ/2,8)_P$.
Its first factor has only a free degree-two generator in
the relevant attaching group, and its second has only the
bottom eta. Both $K$ coordinates are zero, by its degree
and the cancellation in $P+B^2$. This is the quadratic
part of the complete zero-$K_6$ component check in
\textup{D.2.32}; the explicit calculation
\textup{(D.2.33m)} also proves it before collapsing.
Every other lower-factor contribution therefore remains
inside $\Sigma L_8$ and cannot change the $BX$ coefficient
on $p$. Possible $s$-valued obstructions are not discarded.

The inherited functional has value one on $J=AB+Q^4A$,
and zero on $Q^4A$. Therefore the $J$ bottom maps to an
odd multiple of $\rho$. Fix the corresponding odd multiple
of the reference Moore map in \textup{D.2.32}.
All maps with that bottom differ on the Moore top by
$\pi_8W=\ZZ/2\{\lambda\}$. Denote this choice by $\tau$.
Choices on its nine-cell cannot change the $K$ composite,
because the integral degree on that cell is zero.
For the reference top, the Moore/Toda computation of
\textup{(D.2.32j)} sends the eta-coextension to $2n$ in
$\pi_9\Sigma L_8$; an odd multiple gives the same value.
Fill it with $2N'$. All other components of $K$ are null
already in that lower target by \textup{(D.2.33l)--(D.2.33m)}.
The relation \textup{(D.2.33k)} consequently has integral
Hurewicz $4N'$, hence obstruction $p$ modulo $\langle s\rangle$.
Changes of ten-cell fillings add twice an element of
$\pi_{10}W$, which is zero by \textup{D.2.31}.

\medskip\noindent
\textbf{The only change in that direction is one exact Toda value.}
There is an actual quotient $q_M$ from the nine-stage onto
$S^8$, projecting the $J$-Moore top and killing all other
factors. Its nine-cell attaches into the bottom seven-sphere,
so the quotient is well-defined. The other top choice is the
ordinary loop-space sum with $\lambda q_M$.

Put $C\eta_8=S^8\cup_{\eta_8}e^{10}$. Its exact sequence gives
\[
 \pi_9C\eta_8=0,\qquad
 \pi_{10}C\eta_8=\ZZ_{(2)}\{\Gamma\},
                \qquad h_{\ZZ}(\Gamma)=2[e^{10}].
 \tag{D.2.33n}
\]
The pinch is injective with image $2\ZZ_{(2)}$:
multiplication by $\eta$ maps the first stable stem onto
the second, while the following map reduces degree modulo two.
Compose $q_M$ with $S^8\to C\eta_8$ and extend over all
ten-cells. The second-square table gives odd relative top
degree only on $K$, and even degree on the other three cycle
coordinates. Changes by $\Gamma$ adjust these degrees by two,
so they may be chosen as $1$ and $0$, respectively.
The $BX$ attachment then maps to exactly $\Gamma$;
every zero-degree relation maps to zero by
\textup{(D.2.33n)}. This checks the whole relevant coupling,
without assuming a four-cell subcomplex of the source.

Since $\lambda\eta_8=0$ by adjunction from
\textup{(D.2.20h2)}, extend $\lambda$ to $C\eta_8$.
Its value on $\Gamma$ is exactly the bracket in
\textup{(D.2.33a)}. Its full indeterminacy is
\[
 \lambda\circ\pi_{10}S^8+2\pi_{10}W
       =\langle\lambda\eta_8^2\rangle+2\pi_{10}W=0.
 \tag{D.2.33o}
\]
This proves the asserted singleton and the change
$p\mapsto p+\tau\Delta$.

To determine its possible values, lift the extension
$C\eta_8\to W$ through $\Sigma E$ in the same ordinary
range as above. The degree-two pinch of $\Gamma$ makes
the corresponding integral Hurewicz in $H_9E$ twice a
class. Since $\pi_{10}W$ has exponent two, this is either
$0$ or $4N'$: the other summands of
$H_9E=\ZZ/8\oplus(\ZZ/2)^3$ have zero double.
The actual cone comparison
$C\eta_8\cup_\Gamma e^{11}\to C_\Delta$ has relative
top degree one. Its source has no cohomology in degrees
seven or nine, and its top mod-two class is nonzero.
Consequently the suspended detector
$\operatorname{Sq}^4\vartheta_7+
\operatorname{Sq}^2\lambda_9$ has value zero on this cone.
The exact kernel and detector of \textup{D.2.22--D.2.23}
therefore give $\Delta\in\{0,p\}$, with no $s$ addition.

An actual extension on the quadratic eleven-stage must
have zero $BX$ obstruction. The preceding computation
forces $p+\tau\Delta=0$ modulo $\langle s\rangle$.
Since $p$ is nonzero in that quotient, this requires
$\Delta=p$ and $\tau=1$. Applied to the actual
eleven-stage map obtained from \textup{(D.2.33f)}, it proves
the necessary condition for every full ordinary map,
with exactly the scope stated.
\end{proof}

\begin{claim}{D.2.34}{Proved}{The remaining Toda value vanishes}
Keep the spaces and the normalization of \textup{D.2.33}. Then
\[
 \Delta=\langle\lambda,\eta_8,2\rangle
             =\{0\}\subseteq\pi_{10}W.
 \tag{D.2.34a}
\]
Consequently every ordinary pointed map
$h:\Sigma A\to Y$ has $\epsilon=0$ in the normalized
once-suspended functional of \textup{D.2.28}.
The simultaneous first-cone vanishing of \textup{D.2.27}
therefore holds for every mixture of suspension exponents
$0,1,2,3$ and at least $5$. Exponent $4$ and propagation
through later cones are not settled by this theorem.
\end{claim}
\begin{proof}
Write $\rho_T:S^6\to T$ for the actual sphere in
\textup{D.2.17}, and $\rho_Y:S^8\to Y$ for its double
adjoint. With the homology labels used there,
\[
 h_2(\rho_T)=b^2+\{a,b\}+f+ac+a^3,
       \qquad h_2(\rho_Y)=\sigma^2f.
 \tag{D.2.34b}
\]
We first identify an actual representative of the shifted
generator $\lambda_Y\in\pi_9Y\cong\ZZ/2$.
Put
\[
 h_0=\eta_2\wedge\id_B:\Sigma Y\longrightarrow Y,
       \qquad \varphi_0:Y\longrightarrow\Omega Y
 \tag{D.2.34c}
\]
for this map and its once adjoint. The quadratic calculation
of \textup{D.2.13} applies in degree eight and gives
\[
 \varphi_{0*}(\sigma^2f)
       =(\sigma a)(\sigma c)+(\sigma c)(\sigma a).
 \tag{D.2.34d}
\]
Here is the complete degree check. In the cone
$\mathbb CP^2\wedge B$, the product of the two old
classes $tu$ and $tu^2$ is $t^2u^3$. Thus the reduced
diagonal of the corresponding new class has precisely
the two ordered terms giving \textup{(D.2.34d)}.
Equivalently, $\overline\Delta f=a\otimes c+c\otimes a$.
The linear component vanishes because $h_0$ induces zero
on reduced homology. In the James tensor algebra
$H_*(\Omega\Sigma^2B;\FF_2)$, the bottom label
$\sigma a$ has degree three, so tensor length three
starts in degree nine and cannot contribute here.
The relative Adams--Hilton formula for this quadratic
coefficient is the one already audited in
\cite[\S3.1, p.~300; proof of Proposition~7.4,
p.~309]{FGHHalperinLemaire1995}; the tensor-algebra
identification is \cite[\S3; Theorem~5.6 and
Corollary~5.7]{I_James1955}.

The right side of \textup{(D.2.34d)} is the nonzero
homology suspension of $R=\{a,c\}$, which detects the
generator in \textup{D.2.19}. Applying this calculation
to \textup{(D.2.34b)} proves the equality of actual
homotopy classes
\[
                 \lambda_Y=h_0\circ\Sigma\rho_Y.
 \tag{D.2.34e}
\]
No formula on higher nonprimitive labels is used.

The relation $\rho_T\eta_6=0$ is part of
\textup{(D.2.19a)}. Choose an actual extension
\[
 \overline\rho:C\eta_6=S^6\cup_{\eta_6}e^8
                      \longrightarrow T.
\]
The ordinary cofiber exact sequence gives
\[
 \pi_8(C\eta_6)=\ZZ_{(2)}\{\Gamma\},
       \qquad h_{\ZZ}(\Gamma)=2[e^8].
 \tag{D.2.34f}
\]
Indeed the preceding multiplication by $\eta$ from
$\pi_1^S$ to $\pi_2^S$ is an isomorphism, and the
image in the integral top-cell group consists exactly
of the even integers. These are ordinary groups in
the stable range of the five-connected space $C\eta_6$;
the sphere identities are those in
\cite[\S2, (2.1), p.~149]{I_Oka1986}.
The class $\Gamma$ is a coextension of degree two
across $\eta_6$, and its mod-two Hurewicz image is zero.
Consequently
\[
 h_2(\overline\rho\Gamma)=0.
\]
The decisive input is the injectivity of the actual
ordinary Hurewicz map
\[
               \pi_8T\longrightarrow H_8(T;\FF_2)
\]
proved in \textup{D.2.20}. It gives
$\overline\rho\Gamma=0$, an actual zero representative
of $\langle\rho_T,\eta_6,2\rangle$.

Double suspension followed by the double-loop counit
transports this extension, this coextension, and this
nullhomotopy to a zero representative of
$\langle\rho_Y,\eta_8,2\rangle$ in $\pi_{10}Y$.
Suspending once more and postcomposing with $h_0$,
the extension/coextension definition of the Toda bracket
and \textup{(D.2.34e)} give
\[
            0\in\langle\lambda_Y,\eta_9,2\rangle
                    \subseteq\pi_{11}Y.
 \tag{D.2.34g}
\]
This uses a specified zero representative, so no
surjectivity assertion about suspension of bracket sets
is required. All the cone domains here have bottom cell
above degree four; their maps and homotopies lift through
$Z=\fib(Y\to K(\FF_2,4))$ with the needed uniqueness.
Adjunction therefore gives a zero representative of the
bracket in \textup{(D.2.34a)} in $W=\Omega Z$.
Its full indeterminacy is zero by
\textup{(D.2.33o)}, using \textup{D.2.31}; hence that
bracket is the singleton $\{0\}$.

Finally \textup{D.2.33} shows that any ordinary map with
$\epsilon=1$ would force $\Delta=p$, where $p\ne0$.
Equation \textup{(D.2.34a)} excludes this possibility
for every such map, proving the stated uniform conclusion.
The normalization is the same single target class used in
\textup{D.2.27--D.2.28}. The common-factor criterion of
\textup{D.2.26} therefore combines this vanishing with all
the exponents already treated in \textup{D.2.27}, including
arbitrary simultaneous mixtures of them.
\end{proof}

\begin{claim}{D.2.35}{Proved}{Finite Postnikov data cannot give a universal cone-preservation invariant}
Put
\[
 B=K(\FF_2,2),\qquad A=Q\Sigma B,\qquad Y=\Sigma^2B.
\]
Every finite Postnikov section of $Y$ is ordinarily $A$-acyclic:
\[
                    P_A(P_NY)\simeq *\qquad (N<\infty).
                    \tag{D.2.35a}
\]
Consequently there is no class $\mathcal C$ of pointed CW spaces
with all of the following properties, for any fixed finite $N$:
\begin{enumerate}
\item Membership in $\mathcal C$ depends only on the $N$-type;
      $Y\in\mathcal C$ and $*\notin\mathcal C$.
\item If $C\in\mathcal C$ and
      $h:\bigvee_i\Sigma^{k_i}A\to C$, with a small indexing
      set and $k_i\geq0$, then $\cofib(h)\in\mathcal C$.
\item The class is closed under transfinite CW colimits.
\end{enumerate}
The same assertion holds for spaces with compatible marked
classes, if the marks are transported by $N$-equivalences and
by the specified successor and limit constructions.

In particular, $N\geq11$ already preserves the full mod-two
Hopf algebra and Steenrod data of $\Omega Y$ through degree ten.
A criterion involving only those data, including the primitive
operation
\[
 \operatorname{Sq}_*^1\operatorname{Sq}_*^2:
       \operatorname{Prim}H_{10}(\Omega C;\FF_2)
                    \longrightarrow H_7(\Omega C;\FF_2),
                    \tag{D.2.35b}
\]
cannot by itself supply the three properties above. This is
not a claim that $P_NY$ occurs in the particular nullification
sequence beginning at $Y$.
\end{claim}
\begin{proof}
The space $A$ is two-connected and $H_3(A;\ZZ)=\FF_2$.
Bousfield's symmetric-product comparison says that $A$ kills
$SP^\infty A$ \cite[Corollary~3.2(ii), p.~39]{Bousfield1996ULP}.
The bottom factor $K(\FF_2,3)$ is a homotopy retract of this
generalized Eilenberg--MacLane space. Thus
\[
                         P_AK(\FF_2,3)\simeq *.
                         \tag{D.2.35c}
\]
Here the bottom homology calculation follows from the
stabilization unit and the Hurewicz theorem; the factorization
of the infinite symmetric product is the usual Dold--Thom
generalized Eilenberg--MacLane decomposition.

The integral reduced homology of $Y$ is two-primary, and $Y$
is simply connected and of finite type. The Serre-class theorem
therefore makes every $\pi_iY$ a finite two-primary group
\cite[Chapter~III, Theorem~1 and Corollaries~1--2,
pp.~272--274]{I_SerreClasses1953}. In particular, $P_NY$ is a
three-connected Postnikov space with two-primary homotopy
groups. It is a generalized polyGEM by
\cite[Proposition~8.5, p.~1340]{EH5Bousfield1997}.
Apply \cite[Lemma~8.13, p.~1341]{EH5Bousfield1997} with
$G=\FF_2$, $p=2$, and $n=3$. Its extra condition applies
only when $G/pG=0$, so it imposes no condition here. The lemma
gives
\[
                    P_{K(\FF_2,3)}(P_NY)\simeq *.
                    \tag{D.2.35d}
\]
Transitivity of ordinary acyclicity, together with
\textup{(D.2.35c)}, proves \textup{(D.2.35a)}. Equivalently,
every $A$-null target is $K(\FF_2,3)$-null, and hence receives
only a contractible pointed mapping space from $P_NY$.

For completeness, the closure assertion uses actual cone
attachments, rather than an inverse-limit assertion about
nullification. Ordinary $A$-nullification can be constructed
by iteratively attaching the mapping cones of all maps from
$A$ and its nonnegative suspensions, continuing to a
sufficiently long transfinite CW colimit
\cite[Theorem~1.1 and the following construction, p.~34]
     {Bousfield1996ULP}.
At each stage the maps form a set. One may use a regular
cardinal bounding a CW model of $A$ so that maps from every
$\Sigma^kA$ into the final colimit factor through an earlier
stage up to homotopy. The final space is $A$-null, and every
cone attachment is an $A$-equivalence, which verifies the
nullification universal property.

If $\mathcal C$ had the three stated properties, the equality
of $N$-types would put $P_NY$ in $\mathcal C$. Apply this cone
construction to $P_NY$. Successor closure and transfinite
closure would put $P_A(P_NY)$ in $\mathcal C$, contradicting
\textup{(D.2.35a)} and $*\notin\mathcal C$. Compatible marked
classes give the same contradiction when the property requires
their nonvanishing.

Finally, $Y\to P_NY$ is $(N+1)$-connected, and its loop map
is $N$-connected. It consequently induces homology
isomorphisms in every degree strictly below $N$. For $N\geq11$
these identify all products, coproducts, primitives, and
Steenrod operations occurring in \textup{(D.2.35b)}, as well
as the old degree-three and degree-four loop classes and their
degree-seven commutator. Thus a preservation condition based
only on this fixed data has exactly the same initial value
on $Y$ and $P_NY$. Additional global information, or a proof
restricted to the actual intermediate targets reached from
$Y$, is necessary.
\end{proof}

\begin{claim}{D.2.36}{Proved}{Every entire fourth-suspension unit has zero linear image}
Put $B=K(\FF_2,2)$, $X=\Sigma B$, $A=QX$, and
$Y=\Sigma^2B$. Every ordinary pointed map
\[
             h:\Sigma^5B\longrightarrow\Sigma^2B
\]
induces the zero map on degree-eight mod-two homology.
The same is therefore true of every ordinary pointed map
$h:\Sigma^4A\to Y$.

For the double adjunct $G:\Sigma^2A\to T=\Omega^2Y$,
write $B_0=\sigma b\in H_4A$, where $b$ is dual to
$v=\operatorname{Sq}^1u\in H^3B$. Then
\[
                   G_*(\sigma^2B_0)=\epsilon k,
                   \qquad \epsilon\in\FF_2,
                   \quad k=\{a,b\}.
                   \tag{D.2.36a}
\]
In particular, the alternative bottom pattern
$b^2+k+p_3$ cannot occur on an entire unit. The pure
commutator pattern is realized on the entire unit by D.2.29;
its extension over the full infinite-loop source is not
settled by this statement.
\end{claim}
\begin{proof}
All coefficients are $\FF_2$. Remove the suspensions from
the induced cohomology homomorphism to obtain a linear map
\[
 \phi=\sigma^{-5}h^*\sigma^2:
 \widetilde H^dB\longrightarrow\widetilde H^{d-3}B.
                   \tag{D.2.36b}
\]
It commutes with every Steenrod square, by naturality and
suspension compatibility. It is not assumed multiplicative.
Serre's theorem gives
\[
 H^*B=\FF_2[u_2,v_3,w_5,t_9,s_{17},\ldots],\qquad
 v=\operatorname{Sq}^1u,\quad w=\operatorname{Sq}^2v,
 \quad t=\operatorname{Sq}^4w,\quad s=\operatorname{Sq}^8t;
                   \tag{D.2.36c}
\]
the later generator degrees are $2^j+1$. The exact source
and operation conventions are
\cite[\S1, nos.~2.1--2.4; \S2, no.~8, Theorem~2,
and no.~9, pp.~199, 203--204]{I_Serre1953}.

Suppose $\phi(u^3)=v$. The following five identities follow
directly from Cartan and instability:
\[
\begin{aligned}
 \operatorname{Sq}^4(u^2)&=u^4,&
 \operatorname{Sq}^3(uv)&=uv^2+vw,\\
 \operatorname{Sq}^2(u^3)&=u^4+uv^2,&
 \operatorname{Sq}^6(v^2)&=v^4,\\
 \operatorname{Sq}^4(vw)&=vt+v^4.
\end{aligned}
                   \tag{D.2.36d}
\]
Since $\phi(u^2)\in H^1B=0$, the first gives
$\phi(u^4)=0$. The third now gives $\phi(uv^2)=w$.
The second gives $\phi(vw)=w$, since
$\operatorname{Sq}^3\phi(uv)=0$ by instability in degree two.
Likewise $\operatorname{Sq}^6\phi(v^2)=0$ in degree three,
so the last two identities give
\[
                       \phi(v^4)=0,\qquad \phi(vt)=t.
                       \tag{D.2.36e}
\]

Four more explicit source identities finish the contradiction:
\[
\begin{aligned}
 \operatorname{Sq}^{10}(uvw)&=u^2v^2w^2,&
 \operatorname{Sq}^9(ut)&=ut^2+vs,\\
 \operatorname{Sq}^8(vt)&=vs,&
 \operatorname{Sq}^4(uwt)&=ut^2+u^2v^2w^2.
\end{aligned}
                       \tag{D.2.36f}
\]
For clarity, the generator formulas used in Cartan here are
\[
 \operatorname{Sq}^1t=w^2,\quad
 \operatorname{Sq}^8t=s,\quad
 \operatorname{Sq}^9t=t^2,\quad
 \operatorname{Sq}^jt=0\quad(2\leq j\leq7).
                       \tag{D.2.36g}
\]
The first and last follow by Adem applied to $t=\operatorname{Sq}^4w$,
using $\operatorname{Sq}^1w=v^2$ and
$\operatorname{Sq}^2w=\operatorname{Sq}^3w=0$;
the middle two are the generator definition and the top-square
identity. Thus every term in \textup{(D.2.36f)} is specified.

The first identity of \textup{(D.2.36f)} has zero image under
$\phi$, since $\phi(uvw)$ has degree seven. The second has
zero image as a ninth square of a degree-eight class.
The third, together with \textup{(D.2.36e)}, gives
\[
 \phi(u^2v^2w^2)=0,\qquad
 \phi(ut^2)=\phi(vs)=s.
                       \tag{D.2.36h}
\]
The final identity would therefore give
\[
                       \operatorname{Sq}^4\phi(uwt)=s.
                       \tag{D.2.36i}
\]
But $\phi(uwt)$ has degree thirteen, where every class of
$H^*B$ is decomposable: there is no algebra generator in
that degree. Cartan preserves the ideal of decomposables,
whereas $s$ is an indecomposable generator. This contradiction
proves $\phi(u^3)=0$. Also
\[
 \phi(v^2)=\phi(\operatorname{Sq}^2u^2)
          =\operatorname{Sq}^2\phi(u^2)=0.
\]
The degree-six basis is $u^3,v^2$, so duality proves the
degree-eight homology assertion for the unit map.

The unit $X\to QX$ is an isomorphism on degree-four
homology: the first extra stable weight starts in degree six.
After four suspensions it therefore spans all degree-eight
source homology. Restriction to this actual unit proves the
same vanishing for every $\Sigma^4A\to Y$.

Finally, the complete primitive basis of $H_6T$ is
$b^2,k,p_3$, as in D.2.24. Double homology suspension sends
$p_3$ to the linear $f$ label and kills the other two.
The homology vanishing just proved excludes its coefficient.
Since $\operatorname{Sq}_*^2(\sigma^2B_0)=0$, while
$\operatorname{Sq}_*^2b^2=a^2$ and
$\operatorname{Sq}_*^2k=0$, the $b^2$ coefficient is zero
as well. This proves \textup{(D.2.36a)} without imposing
any multiplicativity on the ordinary map.

The exact program \texttt{code/D2\_k4\_linear\_module.py}
independently constructs all source monomials, uses Adem and
Cartan, and allows every coefficient of $\phi$ except the
single prescribed value $\phi(u^3)=v$. Its degree-twenty
inconsistency certificate is precisely the nine identities
above. The range through degree nineteen is consistent as
a truncated module problem; no realization of that truncated
solution is asserted.
\end{proof}

\begin{claim}{D.2.37}{Proved}{Smashing a primary Eilenberg--MacLane generator with $B\mathbb Z/p$ preserves its reversible closure}
Let $p$ be a prime and $n\geq1$, and put
\[
 G=B\mathbb Z/p=K(\FF_p,1),\qquad
 B=K(\FF_p,n),\qquad X=G\wedge B.
\]
Then
\[
 P_X(\Sigma^2B)\simeq *,\qquad
 \min\mathrm{Rev}\{X\to *\}
       =\min\mathrm{Rev}\{B\to *\}.
 \tag{D.2.37a}
\]
In particular, $\RP^2\wedge K(\FF_2,2)$ belongs to the
reversible null class generated by
$B\mathbb Z/2\wedge K(\FF_2,2)$. Tensoring the
$B\mathbb Z/2$, $\RP^2$ example with $K(\FF_2,2)$ therefore
does not provide the proposed strict separation within the
stable part $\mathcal I_2$.
\end{claim}
\begin{proof}
We use pointed derived mapping spaces throughout. Recall that
a space $U$ is ordinarily $V$-acyclic precisely when
$\Map_*(U,Z)\simeq *$ for every $V$-null target $Z$.
Such acyclic spaces are closed under pointed homotopy colimits
and under smashing with an arbitrary pointed space. The latter
assertion follows from
\[
 \Map_*(U\wedge C,Z)
       \simeq\Map_*(C,\Map_*(U,Z)).
 \tag{D.2.37b}
\]
They are also closed under finite products: use the cofibration
$U\vee V\to U\times V\to U\wedge V$ and its mapping-space
fibration. Iterating the bar construction, whose levels are
finite products, consequently proves that every
$K(\FF_p,n)$ with $n\geq1$ is ordinarily $G$-acyclic.

Let $Z$ be $X$-null. Exponential adjunction says that
$\Map_*(B,Z)$ is $G$-null. Since $B$ is $G$-acyclic, it follows
that
\[
 \Map_*(B\wedge B,Z)
    \simeq\Map_*(B,\Map_*(B,Z))\simeq *.
 \tag{D.2.37c}
\]
Thus $B\wedge B$ is $X$-acyclic. The same holds for every
$B^{\wedge j}$ with $j\geq2$, by \textup{(D.2.37b)}.
The usual fat-wedge filtration of $B^r/\bigvee_{i=1}^r B$
has successive quotients wedges of these smash powers. It
therefore gives
\[
          B^r/\bigvee_{i=1}^rB\quad\hbox{is $X$-acyclic}
          \qquad(r\geq1).
          \tag{D.2.37d}
\]
At a filtration step, both the preceding stage and the
cofiber are acyclic, so the mapping-space fibration proves
acyclicity of the new stage.

The components of the next mapping spaces require care.
Replace $\Omega Z$ by a weakly equivalent simplicial group
$H$. This is still $X$-null. Choose the standard simplicial
$\FF_p$-vector-space model for $B$. Restriction is a map of
simplicial groups
\[
 r_r:\Map_*(B^r,H)\longrightarrow
             \Map_*(B,H)^r.
 \tag{D.2.37e}
\]
Its homotopy fiber at the constant tuple is contractible by
\textup{(D.2.37d)}. Every tuple $(f_1,\ldots,f_r)$ extends to
$B^r$ by the ordered product
\[
       (b_1,\ldots,b_r)\longmapsto
                          f_1(b_1)\cdots f_r(b_r).
 \tag{D.2.37f}
\]
Thus $r_r$ is surjective on components, and translation by
an extension identifies any fiber with the fiber at the
constant tuple. Equivalently, the wedge inclusion is a
cofibration of simplicial sets, so \textup{(D.2.37e)} is a
Kan fibration of simplicial groups; all its fibers are
contractible. Hence \textup{(D.2.37e)} is a weak equivalence
on all components. Merely knowing that its zero fiber is
contractible would not justify this conclusion for an
unlooped target.

Let $\mathcal T$ be the opposite of a skeleton of finite
free $\FF_p$-vector spaces. Its object $\FF_p^r$ is the
product of $r$ copies of $\FF_p$ in $\mathcal T$; its
projections are opposite to the coordinate inclusions
$\FF_p\to\FF_p^r$. Matrices over $\FF_p$ act strictly on
the chosen model for $B$. Consequently
\[
 \mathcal T\longrightarrow\mathrm{sSet},\qquad
        \FF_p^r\longmapsto\Map_*(B^r,H)
 \tag{D.2.37g}
\]
is an actual functor, and its product comparison maps are
exactly \textup{(D.2.37e)}. It is a homotopy algebra over
this algebraic theory, with value $*$ at the zero object.
Badzioch's rigidification theorem replaces it by a strict
$\mathcal T$-algebra through an objectwise weak equivalence
\cite[Definitions~1.1--1.2 and Theorem~1.3, pp.~895--897;
Theorem~6.4 and Lemma~6.5, pp.~910--912]{I_Badzioch2002}.
Strict algebras for this theory are precisely simplicial
$\FF_p$-vector spaces. In particular,
\[
             K=\Map_*(B,\Omega Z)
                 \simeq\prod_{i\geq0}K(V_i,i)
 \tag{D.2.37h}
\]
for some $\FF_p$-vector spaces $V_i$. To see the last
description explicitly, apply Dold--Kan and split the
normalized chain complex into its homology and contractible
summands over the field $\FF_p$. The zero operation in
$\mathcal T$ makes this a pointed equivalence. No connectivity
or finite-dimensionality assumption on $K$ is used.

Exponential adjunction also gives
\[
                \Map_*(G,K)
                   \simeq\Map_*(X,\Omega Z)\simeq *.
 \tag{D.2.37i}
\]
For a trivial $\FF_p[\mathbb Z/p]$-module $V$,
$H^j(G;V)=V$ for every $j\geq1$: the periodic free resolution
over $\FF_p[t]/(t^p)$ has alternating differentials $t$ and
$t^{p-1}$, both of which become zero after applying
$\Hom(-,V)$. In particular, for $i\geq1$,
\[
 \pi_{i-1}\Map_*(G,K(V_i,i))
                  =H^1(G;V_i)=V_i.
 \tag{D.2.37j}
\]
Every factor in \textup{(D.2.37h)} is a pointed retract.
Equation \textup{(D.2.37i)} therefore forces $V_i=0$ for
every $i\geq1$. The possibly nonzero discrete factor $V_0$
disappears on looping. We obtain
\[
          \Map_*(\Sigma^2B,Z)
             \simeq\Omega K\simeq *.
 \tag{D.2.37k}
\]
This proves the ordinary acyclicity asserted in
\textup{(D.2.37a)}.

The reversible closure generated by $X\to *$ contains every
ordinary $X$-acyclic space and is closed under desuspension.
It therefore contains $B\to *$. Conversely, $X=B\wedge G$
belongs to the null class generated by $B$, since reversible
null classes are closed under smashing with pointed spaces.
These two inclusions prove the equality of reversible
closures. The same smash closure then includes
$\RP^2\wedge B$ when $p=2$ and $n=2$.
\end{proof}

\begin{claim}{D.2.38}{Proved}{The exact order of the finite fourth-stage source class}
Let $M=M(\mathbb Z/2,3)$, with bottom and top mod-two homology
classes $A_0,B_0$, and let
\[
 D=D_2^{(4)}M=S^3_+\wedge_{C_2}M^{\wedge2}.
 \tag{D.2.38a}
\]
Here the action on $S^3$ is antipodal and the action on the smash
square interchanges the factors.  The class
$Y_{10}=\lambda_3(A_0,B_0)\in H_{10}(D;\mathbb F_2)$
is spherical, and
\[
 \pi_{10}D\cong(\mathbb Z/4)^2\oplus\mathbb Z/2.
 \tag{D.2.38b}
\]
Every element with mod-two Hurewicz image $Y_{10}$ has exact
order four.  In a stable cell model described below, the last
attaching map has the form
\[
 2y+\zeta+2d n+e t,\qquad d,e\in\mathbb F_2,
 \tag{D.2.38c}
\]
where $y$ is a free degree-ten sphere and
$\pi_{10}R=\mathbb Z/2\{\zeta\}\oplus
\mathbb Z/4\{n\}\oplus\mathbb Z/2\{t\}$ is the group of
the other ten-dimensional summand.  The conclusions do not
require the values of $d,e$.  In particular the last cell cannot
be treated as a pure degree-two Moore attachment.
The sphere $S^{10}\to C_4M$ used below to realize $Y_{10}$
has zero image under $C_4M\to QM$.
\end{claim}
\begin{proof}
We work at the prime two.  The complex $D$ is five-connected and
eleven-dimensional.  Freudenthal therefore identifies its ordinary
$\pi_{10}$ with stable homotopy.  Its suspension spectrum is
two-primary: its integral homology, computed below, is two-primary,
and the finite simply connected complex becomes contractible
after inverting two.  Thus the calculation also determines the
unlocalized group in \textup{(D.2.38b)}.

\medskip\noindent
\textbf{The configuration cells and their operations.}
The filtered stable splitting identifies \textup{(D.2.38a)}
with the weight-two summand of $C_4M$; see
\cite[Theorem~1.1, p.~577, and Theorem~3.2, p.~580]{Snaith1974}.
The homology operations and their top-operation corrections used
here are \cite[Chapter~III, Theorems~1.2--1.4,
pp.~216--219]{I_CohenLadaMay1976}.  All unlabeled homology in
the proof has coefficients $\mathbb F_2$.
The complete table needed in degrees six through eleven is
\[
\begin{array}{c|c|c|c|c}
\text{degree}&\text{class}&\beta&\operatorname{Sq}_*^2&
 \operatorname{Sq}_*^4\\ \hline
6&P=A_0^2&0&0&0\\
7&R=A_0B_0&P&0&0\\
7&S=Q^4A_0&P&0&0\\
8&T=B_0^2&0&P&0\\
8&U=Q^5A_0&0&P&0\\
9&V=Q^5B_0&0&S&0\\
9&W=Q^6A_0&U&0&0\\
10&X=Q^6B_0&V&U&P\\
10&Y=\lambda_3(A_0,B_0)&0&0&0\\
11&Z=Q^7B_0&Y&W&S
\end{array}
\tag{D.2.38d}
\]
For completeness, the integral chain calculation uses the Moore
chains $a_3,b_4$, $\partial b=2a$, and the antipodal resolution
cells $e_i$, $0\leq i\leq3$:
\[
 \partial(e_i xy)=e_{i-1}\bigl(xy+(-1)^i\tau(xy)\bigr)
                  +(-1)^i e_i\partial(xy),
 \tag{D.2.38e}
\]
with the first term omitted for $i=0$ and
$\tau(aa)=-aa$, $\tau(ab)=ba$, $\tau(ba)=ab$, $\tau(bb)=bb$.
Cancelling the three unit chain pairs gives elementary boundaries
$2,4,2,2,2$ in successive dimensions
$7\to6,8\to7,9\to8,10\to9,11\to10$.
Thus $H_i(D;\mathbb Z)$ is respectively
$\mathbb Z/2,\mathbb Z/4,\mathbb Z/2,\mathbb Z/2,
\mathbb Z/2$ for $6\leq i\leq10$, and zero in degree eleven.
The unit cancellations can be made in a CW-spectrum model.
We use that stable model below; chain diagonalization alone is
not being used to assert an unstable wedge decomposition.

\medskip\noindent
\textbf{An actual spherical representative of $Y$.}
The generalized Browder map on the bottom sphere and the unit is
an actual map
\[
 \Sigma^3(S^3\wedge M)=M(\mathbb Z/2,9)\longrightarrow C_4M.
 \tag{D.2.38f}
\]
Its restriction to the bottom $S^9$ has fourfold adjunct the
Whitehead square of the bottom $S^7\to\Sigma^4M$.
The octonionic multiplication makes $S^7$ an $H$-space, so
$[\iota_7,\iota_7]=0$; naturality makes that restriction null.
Collapsing the null bottom gives an actual sphere $S^{10}\to C_4M$
whose Hurewicz image is $Y$.  The generalized product/adjunction
interface is \cite[Definitions~2.2--2.3 and
Theorem~2.4]{EH5Arkowitz1962}, and its homology identification is
the Browder operation cited above.  Projecting stably to $D$
retains this weight-two class.  The stable range for $D$ then
gives a sphere in its ordinary $\pi_{10}$ with the same image.

\medskip\noindent
\textbf{The stable eight- and nine-dimensional models.}
Choose integral lifts of the mod-two coordinates
$z_7=S$, $c_7=R+S$, and $T'_8=T+U$, with orientations
chosen in the chain reduction \textup{(D.2.38e)}.
Their cellular boundaries are
$\partial z=2P$, $\partial T'=4c$, $\partial U=0$.
Table \textup{(D.2.38d)} gives
$\operatorname{Sq}_*^2 U=P$ and
$\operatorname{Sq}_*^2 T'=0$.
The only possible degree-seven attaching torsion is the bottom
$\eta$, detected by this second square.  Therefore
\[
 D_8=C\vee M(\mathbb Z/4,7),\qquad
 C=\operatorname{Cof}\bigl((2,\eta_6):S^6\vee S^7\to S^6\bigr).
 \tag{D.2.38g}
\]
Let $i:S^6\to C$ and $q:C\to S^7\vee S^8$ be its structure maps.
The direct Chang-cofiber calculations of D.2.18--D.2.20,
shifted by two, give
\[
 \pi_8 C=\mathbb Z_{(2)}\{\gamma\}\oplus\mathbb Z/2\{\delta\},
 \qquad q\gamma=(0,2),\quad q\delta=(\eta_7,0),
 \qquad \gamma\eta_8=0.
 \tag{D.2.38h}
\]
These calculations use only the cofiber exact sequence and
the low sphere stems and Moore coextension identities in
\cite[Section~2, equations~(2.1)--(2.2), p.~149]{I_Oka1986}.
The degree-nine cells $W,V$ have boundaries $2U,0$ and
second squares $0,z$, respectively.  Their possible Moore-four
components are detected by the absent $c$ coefficient of the
second square.  Their attachments are consequently exactly
$\gamma,\delta$.  Put
\[
 C_b=C\cup_\gamma e^9_W\cup_\delta e^9_V,
 \qquad D_9=C_b\vee M(\mathbb Z/4,7).
 \tag{D.2.38i}
\]
This notation includes the stable unit cancellations of
\textup{(D.2.38e)}; no cancelled nine-cell is retained.

\medskip\noindent
\textbf{The component groups and the decisive composition.}
The same cofiber sequences give
\[
 \pi_9 C_b=\mathbb Z_{(2)}\{\kappa\}\oplus\mathbb Z/2\{v\},
 \quad h_{\mathbb Z}(\kappa)=2V,
 \quad v=i\nu_6,
 \tag{D.2.38j}
\]
where $\kappa$ is normalized to have fourth-square detector zero.
Its second-square detector is $U$.  This is the degree-two
pinch/Bockstein and Adem calculation used for $\kappa'$ in
D.2.20; adding $v$ changes exactly the fourth-square detector.
One degree higher the exact sequence for $C$ gives
\[
 \pi_{10}C=\mathbb Z/4\{\omega\},\qquad
 q\omega=(2\nu_7,\eta_8^2),\qquad
 \delta\eta_8^2=2\omega.
 \tag{D.2.38k}
\]
Since $\gamma\eta=0$ and $\delta\eta\ne0$, attachment of
the two nine-cells gives
\[
 0\longrightarrow\mathbb Z/2\{\bar\omega\}
 \longrightarrow\pi_{10}C_b
 \longrightarrow\mathbb Z/2\{\bar\zeta\}
 \longrightarrow0.
 \tag{D.2.38l}
\]
The right quotient is $\eta$ on the $W$-cell.  Its extension
need not be split at this point.

The specific composition needed below is
\[
 \kappa\eta_9=\bar\omega.
 \tag{D.2.38m}
\]
To check it, first attach only the $\delta$-cell.
In $C'=C\cup_\delta e^9_V$ the group $\pi_{10}C'$ is
$\mathbb Z/2\{\bar\omega\}$.  A lift $\kappa'$ of degree
two on $V$ has second-square detector $U$.
After collapsing the bottom $S^6$, its $U$-component is $\eta_8$;
its further $\eta$ therefore has the nonzero $\eta_8^2$ component
and is $\bar\omega$.  Attaching $\gamma$ preserves this subgroup
because $\gamma\eta^2=0$.  Normalizing by $v$ has no effect,
since $\nu\eta=0$.  The right quotient in
\textup{(D.2.38l)} is detected by the second square to $W$.

For the other wedge summand the Moore coextension sequence gives
\[
\begin{split}
 \pi_9 M(\mathbb Z/4,7)
   &=\mathbb Z/2\{i_7\eta_7^2\}
      \oplus\mathbb Z/2\{\widetilde\eta\},\\
 \pi_{10} M(\mathbb Z/4,7)
   &=\mathbb Z/4\{n=i_7\nu_7\}
      \oplus\mathbb Z/2\{t=\widetilde\eta\eta_9\}.
\end{split}
\tag{D.2.38n}
\]
Here $\widetilde\eta$ is the order-two Moore-four coextension
already used in D.2.19--D.2.22.  In particular
$v\eta_9=0$ and $i_7\eta_7^3=4i_7\nu_7=0$.

\medskip\noindent
\textbf{The two ten-cells and the spherical safeguard.}
The boundary and square vectors in \textup{(D.2.38d)} imply
\[
 a_X=\kappa+v+\epsilon i_7\eta_7^2,
 \qquad \epsilon\in\mathbb F_2.
 \tag{D.2.38o}
\]
Its Moore-four coextension component would have a nonzero
$T'$ second-square coefficient, which is absent.  The displayed
$i\eta^2$ coefficient is invisible to these primary detectors
and is deliberately retained.  For $a_Y$, zero integral boundary
and zero second/fourth squares leave only a possible $i\eta^2$.

The spherical class constructed in \textup{(D.2.38f)} forces
$a_Y=0$, for the following integral reason.  Its stable sphere
lifts to the ten-skeleton.  If its degrees on the two ten-cells
are $(a,b)$, the cofiber sequence requires
$a a_X+b a_Y=0$ in $\pi_9D_9$.
The free $\kappa$ coefficient of $a_X$ forces $a=0$.
Its Hurewicz image forces $b$ odd, and multiplication by such
$b$ is invertible on the remaining two-primary torsion.
Consequently $a_Y=0$ and there is a stable splitting
\[
 D_{10}=R\vee S^{10}_y,\qquad
 R=D_9\cup_{a_X}e^{10}_X.
 \tag{D.2.38p}
\]
This uses the free $\kappa$ coefficient to rule out a correction
by $X$; mod-two sphericity alone would not suffice.

In the cofiber sequence for $a_X$, the map
$\pi_9S^9\to\pi_9D_9$ is injective because of that same free
coefficient.  There is no right-hand kernel term in the
calculation of $\pi_{10}R$.
Equations \textup{(D.2.38m)} and \textup{(D.2.38n)} give
$a_X\eta_9=\bar\omega$, independently of $\epsilon$.
Killing the left subgroup of \textup{(D.2.38l)} therefore gives
the exact direct sum
\[
 \pi_{10}R=\mathbb Z/2\{\zeta\}\oplus
           \mathbb Z/4\{n\}\oplus\mathbb Z/2\{t\}.
 \tag{D.2.38q}
\]
The class $\zeta$ is the image of a lift of $\bar\zeta$.

\medskip\noindent
\textbf{The final attachment.}
Its integral boundary is $2Y$, so it is $2y+z$ with
$z\in\pi_{10}R$.
The second-square vector is $W$.  Among the torsion summands
in \textup{(D.2.38q)}, precisely $\zeta$ has that detector,
so its coefficient is one.
The fourth-square vector is $z_7=S$, with zero $c_7$ coefficient.
The class $\zeta$ comes from $C_b$, so has zero $c_7$ detector.
The class $t$ factors through $\eta$ on $S^9$; naturality from
its cone $C\eta_9$, which has no degree-seven cohomology,
also gives zero $c_7$ detector.
The $c_7$ fourth-square detector of $n=i_7\nu_7$ is one,
by the Hopf invariant of $\nu$.  Its coefficient must thus
be even.  This proves \textup{(D.2.38c)}.

Eliminating $\zeta$ from the resulting presentation gives
$(\mathbb Z/4)^2\oplus\mathbb Z/2$:
the relation $2\zeta=0$ becomes $4y=0$, independently of
$d,e$.  Moreover $2y$ is nonzero, since modulo
$2\pi_{10}R$ its value has the nonzero $\zeta$ coordinate.
Every representative of $Y$ is an odd multiple of $y$ plus
torsion from \textup{(D.2.38q)}.  Its double retains this
nonzero $\zeta$ coordinate, while its fourth multiple vanishes.
It therefore has exact order four, as asserted.

\medskip\noindent
\textbf{The image of the actual sphere factor in $QM$.}
The composite of \textup{(D.2.38f)} with $C_4M\to QM$ is
null: the generalized Browder product becomes null after
stabilization, equivalently its generalized Whitehead adjunct
has zero stable suspension.  The Moore Puppe sequence therefore
places the image of its sphere factor in $2\pi_{10}QM$.
This subgroup is zero in the present degree, as can be checked
without making a general inference about factors of null maps.
Write $M_2=\operatorname{Cof}(2:S\to S)$, so
$\pi_{10}QM=\pi_7M_2$.
The map $2\id_{M_2}$ restricts to zero on the bottom sphere,
so it factors through $q:M_2\to\Sigma S$.
The Moore exact sequence gives
$\pi_1M_2=\mathbb Z/2\{i\eta\}$; hence this factorization
is either zero or $i\eta q$.
For $x\in\pi_7M_2$, its pinch $qx$ lies in
$\pi_6S=\mathbb Z/2\{\nu^2\}$, and
$\eta\nu^2=(\eta\nu)\nu=0$.
The two sphere facts are
\cite[equation~(2.1)(e'),(g), p.~149]{I_Oka1986}.
Consequently $2x=0$ for every $x$, proving the asserted zero
image of this actual sphere $S^{10}\to C_4M\to QM$.
No ordinary map $D\to QM$ is inferred from the stable weight
projection in this argument.

This source calculation leaves the lower attaching coefficients
in \textup{(D.2.38o)} and \textup{(D.2.38c)} uncomputed.
It does not assert that the paired target functional vanishes
under every map $D\to V$, nor that a full fourth-suspension
extension exists or is impossible.
\end{proof}

\begin{claim}{D.2.39}{Proved}{Classifying-space smash preserves every reversible class below its own}
Let $p$ be a prime, let $G=B\mathbb Z/p$, and let $C$ be a
connected pointed CW space which is ordinarily $G$-acyclic.
Put
\[
       M=M(\mathbb Z/p,1)=\cofib(p:S^1\to S^1),
       \qquad X=G\wedge C.
\]
Then
\[
 \begin{gathered}
 P_X\bigl(\Sigma^2(M\wedge C)\bigr)\simeq *,\qquad
 P_X(\Sigma^4C)\simeq *,\\
 \min\mathrm{Rev}\{G\wedge C\to *\}
   =\min\mathrm{Rev}\{M\wedge C\to *\}
   =\min\mathrm{Rev}\{C\to *\}.
 \end{gathered}
 \tag{D.2.39a}
\]
More generally, for every $C$ in the aspherical class of
$\min\mathrm{Rev}(G\to *)$, one has
\[
 \min\mathrm{Rev}\{G\wedge C\to *\}
             =\min\mathrm{Rev}\{C\to *\}.
 \tag{D.2.39a'}
\]
Thus smashing the classifying-space/Moore-space distinction with
any such factor cannot retain that distinction between the
generated reversible classes.
No finiteness assumption on $C$, or compactness assumption on its
suspension spectrum, is required.
\end{claim}
\begin{proof}
We first establish the spectral consequence of
$B\mathbb Z/p$-nullity needed here:
\[
 \begin{gathered}
 K\simeq\Omega^\infty E\text{ for a connective spectrum }E,
 \quad\Map_*(G,K)\simeq *\quad\Longrightarrow\\
 p:\pi_jK\xrightarrow{\cong}\pi_jK
                  \quad\text{for every }j\geq2.
 \end{gathered}
 \tag{D.2.39b}
\]
Set $W=\Sigma^\infty G$ and $W_i=\Sigma^{-i}W$ for $i\geq0$.
The adjunction between suspension spectra and infinite-loop
spaces identifies the hypothesis with spectral $W$-nullity of
$E$, using mapping spaces, so $P_W^{\Sp}E\simeq E$.
The spectrum $W$ is $1$-connective and satisfies the
$1$-supported $p$-torsion condition. Write
\[
       E_i=P_{W_i}^{\Sp}E,
       \qquad E_\infty=\colim_{i\geq0}E_i.
 \tag{D.2.39c}
\]
The natural maps point from $E_i$ to $E_{i+1}$.
Bousfield's exact comparison theorem gives
\[
       \fib(E_i\to E_{i+1})\simeq\Sigma^{-i}HG_i
 \tag{D.2.39d}
\]
for some $p$-primary torsion group $G_i$
\cite[Theorem~2.6, p.~37]{Bousfield1996ULP}.
By extension closure of homotopy-group bounds and exactness of
filtered colimits in spectra, this implies
\[
       \fib(E_i\to E_\infty)\in\Ho(\Sp)_{\leq-i}.
 \tag{D.2.39e}
\]

Here is the required check at the limit. Fix $i\geq0$ and use
$E_{i+1}\to E_\infty$. Its fiber lies in degrees at most
$-i-1$, and its suspended fiber lies in degrees at most $-i$.
On the other hand $W_i\in\Ho(\Sp)_{\geq1-i}$.
The two $t$-structure orthogonalities therefore give an
equivalence of entire mapping spaces
\[
 \Map_{\Sp}(W_i,E_{i+1})
          \xrightarrow{\simeq}\Map_{\Sp}(W_i,E_\infty).
 \tag{D.2.39f}
\]
The left side is contractible, since $E_{i+1}$ is
$W_{i+1}$-null and $W_i=\Sigma W_{i+1}$.
It follows that $E_\infty$ is right orthogonal to every
integer suspension of $W$. The established Kahn--Priddy
calculation, C.3 and its prime-$p$ version in D.2.2, gives
\[
             \loc(W)=\loc(\mathbb S/p).
 \tag{D.2.39g}
\]
Consequently the function spectrum
$F(\mathbb S/p,E_\infty)$ is zero, and multiplication by $p$
on $E_\infty$ is an equivalence. Since $E_0=E$, the case
$i=0$ of \textup{(D.2.39e)} identifies $\pi_jE$ with
$\pi_jE_\infty$ for $j\geq2$. This proves
\textup{(D.2.39b)}. In particular, this argument has not
commuted $\Map_{\Sp}(W,-)$ with a filtered colimit; the
comparison with $E_{i+1}$ controls components as well as
positive homotopy groups.

Now let $Z$ be $X$-null. Then $\Map_*(C,Z)$ is $G$-null.
The assumed ordinary $G$-acyclicity of $C$ gives
\[
 \Map_*(C\wedge C,Z)
       \simeq\Map_*(C,\Map_*(C,Z))\simeq *.
 \tag{D.2.39h}
\]
As in D.2.37, all higher smash powers of $C$ are $X$-acyclic,
and the finite product filtration proves that
$C^r/\bigvee_{i=1}^rC$ is $X$-acyclic for every finite $r$.
Replace $\Omega Z$ by a simplicial group $H$. Restriction
\[
       \Map_*(C^r,H)\longrightarrow\Map_*(C,H)^r
 \tag{D.2.39i}
\]
is a Kan fibration of simplicial groups with contractible
zero fiber. Every tuple extends by ordered multiplication
on its separate coordinates. Translation therefore identifies
every fiber with the zero fiber, proving that
\textup{(D.2.39i)} is a weak equivalence on all components.

For this general $C$ there is no linear algebraic theory.
There is instead an actual functor on finite pointed sets:
\[
 \mathcal F(S)=\Map_*(C^{S\setminus\{*\}},H).
 \tag{D.2.39j}
\]
A pointed map $f:S\to T$ pulls coordinates back to give
$C^{T\setminus\{*\}}\to C^{S\setminus\{*\}}$, with the
basepoint in coordinates mapped to $*$. Precomposition gives
the required covariant map $\mathcal F(S)\to\mathcal F(T)$.
This is a $\Gamma$-space in the convention that
$\Gamma^{\mathrm{op}}$ is the category of finite pointed
sets. Its Segal maps are precisely
\textup{(D.2.39i)}, so it is special. The multiplication on
$\pi_0\mathcal F(\{*,1\})$ is induced by the fold map
$\{*,1,2\}\to\{*,1\}$. On powers of $C$ this is the
diagonal, and, using the ordered-product extension in
\textup{(D.2.39i)}, the induced law is pointwise
multiplication in $H$. Its components therefore form a
group. Thus $\mathcal F$ is very special.

Segal's recognition theorem now supplies a connective
spectrum $E$ and a pointed weak equivalence
\[
             K=\Map_*(C,\Omega Z)\simeq\Omega^\infty E
 \tag{D.2.39k}
\]
\cite[Definitions~1.1--1.3 and Proposition~1.4,
pp.~295--296; proof in Propositions~1.5--1.6 and
Lemmas~1.7--1.8, pp.~296--298]{I_Segal1974}.
One can realize the simplicial mapping spaces before applying
that formulation; their weak equivalences are homotopy
equivalences of CW types. Repeated classifying spaces give
the connective spectrum. This is an infinite-loop-space
conclusion, with no assertion that $K$ is a GEM.

By adjunction,
\[
             \Map_*(G,K)
                 \simeq\Map_*(X,\Omega Z)\simeq *.
 \tag{D.2.39l}
\]
Apply \textup{(D.2.39b)}. Multiplication by $p$ is an
isomorphism on all homotopy groups of $\Omega^2K$, including
its group of components. The $p$-power map on $\Omega^2K$,
induced by the degree-$p$ circle map, is consequently an
equivalence. The defining
Moore cofibration gives
\[
 \begin{aligned}
 \Map_*(M,\Omega K)
    &\simeq\fib\bigl(p:\Omega^2K\to\Omega^2K\bigr)
       \simeq *,\\
 \Map_*(M,\Omega K)
    &\simeq\Map_*\bigl(\Sigma^2(M\wedge C),Z\bigr).
 \end{aligned}
 \tag{D.2.39m}
\]
This proves the first ordinary acyclicity in
\textup{(D.2.39a)} and, by desuspension, puts $M\wedge C$
in the reversible null class generated by $X$.

We finish with an elementary telescope argument which also
recovers $C$ itself. More generally, let $D$ be any connected
space with $p$-primary reduced integral homology, put
$T=\Sigma D$, and take the suspension degree map
\[
 f=(p:S^1\to S^1)\wedge\id_D:T\longrightarrow T.
 \tag{D.2.39n}
\]
It induces multiplication by $p$ on reduced integral homology.
Its mapping telescope is simply connected and has zero reduced
integral homology, since each homology element is killed by a
power of $p$. The telescope is consequently contractible.
Set $D_r=\cofib(f^r)$, for $r\geq1$. The cofiber-of-composition
pushout gives actual homotopy cofiber sequences
\[
            D_r\longrightarrow D_{r+1}\longrightarrow D_1,
            \qquad D_1=M\wedge D.
 \tag{D.2.39o}
\]
Thus every $D_r$ is ordinarily $D_1$-acyclic. The structure
maps in this sequence use the identity on the source copy of
$T$ and $f$ on the target copy. Homotopy colimits preserve
these cofibers, so
\[
 \hocolim_rD_r
     \simeq\cofib\bigl(T\to\operatorname{tel}(f)\bigr)
     \simeq\Sigma T=\Sigma^2D.
 \tag{D.2.39p}
\]
In particular $\Sigma^2D$ is ordinarily $(M\wedge D)$-acyclic,
and
\[
       \min\mathrm{Rev}\{M\wedge D\to *\}
                  =\min\mathrm{Rev}\{D\to *\}.
 \tag{D.2.39q}
\]
The other inclusion in \textup{(D.2.39q)} is just closure
under smash products.

An ordinarily $G$-acyclic space has $p$-primary reduced integral
homology, as proved in D.2.7: every $H\mathbb Z[1/p]$-local
space is $G$-null. Apply the preceding construction with
$D=C$, and suspend every cofiber sequence and its final
homotopy colimit twice. Its building space
$\Sigma^2(M\wedge C)$ is $X$-acyclic by
\textup{(D.2.39m)}. Hence its resulting space $\Sigma^4C$
is $X$-acyclic. Desuspension puts $C$ in the reversible
class generated by $G\wedge C$, while $G\wedge C$ belongs
to the class generated by $C$ by smash closure. Together
with \textup{(D.2.39q)}, this proves all of
\textup{(D.2.39a)}.

Finally, if $C$ lies in the reversible null class generated
by $G$, D.2.7 with $r=1$ says that $\Sigma C$ is ordinarily
$G$-acyclic. In particular $C$ has $p$-primary reduced
homology and is connected. Apply the proved case to
$\Sigma C$. Since $G\wedge\Sigma C\simeq\Sigma(G\wedge C)$,
and a space and its suspension generate the same reversible
class, the resulting equality is exactly
\textup{(D.2.39a')}.
\end{proof}

\begin{claim}{D.2.40}{Proved}{The fourth-loop five-type and the failure of pure descent}
Put $B=K(\FF_2,2)$, $Y=\Sigma^2B$, $Z'=Y\langle7\rangle$,
$E=\Omega^2Z'$ and $U=\Omega^2E$.  Thus $Z'$ is six-connected
and $U$ is two-connected.  Let $\rho_2$ denote reduction modulo two
and $i:\FF_2\to\mathbb Z/4$ the inclusion $i(1)=2$.
There is a space-level equivalence
\[
 P_5U\simeq
 \operatorname{fib}\bigl(i\operatorname{Sq}^2\rho_2:
      K(\mathbb Z/4,3)\longrightarrow K(\mathbb Z/4,5)\bigr)
       \times K(\FF_2,5).
 \tag{D.2.40a}
\]
This statement does not assert an infinite-loop splitting.

Use the free double-loop homology coordinates of D.2.29 and write
\[
 \mu_{11}=\sigma^{*2}\bigl((\xi_1k)^*+(\xi_1f)^*\bigr)
       \in H^{11}(U;\FF_2).
 \tag{D.2.40b}
\]
Here the two ambient functionals on $H_{13}(\Omega^2Y;\FF_2)$
are restricted to $E$ before the two cohomology suspensions, and
$f\in H_6B$ is dual to $u^3$ in Serre's polynomial presentation.
Then
\[
 \mu_{11}\notin
 \operatorname{im}\bigl(H^{11}(P_5U;\FF_2)
                       \longrightarrow H^{11}(U;\FF_2)\bigr).
 \tag{D.2.40c}
\]
This excludes pure descent to the five-type.  It does not determine
the restriction of $\mu_{11}$ to $V=U\langle6\rangle$, or its
evaluation on the finite relative source in D.2.38.
\end{claim}
\begin{proof}
All unlabeled cohomology has coefficients $\FF_2$.
We first determine the two Postnikov invariants and then compute
only the degree-eleven quotient needed for \textup{(D.2.40c)}.

\medskip\noindent
\textbf{The two Postnikov invariants.}
The actual target calculations of D.2.18--D.2.22 give
\[
 \pi_3U=\mathbb Z/4\{\alpha\},\qquad
 \pi_4U=\mathbb Z/4\{\rho\},\qquad
 \alpha\eta_3=2\rho\ne0,\qquad
 \pi_5U=\mathbb Z/2\{\lambda\}.
 \tag{D.2.40d}
\]
Serre's integral calculation gives
\[
 H_4(K(\mathbb Z/4,3);\mathbb Z)=0,
 \qquad H_5(K(\mathbb Z/4,3);\mathbb Z)=\mathbb Z/2.
 \tag{D.2.40e}
\]
Indeed, his stable groups $A_n$ equal $H_{q+n}$ for $q>n$;
for a cyclic group of order at least four and $n>0$ they equal
those for $\mathbb Z/2$, indexed by admissible sequences of
degree $n$ whose first entry is even.  For $n=1$ there is none,
and for $n=2$ only $(2)$; take $q=3$.
See \cite[\S2, nos.~14--15, pp.~207--209,
Theorems~8--9]{I_Serre1953}.
The coefficient universal-coefficient theorem therefore identifies
\[
 H^5(K(\mathbb Z/4,3);\mathbb Z/4)
       =\operatorname{Hom}(\mathbb Z/2,\mathbb Z/4)
       \cong\mathbb Z/2.
 \tag{D.2.40f}
\]
Its nonzero class is $i\operatorname{Sq}^2\rho_2$:
$\operatorname{Sq}^2\rho_2$ is the unique nonzero mod-two
degree-five class, and coefficient inclusion is nonzero on
the displayed $\operatorname{Hom}$ group.
The zero Postnikov invariant would give a product of the two
Eilenberg--Mac Lane spaces and hence $\alpha\eta_3=0$.
Thus \textup{(D.2.40d)} forces the unique nonzero invariant and
identifies the first factor $P=P_4U$ in \textup{(D.2.40a)}.

The actual mod-two Hurewicz image of $\lambda$ is nonzero.
Two homology suspensions carry it to the nonzero class
$R=\{a,c\}\in H_7E$, as identified in D.2.19 and the
$D_0\vee S^7_R$ model of D.2.22.
In the next Postnikov fibration
\[
 K(\FF_2,5)\longrightarrow P_5U\longrightarrow P_4U,
 \tag{D.2.40g}
\]
the base is simply connected and the fiber map is an isomorphism
on $\pi_5$, since $\pi_5P_4U=\pi_6P_4U=0$.
A nonzero invariant in $H^6(P_4U;\FF_2)$ would evaluate
nontrivially on a mod-two six-cycle.  Its homology transgression
would kill the unique fiber $H_5$ class, forcing the Hurewicz
image of that $\pi_5$ generator to be zero in the total space.
This contradicts the value just established.  The invariant is
therefore zero, proving the space-level splitting.

\medskip\noindent
\textbf{The low Serre calculation.}
Write $x_3$ for the mod-two fundamental class of
$K(\mathbb Z/4,3)$, $y_4$ for its higher coefficient connecting
class, and $z_5=\operatorname{Sq}^2x_3$.
The coefficient fibration also gives the exact alternative model
\[
 P\simeq\operatorname{fib}\bigl(
 \operatorname{Sq}^2\rho_2+\operatorname{Sq}^1:
 K(\mathbb Z/4,3)\times K(\FF_2,4)
           \longrightarrow K(\FF_2,5)\bigr).
 \tag{D.2.40h}
\]
In fact $K(\mathbb Z/4,4)\to K(\FF_2,4)
\xrightarrow{\operatorname{Sq}^1}K(\FF_2,5)$ is the coefficient
fibration, classified by $i$; pulling it back gives precisely the
fiber of $i\operatorname{Sq}^2\rho_2$.
Consequently there is an actual class $r_4$ on $P$, restricting
to the fiber's mod-two fundamental class, with
\[
 \operatorname{Sq}^1r_4=z_5.
 \tag{D.2.40i}
\]
This is an actual Bockstein extension, not an inference from
the dimensions of a spectral sequence.

For the required spanning assertion use the Serre sequence of
$K(\mathbb Z/4,4)\to P\to K(\mathbb Z/4,3)$.
Serre's polynomial-generator theorem gives the complete fiber
generator list through degree eleven:
\[
\begin{gathered}
 r_4,\ s_5,\ \operatorname{Sq}^2r_4,\
 \operatorname{Sq}^3r_4,\ \operatorname{Sq}^2s_5,\
 \operatorname{Sq}^3s_5,\ \operatorname{Sq}^4s_5,\\
 \operatorname{Sq}^4\operatorname{Sq}^2r_4,\
 \operatorname{Sq}^5\operatorname{Sq}^2r_4,\
 \operatorname{Sq}^4\operatorname{Sq}^2s_5.
\end{gathered}
 \tag{D.2.40j}
\]
Here $s_5$ is the higher coefficient connecting class for
$\mathbb Z/4$, not $\operatorname{Sq}^1r_4$ in the total space.
The simple multiplicative system also includes $r_4^2,s_5^2$
in degrees eight and ten.
Serre's universal path-fibration construction makes these
generators and their $2^j$ powers transgressive, with suspension
and Steenrod operations compatible with transgression;
see \cite[\S2, no.~7, Theorem~1 and paragraph~7.3,
pp.~201--203; no.~11, Theorem~4 and its proof,
pp.~204--205]{I_Serre1953}.
After pullback a universally transgressive class survives every
earlier differential by naturality.  If its transgression is
zero, it is permanent, since no later outgoing differential can
start on that fiber row.

The degree-five fundamental class of $K(\mathbb Z/4,5)$ pulls
back to zero, while its higher degree-six coefficient class
pulls back to $x_3^2$.  For the latter assertion, pulling
$0\to\FF_2\to\mathbb Z/8\to\mathbb Z/4\to0$ back along
$i$ gives $0\to\FF_2\to\mathbb Z/4\to\FF_2\to0$;
thus the coefficient identity is $\delta_4i=\operatorname{Sq}^1$.
It follows that the complete nonzero differential list in the
needed range is
\[
 d_6(s_5)=x_3^2,\qquad
 d_{10}(\operatorname{Sq}^4s_5)=z_5^2.
 \tag{D.2.40k}
\]
Indeed Cartan and $\operatorname{Sq}^1x_3=0$ give
\[
 \operatorname{Sq}^2(x_3^2)=\operatorname{Sq}^3(x_3^2)
 =\operatorname{Sq}^5(x_3^2)=0,\quad
 \operatorname{Sq}^4(x_3^2)=z_5^2,\quad
 \operatorname{Sq}^4\operatorname{Sq}^2(x_3^2)=0.
 \tag{D.2.40l}
\]
In particular $s_5^2$, a separate transgressive simple-system
generator, survives.  The two nonzero images are a regular
sequence in the relevant polynomial algebra: $x_3^2$ is a
non-zero-divisor, and after quotienting by it the independent
polynomial generator $z_5$ still has non-zero-divisor square.
Thus these differentials remove the odd powers of $s_5$ and
$\operatorname{Sq}^4s_5$, retain $s_5^2$, and quotient by
$(x_3^2,z_5^2)$.  There are no additional annihilator cycles.
Every other generator needed through total degree eleven is
permanent by the universal-transgression argument, and
multiplicativity controls all mixed terms in that range.

The complete base generator list through degree eleven is
\[
\begin{gathered}
 x_3,\ y_4,\ z_5,\ \operatorname{Sq}^2y_4,\
 \operatorname{Sq}^3y_4,\\
 \operatorname{Sq}^4z_5,\
 \operatorname{Sq}^4\operatorname{Sq}^2y_4,\
 \operatorname{Sq}^5\operatorname{Sq}^2y_4.
\end{gathered}
 \tag{D.2.40m}
\]
Its only degree-eleven indecomposable is a positive operation.
On the fiber side the only degree-eleven indecomposables are
$\operatorname{Sq}^5\operatorname{Sq}^2r_4$ and
$\operatorname{Sq}^4\operatorname{Sq}^2s_5$.
The former is an operation of the actual class $r_4$.
For the latter, choose any lift $v_7$ of the permanent fiber
class $\operatorname{Sq}^2s_5$; then
$\operatorname{Sq}^4v_7$ has the required leading filtration term.
The new class $s_5^2$ has degree ten and no degree-one factor
can produce a further degree-eleven indecomposable.
All other filtration terms of degree eleven are products of
positive-degree permanent lifts.  Induction over the Serre
filtration therefore proves the spanning assertion without
requiring every multiplicative or Steenrod extension.
The extra factor $K(\FF_2,5)$ has only the degree-eleven
indecomposables $\operatorname{Sq}^5\operatorname{Sq}^1$ and
$\operatorname{Sq}^4\operatorname{Sq}^2$ of its fundamental
class, so altogether
\[
 H^{11}(P_5U)=
 \sum_{a=1}^{10} H^a(P_5U)H^{11-a}(P_5U)
 +\sum_{j=1}^{11}\operatorname{Sq}^jH^{11-j}(P_5U).
 \tag{D.2.40n}
\]

\medskip\noindent
\textbf{The canonical class detects failure of descent.}
Let $q_4\in H_4U$ be the top homology value of the canonical
Moore map $M(\mathbb Z/2,3)\to U$ from D.2.29.
Its bottom value is $2\alpha$ and has zero mod-two Hurewicz,
whereas two homology suspensions send $q_4$ to $k$.
Thus $q_4$ is primitive and $\beta q_4=0$; all other positive
homology squares vanish by connectivity.
The class $Q^7q_4\in H_{11}U$ is the top $E_4$ operation.
It is primitive, by the diagonal Cartan formula and the
vanishing of its possible mixed terms with units.
The full top Nishida formula, including its Browder correction,
shows that every positive homology square of this class is zero;
see \cite[Chapter~III, Theorems~1.2--1.4,
pp.~216--219]{I_CohenLadaMay1976}.
Explicitly, the potentially nonzero leading coefficients for
indices $1,2,3$ are
$\binom61,\binom52,\binom43$, all even.
Every remaining term involving a positive square of $q_4$,
including the Browder correction, vanishes.  For indices four
and five the remaining operations are unstable; for larger
indices ordinary Steenrod instability on the dual degree gives
zero.  Hence the image of $Q^7q_4$ in $H_{11}(P_5U)$ pairs
trivially with every summand of \textup{(D.2.40n)} and is zero.
On the other hand, homology suspension commutes with the
operations in this range and gives
\[
 \sigma^2(Q^7q_4)=Q^7k=\xi_1k,
 \qquad \langle\mu_{11},Q^7q_4\rangle=1.
 \tag{D.2.40o}
\]
This proves \textup{(D.2.40c)}.
The fiber of $U\to P_5U$ still has mixed Serre terms in the
range at issue.  Failure of descent therefore supplies no
assertion that restriction to that fiber is nonzero, and no
conclusion about the full fourth-suspension attachment follows.
\end{proof}

\begin{claim}{D.2.41}{Proved}{The actual sphere constraint on fourth-stage relative choices}
Let $M=M(\mathbb Z/2,3)$, let $D=D_2^{(4)}M$ be the
finite quadratic source of D.2.38, and let
$y_0:S^{10}\to C_4M$ be the actual spherical representative
constructed there. Its relative stable class has only a
weight-two component $y_D\in\pi_{10}D$. The fourfold adjunct
of $y_0$ is
\[
 \widetilde y_0=\pm l_4\in\pi_{14}P^8(2),
 \qquad |y_0|=|y_D|=4,
 \tag{D.2.41a}
\]
where $P^8(2)=M(\mathbb Z/2,7)$ and $l_4$ is the
attaching map of the top cell of the Stiefel manifold $V_{9,2}$.

Put $Y=\Sigma^2K(\mathbb F_2,2)$,
$Z'=Y\langle7\rangle$, $U=\Omega^4Z'$, and
$V=U\langle6\rangle$. Thus $V$ is five-connected.
Let $m:M\to U$ be the canonical dangerous Moore restriction
of D.2.29, and let $g_4:C_4M\to U$ be its free $E_4$
extension. If an ordinary map $f:QM\to U$ restricts to $m$,
its difference from $g_4$ determines a weight-two relative map
$v_2:D\to V$ satisfying
\[
 (v_2)_*y_D=-(g_4)_*y_0\in2\pi_{10}V,
 \qquad (v_2)_*Y_{10}=0\in H_{10}(V;\mathbb F_2).
 \tag{D.2.41b}
\]
The identification $\pi_{10}V=\pi_{10}U$ is used in this formula.
In particular, for the fixed paired functional
\[
 \mu_{11}=\sigma^{*2}
 \bigl((\xi_1k)^*+(\xi_1f)^*\bigr)\in H^{11}U
 \tag{D.2.41c}
\]
of D.2.40, the target relation
\[
 \left.\mu_{11}\right|_V\in
 \operatorname{Sq}^1H^{10}V+
 \operatorname{Sq}^3H^8V+
 \operatorname{Sq}^5H^6V
 \tag{D.2.41d}
\]
would suffice to exclude such a full extension. All cohomology
in \textup{(D.2.41c)--(D.2.41d)} has coefficients $\mathbb F_2$.
Relation \textup{(D.2.41d)} is a sufficient test, not a proved
target identity. The fourth-suspension extension problem and
Part D.2.4 remain open.
\end{claim}
\begin{proof}
\textbf{The arity-three component.}
Write $A_0,B_0$ for the bottom and top homology classes of $M$.
The free $E_4$ homology basis gives the following complete list
through degree eleven in $E_3=D_3^{(4)}M$:
\[
\begin{array}{c|l}
9&A_0^3\\
10&A_0^2B_0,\ A_0Q^4A_0\\
11&A_0B_0^2,\ A_0Q^5A_0,\ B_0Q^4A_0.
\end{array}
 \tag{D.2.41e}
\]
An indecomposable Dyer--Lashof expression has weight a power
of two. In this range the weight-three classes are therefore
products of weights one and two: the first weight-three
Browder expression has degree at least twelve. This proves
the completeness of the list, using
\cite[Chapter III, Theorems~3.1--3.2]{I_CohenLadaMay1976}.
Cartan and the operation table in D.2.38 give
\[
\begin{split}
 \beta(A_0^2B_0)&=\beta(A_0Q^4A_0)=A_0^3,\\
 \beta(B_0Q^4A_0)&=A_0^2B_0+A_0Q^4A_0,
\end{split}
 \tag{D.2.41f}
\]
with the other two degree-eleven Bocksteins zero. Each of the
three degree-eleven classes has second square $A_0^3$.

Consequently the stable ten-skeleton is
$M(\mathbb Z/2,9)\vee S^{10}_c$, where $c$ reduces to
$A_0^2B_0+A_0Q^4A_0$. Normalize the nonzero eleven-cell
integral boundary to $2c$ by an odd $2$-local unit. The other
two boundaries are divisible by four and can be made zero by
even cell slides. These slides preserve all the mod-two square
vectors. The cell $A_0B_0^2$ then attaches by the sole nonzero
element $i\eta_9\in\pi_{10}M(\mathbb Z/2,9)$, detected by
its second square. It kills that torsion summand; the remaining
relation is $2c$. Thus
\[
 \pi_{10}E_3=\mathbb Z/2,
 \qquad
 h_2:\pi_{10}E_3\longrightarrow H_{10}(E_3;\mathbb F_2)
 \text{ is injective}.
 \tag{D.2.41g}
\]
These are stable cell models, not an unstable cell cancellation
claim. Since $E_3$ is eight-connected, its ordinary group in
this degree agrees with the stable group. Cells above degree
eleven do not affect it. The degree-two boundary normalization
also follows directly from the rank-one first Bockstein in
\textup{(D.2.41f)}; there is no assumption that a nonzero
boundary has an unspecified higher power of two.

The relative Snaith splitting of $C_4M/M$ is compatible with
weights; see \cite[Theorems~1.1 and~3.2]{Snaith1974}.
The actual sphere $y_0$ has homology solely in the weight-two
class $Y_{10}$. Its weight-three projection is therefore zero
by \textup{(D.2.41g)}. The summands of weight at least four
are eleven-connected and contribute nothing to $\pi_{10}$.
Its entire relative stable class is consequently its
weight-two projection $y_D$. D.2.38 gives $|y_D|=4$.

\medskip\noindent
\textbf{The actual sphere before stable projection.}
D.2.38 proves that $y_0$ has zero image in $QM$. Thus its
fourfold adjunct $\widetilde y_0:S^{14}\to P^8(2)$ has zero
stable image. The primary calculation
\cite[Lemmas~2.7 and~2.8(ii), p.~523]{EH5MorisugiMukai2000}
is
\[
\begin{split}
 \pi_{14}P^8(2)&=\mathbb Z/4\{l_4\}
       \oplus\mathbb Z/2\{\widetilde{\nu_7^2}\},\\
 2l_4&=i_8\sigma',\qquad
 \Sigma l_4=i_9[\iota_8,\iota_8].
\end{split}
 \tag{D.2.41h}
\]
The second summand pinches to $\nu_8^2$, whose stable image is
nonzero. The $l_4$ summand stabilizes to zero by the displayed
Whitehead-square formula. Hence the exact stable kernel is
$\mathbb Z/4\{l_4\}$. Since the relative projection of
$\widetilde y_0$ has order four, its coefficient on $l_4$ is
odd. This proves \textup{(D.2.41a)}, including the order of
the actual sphere rather than merely its stable projection.

\medskip\noindent
\textbf{The relative map and its precise sphere value.}
The space $C_4M/M$ is five-connected. Form the pointwise
difference of $f|_{C_4M}$ and $g_4$ in the H-group $U$,
using the chosen homotopy between their restrictions to $M$.
The cofiber universal property produces a map
$d:C_4M/M\to U$, and connectivity gives its unique lift
to $V$. No multiplicativity of $f$ is used.

For a finite source $K$ of dimension at most eleven, stabilization
gives a bijection
\[
 [K,V]\ \xrightarrow{\ \cong\ }\
 [\Sigma^\infty K,\Sigma^\infty V].
 \tag{D.2.41i}
\]
Indeed the five-connectivity of $V$ makes the first suspension
surjective in this range. The unit $V\to\Omega\Sigma V$
has an actual left inverse because $V$ is a loop space, which
gives injectivity. The space $\Sigma V$ is six-connected and
$\dim\Sigma K\leq12$, so every subsequent suspension is
bijective by Freudenthal. Apply \textup{(D.2.41i)} to the
weight-two stable restriction of $d$ to obtain $v_2:D\to V$.
Apply it also to $S^{10}$: the weight calculation above and
the nullity of $y_0$ in $QM$ give the actual equality
\[
 (v_2)_*y_D=d_*y_0=-(g_4)_*y_0.
 \tag{D.2.41j}
\]
The middle notation means composition with the quotient
$C_4M\to C_4M/M$. This equality is proved in $\pi_{10}V$
and then identified with $\pi_{10}U$; it does not require
injectivity of the map of suspension spectra induced by
$V\to U$.

Write $\alpha\in\pi_3U=\mathbb Z/4$ and
$\rho\in\pi_4U=\mathbb Z/4$ for the generators already
computed, so $\alpha\eta_3=2\rho$. The bottom restriction
of $m$ is $-2\alpha$. In the stable range of maps $M\to M$,
the Moore Puppe sequence factors $2\id_M$ through the top
pinch followed by either zero or $i\eta_3$. Therefore
\[
 2m=0,
 \qquad
 \lambda_3(-2\alpha,m)
   =-2\lambda_3(\alpha,m)
   =-\lambda_3(\alpha,2m)=0.
 \tag{D.2.41k}
\]
For the first equality, the possibly nonzero factor gives
$(-2\alpha)\eta_3q=-4\rho q=0$. The second uses genuine
bi-additivity: the desuspensions of the two inputs of the
fourfold-adjoint Whitehead product are $S^6$ and $\Sigma^3M$,
both suspensions. Thus the hypotheses of
\cite[Proposition~3.4, p.~14]{EH5Arkowitz1962} apply.

The Moore map whose top-sphere factor is $(g_4)_*y_0$ is
exactly the generalized Browder product in
\textup{(D.2.41k)}. The actual Puppe sequence for
$M(\mathbb Z/2,9)$ now implies
$(g_4)_*y_0\in2\pi_{10}U$. Since $\pi_{10}V\to\pi_{10}U$
is an isomorphism, its lift is also a double. Equations
\textup{(D.2.41j)--(D.2.41k)} prove \textup{(D.2.41b)}.
They do not assert this doubled element is zero. More precisely,
if $h_M:P^8(2)\to Y$ is the adjoint of $m$, its value is
\[
 (g_4)_*y_0=\pm h_Ml_4\in\pi_{14}Y,
 \qquad 2(h_Ml_4)=(-2\alpha)\sigma'.
 \tag{D.2.41l}
\]
In the final formula $\alpha$ denotes its corresponding
generator of $\pi_7Y$. The composite $h_Ml_4$ is not evaluated.

\medskip\noindent
\textbf{The weaker sufficient target identity.}
The source table in D.2.38 gives
\[
 \beta Z_{11}=Y_{10},\qquad
 \operatorname{Sq}_*^3Z_{11}
   =\operatorname{Sq}_*^5Z_{11}=0.
 \tag{D.2.41m}
\]
Consequently \textup{(D.2.41b)} makes every term on the
right of \textup{(D.2.41d)} evaluate to zero on
$(v_2)_*Z_{11}$. The first-square term uses the specific
sphere constraint; it need not vanish for an arbitrary map
$D\to V$.

The top class $Q^7B_0$ is primitive, so the pointwise difference
adds its two homology images without a convolution term.
For the canonical $g_4$ its $\mu_{11}$ value is one, by
D.2.29 and D.2.40. For a map from $QM$ the value is zero:
$Q^7B_0=\beta Q^8B_0$ in the infinite-loop source, whereas
neither outer restricted-square coefficient defining
\textup{(D.2.41c)} occurs in $\beta H_{14}T$. The latter is
the free restricted-Lie basis argument of D.2.29 and applies
to both $\xi_1k$ and $\xi_1f$. Thus a full extension would
require the relative value one, contradicting
\textup{(D.2.41d)} if that relation were established.

Finally, D.2.36 and the Moore Puppe sequence show that every
dangerous fourth-suspension map has this same Moore restriction:
for fixed bottom map, the only nontrivial top change is
$\rho q$, whose double homology suspension adds
$r=b^2+k+p_3$. The entire-unit linear exclusion of D.2.36
rules out that change. Hence the sufficient test has the
stated relevance to all dangerous fourth-suspension maps,
without fixing their entire infinite unit maps. D.2.42 below
disproves the target membership in this conditional test.
\end{proof}

\begin{claim}{D.2.42}{Proved}{The actual secondary class lies outside the proposed odd-operation test}
Retain the notation of D.2.40, and put
\[
 L=Z'=Y\langle7\rangle,\qquad
 L'=Y\langle10\rangle,\qquad V=\Omega^4L'=U\langle6\rangle.
 \tag{D.2.42a}
\]
Thus $L$ is six-connected, $L'$ nine-connected and $V$
five-connected.  Let $y_8\in H^8L$ be the mod-two reduction
of the integral order-four coefficient connecting class of
the fundamental $\mathbb Z/4$ class in degree seven.
In the established basis it is $D_8^*+E_8^*$.
There is a common normalization of the relative square cocycle
which gives a class $\Theta_{15}\in H^{15}L'$ with
\[
 \sigma^{*4}\Theta_{15}=\mu_{11}|_V.
 \tag{D.2.42b}
\]
Here $\mu_{11}$ is the paired functional of D.2.40 and all
unlabeled cohomology has coefficients $\FF_2$.
The class is nonzero and satisfies
\[
 \mu_{11}|_V\notin
 \operatorname{Sq}^1H^{10}V+
 \operatorname{Sq}^3H^8V+
 \operatorname{Sq}^5H^6V.
 \tag{D.2.42c}
\]
This refutes that particular proposed sufficient operation
identity.  It does not assert that a map from the finite source
$D$ of D.2.38 realizes a nonzero evaluation, nor does it decide
the existence of a full fourth-suspension extension.
\end{claim}
\begin{proof}
We give the actual relative-cocycle construction before computing
its transgression.  This order distinguishes the class from a
formal bar symbol or an unspecified lift in a spectral sequence.

\medskip\noindent
\textbf{The square and its universal cocycle.}
The low double-loop cover basis of D.2.22 is
\[
 H_5E=\FF_2\{u\},\qquad
 H_6E=\FF_2\{D,E_6\},\qquad
 H_7E=\FF_2\{S,R\}.
 \tag{D.2.42d}
\]
The corresponding classes in $H_*(\Omega L)$, denoted
$u_6,D_7,E_7,S_8,R_8$, have images in the actual James
tensor algebra $H_*(\Omega Y)$
\[
 aa,\quad ab+ba,\quad f,\quad bb+g+h,\quad ac+ca,
 \tag{D.2.42e}
\]
respectively.  The tensor letters have degrees
$|a|=3$, $|b|=4$, $|c|=5$, $|f|=7$, $|g|=|h|=8$.
The complete eight ordered products of total degree fourteen
are linearly independent: the following table gives a tensor
word distinguishing each successive product from the others.
\[
\begin{array}{c|cccccccc}
\text{product}&uS&Su&uR&Ru&D^2&DE&ED&E^2\\ \hline
\text{word}&aag&gaa&aaac&caaa&abab&abf&fab&ff
\end{array}
\tag{D.2.42f}
\]
In a minimal Adams--Hilton model of the six-connected space
$L$, the least generator degree is six.  A differential of
degree fourteen can therefore have only these quadratic
terms; cubic terms start in degree eighteen.  Their independent
nonzero homology images exclude any nonzero such differential.
The quadratic-differential/reduced-diagonal identity gives
\[
 H^7L\,H^9L=0,\qquad H^8L\,H^8L=0
       \quad\text{in }H^{16}L.
 \tag{D.2.42g}
\]
The model and its diagonal interpretation are the same exact
ones used in D.2.26; see also
\cite[\S3.1, p.~300]{FGHHalperinLemaire1995}.
In particular $y_8^2=0$.

Choose an actual singular cocycle representing $y_8$ and a
cochain $c_{15}$ with $dc_{15}=y_8\smile y_8$.
For the cover map $q:L'\to L$, choose $P_7$ with
$dP_7=q^*y_8$, possible by nine-connectivity.  Then
\[
 \Theta_{15}=[q^*c_{15}-P_7\smile q^*y_8]
                     \in H^{15}L'
 \tag{D.2.42h}
\]
is closed.  A different $P_7$ changes it by a boundary, because
the difference is closed and $q^*y_8$ is exact.  Changing
$c_{15}$ by a cocycle $\ell_{15}$ changes $\Theta_{15}$ by
$q^*\ell_{15}$.
The actual singular-cochain Mayer--Vietoris construction of
D.2.26 says that its pullback by a lift of $h:S\to L$ is the
relative square class of the cone of $h$ whenever
$H^7S=H^8S=0$.  No abstract cone algebra is substituted for
the space in this assertion.

We now identify the required one common normalization.  Use
the universal evaluation map
\[
 \varepsilon:\Sigma^4V\longrightarrow L'\xrightarrow{q}L.
 \tag{D.2.42i}
\]
Its source starts in dimension ten.  For
$z\in\ker(\varepsilon_*:H_{15}(\Sigma^4V)\to H_{15}L)$,
take its connecting cone cycle and cancel the linear part of
the relative Adams--Hilton differential.  The remaining
degree-fourteen term has exactly the eight quadratic coordinates
in \textup{(D.2.42f)}.  Terms involving an old and a new
generator start in degree $6+10=16$, and old cubic terms
start in degree eighteen.  The quadratic-diagonal identity and
the actual cone pairing therefore identify the relative square
value on $z$ with the sum of the $D_7^2$ and $E_7^2$
coefficients of its once-adjoint image.  The mixed $D_7E_7$
and $E_7D_7$ coefficients are equal by primitivity and cancel;
the $uS,uR$ terms pair trivially with $y_8\otimes y_8$.
This is the per-vector argument of D.2.26j, in its explicitly
verified degree range.

Under \textup{(D.2.42e)}, the two square coordinates become
$(ab+ba)^2$ and $f^2$.  They are precisely the independent
restricted-square coordinates obtained by one homology
suspension from $\xi_1k$ and $\xi_1f$ in the free double-loop
basis.  Ordinary Hall-bracket coordinates and the other
restricted squares cannot change either coefficient; this is
the free basis used in D.2.29, from
\cite[Chapter~III, Theorems~1.2--1.4 and~3.1--3.2]{I_CohenLadaMay1976}.
Thus on the displayed kernel the relative square functional
equals $\mu_{11}|_V$, after removing the four source
suspensions.  Their difference consequently factors through
$\varepsilon_*H_{15}(\Sigma^4V)$.  Extend that functional to
one $\ell_{15}\in H^{15}L$ and replace $c_{15}$ by
$c_{15}+\ell_{15}$.  This proves \textup{(D.2.42b)} with
one normalization for all maps into $V$.
The fourfold cohomology suspension is an isomorphism here:
$L'$ is nine-connected and $V$ five-connected, and the
degree-eleven range is below the first product degree twelve.

\medskip\noindent
\textbf{The shifted Postnikov base.}
The proof of D.2.40, with the lower dimensions shifted by four,
gives the space-level description
\[
 P_9L\simeq P\times K(\FF_2,9),\qquad
 P=\operatorname{fib}\bigl(i\operatorname{Sq}^2\rho_2:
              K(\mathbb Z/4,7)\to K(\mathbb Z/4,9)\bigr).
 \tag{D.2.42j}
\]
For precision, the actual lower groups are
$\mathbb Z/4\{\alpha\}$ in degree seven,
$\mathbb Z/4\{\rho\}$ in degree eight, with
$\alpha\eta_7=2\rho$, and $\mathbb Z/2\{\lambda\}$
in degree nine.  Serre's stable integral calculation again
gives $H_8(K(\mathbb Z/4,7);\mathbb Z)=0$ and $H_9=\mathbb Z/2$,
so the unique nonzero lower invariant is the displayed one.
The homology suspension $H_7E\to H_9L$ is an isomorphism
in this range.  Thus the nonzero Hurewicz value $R$ of
$\lambda$ remains nonzero and excludes a nonzero next
Postnikov invariant, by exactly the transgression argument
of D.2.40g.  This proves \textup{(D.2.42j)} without assuming
an infinite-loop splitting.

Use $x_7,y_8$ for the fundamental and higher coefficient classes
on $K(\mathbb Z/4,7)$, $r_8,s_9$ for those of the fiber
$K(\mathbb Z/4,8)$, and $t_9$ for the fundamental class of
the separate $K(\FF_2,9)$ factor.
There is an actual class $r_8$ on $P$ with
\[
 \operatorname{Sq}^1x_7=\operatorname{Sq}^1y_8=0,
 \qquad\operatorname{Sq}^1r_8=\operatorname{Sq}^2x_7.
 \tag{D.2.42k}
\]
The last identity follows from the actual coefficient-fiber
model D.2.40h, with the dimensions shifted.  The pullback of
$y_8$ to $L$ is the class used above: the integral Moore-four
coordinates of D.2.22 identify its value as $D_8^*+E_8^*$.

The universal transgressive simple system of Serre gives the
following complete nonzero differential list with fiber source
degree at most fifteen:
\[
\begin{array}{c|c|c}
\text{source}&\text{differential}&\text{image}\\ \hline
s_9&d_{10}&\operatorname{Sq}^3x_7\\
\operatorname{Sq}^2s_9&d_{12}&\operatorname{Sq}^5x_7\\
\operatorname{Sq}^4s_9&d_{14}&
             \operatorname{Sq}^5\operatorname{Sq}^2x_7\\
\operatorname{Sq}^6s_9&d_{16}&
             \operatorname{Sq}^6\operatorname{Sq}^3x_7\\
\operatorname{Sq}^4\operatorname{Sq}^2s_9&d_{16}&
             \operatorname{Sq}^7\operatorname{Sq}^2x_7
\end{array}
\tag{D.2.42l}
\]
Here the higher coefficient image is
$\operatorname{Sq}^1\operatorname{Sq}^2x_7
=\operatorname{Sq}^3x_7$, not $x_7^2$.
Adem and $\operatorname{Sq}^1x_7=0$ give the later entries;
for the last also use instability of $\operatorname{Sq}^9x_7$.
The other relevant $s$-generators,
$\operatorname{Sq}^3s_9$ and $\operatorname{Sq}^5s_9$,
have zero transgression and survive in degrees twelve and
fourteen.  All $r$-generators survive as actual operations of
$r_8$.  The five images in \textup{(D.2.42l)} are independent
polynomial base generators and form a regular sequence.
Consequently no additional odd-annihilator cycles occur.
This uses the same precise simple-system theorem as D.2.40;
see \cite[\S2, no.~7, pp.~201--203, and no.~11,
pp.~204--205]{I_Serre1953}.

In particular the complete spanning lists for the three
degrees needed below are as follows; subscripts are suppressed
in the table.
\[
\begin{array}{c|l}
11&\operatorname{Sq}^4x,\ 
    \operatorname{Sq}^3y,\ \operatorname{Sq}^3r\\
13&\operatorname{Sq}^6x,\ \operatorname{Sq}^5y,\
    \operatorname{Sq}^4\operatorname{Sq}^2x,\
    \operatorname{Sq}^5r\\
15&\operatorname{Sq}^7y,\
    \operatorname{Sq}^6\operatorname{Sq}^2x,\
    \operatorname{Sq}^5\operatorname{Sq}^2y,\ xy,\\
  &\operatorname{Sq}^7r,\
    \operatorname{Sq}^5\operatorname{Sq}^2r,\ xr
\end{array}
\tag{D.2.42m}
\]
These are lists for $H^*(P)$, with no unspecified fiber lift
among their entries.  The extra permanent $s$-generators are
even-dimensional and produce no additional class in these
odd degrees: the first possible product involving $s_9$ is
$x_7s_9$ in degree sixteen, and it is nonpermanent; its
square begins in degree eighteen.
The separate $K(\FF_2,9)$ factor adds, respectively,
\[
\begin{array}{c|l}
11&\operatorname{Sq}^2t\\
13&\operatorname{Sq}^4t,\ 
            \operatorname{Sq}^3\operatorname{Sq}^1t\\
15&\operatorname{Sq}^6t,\
            \operatorname{Sq}^5\operatorname{Sq}^1t,\
            \operatorname{Sq}^4\operatorname{Sq}^2t.
\end{array}
\tag{D.2.42n}
\]
Cross-factor products start in degree $7+9=16$.

Apply the three odd squares to these complete lists.
On the $P$ factor, $\operatorname{Sq}^1H^{15}$ has only
the possible surviving value $x\operatorname{Sq}^2x$;
the other possibly nonzero value
$\operatorname{Sq}^7\operatorname{Sq}^2x$ is killed in
\textup{(D.2.42l)}.
Both $\operatorname{Sq}^3H^{13}(P)$ and
$\operatorname{Sq}^5H^{11}(P)$ give zero or that same killed
class.  For example
$\operatorname{Sq}^3\operatorname{Sq}^5r
=\operatorname{Sq}^7\operatorname{Sq}^1r
=\operatorname{Sq}^7\operatorname{Sq}^2x$ and
$\operatorname{Sq}^1(xr)=x\operatorname{Sq}^2x$.
On the separate factor the only resulting values are
$\operatorname{Sq}^7t$ and
$\operatorname{Sq}^5\operatorname{Sq}^2t$.
Therefore
\[
\begin{split}
 &\operatorname{Sq}^1H^{15}(P_9L)
  +\operatorname{Sq}^3H^{13}(P_9L)
  +\operatorname{Sq}^5H^{11}(P_9L)\\
 &\hspace{1cm}\subseteq
 \langle x\operatorname{Sq}^2x,
         \operatorname{Sq}^7t,
         \operatorname{Sq}^5\operatorname{Sq}^2t\rangle.
\end{split}
\tag{D.2.42o}
\]
The class $y_8^2$ is nonzero and independent of this span.
It is a polynomial square in the base of the first factor;
none of the five independent transgressions in
\textup{(D.2.42l)} kills it or relates it to another base
monomial.  Those two degree-fifteen fiber generators are also
the only possible sources of a differential into the scalar
base line in degree sixteen.  Thus its nonzero associated
filtration class survives to $H^{16}(P_9L)$.

\medskip\noindent
\textbf{Transgression of the actual relative class.}
In the actual fibration $L'\to L\to P_9L$, a mixed term
has total degree at least $7+10=17$.  Every fiber class of
degree $j$, $10\le j\le15$, is therefore transgressive:
for $r<7$ its outgoing differential has zero base target,
and for $7\le r\le j$ the target fiber degree is at most
nine and is zero.  The first possible differential is
$d_{j+1}$ to the scalar base line, and no earlier mixed
source can hit that line in this range.  Hence there are
unambiguous transgressions
\[
 \tau:H^jL'\longrightarrow H^{j+1}(P_9L),
             \qquad 10\le j\le15,
 \tag{D.2.42p}
\]
commuting with the Steenrod operations whose source and target
degrees are in this range.  This is ordinary cohomology
transgression compatibility, not a later-page operation with
an unmeasured indeterminacy.

For the particular class constructed above,
\[
                         \tau\Theta_{15}=y_8^2.
 \tag{D.2.42q}
\]
Indeed use a strict fibration model and choose $y_8$ on $L$
as the actual pullback of the base cocycle.  Its restriction
to the fiber is literally zero, so $P_7=0$ in
\textup{(D.2.42h)}.  The cochain $c_{15}$ satisfies
$dc_{15}=p^*(y_8^2)$, and the definition of transgression
gives \textup{(D.2.42q)}.  Adding the common closed
$\ell_{15}$ used to normalize \textup{(D.2.42b)} changes
the fiber class by a total-space restriction, whose
transgression is zero.

If $\Theta_{15}$ belonged to
$\operatorname{Sq}^1H^{14}L'+
\operatorname{Sq}^3H^{12}L'+
\operatorname{Sq}^5H^{10}L'$, operation compatibility in
\textup{(D.2.42p)} would put $y_8^2$ in the left side of
\textup{(D.2.42o)}, contradicting its proved independence.
The fourfold cohomology-suspension isomorphism commutes with
these operations, so \textup{(D.2.42b)} proves
\textup{(D.2.42c)}, including nonvanishing.
The operation condition was only a sufficient test for maps
from $D$.  Its failure does not settle the further attaching-map
and integral constraints on actual relative differences.
\end{proof}

\begin{claim}{D.2.43}{Proved}{The exact quadratic obstruction and its finite-model survival}
Work at the prime two and put
\[
 D_2E=(E\wedge E)_{hC_2},\qquad L_2E=\Sigma^{-1}D_2E,
 \qquad M_q=M(\mathbb Z/q,7).
\]
For a space in the argument of $D_2$ or $L_2$, its suspension
spectrum is understood. Let
\[
 K=(M_4\vee S^8_r)\cup_{i\eta_7+2r}e^9_s,
 \qquad m_0:M_2\xrightarrow{c_{2,4}}M_4\longrightarrow K,
 \tag{D.2.43a}
\]
where the first map has bottom degree two and top degree one.
Then
\[
 \begin{split}
  \pi_{14}L_2M_2&=\mathbb Z/4\{c_2\},\\
  \pi_{14}L_2K&=\mathbb Z/8\{c\}\oplus\mathbb Z/2\{u\},
  \qquad (L_2m_0)_*c_2=2c.
 \end{split}
 \tag{D.2.43b}
\]
The same target quadratic group applies to
$Z'=(\Sigma^2K(\mathbb F_2,2))\langle7\rangle$.
For the actual Moore map $h_M:M_2\to Z'$ fixed in D.2.41,
the natural metastable boundary satisfies
\[
 h_Ml_4=\partial_{Z'}(\pm2c).
 \tag{D.2.43c}
\]
Consequently the precise remaining nullity test is
\[
 h_Ml_4=0
 \quad\Longleftrightarrow\quad
 2c\in\operatorname{im}\left(
  JH_*:\pi_{15}^{s}Z'\longrightarrow
       \mathbb Z/8\{c\}\oplus\mathbb Z/2\{u\}
 \right),
 \tag{D.2.43d}
\]
where $JH$ is the stable James--Hopf map.

For the finite model $K$ itself, the James--Hopf homomorphism
is zero. Its quadratic boundary is therefore injective, and
\[
 |\partial_K(2c)|=4.
 \tag{D.2.43e}
\]
This class retains exact order four after adding the free
nine-sphere of the target model and any cells of dimensions
ten and eleven. The first possible new contribution to the
James--Hopf image comes from a twelve-cell.
No equality between the finite-model image and the image for
the full $Z'$ is asserted. Thus
\textup{(D.2.43d)} is an exact unresolved target calculation;
it does not settle the fourth-suspension extension problem.
\end{claim}
\begin{proof}
\textbf{The quadratic layer and its integral coordinates.}
The interchange in $D_2$ is the usual symmetric-monoidal
interchange. There is no additional sign representation:
\cite[Section~8, p.~704]{FGHGoodwillie2003} identifies the
second derivative of the identity with $S^{-1}$ carrying
trivial $C_2$-action.

\medskip\noindent
\textbf{The finite quadratic calculation, including the eta attachment.}
Write $D_2E=(E\wedge E)_{hC_2}$ and $L_2E=\Sigma^{-1}D_2E$,
with the ordinary interchange action.  Put $M_q=M(\ZZ/q,7)$,
where $q$ is even.  Its integral label chains have generators
$a_7,b_8$ and $\partial b=qa$.  Use the free $C_2$ resolution
with the convention
\[
 \partial(e_i xy)=e_{i-1}\bigl(xy+(-1)^i\tau(xy)\bigr)
                   +(-1)^ie_i\partial(xy).
 \tag{D.2.43q1}
\]
Thus $\tau(aa)=-aa$, $\tau(ab)=ba$, and $\tau(bb)=bb$.
The cells relevant to integral homology through degree fifteen are
\[
\begin{array}{c|l}
14&P=e_0aa\\
15&x=e_0ab,\quad y=e_0ba,\quad z=e_1aa\\
16&T=e_0bb,\quad h=e_1ab,\quad h'=e_1ba,\quad U=e_2aa.
\end{array}
\]
Their boundaries are
\[
\begin{aligned}
 \partial x&=-qP,&\partial y&=qP,&\partial z&=2P,\\
 \partial T&=q(x+y),&\partial h&=x-y+qz,
       &\partial h'&=-\partial h,\qquad\partial U=0.
\end{aligned}
\tag{D.2.43q2}
\]
Cancel the unit incidence in $\partial h$, so that $y=x+qz$,
and set $c_q=x+(q/2)z$.  Then $c_q$ is an integral cycle and
$\partial T=2q c_q$.  It follows that
\[
 H_{14}(D_2M_q;\ZZ)=\ZZ/2\{P\},\qquad
 H_{15}(D_2M_q;\ZZ)=\ZZ/(2q)\{c_q\}.
 \tag{D.2.43q3}
\]
This chain computation applies to spectra: extension of scalars
to $H\ZZ$ is symmetric monoidal, so the derived coinvariant
tensor complex computes the integral homology of the
homotopy-orbit square.  It does not discard the homotopy
information in eta attachments; that information is checked next.

The bottom inclusion gives $D_2S^7\to D_2M_q$.  Its first three
cells are $P_{14},z_{15},U_{16}$, with $\partial z=2P$ and
$\partial U=0$.  The extended-power basis and Nishida relation
in \cite[\S1.4, Lemma~1.4.1 and the following displayed relation,
p.~4]{EH5Behrens2012} identify these mod-two classes as
$Q^7a,Q^8a,Q^9a$ and give
\[
 \operatorname{Sq}_*^2U
  =\operatorname{Sq}_*^2Q^9a
  =\binom72 Q^7a=P.
 \tag{D.2.43q4}
\]
The fifteen-skeleton of $D_2S^7$ is $M(\ZZ/2,14)$;
its fifteenth stable homotopy group is generated by the bottom
eta.  Hence the sixteen-cell $U$ attaches by, and kills, $P\eta$.
The same class remains zero after mapping to any of the quadratic
spectra considered below.

To see that this accounts for the entire possible Hurewicz kernel,
let $E$ be a spectrum with cells in dimensions at least fourteen
and $H_{14}E=\ZZ/2\{P\}$.  The fiber of
$E\to H\ZZ\wedge E$ is $E\wedge\tau_{\ge1}\mathbb S$.
Its lowest possible homotopy group is
\[
 \pi_{15}(E\wedge\tau_{\ge1}\mathbb S)
       =H_{14}E\otimes\pi_1\mathbb S=\ZZ/2,
\]
and its map into $\pi_{15}E$ is the bottom eta action.
Thus the degree-fifteen stable Hurewicz map is surjective with
kernel generated at most by $P\eta$.  Equation~\textup{(D.2.43q4)}
kills that kernel here, proving the natural isomorphisms
\[
 \pi_{15}D_2M_q\xrightarrow{\ \cong\ }
       H_{15}(D_2M_q;\ZZ)=\ZZ/(2q)\{c_q\}.
 \tag{D.2.43q5}
\]
Equivalently, stable unit cancellation gives the relevant
Chang block with attaching matrix $(2,\eta_{14})$, together
with a Moore $2q$ block whose bottom dimension is fifteen.
Any eta component on that Moore block's top attachment is
absorbed by $U$; no hidden homotopy extension is being inferred
from integral chains alone.

The specified comparison $c_{2,4}:M_2\to M_4$ acts on label chains
by $a\mapsto2a$, $b\mapsto b$.  Therefore
\[
 x\longmapsto2x,\qquad z\longmapsto4z,
 \qquad c_2=x+z\longmapsto2x+4z=2c_4.
 \tag{D.2.43q6}
\]
Naturality of the Hurewicz isomorphism computes the map on
$\pi_{15}$ as $\ZZ/4\to\ZZ/8$, $1\mapsto2$.

Now put $K_8=M_4\vee S^8_r$.  Distributing the two smash factors
and taking homotopy orbits gives the actual stable splitting
\[
 D_2K_8\simeq D_2M_4\vee(M_4\wedge S^8_r)\vee D_2S^8_r.
\]
The middle summand is $M(\ZZ/4,15)$, with bottom class $u=ar$;
the last summand starts in dimension sixteen.  Thus
\[
 \pi_{15}D_2K_8=\ZZ/8\{c_4\}\oplus\ZZ/4\{u\}.
 \tag{D.2.43q7}
\]
Attach the nine-cell of
$K=K_8\cup_{i\eta_7+2r}e^9_s$ and filter its square by the
number of factors in the new cell.  The one-new-factor quotient
is $K_8\wedge S^9$; the two-new-factor quotient is $D_2S^9$,
which starts in dimension eighteen.  Through dimension sixteen,
the former quotient has just the cell $as$.  Its attachment is
the orbit image of $a\wedge(i\eta_7+2r)$, hence
$P\eta+2u$.  There is no extra factor of two on the eta term:
the off-diagonal cross term is an induced free $C_2$-orbit,
whose homotopy orbits give one copy, as in the preceding wedge
splitting.  By \textup{(D.2.43q4)} this imposes exactly $2u=0$.
All further cells have dimension at least seventeen and do not
affect $\pi_{15}$.  Consequently
\[
 \begin{aligned}
 \pi_{14}L_2M_2&=\ZZ/4\{c_2\},\\
 \pi_{14}L_2K&=\ZZ/8\{c_4\}\oplus\ZZ/2\{u\},
       & (L_2m_0)_*c_2&=2c_4.
 \end{aligned}
 \tag{D.2.43q8}
\]

Finally, the extra free nine-sphere in the complete target model
adds a free sixteen-cell to the quadratic spectrum, with zero
attachment, and higher-dimensional cells.  It has no effect on
\textup{(D.2.43q8)}.  A target cell of dimension at least ten
first contributes a cross term in dimension at least seventeen,
because the bottom target cell has dimension seven.  Homotopy
orbits add only nonnegative cell dimensions.  Thus the complete
target's quadratic spectrum has the same fifteenth homotopy
group and the same specified map.  This truncation statement
concerns the quadratic layer; it does not truncate the domain
of the target James--Hopf homomorphism.


\medskip\noindent
\textbf{The actual metastable boundary.}
Theorem~4.2.1 of \cite{EH5Behrens2012}, printed p.~37 in
the opened author version, gives the natural fiber sequence
\[
 P_2\operatorname{Id}(X)\longrightarrow QX
       \xrightarrow{JH}\Omega^\infty D_2\Sigma^\infty X.
 \tag{D.2.43f}
\]
Here $JH$ is adjoint to the weight-two projection of the
Snaith splitting. The following comparison statement in that
source gives
$\pi_tX\cong\pi_tP_2\operatorname{Id}(X)$ for
$t\leq3\operatorname{conn}(X)$. All our spaces in this
paragraph are six-connected, so degrees fourteen and fifteen
are within the stated range. The resulting exact sequence is
\[
 \pi_{15}^{s}X\xrightarrow{JH_*}\pi_{15}D_2X
       \xrightarrow{\partial_X}\pi_{14}X
       \longrightarrow\pi_{14}^{s}X.
 \tag{D.2.43g}
\]

D.2.41, using
\cite[Lemmas~2.7--2.8(ii), p.~523]{EH5MorisugiMukai2000},
identifies the exact stable kernel of $M_2$ as
$\mathbb Z/4\{l_4\}$. The quadratic source in
\textup{(D.2.43b)} also has order four. Hence its boundary
onto that kernel is an isomorphism:
\[
 \partial_{M_2}:\pi_{15}D_2M_2
       \xrightarrow{\ \cong\ }
       \ker(\pi_{14}M_2\to\pi_{14}^{s}M_2).
 \tag{D.2.43h}
\]
In particular $JH_*:\pi_{15}^{s}M_2\to\pi_{15}D_2M_2$
is zero. The generator lifting $l_4$ is unique, up to its
sign convention; there is no unrecorded choice of a source
connecting class.

The stable nine-skeleton of $Z'$ is $K\vee S^9_R$,
by twice suspending the model in \textup{(D.2.22t2)}.
In this low range the model and its specified Moore map
are realized as ordinary maps: their source dimensions are
at most nine and the targets are six-connected, so
Freudenthal gives the required bijection with stable maps.
The pure-$k$ Moore restriction of D.2.41 is the
bottom-two/top-one comparison in \textup{(D.2.43a)}.
The quadratic truncation proved above therefore identifies
its quadratic map with \textup{(D.2.43b)}.
Naturality of \textup{(D.2.43g)--(D.2.43h)} gives
\textup{(D.2.43c)}; exactness gives \textup{(D.2.43d)}.
Unlike the quadratic group, the James--Hopf image cannot
be computed by merely replacing its domain with the
nine-skeleton.

\medskip\noindent
\textbf{Exact order four in the finite model.}
The intrinsic Thom-cell calculation used above also gives
\[
 \pi_{15}D_2S^7=0:
 \tag{D.2.43i}
\]
its integral homology in this degree is zero, and the
sixteen-cell kills its only possible Hurewicz-kernel class,
the bottom eta. By naturality every class induced from
the bottom sphere has zero James--Hopf image.

We first specify an actual boundary calculation.
Let $L=M_4$ or $K$, put $B=L/S^7$, and let
$F=\operatorname{hofib}(L\to B)$. The defining homotopy
pushout square has initial maps $S^7\to L$ and $S^7\to *$,
both seven-connected. Blakers--Massey, in the form
\cite[Section~2, p.~309, and Theorem~2.3, p.~310]{FGHGoodwillie1992},
makes the canonical map $i:S^7\to F$ thirteen-connected.
In both cases $B=\Sigma B'$ with $\dim B'\leq8$. The
fibration boundary restricted to the suspension unit,
$B'\to\Omega B\to F$, therefore factors uniquely through $i$.

This factor is determined by stabilization, since maps
from $B'$ into $S^7$ are in the stable range. To identify it,
compare with the stable fiber:
\[
 F\longrightarrow\operatorname{hofib}(QL\to QB)
       \simeq QS^7.
\]
The map of fiber sequences identifies the stabilized factor
with the desuspension of the stable cofiber connecting map.
The resulting \emph{actual} factors are
\[
 \begin{array}{c|c|c}
 L&B'&B'\longrightarrow S^7\\ \hline
 M_4&S^7&4\iota_7\\
 K&S^7_b\vee M(\mathbb Z/2,7)_r&(4\iota_7,\eta_7q).
 \end{array}
 \tag{D.2.43j}
\]
For the second row, collapsing the bottom gives
\[
 K/S^7\simeq S^8_b\vee M(\mathbb Z/2,8)_r,\qquad
 \delta=(4,\eta_8q):K/S^7\longrightarrow S^8.
 \tag{D.2.43k}
\]
Thus the eta component of the nine-cell is retained in both
the stable connecting map and its actual low-unit factor.
This argument uses a homotopy fiber, not an unstable
covariant exact sequence for a cofiber.

Next, \cite[Introduction, p.~515]{EH5MorisugiMukai2000}
states
\[
 \pi_{14}S^7_{(2)}=\mathbb Z/8\{\sigma'\},\qquad
 [\iota_8,\iota_8]=\pm(2\sigma_8-E\sigma').
 \tag{D.2.43l}
\]
Consequently $E^\infty\sigma'=2\sigma$, up to sign.
The actual class $2E\sigma'\in\pi_{15}S^8$ has stable
image $4\sigma$. Its boundary in the first case of
\textup{(D.2.43j)} is
\[
 i\bigl((4\iota_7)\circ2\sigma'\bigr)=8i\sigma'=0.
 \tag{D.2.43m}
\]
There is no omitted Hilton term: $S^7$ is an H-space,
so its degree-four map induces multiplication by four
on homotopy groups. Equivalently the Hilton correction
vanishes because $[\iota_7,\iota_7]=0$; see also
\cite[Lemma~1.1]{EH5MorisugiMukai2000}.

The actual fibration exact sequence therefore supplies
$\widehat a\in\pi_{15}M_4$ with pinch $2E\sigma'$.
Its stable class has zero James--Hopf image and pinch
$4\sigma$. Every stable class of $M_4$ pinches to
$4j\sigma\in\ker(4:\mathbb Z/16\{\sigma\}\to
\mathbb Z/16\{\sigma\})$. After subtracting
$j\widehat a$, it comes from the bottom $S^7$, whose
James--Hopf image is zero by \textup{(D.2.43i)}. Hence
\[
 JH_*:\pi_{15}^{s}M_4\longrightarrow\pi_{15}D_2M_4
 \quad\text{is zero}.
 \tag{D.2.43n}
\]

It remains to handle the additional stable quotient
classes in \textup{(D.2.43k)}. As in D.2.38,
\[
 \pi_{15}^{s}M(\mathbb Z/2,8)=\pi_7(\mathbb S/2),
 \qquad 2\pi_7(\mathbb S/2)=0.
 \tag{D.2.43o}
\]
Indeed $2\id_{\mathbb S/2}$ factors through the top
pinch followed by zero or $i\eta$; the pinch of a seventh
homotopy class is a multiple of $\nu^2$, and
$\eta\nu^2=0$. The low stable stems used here and below
are listed in \cite[Section~2, (2.1), p.~149]{I_Oka1986}.
The Moore exact sequence shows that this group is
generated by a bottom $\sigma$-class and a coextension
of $\nu^2$.

For that coextension use the actual class
\[
 t:S^{14}\longrightarrow M(\mathbb Z/2,7),
 \qquad qt=\nu_8^2,
 \tag{D.2.43p}
\]
provided by
\cite[Lemma~2.8(ii), p.~523]{EH5MorisugiMukai2000}.
Its suspension
$\xi=Et:S^{15}\to M(\mathbb Z/2,8)\subset B$
has the required stable pinch. The second actual factor
in \textup{(D.2.43j)} gives
\[
 \partial\xi=i\eta_7qt=i\eta_7\nu_8^2=0.
 \tag{D.2.43q}
\]
Here $\eta_7\nu_8\in\pi_{11}S^7$ is already in the
stable range, where the fourth stem is zero.
Thus $\xi$ has an actual lift to $\pi_{15}K$, whose
stable class has zero James--Hopf image. The other
Moore quotient generator is covered by the actual map
$r\sigma_8:S^{15}\to S^8_r\to K$, again with zero
James--Hopf image.

Finally take any $x\in\pi_{15}^{s}K$. If its quotient
coordinates are $(a,\xi')$, exactness of the stable
cofiber sequence gives $4a+\eta q\xi'=0$.
The second term vanishes, since $q\xi'$ is a multiple
of $\nu^2$. Thus $a=4j\sigma$. Subtract the old
Moore-four class, a bottom-$r$ sphere class, and the
actual lift from \textup{(D.2.43q)}, according to these
quotient coordinates. The remainder comes from the
bottom $S^7$. Each summand has zero James--Hopf image,
the last by \textup{(D.2.43i)}. We have proved
\[
 JH_*:\pi_{15}^{s}K\longrightarrow\pi_{15}D_2K
 \quad\text{is zero}.
 \tag{D.2.43r}
\]
For the stable bottom remainder, no separate
unstabilization claim is needed. Exactness of
\textup{(D.2.43g)} makes $\partial_K$ injective, and
$2c$ has order four, proving \textup{(D.2.43e)}.

The additional free $S^9_R$ contributes stable classes
through its bottom inclusion, but $D_2S^9$ starts in
dimension eighteen. Their James--Hopf images in degree
fifteen are therefore zero. Next, attaching cells of
dimensions ten and eleven induces a surjection on
fifteenth stable homotopy: the corresponding quotient
terms are sums of $\pi_5^s\mathbb S$ and
$\pi_4^s\mathbb S$, both zero. The quadratic target group
remains unchanged by the dimension argument above, so
naturality prevents any enlargement of the James--Hopf
image at these stages. A twelve-cell is the first
possibility, since $\pi_{15}^sS^{12}=\pi_3^s\mathbb S$.
This proves the final finite-range assertion while
retaining the unresolved full-target test
\textup{(D.2.43d)}.
\end{proof}

\begin{claim}{D.2.44}{Proved}{The universal Moore composite comes from the bottom sphere}
Work at the prime two, write $P_n=M(\mathbb Z/2,n-1)$, and
let $i_n:S^{n-1}\to P_n$ be the bottom inclusion.  Put
\[
 w=[i_5,\operatorname{id}_{P_5}]:P_8\longrightarrow P_5,
 \qquad h_M=jw:P_8\longrightarrow Y=\Sigma^2K(\FF_2,2),
 \tag{D.2.44a}
\]
where $j:P_5\to Y$ is the bottom Moore map of D.2.41.
For the Stiefel attachment $l_4\in\pi_{14}P_8$ of that
claim, the following three assertions hold.

First, the actual universal composite satisfies
\[
 wl_4\in\operatorname{im}\bigl(
    i_{5*}:\pi_{14}S^4_{(2)}\longrightarrow\pi_{14}P_5\bigr).
 \tag{D.2.44b}
\]
Second, with $u_4=ji_5$ and $\alpha=u_4\nu_4$, its target
image has the form
\[
 h_Ml_4=a\,\alpha\sigma'+b\,u_4E\epsilon',
 \qquad a\in\mathbb Z,\quad b\in\{0,1\},
 \qquad 2u_4E\epsilon'=0.
 \tag{D.2.44c}
\]
Neither coefficient is evaluated here.  Third, the two bottom
images that remain in this formula already suspend to zero
separately:
\[
 E(i_5\nu_4\sigma')=0,
 \qquad E(i_5E\epsilon')=0
                 \quad\text{in }\pi_{15}P_6.
 \tag{D.2.44d}
\]
Consequently the known equality $E(wl_4)=0$ imposes no further
relation between these two coefficients.  These assertions do
not evaluate $h_Ml_4$ or decide the James--Hopf image question
of D.2.43.
\end{claim}
\begin{proof}
All Whitehead products below are actual products of spaces.
The generalized product in (D.2.44a) has domain
$\Sigma(S^3\wedge P_4)=P_8$, and its bottom restriction is
\[
 wi_8=i_5\omega_4,\qquad
 \omega_4=[\iota_4,\iota_4].
 \tag{D.2.44e}
\]
This is the generalized-Whitehead/commutator convention of
D.2.29 and D.2.41; see
\cite[Definitions 2.2--2.3 and Theorem 2.4,
pp.~10--13]{EH5Arkowitz1962}.
Morisugi--Mukai give
\[
 \pi_{14}P_8=\mathbb Z/4\{l_4\}
       \oplus\mathbb Z/2\{\widetilde{\nu_7^2}\},
 \quad 2l_4=i_8\sigma',\quad
 j_*l_4=[\gamma_8,\iota_7],
 \tag{D.2.44f}
\]
where the last $j_*$ is the map to the relative group
of $(P_8,S^7)$ and $\gamma_8$ is its characteristic
relative eight-cell
\cite[Lemmas 2.7--2.8, p.~523]{EH5MorisugiMukai2000}.
In particular the pinch $p_8:P_8\to S^8$ kills $l_4$:
in the relative product it kills the absolute bottom entry.

\medskip\noindent
\textbf{The actual third cell of the pinch fiber.}
Let $F_5$ be the homotopy fiber of $p_5:P_5\to S^5$.
Its cells begin in dimensions $4,8,12,16$, and the first
two nontrivial skeleta are
\[
 A_5=S^4\cup_{2\omega_4}e^8,
 \qquad B_5=A_5\cup_{2\tau_5}e^{12}.
 \tag{D.2.44g}
\]
The second formula is an exact two-local statement.
Indeed Mukai computes
\[
 \pi_{11}A_5=\mathbb Z\{\tau_5\}
   \oplus\mathbb Z/4\{\widetilde{2\nu_7}\}
   \oplus\mathbb Z/15\{i\phi\},
 \qquad j_*\tau_5=3[\gamma_8,\iota_4],
 \tag{D.2.44h}
\]
and proves that the actual twelve-cell attachment is
$2\tau_5+a i\phi$, with $\phi\in\pi_{11}S^4$ of
odd order.  In particular it has no component in the displayed
Moore coextension summand
\cite[Lemma 4.1(iii), p.~69, and Lemma 5.1(i), p.~70]
{FGHMukai2001}.  Thus (D.2.44g) does not discard a
two-primary indeterminacy in the attachment.

Choose the map of cofibers $c:P_8\to A_5$ induced by the
square with top map $2:S^7\to S^7$, bottom map
$2\omega_4:S^7\to S^4$, left vertical map the identity,
and right vertical map $\omega_4$.  Its bottom map is
$\omega_4$ and its relative top degree is one.  Write
$f:P_8\xrightarrow{c}A_5\to P_5$ for its actual composite.
The restrictions of $f$ and $w$ to $S^7$ agree, so exactness
of the Moore Puppe sequence gives
\[
                    w=f+h p_8
 \quad\text{for some }h:S^8\longrightarrow P_5.
 \tag{D.2.44i}
\]

\medskip\noindent
\textbf{The necessary Hopf--Hilton correction.}
Although $p_8l_4=0$, precomposition by $l_4$ need not
preserve addition in $[P_8,P_5]$.  We compute the correction
in (D.2.44i) before using it.  Write $P_8=\Sigma P_7$
and apply Hilton--Milnor to
$\Omega\Sigma(P_7\vee P_7)$ through degree fourteen.
Besides the two singleton factors, the only possible factor
is $\Sigma(P_7\wedge P_7)$, with bottom cell thirteen.
Word length three begins in dimension nineteen and cannot
contribute.  This is the usual two-letter Hilton--Milnor
decomposition, with the range obtained from the indicated
bottom dimensions; compare
\cite[Proposition 2.2 and its proof]{I_CMN1987}.

Since the second map is $h p_8$, the mixed factor further
maps through
\[
       \Sigma(P_7\wedge S^7)\simeq P_{15}.
 \tag{D.2.44j}
\]
Its fourteenth homotopy group is its bottom $\mathbb Z/2$
by Hurewicz.  Naturality of the Whitehead map therefore
identifies the entire possible correction with
$\varepsilon[\alpha_0,h]$, where
$\alpha_0=i_5\omega_4=[i_5,i_5]$ and
$\varepsilon\in\FF_2$.  There is no omitted top--top
or higher-word contribution in this range.

Put $u=i_5$.  Its shifted degree is three, hence odd.
Graded Jacobi and bilinearity give
\[
 [[u,u],h]=\pm2[u,[u,h]]
          =\pm[2u,[u,h]]=0,
 \tag{D.2.44k}
\]
since $2u=0$ in the actual bottom Moore homotopy group.
Thus the correction is zero and the required equality is
valid:
\[
                 wl_4=(A_5\longrightarrow P_5)cl_4.
 \tag{D.2.44l}
\]
This argument is essential: the pinch equation alone would
not justify (D.2.44l).

\medskip\noindent
\textbf{Relative Jacobi and the bottom image.}
Naturality of the relative Whitehead product gives
\[
 j_*(cl_4)=[\gamma_8,\omega_4]
          =2[[\gamma_8,\iota_4],\iota_4]
 \tag{D.2.44m}
\]
up to the common orientation convention.  Only one relative
class occurs here.  The mixed relative Whitehead module
satisfies the Jacobi derivation law; the shifted degrees of
$\gamma_8$ and $\iota_4$ are seven and three, respectively,
so its two terms have the same sign.  James proves precisely
this mixed relative Samelson identity, and its equivalent
left-action form, in
\cite[\S4, equations (4.1)--(4.2), p.~337]{FGHJames1971}.
The loop-group model of a CW pair identifies it with the
relative Whitehead module identity used here.  No product
of two arbitrary relative classes or metastable exact
sequence is being assumed.

Combining (D.2.44h) and (D.2.44m) yields
\[
           3j_*(cl_4)=j_*[2\tau_5,\iota_4].
 \tag{D.2.44n}
\]
In $B_5$ the absolute class $2\tau_5$ is zero by the
actual twelve-cell attachment.  Therefore the image of
$3j_*(cl_4)$ in $\pi_{14}(B_5,S^4)$ is zero.  Multiplication
by three is invertible at the prime two.  Exactness of this
pair puts the image of $cl_4$ in $B_5$ in the bottom-sphere
image.  Composing $B_5\to F_5\to P_5$ and using
(D.2.44l) proves (D.2.44b).

\medskip\noindent
\textbf{The two remaining target terms.}
The sphere table and composition relation needed here are
\[
\begin{split}
 \pi_{14}S^4_{(2)}
    &=\mathbb Z/8\{\nu_4\sigma'\}
      \oplus\mathbb Z/4\{E\epsilon'\}
      \oplus\mathbb Z/2\{\eta_4\mu_5\},\\
 (E\nu')\sigma'&=2E\epsilon'.
\end{split}
 \tag{D.2.44o}
\]
They are displayed and proved, respectively, in
\cite[Propositions 7 and 13, pp.~7--9]{FGHStricklandToda}.
For clarity about the source's Hopf-formula warning, the
left factor $E\nu'$ is a suspension.  For the comparison
$H(\nu_4\sigma')=\sigma'$ one may use the formula after
one suspension and the injectivity of
$E:\pi_{14}S^7\to\pi_{15}S^8$, since both $\sigma'$
and $E\sigma'$ have order eight.  The remaining correction
$\nu_7\nu_{10}\eta_{13}$ is zero.  Thus this instance
does not use the unqualified general composite formula
about which that source warns.

The established target relations are
$u_4\eta_4=0$, $\alpha=u_4\nu_4$, $4\alpha=0$, and
$[u_4,u_4]=-2\alpha$ (D.2.17--D.2.18 and D.2.29).
Using $\omega_4=2\nu_4-E\nu'$ with compatible signs gives
\[
 u_4E\nu'=2\alpha-[u_4,u_4]=4\alpha=0,
 \qquad
 2u_4E\epsilon'=u_4(E\nu')\sigma'=0.
 \tag{D.2.44p}
\]
Postcomposition by $u_4$ is a homomorphism on the sphere
homotopy group.  It kills the third summand in (D.2.44o),
and (D.2.44p) reduces the coefficient of the second to
zero or one.  The bottom-image assertion now proves
(D.2.44c).  The exact double from D.2.41 remains
$2(h_Ml_4)=(-2\alpha)\sigma'$; it is not an evaluation
of either coefficient.

\medskip\noindent
\textbf{Why suspension supplies no additional coefficient test.}
The same primary table gives
\[
 E(\nu_4\sigma')=\pm E^2\epsilon'
                 =2\nu_5\sigma_8.
 \tag{D.2.44q}
\]
We show that the bottom inclusion $i_6$ kills this common
class.  Since $\nu_5$ is a suspension, the degree-two map
commutes with it.  The Hopf invariant of $\sigma_8$ is one,
so the actual Hopf--Hilton formula gives
\[
 (2\iota_5)\nu_5\sigma_8
 =\nu_5(2\iota_8)\sigma_8
 =2\nu_5\sigma_8+[\nu_5,\nu_5].
 \tag{D.2.44r}
\]
There are no higher Hilton terms in the dimension of
$\sigma_8$.  Naturality and the classical identity
$[\iota_5,\iota_5]=\nu_5\eta_8$ imply
\[
 [\nu_5,\nu_5]
 =\pm[\iota_5,\iota_5]\,\Sigma(\nu_4\wedge\nu_4)
 =\pm\nu_5\eta_8\nu_9\nu_{12}=0.
 \tag{D.2.44s}
\]
Here $\eta_8\nu_9=0$ because it belongs to the stable
fourth stem, which is zero.  The Whitehead identity is
recorded, with its exact Toda equations (5.10) attribution,
in the accompanying author program of
\cite{FGHStricklandToda}.  Finally $i_6(2\iota_5)=0$
by the Moore attachment, so (D.2.44r)--(D.2.44s) give
$i_6(2\nu_5\sigma_8)=0$.  Equation (D.2.44q) proves
(D.2.44d).

Thus the suspension test vanishes on each of the two
candidate bottom images separately.  In particular it
cannot replace $E(i_5\delta)=0$ by the stronger assertion
$E\delta=0$ for a bottom representative $\delta$.
The actual coefficient problem and the full target
James--Hopf image remain as stated in D.2.43.
\end{proof}

\begin{claim}{D.2.45}{Proved}{The obstruction survives the fourteen-skeleton; the exact integral test}
Put $L=Y\langle7\rangle$ and use the quadratic coordinates
$\pi_{15}D_2L=\mathbb Z/8\{c\}\oplus\mathbb Z/2\{u\}$
of D.2.43.  Let $P_j$ be an actual finite-type CW skeleton of
$L$, with the specified nine-skeleton and no cells below seven.
The following statements hold.

The element $\partial(2c)$, which represents $h_Ml_4$,
is nonzero in every $P_j$ for $9\leq j\leq14$.
Its order there is either four or two.  Thus cells of dimensions
twelve through fourteen can at most reduce its order; only
fifteen-cells can first kill it.  Moreover the full-target
question is already determined on the actual fifteen-skeleton:
\[
 \pi_{14}P_{15}\xrightarrow{\ \cong\ }\pi_{14}L,
 \qquad
 \operatorname{im}JH_{P_{15}}=\operatorname{im}JH_L
             \quad\text{in }\pi_{15}D_2L.
 \tag{D.2.45a}
\]

For the full target, set
\[
 a\in H^7(L;\mathbb Z/4),\qquad
 x=\beta_4a\in H^8(L;\mathbb Z),\qquad b=\rho_4x,
 \tag{D.2.45b}
\]
where $a$ is the fundamental coefficient class and $\beta_4$
is the integral connecting homomorphism for multiplication
by four.  There is a well-defined coset
\[
 \theta\in H^{15}(L;\mathbb F_2)/
                      \operatorname{Sq}^1H^{14}(L;\mathbb F_2)
\]
such that, with $i:\mathbb F_2\to\mathbb Z/4$, $i(1)=2$,
\[
 a b=i\theta,\qquad \beta_{\mathbb Z}\theta=x^2.
 \tag{D.2.45c}
\]
Here $\beta_{\mathbb Z}$ is the integral first Bockstein.
If $JH(f)=2c$ for a stable map $f:S^{15}\to L$, then
\[
 \left\langle a b,\rho_4h_{\mathbb Z}(f)\right\rangle=2,
 \qquad
 \left\langle\theta,\rho_2h_{\mathbb Z}(f)\right\rangle=1.
 \tag{D.2.45d}
\]
The latter pairing is independent of the indicated ambiguity.
In particular $x^2=0$ would exclude $2c$ from the full
James--Hopf image. The square is subsequently evaluated in
D.2.50. For a specified sphere, (D.2.45d) does not fix its exact
James--Hopf value; D.2.51 strengthens the existential pairing
condition to an equivalence. D.2.4 remains open.
\end{claim}
\begin{proof}
\textbf{The integral transfer and its complete kernel.}
For a stable map between suspension spectra, the stable Hopf
invariant satisfies
\[
 \operatorname{tr}JH(f)
    =(f\wedge f)\Delta_{\mathrm{source}}-\Delta_{\mathrm{target}}f.
 \tag{D.2.45e}
\]
This is \cite[Theorem A(iii), p.~3, with proof pp.~12--13]
{Klein2026Hopf}; its Theorem B, p.~4, and full proof
pp.~13--16 identify the invariant with the Segal--Snaith
projection used in D.2.43.  The comparison checks weights
one and two and eliminates higher weights.  Thus the operation
in (D.2.45e) is the specified James--Hopf map, not an
unidentified operation with a similar obstruction property.
For a source sphere the reduced stable diagonal is zero,
so the right side is $-\Delta_L f$.

The integral label complex of the nine-skeleton has
$a_7,b_8,r_8,s_9$, with $db=4a$ and $ds=2r$.
In its ordinary tensor square the degree-fifteen cycles are
\[
 t=ab+ba,\qquad v=ar,\qquad w=ra,
\]
with relations $4t=2v=2w=0$.  Explicitly the degree-fifteen
boundary is $d(ab)=-4aa$, $d(ba)=4aa$; the degree-sixteen
boundaries are
\[
 \begin{aligned}
 d(bb)&=4(ab+ba),& d(as)&=-2ar,& d(sa)&=2ra,\\
 d(br)&=4ar,& d(rb)&=4ra,& d(rr)&=0.
 \end{aligned}
\]
Consequently
\[
 H_{15}(L\wedge L;\mathbb Z)
   =\mathbb Z/4\{t\}\oplus\mathbb Z/2\{v,w\}.
 \tag{D.2.45f}
\]
Higher target cells first contribute smash cells in dimension
seventeen, so this group and its named cycles are unchanged.

In the free $C_2$ resolution from D.2.43 the transfer is the
norm in resolution degree zero and is zero in positive
resolution degrees.  It is a chain map, since the norm
annihilates the first resolution differential.  Therefore
\[
 c=e_0ab+2e_1aa\longmapsto t,
 \qquad u=e_0ar\longmapsto v+w.
 \tag{D.2.45g}
\]
D.2.43 identifies the quadratic degree-fifteen homotopy group
with its integral homology.  Thus the composite of transfer
with integral Hurewicz has exactly the kernel
\[
       \ker\bigl(\pi_{15}D_2L\longrightarrow
                      H_{15}(L\wedge L;\mathbb Z)\bigr)
                    =\langle4c\rangle.
 \tag{D.2.45h}
\]
The independent finite certificate
\texttt{code/D45\_integral\_transfer.py} checks the tensor
boundaries, their integral quotient, and every element of
this transfer homomorphism.

\medskip\noindent
\textbf{The sharp remaining skeleton dimension.}
For a target $P_j$ of dimension at most fourteen,
$H_{15}(P_j;\mathbb Z)=0$.  Thus (D.2.45e), followed by
integral Hurewicz, gives
\[
             \operatorname{im}JH_{P_j}\subseteq\langle4c\rangle.
 \tag{D.2.45i}
\]
The quadratic group has already stabilized at the
nine-skeleton, as proved in D.2.43.  The same claim's
natural Goodwillie comparison applies to these six-connected
spaces and identifies $\ker\partial$ with
$\operatorname{im}JH$.  Since $2c\notin\langle4c\rangle$,
its actual boundary is nonzero.  Its original order is four,
so (D.2.45i) allows precisely orders four and two.  This proves
the assertion for all choices of intervening attaching maps;
no calculation of the individual twelve-cell compositions
is needed for nonvanishing.

Cells of dimension at least sixteen do not change
$\pi_{14}$, by cellular approximation applied to maps and
homotopies.  This proves the first part of (D.2.45a).
The cofiber spectrum of $P_{15}\to L$ has no cells below
sixteen.  Hence $\pi_{15}^sP_{15}\to\pi_{15}^sL$ is
surjective.  The map on the quadratic degree-fifteen groups
is an isomorphism, again by D.2.43.  Naturality of $JH$
then proves equality of images in (D.2.45a).  Sixteen-cells
can impose relations in its domain, but cannot add values
to its image.

\medskip\noindent
\textbf{The full-target coefficient pairing.}
The classes $a,b$ in (D.2.45b) are represented in the low
cellular model by the duals of $a_7,b_8$, respectively;
$b$ is zero on $r_8$.  Consequently $a\otimes b$ evaluates
as one on $t$ and zero on $v,w$.  If $JH(f)=2c$, equations
(D.2.45e)--(D.2.45g) give
\[
                  \Delta_*h_{\mathbb Z}(f)=-2t.
\]
Cup and diagonal compatibility gives the first equality
in (D.2.45d), since $-2=2$ modulo four. This pairing alone does
not specify the James--Hopf value: the transfer kernel is
$\langle4c\rangle$, and the quadratic group also has the separate
$u$ coordinate. D.2.51 uses all the mod-two products below to
remove the latter coordinate and proves an existential
sufficiency statement.

The required mod-two product vanishes on the actual $L$.
Here is the low Adams--Hilton check.  The degree-six and
degree-seven loop classes have images in the James tensor
algebra of $\Omega Y$ given by
\[
 u_6\longmapsto aa,\qquad D_7\longmapsto ab+ba,
                   \qquad E_7\longmapsto f.
\]
The four degree-thirteen products $uD$, $Du$, $uE$, $Eu$
have distinctive words $aaab$, $abaa$, $aaf$, $faa$,
respectively.  Their images are
linearly independent.  These are all possible quadratic
terms of degree thirteen in a minimal mod-two
Adams--Hilton model of a six-connected space; cubic terms
start in degree eighteen.  The quadratic differential,
dual to the reduced diagonal, must therefore be zero in
these four directions.  It follows that
\[
       H^7(L;\mathbb F_2)\,H^8(L;\mathbb F_2)=0.
 \tag{D.2.45j}
\]
This uses the actual model and comparison established in
D.2.26 and D.2.42, including its opened
\cite[\S3.1, p.~300]{FGHHalperinLemaire1995} interface.
It does not promote an integral word projection to an
integral diagonal computation.

The coefficient sequence
$0\to\mathbb F_2\xrightarrow{i}\mathbb Z/4\to\mathbb F_2\to0$
and (D.2.45j) now give a class $\theta$ with $ab=i\theta$,
unique modulo $\operatorname{Sq}^1H^{14}L$.
For its integral Bockstein, choose an integral cochain
$A$ representing $a$, with $dA=4B$; then $[B]=x$.
Choose an integral $C^{14}$ with
$dC\equiv A\smile B\pmod2$ and put
\[
 E^{15}=\frac{A\smile B-dC}{2}.
\]
This divides an even integral cochain.  Leibniz gives
$dE=2B\smile B$, so $\theta=[E\bmod2]$ and
$\beta_{\mathbb Z}\theta=x^2$, proving (D.2.45c).

For any integral cycle $z$ the same coefficient sequence gives
\[
                  \langle ab,\rho_4z\rangle
                      =2\langle\theta,\rho_2z\rangle.
 \tag{D.2.45k}
\]
Its value is unchanged on adding a first square: the
homological first Bockstein vanishes on the reduction of
an integral cycle.  Taking $z=h_{\mathbb Z}(f)$ proves
the second equality in (D.2.45d).

Finally $L$ is rationally acyclic and of finite type:
$Y=\Sigma^2K(\mathbb F_2,2)$ has these properties, and
its finite Postnikov base has only finite homotopy groups.
Thus its integral homology in each positive degree is finite.
If $x^2=0$, the coefficient exact sequence gives an
integral lift of $\theta$.  Integral cohomology pairs as
zero with every torsion integral homology class.  Equation
(D.2.45k) then excludes the required value.  Conversely,
the integral universal coefficient sequence shows that
vanishing of this pairing on \emph{all} integral
fifteen-cycles is equivalent to $\beta_{\mathbb Z}\theta=0$:
the classes annihilating reductions of integral cycles
are precisely the reductions of integral degree-fifteen
classes.  Indeed this annihilator is
$\operatorname{Ext}(H_{14}L,\mathbb F_2)$, and
$\operatorname{Ext}(H_{14}L,\mathbb Z)\to
\operatorname{Ext}(H_{14}L,\mathbb F_2)$ is surjective,
since $\operatorname{Ext}_{\mathbb Z}^2=0$.
This does not identify the smaller subgroup
of stable spherical cycles.

The ordinary Postnikov square does not determine the even
coefficient in (D.2.45k).  Its coefficient map here is
$\mathbb Z/4\to\mathbb Z/8$, $1\mapsto4$, and its
cochain is $A\smile dA=4A\smile B\bmod8$;
see \cite[\S2, formula (2.1), p.~899]{Thomas1959Pontrjagin}.
It kills the value two modulo four.  The suspension identity
of that source's Theorem I does not recover this lost
coefficient. The integral square is determined in D.2.50;
the pairing on the stable Hurewicz image remains unresolved.
\end{proof}

\begin{claim}{D.2.46}{Proved}{The universal Moore composite lies in one cyclic subgroup}
Retain the two-local notation of D.2.44, and put
\[
 h=h_Ml_4\in\pi_{14}Y,\qquad
 e=u_4E\epsilon',\qquad
 q=\alpha\sigma'+e,\qquad \alpha=u_4\nu_4.
 \tag{D.2.46a}
\]
Then
\[
             4q=0,\qquad h\in\langle q\rangle,
             \qquad 2h=2q.
 \tag{D.2.46b}
\]
If $q$ has order four, then $h=q$ or $h=-q$.  If $q$ has
order at most two, then $h=0$ or $h=q$.  More precisely, with
compatible sphere-generator signs, the universal composite itself
satisfies
\[
 wl_4\in
 \left\langle i_{5*}\bigl(\nu_4\sigma'
                    -\varepsilon E\epsilon'\bigr)\right\rangle,
 \qquad \varepsilon\in\{1,-1\}.
 \tag{D.2.46c}
\]
This sharpens the two-component span in D.2.44.  It does not
determine the order of $q$ in the full target or choose between
the possibilities in (D.2.46b).
\end{claim}
\begin{proof}
The stronger suspension constraint is obtained in the actual
two-cell complex $A_5$, before mapping to $P_5$.  This order of
the argument is essential: D.2.44 proves that suspension in
$P_5$ alone cannot distinguish the two candidate components.

Use the actual cofiber map $c:P_8\to A_5$ and the characteristic
class $\tau_5$ from (D.2.44g)--(D.2.44n).  In particular,
\[
 A_5=S^4\cup_{2\omega_4}e^8,
 \qquad \omega_4=[\iota_4,\iota_4],
 \qquad c|_{S^7}=\omega_4.
 \tag{D.2.46d}
\]
Write $i:S^4\to A_5$ for its bottom inclusion and $j_*$ for
the map to the relative group of $(A_5,S^4)$.  The mixed relative
Jacobi calculation of (D.2.44n), in this pair itself, gives
\[
 j_*\bigl(3cl_4-2[\tau_5,i]\bigr)=0.
\]
Exactness therefore supplies a sphere class
$\delta_0\in\pi_{14}S^4_{(2)}$ with
\[
             i_*\delta_0=3cl_4-2[\tau_5,i]
                       \quad\text{in }\pi_{14}A_5.
 \tag{D.2.46e}
\]
There is no inference here from an equality only after the
twelve-cell has been attached.

Both terms on the right of (D.2.46e) suspend to zero.  This is
the usual suspension property for the Whitehead product.  For
the other term, the bottom restriction of
$\Sigma c:P_9\to\Sigma A_5$ is $E\omega_4=0$.
The Moore Puppe sequence makes $\Sigma c$ factor through the
actual pinch $p_9:P_9\to S^9$.  The Stiefel attachment has
\[
                 El_4=i_9[\iota_8,\iota_8],
 \tag{D.2.46f}
\]
as in \cite[Lemmas 2.7--2.8, p.~523]{EH5MorisugiMukai2000}.
Its composite with $p_9$ is zero, so $E(cl_4)=0$.
Finally the suspended attaching map of $A_5$ is zero, giving
\[
          \Sigma A_5\simeq S^5\vee S^9
          \quad\text{under the bottom }S^5.
\]
The bottom inclusion thus has a retraction.  Suspending
(D.2.46e) and applying this retraction proves the stronger
sphere-level equality
\[
                 E\delta_0=0\quad\text{in }\pi_{15}S^5_{(2)}.
 \tag{D.2.46g}
\]

For completeness, the entire sphere suspension kernel is explicit.
The primary tables and suspension formulas reproduced in
\cite[Proposition 7 and its proof, p.~7]{FGHStricklandToda}
give
\[
\begin{aligned}
 \pi_{14}S^4_{(2)}
  &=\mathbb Z/8\{a=\nu_4\sigma'\}
    \oplus\mathbb Z/4\{b=E\epsilon'\}
    \oplus\mathbb Z/2\{c_0=\eta_4\mu_5\},\\
 \pi_{15}S^5_{(2)}
  &=\mathbb Z/8\{d=\nu_5\sigma_8\}
    \oplus\mathbb Z/2\{e_0=\eta_5\mu_6\},\\
 Ea&=2d,\qquad Eb=\varepsilon\,2d,\qquad Ec_0=e_0,
 \qquad \varepsilon\in\{1,-1\}.
\end{aligned}
 \tag{D.2.46h}
\]
These are the stated Toda Theorem 7.3 and equations (7.10),
(7.16) inputs in that proof; no general Hopf-composition formula
is needed for this kernel calculation.  A class $ma+nb+tc_0$
belongs to the kernel precisely when $t=0$ and
$m+\varepsilon n\equiv0\pmod4$.  Hence
\[
 \ker\bigl(E:\pi_{14}S^4_{(2)}\to\pi_{15}S^5_{(2)}\bigr)
       =\mathbb Z/8\{a-\varepsilon b\}.
 \tag{D.2.46i}
\]

The actual third cell of the pinch fiber kills $2\tau_5$.
Moreover the full Hopf--Hilton comparison in D.2.44 proves
$wl_4=(A_5\to P_5)cl_4$, with its mixed correction retained
and shown to vanish.  Mapping (D.2.46e) to $P_5$ therefore gives
\[
                        3wl_4=i_{5*}\delta_0.
 \tag{D.2.46j}
\]
Equations (D.2.46g)--(D.2.46i) and invertibility of three on a
two-primary cyclic group prove (D.2.46c).

In $Y$, D.2.44 gives $2e=0$, so the sign of the $e$ term
does not affect the generated subgroup.  It remains to justify
the exponent and the double without discarding a Hilton term.
The sphere $S^7$ is an H-space, and consequently
\[
 [\alpha,\alpha]=\alpha_*[\iota_7,\iota_7]=0.
\]
For precomposition by $\sigma':S^{14}\to S^7$, the only
possible nonlinear Hilton term is a multiple of this
Whitehead square: the two-letter factor begins in dimension
thirteen, while every three-letter factor begins in dimension
nineteen.  Thus precomposition by $\sigma'$ is additive on
the cyclic subgroup generated by $\alpha$.  The already
proved order-four relation for $\alpha$ in D.2.18 yields
\[
              4(\alpha\sigma')=(4\alpha)\sigma'=0,
              \qquad 4q=0.
 \tag{D.2.46k}
\]
Finally the exact Stiefel double from D.2.41 is
\[
              2h=(-2\alpha)\sigma'
                 =-2(\alpha\sigma')=2q.
 \tag{D.2.46l}
\]
The last equality uses (D.2.46k) and $2e=0$.
Together with $h\in\langle q\rangle$, this proves
(D.2.46b) and the two stated order-dependent alternatives.
In particular no unconditional odd-coefficient assertion is
being made when the double vanishes.
\end{proof}

\begin{claim}{D.2.47}{Proved}{The integral half-class has a nonzero mixed differential}
Put $B=K(\FF_2,2)$, $Y=\Sigma^2B$, $Z=Y\langle5\rangle$,
and $L=Y\langle7\rangle$. Write $U_4=\Sigma^2u_2\in H^4(Y;\FF_2)$.
The actual six-type and its loop space are
\[
\begin{aligned}
 P_6Y&\simeq\fib\bigl(\operatorname{Sq}^3:
       K(\FF_2,4)\longrightarrow K(\FF_2,7)\bigr),
 \\
 H=\Omega P_6Y&\simeq\fib\bigl(v_3^2:
       K(\FF_2,3)\longrightarrow K(\FF_2,6)\bigr).
\end{aligned}
 \tag{D.2.47a}
\]
The fundamental class $a_7\in H^7(L;\ZZ/4)$ of D.2.45 is the
restriction of the actual secondary Pontryagin class on $Z$
obtained from a nullhomotopy of $\mathfrak P_2(U_4)$.
This restriction is independent of that nullhomotopy, up to the
chosen sign of the fundamental generator.

Let $b_8=\beta_4a_7\in H^8(L;\ZZ)$, where $\beta_4$ is the
integral coefficient connecting map, and put
$V_4=\beta_{\ZZ}v_3\in H^4(K(\FF_2,3);\ZZ)$.
For the actual fibration $H\longrightarrow L\longrightarrow Y$,
there is a generator $z_8$ such that
\[
 H^8(H;\ZZ)=\ZZ/4\{z_8\},\qquad
 b_8|_H=V_4^2=2z_8\ne0,\qquad
 (b_8^2)|_H=0.
 \tag{D.2.47b}
\]
Its integral Serre spectral sequence has the nonzero mixed
differential
\[
 d_5(z_8)=\operatorname{Sq}^1U_4\otimes V_4
       \in E_5^{5,4}\cong\FF_2.
 \tag{D.2.47c}
\]
Here the base factor has mod-two coefficients because the
fiber coefficient group $H^4(H;\ZZ)$ is $\ZZ/2$.
Thus the half $z_8$ does not extend to $L$, although its
double does. The value of $b_8^2$ on $L$, and the stable-spherical
pairing in D.2.45, remain undecided by this calculation.
\end{claim}
\begin{proof}
\textbf{The six-type and the actual secondary class.}
The already established low groups are
$\pi_4Y=\FF_2$, $\pi_5Y=0$, and $\pi_6Y=\FF_2$.
The possible invariant lies in
\[
 H^7(K(\FF_2,4);\FF_2)
 =\FF_2\{\operatorname{Sq}^3u_4,
                    \operatorname{Sq}^2\operatorname{Sq}^1u_4\}.
 \tag{D.2.47d}
\]
It is nonzero: otherwise the six-type would be a product,
with two independent degree-six cohomology classes, whereas
$H^6Y$ has dimension one. The map $Y\to P_6Y$ is a homology
isomorphism through degree six. Moreover
\[
 \operatorname{Sq}^3U_4=0,
 \qquad
 \operatorname{Sq}^2\operatorname{Sq}^1U_4
      =\Sigma^2(\operatorname{Sq}^2\operatorname{Sq}^1u_2)\ne0.
\]
This identifies the invariant and proves the first equivalence
in \textup{(D.2.47a)}. Looping gives the second, since
$\operatorname{Sq}^3v_3=v_3^2$.
Equivalently, $P_6Y$ is the pullback of
$K(\ZZ/4,6)\to K(\FF_2,6)$ along $\operatorname{Sq}^2u_4$;
the coefficient connecting map is $\operatorname{Sq}^1$.

Thomas's suspension formula and its second-suspension vanishing give
\[
 S(\mathfrak P_2)=p,\qquad S(p)=0,
 \qquad\mathfrak P_2(U_4)=0
 \tag{D.2.47e}
\]
\cite[Theorem~I, pp.~897--898, and Proposition~6.9,
p.~908]{Thomas1959Pontrjagin}.
Choose an actual nullhomotopy, equivalently a lift of $Y$ to
the fiber of $\mathfrak P_2:K(\FF_2,4)\to K(\ZZ/4,8)$.
Passing to fibers over $K(\FF_2,4)$ gives a map
$a_Z:Z\to K(\ZZ/4,7)$.

To check its coefficient, use the generator
$\alpha=u_4\nu_4:S^7\to Y$, which lifts to $L$ by D.2.18.
There is an actual cone map
\[
 \mathbb H P^2=S^4\cup_{\nu_4}e^8
              \longrightarrow \cofib(\alpha)
 \tag{D.2.47f}
\]
with top-cell degree one. The pulled-back bottom mod-two class
has an integral lift on $\mathbb H P^2$, so its Pontryagin
square is the mod-four reduction of the integral square and
evaluates to one on the top cell. The defining relative
pairing for $a_Z$ therefore evaluates oddly on $\alpha$.
Since $L$ is six-connected and $H_7(L;\ZZ)=\ZZ/4\{\alpha\}$,
this is the required fundamental class, after fixing its sign.
Changing the chosen nullhomotopy adds the pullback of a class
in $H^7(Y;\ZZ/4)$. Every such class evaluates to zero on
$\alpha$, which factors through $S^4$ and has zero ordinary
Hurewicz image in $Y$. The coefficient universal-coefficient
theorem, using $H_6L=0$, shows that its entire restriction
to $L$ is zero. This proves the stated choice-independence.

On the $K(\FF_2,3)$ fiber of $Z\to Y$, the class $a_Z$
restricts to $p(v_3)=2v_3V_4$ with coefficients $\ZZ/4$.
The coefficient is exactly that of Thomas's formula (2.1),
not a replacement coefficient map. If $v$ is an integral
cochain lifting $v_3$, with $dv=2V$, then
$d(2vV)=4V^2$. Consequently
\[
                  (\beta_4a_Z)|_{K(\FF_2,3)}=V_4^2.
                  \tag{D.2.47g}
\]
Restricting further to $H$ identifies $b_8|_H$.

\medskip\noindent
\textbf{The integral degree-eight extension on $H$.}
We record the required low Eilenberg--MacLane groups, including
their coefficient normalizations:
\[
\begin{array}{c|c|c}
\text{space}&\text{degree}&\text{integral cohomology}\\ \hline
K(\FF_2,3)&8&\ZZ/2\{V_4^2\}\\
K(\FF_2,5)&6,7,8&
 \ZZ/2\{\beta_{\ZZ}s_5\},\ 0,
 \ZZ/2\{\beta_{\ZZ}\operatorname{Sq}^2s_5\}\\
K(\FF_2,4)&9&\ZZ/4\{\beta_4\mathfrak P_2u_4\}\\
K(\FF_2,6)&7,8,9&
 \ZZ/2\{\beta_{\ZZ}w_6\},\ 0,
 \ZZ/2\{\beta_{\ZZ}\operatorname{Sq}^2w_6\}.
\end{array}
\tag{D.2.47h}
\]
These follow directly from Serre's polynomial generators and
the coefficient Bockstein calculation
\cite[\S2, nos.~7--11, pp.~201--205]{I_Serre1953}.
Here are enough details to fix the higher-order entry.
For $K(\FF_2,3)$, the relevant first Bockstein sends
$v_3\operatorname{Sq}^1v_3$ to
$(\operatorname{Sq}^1v_3)^2$; all other degree-eight
classes are either noncycles or first-Bockstein boundaries.
This leaves precisely the displayed order-two integral class.
The next Bockstein page is zero in degrees seven and eight,
so there is no higher-order summand in that degree.
The degree-seven vanishing and the degree-eight generator for
$K(\FF_2,5)$ are in the stable range: the only appropriate
nonzero first Bockstein is
$\operatorname{Sq}^1\operatorname{Sq}^2s_5
=\operatorname{Sq}^3s_5$. The same argument gives the last row.
For $K(\FF_2,4)$ the two surviving first-Bockstein-page classes
in degrees eight and nine are
\[
 u_4^2,\qquad
 u_4\operatorname{Sq}^1u_4+
       \operatorname{Sq}^4\operatorname{Sq}^1u_4.
\]
The Pontryagin-square cochain formula gives the second
Bockstein from the first to the second; equivalently its
integral connecting class is
$\beta_4\mathfrak P_2u_4$, of exact order four. Thus no
additional integral degree-nine summand is omitted.
Explicitly, for an integral lift $U$ with $dU=2A$, use
$U\smile U+U\smile_1dU$ modulo four. Expanding its
coboundary, dividing by four, and reducing modulo two gives
$UA+A\smile_1A$. Its class is
$u_4\operatorname{Sq}^1u_4+
\operatorname{Sq}^4\operatorname{Sq}^1u_4$, the stated
nonzero second-Bockstein target.

Apply the integral Serre sequence to
$K(\FF_2,5)\to H\to K(\FF_2,3)$.
The only pieces in total degree eight are the first and the
degree-eight fiber entries of \textup{(D.2.47h)}.
The base class cannot be hit: its only possible incoming
sources are $(1,6)$ and $(0,7)$, both zero.
The fiber class is universally transgressive, being the
cohomology suspension of
$\beta_{\ZZ}\operatorname{Sq}^2w_6$ in the path fibration
over $K(\FF_2,6)$. Naturality therefore excludes earlier
mixed differentials, including the dimensionally possible
$d_3$ into $(3,6)$. Its transgression is
\[
 \beta_{\ZZ}\operatorname{Sq}^2(v_3^2)
       =\beta_{\ZZ}\bigl((\operatorname{Sq}^1v_3)^2\bigr)=0,
 \tag{D.2.47i}
\]
since the squared class has the integral lift $V_4^2$.
Thus the integral degree-eight group is an extension of
two copies of $\ZZ/2$.

On $H$, Cartan and the nullity of $v_3^2$ give
$\rho_2(V_4^2)=\operatorname{Sq}^2(v_3^2)=0$.
The integral coefficient sequence makes $V_4^2$ divisible
by two. It remains nonzero by the base-class argument above,
so the extension is necessarily $\ZZ/4$ with $V_4^2=2z_8$.
In particular $4z_8=0$, and
$(V_4^2)^2=4z_8^2=0$. This proves \textup{(D.2.47b)}.
One may choose $z_8=\beta_4c_7$, where $c_7$ is a mod-four
lift of $v_3\operatorname{Sq}^1v_3$ on $H$; its first
Bockstein is zero there. This choice is not asserted to
extend to $L$.

\medskip\noindent
\textbf{The mixed differential.}
First use the universal path fibration over $P_6Y$.
The integral sequence for
$K(\FF_2,6)\to P_6Y\to K(\FF_2,4)$ gives
\[
 H^9(P_6Y;\ZZ)=\ZZ/4\{B_9\},
 \qquad B_9=\beta_4\mathfrak P_2u_4.
 \tag{D.2.47j}
\]
Indeed the base degree-nine group in \textup{(D.2.47h)}
has no possible incoming differential. The only other
degree-nine candidate is the universally transgressive fiber
class $\beta_{\ZZ}\operatorname{Sq}^2w_6$. It dies, since
its transgression is
\[
 \beta_{\ZZ}\operatorname{Sq}^2\operatorname{Sq}^3u_4
   =\beta_{\ZZ}\operatorname{Sq}^4\operatorname{Sq}^1u_4\ne0.
 \tag{D.2.47k}
\]
Its mod-two reduction is $(\operatorname{Sq}^1u_4)^2$,
nonzero on the base and unaffected by any earlier incoming
differential. There are no mixed total-degree-nine terms.

Cohomology suspension, the coefficient connecting map, and
Thomas's formula commute. Hence in the universal path fibration
$H\to P(P_6Y)\to P_6Y$ the suspension image is
\[
 \sigma^*B_9=\beta_4p(v_3)=V_4^2=2z_8.
 \tag{D.2.47l}
\]
For a path fibration the transgressive part of fiber cohomology
is exactly this suspension image. Thus $z_8$ must support an
earlier differential. The base is three-connected, and
$H^3(H;\ZZ)=H^5(H;\ZZ)=0$, whereas
$H^4(H;\ZZ)=\ZZ/2\{V_4\}$. The only possible earlier
target is consequently
\[
 E_5^{5,4}=H^5(P_6Y;\ZZ/2)\otimes\{V_4\}
          =\FF_2\{\operatorname{Sq}^1u_4\otimes V_4\}.
\]
No earlier differential hits this group. Necessarily $d_5z_8$
is its nonzero generator.

Finally $H\to L\to Y$ is the pullback of that path fibration
along the actual Postnikov projection. The map carries
$\operatorname{Sq}^1u_4$ to
$\operatorname{Sq}^1U_4=\Sigma^2\operatorname{Sq}^1u_2\ne0$.
Again no earlier differential can hit bidegree $(5,4)$.
Naturality proves \textup{(D.2.47c)}. The square $b_8^2$
has zero restriction to the fiber, but can still be a
nonzero class of positive base filtration; no multiplicative
extension is suppressed in this conclusion.
\end{proof}

\begin{claim}{D.2.48}{Proved}{Twice the integral bottom class has mixed filtration four}
Retain $H\to L\xrightarrow{\pi}Y$ and the coefficient classes
of D.2.47. Let $F^\bullet$ be the decreasing integral Serre
filtration on $H^8(L;\ZZ)$. Then
\[
\begin{gathered}
 H^8(L;\ZZ)=\ZZ/4\{b_8\},\qquad F^5H^8(L;\ZZ)=0,\\
 F^4H^8(L;\ZZ)=\langle2b_8\rangle\cong\ZZ/2,\\
 \operatorname{gr}^4(2b_8)=U_4\otimes V_4,
 \qquad H^8(L;\ZZ)/F^4\cong\langle2z_8\rangle\cong\ZZ/2.
\end{gathered}
\tag{D.2.48a}
\]
The entire pullback $\pi^*:H^8(Y;\ZZ)\to H^8(L;\ZZ)$
is zero. Thus $2b_8$ is a nonzero mixed class, rather than
the pullback of the degree-eight class on $Y$.
These conclusions do not determine $b_8^2$ in degree sixteen.
\end{claim}
\begin{proof}
Put $P=P_6Y$ and first consider the universal path fibration
$H\to P(P)\to P$. In addition to D.2.47c, its integral
spectral sequence has
\[
 d_5(u_4\otimes V_4)=2B_9,
 \qquad B_9=\beta_4\mathfrak P_2u_4
                 \in H^9(P;\ZZ)=\ZZ/4.
 \tag{D.2.48b}
\]
Here and below $u_4\otimes V_4$ denotes the class in
$H^4(P;H^4(H;\ZZ))=H^4(P;\ZZ/2)$; no integral lift of
$u_4$ is assumed.

We justify the coefficient and every possible competing
differential. The low fiber groups are
\[
\begin{gathered}
 H^3(H;\ZZ)=H^5(H;\ZZ)=0,
 \qquad H^4(H;\ZZ)=\ZZ/2\{V_4\},\\
 H^7(H;\ZZ)=\ZZ/2\{e_7\},
 \qquad e_7=\beta_{\ZZ}\operatorname{Sq}^2
                         \operatorname{Sq}^1v_3.
\end{gathered}
 \tag{D.2.48c}
\]
For the last group, the sequence
$K(\FF_2,5)\to H\to K(\FF_2,3)$ has just the displayed
base class in total degree seven. Its only possible incoming
transgression is $\beta_{\ZZ}(v_3^2)=0$, because
$v_3^2=\operatorname{Sq}^1\operatorname{Sq}^2v_3$.
The other vanishings are in the same low range of Serre's
generators used in D.2.47h. Moreover $e_7$ is the suspension
of the actual class
\[
 e_8=\beta_{\ZZ}\operatorname{Sq}^2
                      \operatorname{Sq}^1u_4\in H^8(P;\ZZ).
 \tag{D.2.48d}
\]
This base class survives in $P$: the only possible incoming
fiber transgression is
$\beta_{\ZZ}\operatorname{Sq}^3u_4=0$.
Consequently $e_7$ is universally transgressive. Its earlier
$d_4$ into bidegree $(4,4)$ is therefore zero.

The group $E_5^{4,4}=\FF_2\{u_4\otimes V_4\}$ has
no other possible earlier incoming differential. By D.2.47,
$d_5z_8$ is nonzero, so the remaining first-column group
of degree eight is $\langle2z_8\rangle\cong\ZZ/2$.
The only incoming differentials which can hit base degree
nine are $d_5$ from $(4,4)$ and $d_9$ from that remaining
first-column group. All other possibilities are excluded by
three-connectivity of $P$ and \textup{(D.2.48c)}.
Since the path total space is contractible, a homomorphism
from the remaining group of order two cannot by itself kill
$H^9(P;\ZZ)=\ZZ/4$. Thus $d_5$ necessarily hits its unique
order-two subgroup, proving \textup{(D.2.48b)}; the later
$d_9$ hits the quotient.

Pull back along the actual projection $Y\to P$. The class
$B_9$ pulls to $\beta_4\mathfrak P_2(U_4)=0$ by D.2.47e.
Thus the actual differential out of $E_5^{4,4}$ is zero:
\[
                  d_5(U_4\otimes V_4)=0.
                  \tag{D.2.48e}
\]
No later outgoing differential from that bidegree is possible.
The only possible incoming one was the earlier $d_4(e_7)$,
which is zero by universal transgressivity. Hence this mixed
class survives nontrivially to $E_\infty$.

The same universally transgressive fiber class instead has
the actual differential
\[
 d_8(e_7)=\beta_{\ZZ}\operatorname{Sq}^2
                       \operatorname{Sq}^1U_4
       =\Sigma^2\beta_{\ZZ}w_5\ne0,
 \qquad w_5=\operatorname{Sq}^2\operatorname{Sq}^1u_2.
 \tag{D.2.48f}
\]
It kills the whole base group $H^8(Y;\ZZ)=\ZZ/2$.
That group has no earlier possible incoming differential:
the candidates use base degrees one through three or the
vanishing fiber groups in \textup{(D.2.48c)}.
Finally $2z_8$ survives in the actual first column, either
by naturality of $d_9$ or by the existing global class $b_8$
with restriction $2z_8$.

The only possible total-degree-eight terms were $(0,8)$,
$(4,4)$ and $(8,0)$. The last is killed as above. Independently,
six-connectivity, rational acyclicity, and $H_7(L;\ZZ)=\ZZ/4$ give
$H^8(L;\ZZ)=\ZZ/4\{b_8\}$ by the universal-coefficient
theorem. Its restriction has image $\langle2z_8\rangle$,
so its nonzero kernel is exactly $\langle2b_8\rangle$.
This is the surviving filtration-four group, and proves
\textup{(D.2.48a)} and the asserted zero base pullback.
\end{proof}

\begin{claim}{D.2.49}{Proved}{The universal composite has zero integral loop Hurewicz}
Put $X_0=\Sigma B$, $Y=\Sigma X_0$, and $L=Y\langle7\rangle$.
Retain $e=u_4E\epsilon'$ and $h=h_Ml_4$ from D.2.46.
There is an actual class $t\in\pi_{10}X_0$ such that
\[
                    e=(Et)\nu_{11}\quad\text{in }\pi_{14}Y.
                    \tag{D.2.49a}
\]
The map $Et:S^{11}\to Y$ lifts uniquely through $L$ up to
homotopy, and this factorization gives the unique lift of $e$.
For the unique lift $\widetilde h\in\pi_{14}L$ one has
\[
 h_{\ZZ}\bigl(\operatorname{adj}(\widetilde h)\bigr)=0
                  \quad\text{in }H_{13}(\Omega L;\ZZ).
 \tag{D.2.49b}
\]
Consequently, if $a_6$ represents the specified integral bottom
generator in a Pontryagin chain model for $\Omega L$ and
$db_7=4a_6$, then
\[
 z_L=[a_6b_7-b_7a_6]\in H_{13}(\Omega L;\ZZ)
                         \quad\text{satisfies}\quad 2z_L=0.
 \tag{D.2.49c}
\]
The doubled class is independent of these chain choices.
These statements do not assert $h=0$ or evaluate the order of $q$
in D.2.46. The integral cohomology consequence is derived in D.2.50.
\end{claim}
\begin{proof}
\textbf{An actual injective low suspension.}
Write $u_3:S^3\to X_0$ for the bottom map, so $Eu_3=u_4$.
We first verify
\[
 \pi_6X_0=\ZZ/2\{\beta_0\},\qquad
 E\beta_0=2\alpha\in\pi_7Y=\ZZ/4\{\alpha\}.
 \tag{D.2.49d}
\]
The first Ganea evaluation construction for
$B=\Omega K(\FF_2,3)$ gives the actual fibration
\[
 G=\Sigma(B\wedge B)\longrightarrow\Sigma B
                    \longrightarrow K(\FF_2,3).
\]
Its fiber is the join $B*B$: pulling the two cone pieces
back to the path-space model gives the homotopy pushout of
the two projections $B\times B\rightrightarrows B$.
Thus $\pi_6G\to\pi_6X_0$ is an isomorphism.

Here is the complete low attaching calculation for $G$.
Use the standard classes $a_2,b_3,c_4$ of $B$, with
$\operatorname{Sq}_*^1b=a$ and
$\operatorname{Sq}_*^2c=a$.
The mod-two groups of $G$ in degrees five, six and seven
have respective bases
\[
                aa;\qquad ab,ba;\qquad ac,ca,bb.
\]
The first Bockstein takes both $ab,ba$ to $aa$, and takes
$bb$ to $ab+ba$.  The second homology square takes each
of $ac,ca,bb$ to $aa$.
In integral cellular coordinates put $B_6=ab$ and
$C_6=ba-ab$, so $dB_6=2aa$ and $dC_6=0$.
The six-skeleton is $M(\ZZ/2,5)\vee S^6_C$ and
\[
 \pi_6(G^{(6)})=\ZZ/2\{i\eta_5\}\oplus\ZZ\{C_6\}.
 \tag{D.2.49e}
\]
The actual integral cycles $ac,ca$ have zero $C_6$ degree;
their square detector therefore attaches each by $i\eta_5$.
The $bb$ cell attaches by $\pm2C_6+i\eta_5$.
These are all components of (D.2.49e), which is already in
the stable range.  Higher cells cannot change $\pi_6$.
It follows that $\pi_6G=\ZZ/2$, and its integral Hurewicz
map to $H_6G=\ZZ/2$ is an isomorphism.

In the first James pair $(JX_0,X_0)$, relative Hurewicz gives
\[
 \pi_6(JX_0,X_0)=H_6(JX_0,X_0;\ZZ)
       =H_3(X_0;\ZZ)\otimes H_3(X_0;\ZZ)=\ZZ/2.
 \tag{D.2.49f}
\]
This is the first relative word of length two in the reduced
product filtration \cite[Sections 3 and 7]{I_James1955}.
The class $\alpha=u_4\nu_4$ of D.2.18 maps to its nonzero
element: naturality from the bottom $S^3$ identifies the
relative map with reduction of the Hopf invariant one of
$\nu_4$.  Exactness, together with
$\pi_7Y=\ZZ/4\{\alpha\}$, makes the image of suspension
exactly $\langle2\alpha\rangle$.
The preceding computation of its source proves (D.2.49d).
In particular this suspension is injective.

\textbf{The full Toda indeterminacy supplies a sphere factor.}
D.2.44 proved the actual relation $u_4E\nu'=0$.
Equation (D.2.49d) now implies
\[
                         u_3\nu'=0\quad\text{in }\pi_6X_0.
                         \tag{D.2.49g}
\]
The exact sphere input is
\[
 \epsilon'\in\{\nu',2\nu_6,\nu_9\}_3
                  \subseteq\{\nu',2\nu_6,\nu_9\}.
 \tag{D.2.49h}
\]
The indexed definition, its convention, and the same
zero-first-entry argument are explicit in the opened author
source \cite[Section 2.5, p.~4, and Proposition 2, p.~6]
{FGHStricklandToda}; its attribution is Toda, Section VI(iv).
Only membership in the ordinary bracket is needed here.
Postcomposition and (D.2.49g) give
\[
 u_3\epsilon'\in\{0,2\nu_6,\nu_9\}
          =\pi_{10}X_0\circ\nu_{10}.
 \tag{D.2.49i}
\]
The equality includes every choice of nullhomotopy.
With zero first map, the bracket contains zero, its first
indeterminacy summand is zero, and its entire other summand
is $[S^{10},X_0]\circ\Sigma\nu_9$.
Thus there is an actual $t$ with $u_3\epsilon'=t\nu_{10}$;
suspending proves (D.2.49a).  No numerical value of this
remaining composition is inferred.

\textbf{Hurewicz vanishing in the connected cover.}
For $j>6$, the map $\pi_jL\to\pi_jY$ is an isomorphism.
In particular $Et$ and its homotopies lift uniquely through
$L$, and the lifted factorization in (D.2.49a) is the unique
lift of $e$.  Its once-adjoint factors through $\Omega S^{11}$.
The once-adjoint of the lift of $\alpha\sigma'$ similarly
factors through $\Omega S^7$.
The integral James tensor calculation gives
\[
 H_{13}(\Omega S^{11};\ZZ)=0,
                 \qquad H_{13}(\Omega S^7;\ZZ)=0,
 \tag{D.2.49j}
\]
since the respective tensor generators have degrees ten and
six.  Hence both summands of the lift of
$q=\alpha\sigma'+e$ have zero integral loop Hurewicz.
D.2.46 places $h$ in $\langle q\rangle$; the isomorphism
on $\pi_{14}$ and the Hurewicz homomorphism therefore prove
(D.2.49b).  The factor through $S^{11}$ is essential:
factoring $e$ through $S^4$ only gives a conclusion in
$\Omega Y$, because its bottom map does not lift to $L$.

\medskip\noindent
\textbf{The intrinsic doubled commutator.}
Let $M_4=M(\mathbb Z/4,7)$ and use a Pontryagin differential
graded algebra for the actual loop space $\Omega L$.  Choose a
cycle $a$ of degree six representing the specified bottom
generator and a chain $b$ of degree seven with $db=4a$.
The low homology calculation in D.2.22, after the two homology
suspensions to $L$, gives
\[
 H_6(\Omega L;\mathbb Z)=\mathbb Z/4,
 \qquad H_7(\Omega L;\mathbb Z)=\mathbb Z/2.
 \tag{D.2.49e1}
\]
These are homology groups; the integral cohomology group
$H^8(L;\mathbb Z)$ in D.2.47--48 is instead $\mathbb Z/4$.
Put $z_L=[ab-ba]\in H_{13}(\Omega L;\mathbb Z)$.  We prove
the exact comparison
\[
 2z_L=h_{\mathbb Z}\bigl(\operatorname{adj}(h_Ml_4)\bigr).
 \tag{D.2.49e2}
\]

First this doubled class is intrinsic.  The identity
$d(b^2)=4(ab-ba)$ gives $4z_L=0$.  If $a'=a+dv$, with
$|v|=7$, take $b'=b+4v$.  The graded Leibniz rule gives the
literal identity
\[
 a'b'-b'a'-(ab-ba)=d(vb+bv+4v^2).
 \tag{D.2.49e3}
\]
Any further change of $b'$ is a cycle $r$ of degree seven.
By \textup{(D.2.49e1)} write $2r=dR$.  Then
\[
                 2(a'r-ra')=d(a'R-Ra').
 \tag{D.2.49e4}
\]
Thus $2z_L$ is independent of all representatives and top-chain
choices.  Odd orientation changes also leave it fixed because
its order is at most two.

We next compare with the specified actual Moore map, retaining
its entire top cell.  Put $P_5=P^5(2)$ and
$\alpha_P=i_4\nu_4:S^7\to P_5$.  The Hopf--Hilton argument
of D.2.17 gives $4\alpha_P=0$ already in $P_5$, so choose an
extension $\mu_0:M_4\to P_5$.  Let
$c_{2,4}:P^8(2)\to M_4$ be the suspension cofiber map with
bottom degree two and top degree one.  The maps
$\mu_0c_{2,4}$ and $-w=-[i_4,\mathrm{id}_{P_5}]$ have the
same bottom restriction $2\alpha_P$.  Their difference factors
through the top pinch $q_2:P^8(2)\to S^8$ by the Moore
Puppe sequence.  Since $q_4c_{2,4}=q_2$, a top-sphere change
of $\mu_0$ makes
\[
                         \mu_0c_{2,4}=-w
 \tag{D.2.49e5}
\]
as actual homotopy classes.  Here $c_{2,4}$ is a suspension
map, so the comparison respects the coaction without a Hilton
correction.  Composing with $P_5\to Y$ and taking the unique
lift through the six-type quotient defines an actual map
$\mu:M_4\to L$.

The exact natural boundary calculation of D.2.43 gives
\[
 c_{2,4}l_4=\pm\partial(2c_4),\qquad
 \ker(\pi_{14}M_4\to\pi_{14}^{s}M_4)
           =\mathbb Z/8\{\partial c_4\}.
 \tag{D.2.49e6}
\]
The source normalization in this assertion can also be checked
integrally.  The James chain algebra of $\Omega M_4$ is
$T(A_6,B_7)$ with $dB=4A$.  In degree thirteen its only words
are $AB,BA$, and the only degree-fourteen boundary is $d(B^2)$.
Consequently
\[
 H_{13}(\Omega M_4;\mathbb Z)
                    =\mathbb Z/4\{AB-BA\}.
 \tag{D.2.49e7}
\]
There is a stably trivial relative Whitehead generator of
$\pi_{14}M_4$ whose mod-two loop Hurewicz image is the nonzero
commutator.  This is precisely the modulus-four,
bottom-dimension-seven case of
\cite[Theorem~1.7(ii)--(iii), Corollary~1.8, and proof,
pp.~6--9]{EH5ZhuPanComplexes2026}; its order-eight
pinch-kernel presentation is also given in
\cite[Lemma~2.2, p.~6]{EH5ZhuPanClassification2026}.
By \textup{(D.2.49e6)} it is an odd multiple of $\partial c_4$.
Equation \textup{(D.2.49e7)} therefore gives
\[
 h_{\mathbb Z}(\operatorname{adj}\partial c_4)
                  =t(AB-BA),\qquad t\text{ odd}.
 \tag{D.2.49e8}
\]
Only this parity assertion is needed; no numerical choice of
the relative generator is imposed.

The actual loop map $\Omega\mu$ sends these chains to an
allowed pair $a,b$ with $db=4a$.  Naturality,
\textup{(D.2.49e5)}--\textup{(D.2.49e8)} identify the
right side of \textup{(D.2.49e2)} with $\pm2t z_L=2z_L$.
The chain identities \textup{(D.2.49e3)}--\textup{(D.2.49e4)}
make this the same doubled class for every allowed target
choice.  The loop-Hurewicz vanishing just proved now gives
\[
                              2z_L=0.
 \tag{D.2.49e9}
\]
This supplies a chain $W_{14}$ with $dW=2(ab-ba)$ in the
actual loop-chain model.  It does not specify the compatible
higher relation on $W$ or the degree-fifteen Adams--Hilton
differential.  Thus no vanishing of $b_8^2$, of its half-product
pairing on the stable Hurewicz image, or of $h_Ml_4$ in
homotopy is inferred from this conclusion.

\end{proof}

\begin{claim}{D.2.50}{Proved}{The integral bottom square is nonzero}
Let $L=Y\langle7\rangle$, and retain the specified classes
$a_7\in H^7(L;\ZZ/4)$ and $b_8=\beta_4a_7\in H^8(L;\ZZ)$.
Then
\[
              0\ne b_8^2\in H^{16}(L;\ZZ),
              \qquad 2b_8^2=0.
              \tag{D.2.50a}
\]
There is an integral class $\omega\in H_{15}(L;\ZZ)$ such that
\[
 \big\langle a_7\smile\operatorname{red}_4b_8,
                     \operatorname{red}_4\omega\big\rangle=2,
 \qquad
 \big\langle\theta,\operatorname{red}_2\omega\big\rangle=1.
 \tag{D.2.50b}
\]
Here $\theta$ is the half-product class of D.2.45, with its
stated $\operatorname{Sq}^1$ ambiguity. The class $\omega$ is
torsion of order at least four. The argument does not show that
$\omega$ is in the stable Hurewicz image; consequently it does
not settle the stable-spherical pairing or the actual composite
$h_Ml_4$.
\end{claim}
\begin{proof}
\textbf{An actual integral cobar model.}
Choose a six-reduced simplicial model $K$ of the six-connected
space $L$. Such a model can be obtained from a minimal Kan
model: minimality and the vanishing homotopy groups through
degree six make all simplices in those degrees degenerate.
The same model-existence statement is recorded explicitly in
\cite[proof of Lemma 5.4.2, p.~43 in the opened author version]
{I_ABFJ2022}. Let $C=C_*(K;\ZZ)$ be normalized chains, with the
Alexander--Whitney coproduct. Thus $\overline C_i=0$ for
$i\le6$. The canonical integral cobar construction is
\[
 \Omega C=T(s^{-1}\overline C),\qquad
 d_\Omega=d_1+d_2,
 \quad d_1=-s^{-1}d_Cs,
 \quad d_2=(s^{-1}\otimes s^{-1})\overline\Delta s.
 \tag{D.2.50c}
\]
It is naturally quasi-isomorphic as a chain algebra to the
integral chains of the Kan loop group, hence computes the
actual Pontryagin algebra of $\Omega L$; see
\cite[pp.~3 and 5]{FGHHessParentScottTonks2005}.
The coproduct in \textup{(D.2.50c)} is the ordinary
Alexander--Whitney coproduct. Its dual therefore computes the
ordinary cup product on $L$.

Choose integral cochains $A\in C^7$ and $B\in C^8$ with
\[
       A\bmod4\text{ representing }a_7,
       \qquad dA=4B,
       \qquad dB=0.
       \tag{D.2.50d}
\]
Then $B$ represents $b_8$. Choose a cycle $x\in C_7$
representing a generator of $H_7(L;\ZZ)=\ZZ/4$, and
$y\in C_8$ with $d_Cy=4x$. The integer $A(x)$ is odd, and
\[
                         B(y)=A(x).
                         \tag{D.2.50e}
\]
Both $x$ and $y$ have zero reduced coproduct by degree.
Put $u=s^{-1}x$ and $v=-s^{-1}y$. Then
$d_\Omega u=0$ and $d_\Omega v=4u$. D.2.49 and its
chain-choice independence imply
\[
       2(uv-vu)=d_\Omega W
                    \quad\text{for some }W\in(\Omega C)_{14}.
       \tag{D.2.50f}
\]
In particular this is a boundary in the canonical model,
not only after projecting to another loop space.

\textbf{The nullchain gives an integral fifteen-cycle.}
Every positive generator of $\Omega C$ has degree at least
six. Thus $W$ has only word lengths one and two:
\[
                    W=s^{-1}w+W_2,
             \qquad w\in C_{15},\quad W_2\in T^2(s^{-1}\overline C).
             \tag{D.2.50g}
\]
The word-length-one part of \textup{(D.2.50f)} gives
$d_Cw=0$. Hence $w$ determines an actual integral homology
class $\omega=[w]\in H_{15}(L;\ZZ)$.

Define an integer-valued functional $\lambda$ on word-length
two and total cobar degree thirteen by
\[
 \lambda\big((s^{-1}c)(s^{-1}d)\big)=
 \begin{cases}
 A(c)B(d),& |c|=7,\ |d|=8,\\
 0,&\text{otherwise}.
 \end{cases}
 \tag{D.2.50h}
\]
Equations \textup{(D.2.50e)--(D.2.50f)} give
\[
                \lambda\big(2(uv-vu)\big)
                  =-2A(x)^2\equiv2\pmod4.
                  \tag{D.2.50i}
\]
The possible unshifted degree pairs in $W_2$ are exactly
$(7,9)$, $(8,8)$ and $(9,7)$. For the first pair,
$\lambda(d_1W_2)=0$ because $B(d_Cd)=(dB)(d)=0$.
For the middle pair its only possibly nonzero contribution
is $-A(d_Cc)B(d)=-4B(c)B(d)$, by \textup{(D.2.50d)}.
The last pair contributes zero for degree reasons. Therefore
\[
                            \lambda(d_1W_2)\equiv0\pmod4.
                            \tag{D.2.50j}
\]
There are no words of length three in the relevant degrees.
The length-two part of \textup{(D.2.50f)} is consequently
$d_2(s^{-1}w)+d_1W_2=2(uv-vu)$.
In the sign convention of \textup{(D.2.50c)}, desuspending
the first factor of degree seven gives
\[
 \lambda\big(d_2(s^{-1}w)\big)
             =-\langle A\smile B,w\rangle.
 \tag{D.2.50k}
\]
Combining \textup{(D.2.50i)--(D.2.50k)} proves
$\langle A\smile B,w\rangle\equiv2\pmod4$.
This is the first pairing in \textup{(D.2.50b)}. The sign
does not affect the value two modulo four.

\textbf{The half-product and its integral Bockstein.}
D.2.45 constructs $\theta$ with
\[
 a_7\smile\operatorname{red}_4b_8=i\theta,
 \qquad \beta_{\ZZ}\theta=b_8^2,
 \qquad i:\FF_2\longrightarrow\ZZ/4,\quad i(1)=2.
 \tag{D.2.50l}
\]
Its $\operatorname{Sq}^1H^{14}$ ambiguity evaluates to zero
on the reduction of every integral cycle. Thus the first
pairing implies the second pairing in \textup{(D.2.50b)}.
If $b_8^2=0$, the coefficient exact sequence would give an
integral lift of $\theta$. Every integral degree-fifteen
cohomology class on $L$ is torsion, since $L$ is rationally
acyclic and of finite type. Such a class evaluates to zero
on every integral cycle, contrary to
\textup{(D.2.50b)}. Hence $b_8^2\ne0$.
It has exact order two because it is the integral
Bockstein in \textup{(D.2.50l)}.

Finally, $H_{15}(L;\ZZ)$ is a finite torsion group.
Here is the needed mod-four lift of $\theta$ explicitly.
The half-product construction in D.2.45 supplies an integral
cochain $E\in C^{15}$ with $dE=2B^2$ and $E\bmod2$
representing $\theta$. The mod-two square vanishing proved
in D.2.42 allows an integral cochain $J\in C^{15}$ with
$dJ=B^2+2K$ for an integral cochain $K\in C^{16}$.
Consequently
\[
 \theta_4=[E-2J\bmod4]\in H^{15}(L;\ZZ/4),
 \qquad\operatorname{red}_2\theta_4=\theta,
 \quad d(E-2J)=-4K.
 \tag{D.2.50m}
\]
Thus $\theta$ evaluates to zero on every integral class
of order at most two: the evaluation of $\theta_4$ on such
a class lies in $\{0,2\}$ and reduces to zero modulo two.
The class $\omega$ from \textup{(D.2.50b)} therefore has
order at least four. The construction made $w$ an integral
cycle, with no claim that it is represented by a stable map
from a sphere.
\end{proof}

\begin{claim}{D.2.51}{Proved}{The exact stable-spherical criterion and the three remaining orders}
Put $L=Y\langle7\rangle$ and retain the quadratic coordinates
\[
 \pi_{15}D_2L=\ZZ/8\{c\}\oplus\ZZ/2\{u\},
 \qquad h=h_Ml_4=\partial(2c),
 \tag{D.2.51a}
\]
with the compatible orientation of D.2.43.  For the actual
James--Hopf map of that claim,
\[
          \operatorname{im}\bigl(JH:\pi_{15}^sL
                      \longrightarrow\pi_{15}D_2L\bigr)
                         \subseteq\langle2c\rangle.
          \tag{D.2.51b}
\]
Let $a_7\in H^7(L;\ZZ/4)$, $b_8=\beta_4a_7$ and $\theta$
be the half-product class of D.2.45.  The following conditions
are equivalent:
\[
\begin{aligned}
 &h=0\text{ in }\pi_{14}L;\\
 &\text{there is }f\in\pi_{15}^sL\text{ with }
       \langle a_7\smile\operatorname{red}_4b_8,
                    \operatorname{red}_4h_{\ZZ}(f)\rangle=2;\\
 &\text{there is }f\in\pi_{15}^sL\text{ with }
       \langle\theta,\operatorname{red}_2h_{\ZZ}(f)\rangle=1.
\end{aligned}
\tag{D.2.51c}
\]
The last pairing is independent of the
$\operatorname{Sq}^1H^{14}(L;\FF_2)$ ambiguity of $\theta$.
More precisely, the possibilities are exactly
\[
\begin{array}{c|c|c}
 \operatorname{im}JH&\text{order of }h&
     \theta\text{ on the stable Hurewicz image}\\ \hline
 0&4&0\\
 \langle4c\rangle&2&0\\
 \langle2c\rangle&1&\text{nonzero}.
\end{array}
\tag{D.2.51d}
\]
Here order one means that $h$ is zero.  The integral class
constructed in D.2.50 satisfies the indicated pairing, but
its stable sphericity remains unproved.  Thus this criterion
does not choose a row of \textup{(D.2.51d)} or settle D.2.4.
\end{claim}
\begin{proof}
\textbf{The mod-two diagonal removes the independent coordinate.}
We use the full product vanishing in D.2.45j,
\[
                 H^7(L;\FF_2)\,H^8(L;\FF_2)=0,
                 \tag{D.2.51e}
\]
not only the vanishing of the particular reduction of
$a_7\operatorname{red}_4b_8$.  Since $L$ is six-connected,
the only summands of $H_{15}(L\wedge L;\FF_2)$ have
bidegrees $(7,8)$ and $(8,7)$.  Cup--diagonal duality and
\textup{(D.2.51e)} therefore imply
\[
 \overline\Delta_*\operatorname{red}_2z=0
             \quad\text{for every }z\in H_{15}(L;\ZZ).
 \tag{D.2.51f}
\]
All groups in these degrees are finite-dimensional over
$\FF_2$, so testing against the indicated cup products
detects the entire reduced diagonal.

To apply this to the actual stable operation, recall the
complete transfer computation in D.2.45f--h.  Its low
integral label chains are $a_7,b_8,r_8,s_9$, with
$db=4a$ and $ds=2r$.  They give
\[
 \begin{aligned}
 H_{15}(L\wedge L;\ZZ)
   &=\ZZ/4\{t=ab+ba\}\oplus\ZZ/2\{v=ar,w=ra\},\\
 \operatorname{tr}(c)&=t,\qquad
 \operatorname{tr}(u)=v+w.
 \end{aligned}
 \tag{D.2.51g}
\]
Here and below the transfer is followed by integral
Hurewicz.  Write $\bar a,\bar b,\bar r$ for the corresponding
mod-two homology classes.  The classes $\bar b,\bar r$
are independent in $H_8(L;\FF_2)$: the former is the
Moore-four Tor class and the latter is the reduction of
the integral order-two class.  Consequently the four
ordered tensors
\[
 \bar a\otimes\bar b,\quad \bar b\otimes\bar a,
 \quad \bar a\otimes\bar r,\quad\bar r\otimes\bar a
 \tag{D.2.51h}
\]
are linearly independent by the field K\"unneth isomorphism.
For $k\in\ZZ/8$ and $\epsilon\in\ZZ/2$, reduction of
\textup{(D.2.51g)} is zero precisely when
\[
             k\equiv0\pmod2,\qquad\epsilon=0.
             \tag{D.2.51i}
\]

For any stable sphere $f$, the transfer identity D.2.45e
gives
\[
       \operatorname{tr}JH(f)=-\overline\Delta_*h_{\ZZ}(f).
       \tag{D.2.51j}
\]
The sphere's reduced stable diagonal is zero.  The cited
identity is the actual Segal--Snaith James--Hopf operation,
with its comparison already checked in D.2.45 against
\cite[Theorems A(iii) and B]{Klein2026Hopf}.
Equations \textup{(D.2.51f)--(D.2.51j)} show that every
value $JH(f)=kc+\epsilon u$ has even $k$ and zero
$\epsilon$.  This proves \textup{(D.2.51b)} without
assuming anything about the stable Hurewicz image itself.

\textbf{An existential test, with the pointwise ambiguity retained.}
The coefficient functional
$a_7\otimes\operatorname{red}_4b_8$ evaluates as one on
$t$ and as zero on $v,w$.  If $JH(f)=kc$ with $k$ even,
\textup{(D.2.51j)} and D.2.45k give
\[
 \begin{aligned}
 \langle a_7\smile\operatorname{red}_4b_8,
                   \operatorname{red}_4h_{\ZZ}(f)\rangle
        &=-k\pmod4,\\
 2\langle\theta,\operatorname{red}_2h_{\ZZ}(f)\rangle
        &=-k\pmod4.
 \end{aligned}
 \tag{D.2.51k}
\]
The ambiguity of $\theta$ pairs to zero on an integral
reduction, since its homological first Bockstein is zero.
Thus the two pairing conditions in \textup{(D.2.51c)}
are equivalent, and each says exactly that
$JH(f)$ is $2c$ or $6c$.

The map $JH$ induces a homomorphism on the displayed
homotopy groups.  If the value is $6c$, replacing $f$ by
$3f$ changes it to $18c=2c$.  Hence existence of either
pairing is equivalent to $2c\in\operatorname{im}JH$.
The natural exact sequence of D.2.43 identifies the latter
condition with $\partial(2c)=0$, proving
\textup{(D.2.51c)}.  For a specified $f$, the pairing
still does not distinguish $2c$ from $6c$; the transfer
kernel $\langle4c\rangle$ from D.2.45 has not disappeared.
The new sufficiency is existential and uses the odd
rescaling of the sphere class.

Finally the image is a subgroup of the cyclic group
$\langle2c\rangle\cong\ZZ/4$, so it is one of
$0$, $\langle4c\rangle$ and $\langle2c\rangle$.
Exactness computes the order of the image of $2c$ in
each quotient as four, two and one, respectively.
Equation \textup{(D.2.51k)} computes the last column of
\textup{(D.2.51d)}.  This proves the complete table.
The exact finite certificate
\texttt{code/D51\_stable\_transfer\_range.py} checks all sixteen
quadratic-group elements, the mod-two range restriction, both
pairing-two values and their odd rescaling, and the three quotient
orders. It selects no actual James--Hopf image.
\end{proof}

\begin{claim}{D.2.52}{Proved}{The half-product is not a positive primary operation}
For $L=Y\langle7\rangle$, there is an integral class
$\omega\in H_{15}(L;\ZZ)$ such that
\[
 \langle\theta,\operatorname{red}_2\omega\rangle=1,
 \qquad
 \operatorname{Sq}_*^i(\operatorname{red}_2\omega)=0
                           \quad(i>0).
 \tag{D.2.52a}
\]
In particular the half-product of D.2.45--50 does not belong to
the positive Steenrod-operation image:
\[
 \theta\notin
 \operatorname{Sq}^1H^{14}(L;\FF_2)
 +\operatorname{Sq}^2H^{13}(L;\FF_2)
 +\operatorname{Sq}^4H^{11}(L;\FF_2)
       =\bigl(\mathcal A^{>0}H^*(L;\FF_2)\bigr)^{15}.
 \tag{D.2.52b}
\]
The class $\omega$ is torsion of order at least four.
It is obtained from the top integral cycle of an explicit
finite fifteen-dimensional space mapping to $L$.
This does not assert that $\omega$ is stably spherical:
higher differentials and integral coefficient obstructions
in the stable Hurewicz spectral sequence remain possible.
\end{claim}
\begin{proof}
\textbf{The suspended sphere factor has zero mod-two Hurewicz.}
Put $X_0=\Sigma B$, so $Y=\Sigma X_0$. In the notation
of D.2.22 and D.2.31,
\[
 \pi_{11}Y=\FF_2\{p,s,x,y\},\qquad
 \operatorname{im}\big(E:\pi_{10}X_0\to\pi_{11}Y\big)
                                  \subseteq\langle p,s\rangle.
 \tag{D.2.52c}
\]
Here is the exact argument for the containment. Under once
adjunction to $\Omega Y$, the classes $p,s$ have zero
mod-two Hurewicz image, while the images of $x,y$ are
respectively
\[
              af+fa,\qquad aab+baa
       \quad\text{in }H_{10}(\Omega\Sigma X_0;\FF_2).
       \tag{D.2.52d}
\]
These are the two independent expressions in D.2.31i;
the letters in this display are the suspended James labels.
Their word lengths are two and three. The once-adjoint of
a suspended map $Et$ is the composite of $t$ with the James
unit $X_0\to\Omega\Sigma X_0$. Its homology image has
word length one. Thus neither coefficient of $x,y$ can
occur, proving \textup{(D.2.52c)}.

The connected-cover interface is needed here as well.
For $E_L=\Omega^2L$, the actual calculation D.2.22 gives
\[
                  h_{\ZZ}(p)=4N',\qquad h_{\ZZ}(s)=0
                          \quad\text{in }H_9(E_L;\ZZ),
                  \tag{D.2.52e}
\]
where $N'$ has order eight. Both classes therefore have
zero mod-two Hurewicz already in $E_L$, not merely after
mapping to $\Omega^2Y$. Naturality of the two homology
suspensions proves that their unique lifts to $L$ have
zero ordinary mod-two Hurewicz in degree eleven.
In particular the lift
$\widetilde{Et}:S^{11}\to L$ of the factor in D.2.49a
satisfies
\[
                        h_{\FF_2}(\widetilde{Et})=0.
                        \tag{D.2.52f}
\]
No value for the integral homology suspension of $4N'$
is inferred in this step.

\textbf{An actual finite cone representing a detecting cycle.}
Let $M_4=M(\ZZ/4,7)$, let
$\delta=\partial(2c_4):S^{14}\to M_4$, and let
$\mu:M_4\to L$ be the actual Moore map of D.2.49.
D.2.46 and D.2.49 supply an integer $n$, with orientations
chosen consistently, for which
\[
 \mu\delta=n\bigl(\widetilde\alpha\sigma'
                        +\widetilde{Et}\nu_{11}\bigr)
                               \quad\text{in }\pi_{14}L.
 \tag{D.2.52g}
\]
Only the existence of $n$ is used; its value and the
homotopy class on either side remain undetermined.
Form
\[
\begin{split}
 P_0&=M_4\vee S^7\vee S^{11},\\
 r&=\iota_M\delta-n\iota_7\sigma'-n\iota_{11}\nu_{11},\\
 P&=P_0\cup_r e^{15}.
\end{split}
 \tag{D.2.52h}
\]
The maps $\mu,\widetilde\alpha,\widetilde{Et}$ on the
three summands extend to an actual map $g:P\to L$ by
\textup{(D.2.52g)}. Let $w\in H_{15}(P;\ZZ)=\ZZ$
be the oriented top cellular class, and put $\omega=g_*w$.

We verify its mixed cup coefficient by an actual cofiber
map. Express $r$ as the pinch $S^{14}\to\bigvee_3S^{14}$
followed by the three separate component maps. The induced
cofiber square gives
\[
 P\longrightarrow
 C(\delta)\vee C(-n\sigma')\vee C(-n\nu_{11}),
 \tag{D.2.52i}
\]
with degree one on each top relative cell. The low classes
$g^*a_7$ and $g^*\operatorname{red}_4b_8$ are the pullbacks
of $a_M+a_S$ and $b_M$, respectively. The sphere $S^7$
contributes to $a_7$, but its integral coefficient Bockstein
is zero; the $S^{11}$ summand contributes to neither class.
All cross products between wedge summands in
\textup{(D.2.52i)} vanish.

The Moore-cone coefficient is exactly
\[
       \langle a_M\smile b_M,[C(\delta)]_{15}\rangle=2
                                      \quad\text{in }\ZZ/4.
       \tag{D.2.52j}
\]
This is an actual Stiefel-cone comparison, before composing
the Moore map into the target. The actual
map $c_{2,4}:M(\ZZ/2,7)\to M_4$ has bottom degree two
and top degree one, and
$c_{2,4}l_4=\pm\delta$ by D.2.43.
The source cone is $C(l_4)=V_{9,2}$ by
\cite[paragraph before Lemma 2.7 and Lemma 2.8(ii), p.~523]
{EH5MorisugiMukai2000}. The resulting cone map has top degree one.
The fundamental mod-four bottom class pulls to twice the
mod-two degree-seven class, while the integral Bockstein
reduces to its degree-eight partner. For the sphere bundle
$S^7\to V_{9,2}\to S^8$, the Euler number is two.
Its mod-two Serre sequence therefore gives the two classes
$x_7,y_8$ with $\operatorname{Sq}^1x_7=y_8$ and
$x_7y_8$ the top class. The pulled-back mod-four product
is consequently twice that top class, giving
\textup{(D.2.52j)}.
Equations \textup{(D.2.52i)--(D.2.52j)} now give
\[
 \langle a_7\smile\operatorname{red}_4b_8,
                       \operatorname{red}_4\omega\rangle=2.
 \tag{D.2.52k}
\]
The half-product identity implies the first assertion
of \textup{(D.2.52a)}. D.2.50m then shows that this integral
class has order at least four.

\textbf{Every positive primary square annihilates the class.}
The finite source in \textup{(D.2.52h)} has
\[
 H_{13}(P;\FF_2)=H_{14}(P;\FF_2)=0,
 \qquad H_{11}(P;\FF_2)=\FF_2\{S^{11}\}.
 \tag{D.2.52l}
\]
Thus the first and second homology squares of its top class
are zero. The fourth square lies in the displayed
$S^{11}$ line, whose image under $g$ is zero by
\textup{(D.2.52f)}. The eighth square is zero by instability:
its dual is $\operatorname{Sq}^8$ on degree-seven
cohomology. Larger powers of two cannot contribute.
Naturality therefore annihilates
$\operatorname{red}_2\omega$ under
$\operatorname{Sq}_*^1,\operatorname{Sq}_*^2,
\operatorname{Sq}_*^4,\operatorname{Sq}_*^8$.
Since the Steenrod algebra is generated by the squares
whose indices are powers of two, this proves the full
annihilation in \textup{(D.2.52a)}.

Finally, six-connectivity and instability show that the
entire degree-fifteen positive-operation image is exactly
the sum in \textup{(D.2.52b)}. Every element of that sum
pairs to zero with $\operatorname{red}_2\omega$, whereas
$\theta$ pairs to one. This proves
\textup{(D.2.52b)} without asserting a stable representative
for the top integral cycle of $P$.
\end{proof}

\begin{claim}{D.2.53}{Proved}{The epsilon-prime target term is zero}
Let $X_0=\Sigma K(\FF_2,2)$, $Y=\Sigma X_0$, and
$u_3:S^3\to X_0$, $u_4=Eu_3:S^4\to Y$ be the bottom maps.
Then
\[
                         u_4E\epsilon'=0\quad\text{in }\pi_{14}Y.
                         \tag{D.2.53a}
\]
More precisely, the factor $t$ in D.2.49 can be chosen with
the compatible Toda extension and nullhomotopies so that
\[
 u_3\epsilon'=\pm t\nu_{10},\qquad
 t\in\{u_3,\nu',2\nu_6\},\qquad Et=0.
 \tag{D.2.53b}
\]
This identifies a chosen factorization; it does not assert
that every map $t$ satisfying the first equality is zero
after suspension.

Consequently the generator and the actual Moore obstruction
of D.2.46 have the sharper form
\[
 \begin{gathered}
 q=\alpha\sigma',\qquad 4q=0,\qquad
 h=h_Ml_4=n\alpha\sigma'\text{ for some integer }n,\\
                         2h=2\alpha\sigma'.
 \end{gathered}
 \tag{D.2.53c}
\]
The same zero holds for the unique lift of the term in
\textup{(D.2.53a)} to $L=Y\langle7\rangle$.
Neither the value nor the order of $\alpha\sigma'$ in
the full target is determined.  If it has order four,
$h=\pm\alpha\sigma'$; if it has order at most two,
$h$ may be zero or $\alpha\sigma'$.
\end{claim}
\begin{proof}
All Toda brackets below are ordinary brackets in the
two-local pointed homotopy category.  The actual homotopy
groups of $Y$ in question are two-primary, so this does
not lose any part of the target assertion.
We use the actual identities
\[
             u_3\nu'=0,\qquad u_4E\nu'=0,
                         \qquad 2\pi_{11}Y=0,
             \tag{D.2.53d}
\]
proved in D.2.49, D.2.44 and D.2.31, respectively.

\textbf{A compatible representative, not an arbitrary factor.}
The indexed definition of $\epsilon'$ used in D.2.49 lies
in the ordinary bracket
\[
                    \epsilon'\in\{\nu',2\nu_6,\nu_9\}.
                    \tag{D.2.53e}
\]
This is the definition and indexed-bracket convention in
\cite[Section 2.5, p.~4, and the accompanying author program]
{FGHStricklandToda}; the original Toda attribution is
Section VI(iv).  Choose an extension and a coextension
realizing this specified value:
\[
 \begin{gathered}
 C=C(2\nu_6)=S^6\cup_{2\nu_6}e^{10},\qquad
 \bar\nu':C\longrightarrow S^3,\qquad
 \widetilde\nu_9:S^{13}\longrightarrow C,\\
 \bar\nu'|_{S^6}=\nu',\qquad
 \bar\nu'\widetilde\nu_9=\epsilon',\qquad
 p\widetilde\nu_9=\nu_{10},
 \quad p:C\longrightarrow S^{10}.
 \end{gathered}
 \tag{D.2.53f}
\]
The harmless common sign depends on the cone convention.
Use a nullhomotopy of $u_3\nu'$ in
\textup{(D.2.53d)} to factor the actual composite as
\[
                         u_3\bar\nu'\simeq t p.
                         \tag{D.2.53g}
\]
The factor obtained from these two compatible
nullhomotopies is precisely, up to that common sign,
a value of $\{u_3,\nu',2\nu_6\}$.  This is the
extension--coextension definition of the ordinary bracket:
the nullhomotopy of $\nu'(2\nu_6)$ defining $\bar\nu'$
and the chosen nullhomotopy of $u_3\nu'$ are glued along
their common face.  Thus this conclusion does not follow
merely from the existence of some pinch factor in a
Puppe sequence.  Composing \textup{(D.2.53g)} with the
\emph{same} $\widetilde\nu_9$ gives
$u_3\epsilon'=\pm t\nu_{10}$.
Suspension of this chosen cone diagram gives
\[
                    Et\in\mathcal B_2
                         :=\{u_4,E\nu',2\nu_7\}
                              \subseteq\pi_{11}Y.
                    \tag{D.2.53h}
\]
No inverse suspension or identification of unrelated
nullhomotopy choices is used.

\textbf{The suspended bracket is the singleton zero.}
The exact two-primary sphere groups needed here are
\[
                 \pi_9(S^3)_{(2)}=0,\qquad
                 \pi_{10}(S^3)_{(2)}=0,\qquad
                 \pi_{11}(S^4)_{(2)}=0.
                 \tag{D.2.53i}
\]
They are explicitly tabulated in the opened author
program accompanying \cite{FGHStricklandToda}, in its
sixth- and seventh-stem blocks, with the attributions
to Toda Proposition 5.11, p.~46, and Proposition 5.15,
p.~50.  The last equality also follows from the first
relevant terms of the Hopf-fibration sequence
$S^3\to S^7\to S^4$, since
$\pi_{11}S^7=\pi_4^s\mathbb S=0$ and
$\pi_{10}(S^3)_{(2)}=0$.

Since $\nu_7=E\nu_6$, the first equality in
\textup{(D.2.53i)} gives
\[
                    (E\nu')\nu_7=E(\nu'\nu_6)=0.
\]
Hence the smaller bracket
$\mathcal B_1=\{u_4,E\nu',\nu_7\}$ is defined.
Precompose its last sphere by the degree-two map
$2\iota_{10}$.  Naturality of the ordinary cone diagram
gives
\[
                    \mathcal B_1\circ(2\iota_{11})
                               \subseteq\mathcal B_2.
                    \tag{D.2.53j}
\]
This is precomposition on a sphere, not an assumed
additivity formula for the middle bracket entry.
Every value on the left is zero by $2\pi_{11}Y=0$.
Thus $0\in\mathcal B_2$.

The complete ordinary indeterminacy of $\mathcal B_2$ is
\[
                 u_4\circ\pi_{11}S^4
                     +\pi_8Y\circ(2\nu_8).
                 \tag{D.2.53k}
\]
The first term is zero by \textup{(D.2.53i)}.
For every $\gamma:S^8\to Y$, the second term is zero
by the exact right-composition identity
\[
              \gamma(2\nu_8)
                  =\gamma\nu_8(2\iota_{11})
                  =2(\gamma\nu_8)=0.
\]
Again only the exponent of $\pi_{11}Y$ is used; no
distributivity in the variable $\gamma$ is needed.
The ordinary bracket is a coset of this zero
indeterminacy, so
\[
                            \mathcal B_2=\{0\}.
                            \tag{D.2.53l}
\]
Equations \textup{(D.2.53h)--(D.2.53l)} prove $Et=0$
for the compatible factor selected above.  Suspending
$u_3\epsilon'=\pm t\nu_{10}$ now proves
$u_4E\epsilon'=\pm(Et)\nu_{11}=0$.

The connected-cover map $L\to Y$ is an isomorphism on
homotopy groups above degree six; the unique lifted
term is consequently zero as well.  Substituting this
value into D.2.46 proves \textup{(D.2.53c)} and its
two order-dependent alternatives.  It does not
evaluate the remaining $\alpha\sigma'$ composite.
\end{proof}

\begin{claim}{D.2.54}{Proved}{A detecting class survives to the eighth stable differential}
Put $L=Y\langle7\rangle$ and use the reduced two-local
stable homotopy Atiyah--Hirzebruch spectral sequence
\[
 E^2_{p,q}=H_p(L;\pi_q^s\mathbb S)
       \Longrightarrow \pi_{p+q}^sL,
 \qquad d_r:E^r_{p,q}\longrightarrow E^r_{p-r,q+r-1}.
 \tag{D.2.54a}
\]
There is an integral class $\omega\in H_{15}(L;\ZZ)$ with
\[
 \langle\theta,\operatorname{red}_2\omega\rangle=1,
 \qquad \omega\in E^8_{15,0}.
 \tag{D.2.54b}
\]
For the integer $n$ in the actual equality
$h_Ml_4=n\widetilde\alpha\sigma'$ supplied by D.2.53,
it can be chosen so that
\[
 d_8\omega=[2n\bar\alpha]\in E^8_{7,7},
 \qquad E^2_{7,7}=H_7(L;\ZZ/16)=\ZZ/4\{\bar\alpha\}.
 \tag{D.2.54c}
\]
Here $\widetilde\alpha:S^7\to L$ is the lift of $\alpha$,
and $\bar\alpha$ denotes its integral bottom Hurewicz class
with the indicated coefficients.  The brackets in
\textup{(D.2.54c)} mean its image in the actual page-eight
quotient; they do not assert that this image is nonzero.

The only possibly nonzero incoming differentials that
form this quotient are
\[
 d_5:E^5_{12,3}\longrightarrow E^5_{7,7},\qquad
 d_6:E^6_{13,2}\longrightarrow E^6_{7,7},\qquad
 d_7:E^7_{14,1}\longrightarrow E^7_{7,7}.
 \tag{D.2.54d}
\]
Each source is its actual page group, with all earlier
differentials taken.  In particular $E^8_{7,7}$ is a
quotient of $\ZZ/4$, and the incoming $d_2$ is zero.

Let
\[
 H=E^8_{15,0}\subseteq H_{15}(L;\ZZ),\qquad
 H_0=\ker\bigl(\langle\theta,\operatorname{red}_2(-)\rangle:
                         H\longrightarrow\FF_2\bigr).
 \tag{D.2.54e}
\]
Then the complete existence test, including a change of
detecting cycle, is
\[
             h_Ml_4=0
       \quad\Longleftrightarrow\quad
             [2n\bar\alpha]\in d_8(H_0).
 \tag{D.2.54f}
\]
Consequently a proper quotient $E^8_{7,7}$ of $\ZZ/4$
would imply $h_Ml_4=0$.  A nonzero $h_Ml_4$ requires
$E^8_{7,7}=\ZZ/4$, odd $n$, and $d_8(H_0)=0$.
Neither the three incoming images in \textup{(D.2.54d)}
nor $d_8(H_0)$ is evaluated here, so D.2.4 remains open.
\end{claim}
\begin{proof}
\textbf{An actual source with only the final possible obstruction.}
Let $M_4=M(\ZZ/4,7)$ and
$\delta=\partial(2c_4):S^{14}\to M_4$.
The actual Moore map $\mu:M_4\to L$ of D.2.49,
together with D.2.53, gives
\[
                         \mu\delta=n\widetilde\alpha\sigma'.
                         \tag{D.2.54g}
\]
Choose compatible orientations and form
\[
 P=(M_4\vee S^7)\cup_{\iota_M\delta-n\iota_7\sigma'}e^{15}.
 \tag{D.2.54h}
\]
The maps $\mu$ and $\widetilde\alpha$ extend to an actual
map $g:P\to L$ by \textup{(D.2.54g)}.  If $w$ is the
oriented top integral cellular cycle, set $\omega=g_*w$.

The same actual cone comparison as D.2.52 applies with
the $S^{11}$ summand omitted.  More explicitly, pinching
the attaching sphere gives a cofiber map
\[
                     P\longrightarrow C(\delta)\vee C(-n\sigma')
                     \tag{D.2.54i}
\]
of degree one on each top relative cell.  The bottom
mod-four class can have a component on the extra $S^7$,
but its integral Bockstein has none there.  All cross
products of the two wedge summands vanish.  The actual
Stiefel-cone map $C(l_4)=V_{9,2}\to C(\delta)$ has top
degree one and bottom Moore degrees two and one.
The cup computation D.2.52j therefore gives
\[
 \langle a_7\smile\operatorname{red}_4b_8,
                  \operatorname{red}_4\omega\rangle=2,
 \qquad
 \langle\theta,\operatorname{red}_2\omega\rangle=1.
 \tag{D.2.54j}
\]
This is a cup calculation on actual spaces and their
actual cone maps; no stable splitting is used for it.

For the stable differential calculation, D.2.43 places
$\delta$ in the kernel of stabilization.  The opened
sphere convention gives $E^\infty\sigma'=2\sigma$
up to a harmless odd unit
\cite[introduction, p.~515]{EH5MorisugiMukai2000}.
Choose a stable nullhomotopy of the $M_4$ component of
the attachment in \textup{(D.2.54h)}.  Homotopy invariance
of the stable cofiber gives a filtered equivalence
\[
 \Sigma^\infty P\simeq
 \Sigma^\infty M_4\vee C(2n\sigma:S^{14}\to S^7),
 \tag{D.2.54k}
\]
identical on the old skeleton and of degree one on the
top quotient.  Thus it preserves the specified top edge
class and the cellular filtration relevant to its
differentials.  In the two-cell summand there is no cell
between dimensions seven and fifteen.  The top class
survives to page eight, and its $d_8$ is $2n\sigma$ on
the bottom seven-cell.  The separate Moore summand
contributes no differential on this top class.
Naturality for the actual map $g$ proves
\textup{(D.2.54b)--(D.2.54c)}.

We use the cellular exact-couple convention of
\cite[Part III, Section 7]{EH5Adams1974AHSS}:
$E^1$ is the cellular chain group with sphere-stem
coefficients and $d_1$ is its ordinary cellular
boundary.  There is therefore no additional unrecorded
$d_1$ in \textup{(D.2.54c)}.  The space is simply
connected, so the coefficient systems are constant.

\textbf{Every incoming source at the final target.}
Six-connectivity makes all outgoing differentials from
$(7,7)$ zero.  An incoming $d_r$, for $2\le r<8$,
would have source $(7+r,8-r)$.  The relevant coefficient
groups are
\[
 \pi_6^s\mathbb S=\ZZ/2,\quad
 \pi_5^s\mathbb S=\pi_4^s\mathbb S=0,\quad
 \pi_3^s\mathbb S=\ZZ/8,\quad
 \pi_2^s\mathbb S=\pi_1^s\mathbb S=\ZZ/2.
 \tag{D.2.54l}
\]
The zero fourth and fifth stems eliminate $d_3$ and
$d_4$ at this target.  The remaining possible $d_2$
has source $H_9(L;\ZZ/2)$, and it is zero for the
following exact naturality reason.

Let $j$ be the connective image-of-$J$ spectrum at two.
The defining fiber sequence
$j\to ko\to\Sigma^4ksp$ and the opened computation
\cite[Theorem 1.2 and Remark 1.3, pp.~2--3]
{EH5BrunerRognes2022} give
\[
                 \pi_6j=0,
       \qquad \pi_7^s\mathbb S\xrightarrow{\cong}\pi_7j=\ZZ/16
       \tag{D.2.54m}
\]
under the unit.  In the cited computation the latter
isomorphism is the four surviving $h_0$ multiples of
the image of $h_3$.  Applying the unit to the cellular
AHSS maps the source of $d_2:E^2_{9,6}\to E^2_{7,7}$
to zero and maps its target isomorphically to the
corresponding $j$-homology target.  The commuting square
therefore forces this $d_2$ to vanish.  The two-complete
convention of that source does not alter either finite
coefficient group in \textup{(D.2.54m)}.

Only the three differentials in
\textup{(D.2.54d)} remain.  This argument has not
identified their later-page sources with raw homology
groups, nor computed their images.  Successively
quotienting the target by those images proves the
assertion about $E^8_{7,7}$.

\textbf{Changing the detecting cycle and convergence.}
The sphere spectrum is connective.  Hence the bottom
coefficient row has no incoming differentials, and
$E^8_{15,0}$ is canonically a subgroup of integral
homology, as in \textup{(D.2.54e)}.  The ambiguity of
$\theta$ by $\operatorname{Sq}^1H^{14}$ evaluates to
zero on reductions of integral cycles.  Since
\textup{(D.2.54j)} gives one detecting member $\omega$,
all detecting classes on this page are exactly
$\omega+H_0$.

For $r\ge9$, a differential from $(15,0)$ would land
in filtration at most six and is zero by connectivity.
Thus the only remaining obstruction for a member of
$E^8_{15,0}$ is its $d_8$.  The connective AHSS has a
finite filtration in total degree fifteen, and its
edge image is precisely the image of
$\pi_{15}^sL\to H_{15}(L;\ZZ)$; this is the finite-range
convergence in the same cellular exact-couple construction
\cite[Part III, Section 7]{EH5Adams1974AHSS}.
It follows that a detecting stable Hurewicz class exists
if and only if
\[
                      d_8\omega\in d_8(H_0).
                      \tag{D.2.54n}
\]
Here a minus sign makes no difference to the subgroup
condition.  D.2.51 identifies this existence assertion
with $h_Ml_4=0$, proving \textup{(D.2.54f)}.

Every proper quotient of $\ZZ/4$ kills $2\bar\alpha$.
In that case the particular $\omega$ already has zero
$d_8$, and \textup{(D.2.54f)} implies $h_Ml_4=0$.
If instead $h_Ml_4\ne0$, the target must be $\ZZ/4$
and $n$ must be odd.  A nonzero subgroup of $\ZZ/4$
contains $2\bar\alpha$, so the same criterion then
forces $d_8(H_0)=0$.  These are necessary conditions,
not a computation selecting their truth values.
\end{proof}

\begin{claim}{D.2.55}{Proved}{The incoming fifth differential does not change the bottom sigma group}
In the stable Atiyah--Hirzebruch spectral sequence of D.2.54,
\[
 d_5:E^5_{12,3}\longrightarrow E^5_{7,7}=\ZZ/4
       \quad\hbox{is zero},\qquad E^6_{7,7}=\ZZ/4.
 \tag{D.2.55a}
\]
More precisely, the bottom stable subgroup $\ZZ/4$ in the
eleven-stage injects after all twelve-cells are attached.
The doubled free $N'$-lift $p$ of D.2.22 satisfies the actual
identity
\[
                         (2p)\nu_{11}=0
                         \quad\hbox{in }\pi_{14}P_{11}.
                         \tag{D.2.55b}
\]
Here $P_{11}$ is the actual twice-shifted finite model of
D.2.22, before the twelve-cell relations are imposed.
Consequently only incoming $d_6$ and $d_7$ can change the
target in D.2.54 before its final $d_8$.
Neither those two images nor $d_8(H_0)$ is computed here;
no value of $h_Ml_4$ is asserted.
\end{claim}
\begin{proof}
\textbf{The finite stage and its complete coordinates.}
The evaluation $\Sigma^2E\to L=Y\langle7\rangle$ is
twelve-connected: its factorization through
$\Sigma\Omega L$ has successive connectivities twelve
and thirteen, as in the double-counit argument of D.2.43.
Thus it has a relative CW model with cells only in
dimensions at least thirteen.  The twice-shifted stable
model of D.2.22 through dimension twelve realizes
uniquely in the ordinary stable range and provides the
required stages $P_j$.  A later thirteen-cell can kill
the suspended $u_5^2$ sphere; it cannot alter the
twelve-stage calculation below.

All homotopy groups in the proof are stable unless an
ordinary cone calculation is specified.  The groups
$\pi_{11}P_{11}$ are also ordinary groups, since $P_{11}$
is six-connected.  Put
\[
 \begin{split}
 K&=(M_4\vee S^8_r)\cup_{i_7\eta_7+2r}e^9,
       \qquad M_4=M(\ZZ/4,7),\\
 K/S^7&=S^8_b\vee M(\ZZ/2,8)_r.
 \end{split}
 \tag{D.2.55c}
\]
This is the shifted core of D.2.22t2; the separate
$S^9_R$ has zero stable $\pi_{14}$.
The exact shifted component calculation gives
\[
 \pi_{11}K=\ZZ/8\{s_0\}\oplus\ZZ/2\{z\},\qquad
 2s_0=m_4,\quad z=r\nu_8,\quad
 \operatorname{pinch}(m_4)=2\nu_8.
 \tag{D.2.55d}
\]
After the ten- and eleven-cells, the complete old torsion
subgroup and free Hurewicz lattice are
\[
 \begin{split}
 V&=(\ZZ/2)^5\{s,z,\tau,w\eta_{10},j\eta_{10}\},\\
 h(p)&=4N',\qquad h(q)=2Q,\qquad
 h(x)=K_1,\qquad h(y)=K_5.
 \end{split}
 \tag{D.2.55e}
\]
These are D.2.22t9--t12, using the actual $a=b=0$
proved in D.2.31.  In particular $K_1,K_5$ attach
trivially and can be taken as free eleven-spheres.
Changing $x,y$ with the same full Hurewicz coordinates
adds elements of $V$ and disappears on doubling.

The bottom subgroup is
\[
 \operatorname{im}(\pi_{14}^sS^7\to\pi_{14}^sK)
                     =\ZZ/4\{i_7\sigma\}.
                     \tag{D.2.55f}
\]
The Moore attachment quotients the seventh stable stem
$\ZZ/16$ by four.  The additional nine-cell does not
kill anything further in this group: its possible
$\pi_{14}$ composite is
$(i_7\eta_7+2r)\nu_8^2=0$, since
$\eta\nu^2=0$ and $2\nu^2=0$.
The map $\pi_{14}^sK\to\pi_{14}^sP_{11}$ is injective:
ten- and eleven-cell attaching spheres have zero
$\pi_{14}^s$, by $\pi_5^s\mathbb S=\pi_4^s\mathbb S=0$.
The corresponding stable cofiber exact sequences justify
this old-subgroup injection without assumptions about
the later quotient $P_{11}/S^7$.

\textbf{The doubled $N'$-lift.}
Orient the $N'$ eleven-cell so its attaching map is
$\alpha\nu_7$, where $\alpha:S^7\to P_{11}$ is the
bottom Moore map.  It supplies an actual map
\[
             j_\nu:C\nu_7=S^7\cup_{\nu_7}e^{11}
                         \longrightarrow P_{11}.
                         \tag{D.2.55g}
\]
There is a unique $v_8\in\pi_{11}C\nu_7$ of top
degree eight: the pinch image is $8\ZZ$ and
$\pi_{11}S^7=0$.  Both $2p$ and $j_\nu v_8$ have
integral Hurewicz $8N'$.  Their difference lies in $V$.
Its second- and fourth-square cone detectors vanish.
For $2p$ this follows from the degree-two cone map;
for $j_\nu v_8$ it follows by naturality, since
$C\nu_7$ has no cohomology in degrees eight or ten.
The five detector vectors of D.2.22t10 are injective
on $V$.  Therefore
\[
                            2p=j_\nu v_8.
                            \tag{D.2.55h}
\]
This fixes the coextension, rather than discarding an
unmeasured old-cell difference.

The opened ordinary sphere convention gives
\[
       \{\nu_7,8\iota_{10},\nu_{10}\}=\{4\sigma'\}.
       \tag{D.2.55i}
\]
Indeed the author's program in \cite{FGHStricklandToda}
declares $\sigma'''=\{\nu_5,8\iota_8,\nu_8\}$ and
$E\sigma'''=2\sigma''$, $E\sigma''=2\sigma'$,
attributing them to Toda Lemmas 5.13--5.14.
Twice suspending the same ordinary cone choices gives
\textup{(D.2.55i)}; its indeterminacy
$\nu_7\pi_{14}S^{10}+\pi_{11}S^7\nu_{11}$ is zero.
Thus the actual cone composition is
\[
 (2p)\nu_{11}=j_\nu(v_8\nu_{11})
             =\alpha(4\sigma')=0.
             \tag{D.2.55j}
\]
For the last equality, $\alpha$ factors through $M_4$,
and $(4\iota_7)\sigma'=4\sigma'$ because $S^7$ is an
$H$-space.  Its Whitehead-square Hilton term is zero.
This proves the actual assertion \textup{(D.2.55b)}
before any twelve-cell imposing $2p=0$ is attached.

We also need the pure Moore composite
\[
                            m_4\nu_{11}=0.
                            \tag{D.2.55k}
\]
The cone formula places it in the bottom image of
$\{4\iota_7,2\nu_7,\nu_{10}\}$.
The defined bracket $\{4\iota_3,\nu',\nu_6\}$ lies
in $\pi_{10}(S^3)_{(2)}=0$; four suspensions put zero
in the former bracket.  Its only indeterminacy is
$\langle4\sigma'\rangle$, whose bottom image in $M_4$
is zero by the same $S^7$ argument.
The exact sphere vanishings and $E^4\nu'=2\nu_7$
are the opened formulas in \cite{FGHStricklandToda},
also used in D.2.53.

\textbf{All remaining orders, with the even normalization retained.}
Equations \textup{(D.2.55d), (D.2.55k)} imply that
$s_0\nu_{11}$ and $z\nu_{11}$ have order at most two.
Their images under the original quotient in
\textup{(D.2.55c)} have coordinates
\[
  (\nu_8^2,*),\qquad(0,i_8\nu_8^2),
  \qquad
  \pi_{14}^sM(\ZZ/2,8)=\ZZ/2\{i_8\nu_8^2\}.
  \tag{D.2.55l}
\]
The unspecified coordinate is immaterial: these two
vectors are independent.  Their span therefore misses
the bottom subgroup \textup{(D.2.55f)}.  This conclusion
passes to $P_{11}$ by the old-$K$ injection already
proved, without asserting that $P_{11}/S^7$ retains
the original quotient's group.
Also
\[
 \tau\nu_{11}=w\eta_{10}\nu_{11}
              =j\eta_{10}\nu_{11}=0,
 \tag{D.2.55m}
\]
since $\pi_{14}^sM(\ZZ/2,9)=\pi_5^s(\mathbb S/2)=0$
and $\eta\nu=0$.

For $q$ we must retain the even coordinate change
allowed in D.2.22t11.  Write the original $Q$ attachment
as $R\eta_9+2c n$, where $c\in\{0,1\}$ and
$n=\alpha\nu_7$.  Sliding $Q$ by $2cN'$ gives
$Q'=Q-2cN'$ attached by $R\eta_9$.  The resulting
actual $C\eta_9$ map provides $q_0$ with
\[
                  h(q_0)=2Q-4cN',
       \qquad q-q_0-cp\in V.
       \tag{D.2.55n}
\]
Pinch gives an isomorphism
$\pi_{14}C\eta_9\to\pi_{14}S^{11}=\ZZ/8$:
the surrounding fourth and fifth stable stems vanish.
Thus $q_0\nu_{11}$ has order four, since its pinch is
$2\nu_{11}$.  Equations \textup{(D.2.55e), (D.2.55j)}
now imply
\[
                            4q\nu_{11}=0.
                            \tag{D.2.55o}
\]
This accounts for the possible even $n$ correction.
The original $Q$ pinch of $q\nu_{11}$ is still
$2\nu_{11}$ and its other eleven-cell coordinates
are zero.  Likewise $2x\nu_{11},2y\nu_{11}$ have
order four and independent coordinates $2\nu_{11}$
on the free $K_1,K_5$ eleven-spheres.

\textbf{Every twelve-cell and the bottom intersection.}
The entire shifted attachment list D.2.22t13--t15,
with $a=b=0$, now has $\nu_{11}$-images generated by
\[
 s\nu_{11},\quad s\nu_{11}+z\nu_{11},\quad
 q\nu_{11}+s\nu_{11},\quad
 2x\nu_{11},\quad2y\nu_{11}.
 \tag{D.2.55p}
\]
The respective nonrepeated cells are $B_5,T_{10},B_4,
B_3,B_7$ in the original dimension-ten notation.
The $B_6,Z_2,B_2$ composites are zero; the other cells
are zero or repeat this list.  In particular all free
integral directions have been retained.

Suppose a linear combination of
\textup{(D.2.55p)} is in the bottom subgroup
\textup{(D.2.55f)}.  Collapse the lower skeleton and
project onto the $Q,K_1,K_5$ eleven-spheres.
The last three coefficients must each be divisible
by four, since their respective coordinates are
$2\nu\in\ZZ/8$.
Those whole terms vanish by
\textup{(D.2.55o)} and the order bounds above.
What remains is a combination of $s\nu$ and $z\nu$.
It comes from $K$, which injects on old $\pi_{14}^s$,
and its two quotient coordinates in
\textup{(D.2.55l)} force both coefficients even.
It too is zero.  Therefore
\[
 \langle\hbox{all twelve-cell }\nu\hbox{-composites}\rangle
   \quad\cap\quad
 \operatorname{im}(\pi_{14}^sS^7\to\pi_{14}^sP_{11})=0.
 \tag{D.2.55q}
\]

The stable cofiber sequence for the simultaneous
twelve-cell attachment identifies its kernel on
$\pi_{14}^s$ with exactly the subgroup on the left:
$\pi_{14}^s(\bigvee S^{11})$ is the direct sum of
the groups $\ZZ/8\{\nu_{11}\}$.
Thus the complete bottom $\ZZ/4$ injects into
$\pi_{14}^sP_{12}$.
In the cellular exact couple an incoming $d_5$ at
$(7,7)$ is a bottom stable class killed at this
twelve-stage after its lower-filtration components
have been removed
\cite[Part III, Section 7]{EH5Adams1974AHSS}.
The stronger injection just proved excludes every
such class.  This proves \textup{(D.2.55a)} without
identifying the page-five source with raw homology.
Thirteen-cell cellular boundaries may further
restrict that source and do not change the conclusion.
D.2.54 then leaves only incoming $d_6,d_7$, as stated.
\end{proof}

\begin{claim}{D.2.56}{Proved}{The next two homology groups of the connected cover}
Let $L=Y\langle7\rangle$, $E=\Omega^2L$, and retain the
double-loop fiber $F$ of D.2.18.  All groups in this claim have
coefficients $\FF_2$.  The edge map for
$E\to F\xrightarrow{p}K(\FF_2,4)$ identifies
\[
 H^{\le12}E=
 \left(H^*F/(p^*\widetilde H^*K(\FF_2,4))H^*F\right)^{\le12}.
 \tag{D.2.56a}
\]
In particular $\dim H_{11}E=20$ and $\dim H_{12}E=32$.
Use the previously specified classes
\[
 \begin{gathered}
 u=\xi a\in H_5E,\quad D=\{a,b\},\quad E_6=b^2+p_3\in H_6E,\\
 S=\xi b+P+Q,\quad R=\{a,c\}\in H_7E.
 \end{gathered}
 \tag{D.2.56b}
\]
Actual double homology suspension induces the quotient isomorphisms
\[
 \begin{aligned}
 H_{13}L&\cong H_{11}E/\langle uD,uE_6,\xi u\rangle,
                 &\dim H_{13}L&=17,\\
 H_{14}L&\cong H_{12}E/
 \langle uS,uR,D^2,DE_6,E_6^2,\{u,D\},\{u,E_6\}\rangle,
                 &\dim H_{14}L&=25.
 \end{aligned}
 \tag{D.2.56c}
\]
These descriptions commute with Steenrod operations whenever both
degrees are in the displayed range.  They identify the raw homology
inputs for the potential incoming sixth and seventh differentials
in D.2.54, but do not compute their later-page subquotients, higher
Bocksteins, or stable attaching maps.
\end{claim}
\begin{proof}
\textbf{A bounded Hopf-algebra argument.}
Put $A=H^*K(\FF_2,4)$, $C=H^*F$, and $A'=\operatorname{im}(p^*)$.
The map $p$ is the double loop of the degree-six Postnikov map of
the five-connected cover of $Y$.  It is therefore an H-map, so
$A'$ is a connected Hopf subalgebra of $C$.  Since cup multiplication
is commutative, Milnor--Moore's freeness theorem applies:
$C$ is a graded free $A'$-module
\cite[Proposition 4.9 and the following paragraph, pp.~229--230]
{I_MilnorMoore1965}.  The field hypothesis supplies all the required
graded vector-space splittings.

We will prove that $p^*$ is injective through degree thirteen.
Its kernel then has least possible degree fourteen.  A minimal
graded $A$-free resolution of $A'$ has first positive generators
in internal degree at least fourteen.  Subsequent minimal
differentials have entries of degree at least four, the least
positive degree of $A$.  Consequently
\[
 \operatorname{Tor}_s^A(A',\FF_2)_t=0
       \quad(s>0,\ t<14+4(s-1)).
 \tag{D.2.56d}
\]
The same bound holds for $\operatorname{Tor}_s^A(C,\FF_2)$,
because $C$ is a direct sum of nonnegative shifts of $A'$.
Positive Tor has total degree at least $13+3(s-1)$.
The convergent Eilenberg--Moore sequence for this simply connected
base thus has only filtration zero through total degree twelve,
and no incoming differential can affect that range.  Its edge
map is restriction, proving \textup{(D.2.56a)} once the promised
injectivity is established.  This does not assume global
injectivity of $p^*$.

\textbf{The three new primitive tests.}
The proof of the cover calculation in D.2.22 establishes
injectivity through degree eleven.  A nonzero element of least
degree in the kernel of a connected Hopf map is primitive.
The Serre polynomial generators of $A$ are primitive, and the
other primitives are their iterated squares.  The only new
primitive directions to test are
\[
 \begin{array}{c|l}
 12&(\operatorname{Sq}^2x_4)^2,
             \quad\operatorname{Sq}^5\operatorname{Sq}^2
                       \operatorname{Sq}^1x_4,\\
 13&\operatorname{Sq}^6\operatorname{Sq}^3x_4,
             \quad\operatorname{Sq}^6\operatorname{Sq}^2
                       \operatorname{Sq}^1x_4.
 \end{array}
 \tag{D.2.56e}
\]
Here $x_4$ is the fundamental class of $K(\FF_2,4)$.

Write $u_B,v_B,w_B,t_B$ for the degree $2,3,5,9$ Serre
generators of $H^*B$, and $[m]$ for their dual homology labels
in $T=\Omega^2\Sigma^2B$.  The formal label series
$\mathcal G=\sum_m[m]m$ is group-like.  Hence coefficients of
$\mathcal G^{-1}t_B\partial\mathcal G/\partial t_B$,
and the analogous $w_B$ logarithmic derivative, are primitive.
Let
\[
 \begin{aligned}
 Z_{11}&=[u_Bt_B]+[u_B][t_B],\\
 Z_{12}&=[v_Bt_B]+[v_B][t_B],\\
 Z_{13}^{(1)}&=[u_B^4w_B]
       (\mathcal G^{-1}w_B\partial\mathcal G/\partial w_B),\\
 Z_{13}^{(2)}&=[u_B^2t_B]
       (\mathcal G^{-1}t_B\partial\mathcal G/\partial t_B).
 \end{aligned}
 \tag{D.2.56f}
\]
In the last two lines the initial bracket means coefficient
extraction in the formal variables, not multiplication by a
homology label.  All four classes lie in $H_*F$: they are
primitive in $H_*T$, their images in $H_*B$ are zero since
$B$ has no primitive homology in degrees $11,12,13$, and
the actual coalgebra splitting $H_*T=H_*B\otimes H_*F$
splits primitives into those of the two factors.

Cartan and the Serre operations give
\[
 \operatorname{Sq}_*^1Z_{12}=Z_{11},\qquad
 \operatorname{Sq}_*^6Z_{13}^{(1)}=P,\qquad
 \operatorname{Sq}_*^6Z_{13}^{(2)}=Q.
 \tag{D.2.56g}
\]
For the last two identities the linear-label coefficients are
exactly those in
\[
 \operatorname{Sq}^6(u_B^2v_B)=u_B^4w_B,
 \qquad
 \operatorname{Sq}^6(u_Bw_B)=v_Bw_B^2+u_B^2t_B.
 \tag{D.2.56h}
\]
Primitivity then fixes the remaining terms, since the polynomial
primitive degree-seven subspace has basis $P,Q$.  The existing
degree-eleven test pairs $Z_{11}$ to one with
$p^*\operatorname{Sq}^4\operatorname{Sq}^2\operatorname{Sq}^1x_4$.
Thus $Z_{12}$ detects the new degree-twelve generator and pairs
to zero with the square in \textup{(D.2.56e)}.

The remaining square is detected by $[v_B^2]^2$.
Its reduced coproduct is $[v_B]^2\otimes[v_B]^2$.
Both factors, and $[v_B^2]^2$ itself, map to zero under
$T\to B$ by divided-power multiplication in $H_*B$.
Its full right $H_*B$-coaction is therefore
$[v_B^2]^2\otimes1$.  In the actual coalgebra splitting
$H_*B\otimes H_*F$, these coinvariants are precisely
$1\otimes H_*F$; applying the $B$-counit to the first
factor of the coinvariant identity proves this directly.
This justifies its fiber inclusion.  Since
$\operatorname{Sq}_*^2([v_B]^2)=a^2$, it pairs to one
with $(p^*\operatorname{Sq}^2x_4)^2$.

Finally $p^*\operatorname{Sq}^3x_4$ and
$p^*\operatorname{Sq}^2\operatorname{Sq}^1x_4$ pair on
$(P,Q)$ as $(1,0)$ and $(1,1)$.  Equations
\textup{(D.2.56g)--(D.2.56h)} give an invertible matrix
for the two degree-thirteen directions.  This proves
injectivity through thirteen and completes the edge argument.

\textbf{Dimensions and low products.}
The free restricted-Lie description of the double-loop homology
\cite[Chapter III, Theorems 3.1--3.2 and Lemma 3.8]
{I_CohenLadaMay1976}, the space splitting $T\simeq B\times F$,
and Hopf freeness give the following finite Hilbert-series counts:
\[
 \begin{array}{c|rrrrrrrr}
 i&5&6&7&8&9&10&11&12\\\hline
 \dim H_iE&1&2&2&5&8&12&20&32.
 \end{array}
 \tag{D.2.56i}
\]
For reproducibility, the ordinary free-Lie counts in shifted
degrees $3$ through $14$ are
$1,1,1,2,3,3,6,9,12,17,28,37$; adjoining their restricted
squares and taking the polynomial algebra gives
$\dim H_{11}T=69$, $\dim H_{12}T=112$.
Dividing first by the series of $H_*B$ gives $30,49$
for $F$, and then by that of $H^*K(\FF_2,4)$ gives
$20,32$ for $E$.  The exact finite certificate
\texttt{code/I\_cover\_extension.py} checks these integer
series, all four primitive coproducts and
\textup{(D.2.56g)}.  No homotopy calculation is inferred
from this certificate.

In the polynomial free-E2 homology of $T$, the two products
$uD,uE_6$, the five products
$uS,uR,D^2,DE_6,E_6^2$, and the additional generator
$\xi u$ have the claimed independence in their respective
degrees.  Indeed $u,D,R$ are distinct free generators,
$S$ has the distinct $\xi b$ term, and $E_6$ has the
distinct linear-label $f$ term.  Their triangular expressions
generate a polynomial subalgebra.  The class $\xi u$ is
a separate free generator.

\textbf{The two bar filtrations.}
Use the actual reduced bar constructions $BE\simeq\Omega L$
and $B(\Omega L)\simeq L$.  Their field-coefficient
filtration-one edges are homology suspension, and their
$E^1$ differentials multiply adjacent factors
\cite[Chapter 1, Theorem 1.5 and its footnote, p.~7]
{I_MadsenMilgram1979}.  In $BE$, factors have degree at
least five.  Length three starts in total degree eighteen.
In total degree twelve the only length-two term is $u|u$;
its differential is the nonzero class $u^2$.
The length-two boundaries into the degree-twelve
filtration-one term are $uD,uE_6$.  Therefore
\[
 H_{12}(\Omega L)=H_{11}E/\langle uD,uE_6\rangle,
                  \qquad\dim=18.
 \tag{D.2.56j}
\]
In the second bar construction the term
$\sigma u|\sigma u$ has differential $(\sigma u)^2$.
The top-operation suspension formula identifies this
nonzero class with $\sigma(\xi u)$
\cite[Chapter III, Theorem 1.4(2), with Theorem 1.1(2)]
{I_CohenLadaMay1976}.  This proves the first quotient in
\textup{(D.2.56c)}.

In total degree thirteen of $BE$, filtration one has
dimension $32-5=27$, after quotienting by the five
displayed degree-twelve products.  There are also the
two surviving length-two cycles
\[
 u|D+D|u,\qquad u|E_6+E_6|u.
 \tag{D.2.56k}
\]
They are the shuffle-leading terms of
$(\sigma u)(\sigma D)$ and $(\sigma u)(\sigma E_6)$.
The two reversed products must be retained: $\Omega L$
is only an E1 space.  Their differences have filtration
one and are exactly
\[
 \sigma\{u,D\},\qquad\sigma\{u,E_6\},
 \tag{D.2.56l}
\]
by the Browder suspension and commutator formulas
\cite[Chapter III, Theorem 1.4(3) and Theorem 1.2(2)]
{I_CohenLadaMay1976}.  They are independent modulo the
five products: their free-Lie indecomposable parts have
the respective nonzero words $[a,[a,[a,b]]]$ and
$[a,[a,f]]$, with different linear labels and arities.
Thus the four ordered products have rank four, with
two independent shuffle-leading terms and two independent
lower commutators.

The second bar differential kills all four products.
There are no other terms or higher incoming differentials
in this degree, since its factors have degree at least six
and length three starts in total degree twenty-one.
It follows that double suspension is onto in degree
fourteen and its kernel is exactly the seven-dimensional
subspace in \textup{(D.2.56c)}.  Its target has dimension
$27+2-4=25$.  Naturality of actual homology suspension
proves the stated operation compatibility.  The computation
ends with these ordinary homology quotients; no value of
a later stable differential has been used.
\end{proof}

\begin{claim}{D.2.57}{Proved}{The bottom sigma group survives through the fourteenth stage}
The two remaining incoming differentials in D.2.55 vanish:
\[
 d_6:E^6_{13,2}\longrightarrow E^6_{7,7},\qquad
 d_7:E^7_{14,1}\longrightarrow E^7_{7,7},\qquad
                         E^8_{7,7}=\ZZ/4.
 \tag{D.2.57a}
\]
In an actual fourteen-skeleton $P_{14}$ of
$L=Y\langle7\rangle$, the image of the bottom stable
element $2\sigma$ is nonzero.  It is detected in Adams
filtration two by the image of $h_0h_3$.

Consequently the detecting class of D.2.54 has the exact
remaining obstruction
\[
       d_8\omega=2n\bar\alpha\in\ZZ/4,
       \qquad
       h_Ml_4=0\quad\Longleftrightarrow\quad
                   2n\bar\alpha\in d_8(H_0).
       \tag{D.2.57b}
\]
The integer $n$ and the source-changing image $d_8(H_0)$
are not evaluated here.  In particular this finite-stage
survival result does not decide the actual full-target
obstruction or D.2.4.
\end{claim}
\begin{proof}
\textbf{A four-dimensional Steenrod submodule.}
Let $a_7\in H^7(L;\FF_2)$ be the nonzero bottom class, and put
\[
 s_9=\operatorname{Sq}^2a_7,\qquad
 n_{11}=\operatorname{Sq}^4a_7,\qquad
 z_{13}=\operatorname{Sq}^4\operatorname{Sq}^2a_7.
 \tag{D.2.57h1}
\]
Through degree fourteen these four nonzero classes span the entire
Steenrod cyclic submodule generated by $a_7$. Its only nonzero positive
operations in this range are the three in \textup{(D.2.57h1)}.

Here is a direct verification in the actual spaces. By D.2.56,
double cohomology suspension embeds $H^{\le14}L$ into $H^{\le12}E$.
Under this embedding, the four classes in
\textup{(D.2.57h1)} are the coefficient functionals of the outer
restricted generators
\[
 \xi a,\qquad \xi b,\qquad \xi c,\qquad \xi d
 \quad\hbox{in degrees }5,7,9,11\hbox{ of }H_*T,
 \tag{D.2.57h2}
\]
restricted to $H_*E$. Here $a,b,c,d$ are the homology labels dual
to $u_B,v_B,u_B^2,w_B$.
The top Nishida formulas
\cite[Chapter III, Theorem 1.3(3), pp.~217--218]{I_CohenLadaMay1976}
give these identifications. In particular the only linear-label
inputs relevant to the indicated even operations are
$\operatorname{Sq}_*^1b=a$,
$\operatorname{Sq}_*^2c=a$, and
$\operatorname{Sq}_*^2d=b$.
The other degree-five inputs, the label dual to $u_Bv_B$ and
$\xi a$, have zero second square.

Every odd homology square has zero coefficient on an outer restricted
generator. Indeed, for a restricted generator the odd top Nishida
formula has only Pontryagin-product and Browder-bracket terms.
The restricted adjoint identity expresses brackets of restricted
words as ordinary Lie words, so these have no outer restricted-power
component. On labels, ordinary Hall words, and Pontryagin products,
the same assertion follows from naturality and the two Cartan rules.
Consequently
\[
 \operatorname{Sq}^1a_7=\operatorname{Sq}^1s_9
 =\operatorname{Sq}^1n_{11}
 =\operatorname{Sq}^1z_{13}=0.
 \tag{D.2.57h3}
\]
Adem gives $\operatorname{Sq}^2s_9=0$ and
\[
 \operatorname{Sq}^2n_{11}
 =\operatorname{Sq}^6a_7+\operatorname{Sq}^5
       \operatorname{Sq}^1a_7=0.
 \tag{D.2.57h4}
\]
For the last equality, $\operatorname{Sq}^6a_7=0$ follows after
double suspension from instability on the degree-five class; the
injectivity in degree thirteen is part of D.2.56.
Adem, \textup{(D.2.57h3)--(D.2.57h4)}, and instability now give
all the remaining vanishings through degree fourteen.

The established low-degree tables show that $a_7,s_9,n_{11}$ are
nonzero. To check $z_{13}$ without a full high-degree table, let
$P_{4v}$ be the coefficient of $u_B^4v_B$ in
$\mathcal G^{-1}v_B\partial\mathcal G/\partial v_B$, where
$\mathcal G=\sum_m[m]m$ is the group-like label series. This is a
primitive homology class. Set
\[
 w_{11}=\xi d+Z_{11}+P_{4v},\qquad
 Z_{11}=[u_Bt_B]+a[t_B].
 \tag{D.2.57h5}
\]
All three summands are primitive classes of the fiber $F$. The only
polynomial generators of $H^*K(\FF_2,4)$ in degree eleven are
$x_{11}=\operatorname{Sq}^5\operatorname{Sq}^2x_4$ and
$y_{11}=\operatorname{Sq}^4\operatorname{Sq}^2
\operatorname{Sq}^1x_4$. Their pairing vectors on the three
summands in \textup{(D.2.57h5)} are
\[
 (1,0),\qquad(0,1),\qquad(1,1),
 \tag{D.2.57h6}
\]
respectively. For the last pair, Cartan on the finite label series
gives $\beta P_{4v}=P_p$ and
$\operatorname{Sq}_*^4P_{4v}=P$, with the previously specified
degree-ten Newton class $P_p$ and degree-seven class $P$.
The earlier tables give the first two pairs. Thus $w_{11}$
annihilates both degree-eleven base generators. Primitivity also
annihilates every positive cohomology product, so the exact edge
description in D.2.56 places $w_{11}$ in $H_{11}E$. Finally,
\[
 \operatorname{Sq}_*^4w_{11}=\xi b+P+Q=S,
 \qquad\operatorname{Sq}_*^2S=\xi a.
 \tag{D.2.57h7}
\]
Hence $z_{13}$ evaluates to one on the actual double suspension
of $w_{11}$, proving its nonvanishing.

Restricting to the ordinary fourteen-skeleton $P_{14}$ preserves this
submodule: cohomology restriction is an isomorphism through degree
thirteen and injective in degree fourteen. No minimal CW model is
required. Dually, $H_*P_{14}$ has a four-dimensional quotient
comodule, with basis denoted $a,s,n,z$ in degrees $7,9,11,13$ and
reduced left Milnor coaction
\[
 \begin{aligned}
 \overline\psi(a)&=0,&
 \overline\psi(s)&=\xi_1^2\otimes a,\\
 \overline\psi(n)&=\xi_1^4\otimes a,&
 \overline\psi(z)&=\xi_1^4\otimes s+\xi_2^2\otimes a.
 \end{aligned}
 \tag{D.2.57h8}
\]
The conventions are
$\Delta\xi_j=\sum_i\xi_{j-i}^{2^i}\otimes\xi_i$ and the ordinary
left coaction dual to Steenrod operations.


\medskip\noindent
\textbf{Restriction to an actual finite skeleton.}
Choose the ordinary finite fourteen-skeleton $P_{14}$
as above, under its specified bottom sphere.  The restriction
\[
 H^i(L;\FF_2)\longrightarrow H^i(P_{14};\FF_2)
 \tag{D.2.57c}
\]
is an isomorphism for $i\le13$ and is injective for
$i=14$.  As explained in \textup{(D.2.57h8)}, this
gives the required actual quotient without identifying
all of $H^{14}(P_{14})$ with $H^{14}(L)$.

Let $\mathcal A$ be the mod-two Steenrod algebra.
Dualizing the four-class submodule gives a left
$\mathcal A_*$-comodule quotient of $H_*P_{14}$,
with basis $A,S,N,Z$ in degrees $7,9,11,13$.
Its reduced coaction is
\[
 \overline\psi A=0,\qquad
 \overline\psi S=\xi_1^2\otimes A,\qquad
 \overline\psi N=\xi_1^4\otimes A,\qquad
 \overline\psi Z=\xi_1^4\otimes S+\xi_2^2\otimes A.
 \tag{D.2.57d}
\]
Here $\xi_i$ are the standard Milnor generators, with
\[
 \Delta\xi_k=\sum_{i=0}^k\xi_{k-i}^{2^i}\otimes\xi_i,
                 \qquad\xi_0=1;
 \tag{D.2.57e}
\]
see \cite[Theorem 3, p.~162, and Appendix 1, p.~171]
{EH5Milnor1958Steenrod}.
For example $\xi_2^2$ in \textup{(D.2.57d)} is
forced by the $\operatorname{Sq}^4\operatorname{Sq}^2$
value and coassociativity; the possible $\xi_1^6$
coefficient is zero because $\operatorname{Sq}^6a_7=0$.

\medskip\noindent
\textbf{A complete finite cobar certificate.}
Use the normalized cobar complex
$\overline{\mathcal A}_*^{\otimes s}\otimes
\langle A,S,N,Z\rangle$.  Its differential splits
each algebra tensor by the reduced coproduct and
appends a reduced comodule term on the right.
This convention corresponds to the left coaction in
\textup{(D.2.57d)}.  Write $[u|v]M$ for a basis
element of $C^{2,16}$.  Then
\[
                         t=[\xi_1|\xi_1^8]A
                         \tag{D.2.57f}
\]
is a cycle: both algebra factors are primitive and
$A$ has zero reduced coaction.

Define a functional $\Phi:C^{2,16}\to\FF_2$ to be
one precisely on the following eleven basis elements,
and zero on every other basis element:
\[
\begin{array}{c|cc}
 M&u&v\\ \hline
 A&\xi_1&\xi_1^8\\
 A&\xi_3&\xi_1^2\\
 A&\xi_1^4\xi_2&\xi_1^2\\
 A&\xi_1^7&\xi_1^2\\
 A&\xi_1^2\xi_2^2&\xi_1\\
 A&\xi_1^8&\xi_1\\
 S&\xi_1^3&\xi_1^4\\
 S&\xi_1^5&\xi_1^2\\
 S&\xi_2^2&\xi_1\\
 S&\xi_1^6&\xi_1\\
 Z&\xi_1&\xi_1^2
\end{array}
 \tag{D.2.57g}
\]
There are exactly thirteen incoming basis elements in
$C^{1,16}$.  The next table lists all of them, together
with the number of terms of their differential on which
$\Phi$ is one.  Every number is even.
\[
\begin{array}{c|c|r}
 M&u\text{ in }[u]M&\text{number}\\ \hline
 A&\xi_2^3&2\\
 A&\xi_1^2\xi_3&2\\
 A&\xi_1^3\xi_2^2&2\\
 A&\xi_1^6\xi_2&2\\
 A&\xi_1^9&2\\
 S&\xi_3&2\\
 S&\xi_1\xi_2^2&2\\
 S&\xi_1^4\xi_2&2\\
 S&\xi_1^7&4\\
 N&\xi_1^2\xi_2&0\\
 N&\xi_1^5&0\\
 Z&\xi_2&0\\
 Z&\xi_1^3&2
\end{array}
 \tag{D.2.57h}
\]
This is a direct expansion of
\textup{(D.2.57d)--(D.2.57e)}.  For instance the two
supported terms in $d[\xi_3]S$ are
$[\xi_2^2|\xi_1]S$ and $[\xi_3|\xi_1^2]A$;
those in $d[\xi_1^3]Z$ are
$[\xi_1|\xi_1^2]Z$ and $[\xi_1^3|\xi_1^4]S$.
Thus $\Phi$ annihilates the entire incoming boundary
space, whereas $\Phi(t)=1$.  In particular
\[
       [t]\ne0\quad\hbox{in the quotient-comodule
                     cobar group }\operatorname{Ext}^{2,16}.
       \tag{D.2.57i}
\]
The comodule quotient preserves the bottom class.
Therefore the image of $h_0h_3$ in
$\operatorname{Ext}_{\mathcal A}^{2,16}
(H^*P_{14},\FF_2)$ is nonzero as well.

The authored exact calculation
\texttt{code/EH5\_bottom\_sigma\_cobar.py} checks
the same thirteen columns, the left coaction identity,
and $d^2=0$; its four-class input and explicit dual
certificate are saved in the corresponding
\texttt{code/outputs/EH5\_bottom\_sigma\_four\_class\_*}
files.  The tables above provide the complete
nonboundary proof independently of that implementation.
The larger operation-matrix calculation is a separate
check and is unnecessary for this implication.

\medskip\noindent
\textbf{Adams permanence and the two incoming differentials.}
The Hopf-invariant-one element $\sigma$ is detected
by $h_3$: its cone has nonzero eighth square between
the two cells.  Multiplication by two is detected by
$h_0$, and the nonzero product $h_0h_3$ detects
$2\sigma$.  The survival of this product is also
certified by its nonzero image under the sphere-to-$j$
unit in \cite[Theorem 1.2 and Remark 1.3, p.~3]
{EH5BrunerRognes2022}.
The bottom sphere map sends its cobar representative to
\textup{(D.2.57f)}.  We use the natural Adams spectral
sequence and its comodule-resolution identification in
\cite[Part III, Theorem 15.1, pp.~370--371, and
pp.~372--377]{EH5Adams1974AHSS}.
Naturality from the permanent sphere class prevents
the nonzero image in \textup{(D.2.57i)} from supporting
an Adams differential.  It cannot be hit either.
The only possible incoming differential to bidegree
$(2,16)$ is $d_2$ from $(0,15)$, whose group is zero
because $H_{15}(P_{14};\FF_2)=0$.
Higher incoming differentials have negative source
filtration.

Adams convergence for this finite connective spectrum
therefore detects a nonzero completed stable image
of the specified $2\sigma$.  Its two-local stable
image is consequently nonzero as well: a zero element
would have zero completed image.  Thus
\[
               2i_7\sigma\ne0\quad\hbox{in }\pi_{14}^sP_{14}.
               \tag{D.2.57j}
\]
By D.2.55 the bottom subgroup through dimension twelve
is exactly $\ZZ/4$.  Every nonzero subgroup of
$\ZZ/4$ contains twice its generator.  Consequently
\textup{(D.2.57j)} forces this entire bottom subgroup
to inject through dimensions thirteen and fourteen.
In the cellular exact couple this says precisely that
the incoming $d_6$ and $d_7$ at $(7,7)$ are zero.
Together with D.2.54--D.2.55 it proves
\textup{(D.2.57a)}.  Substituting the now exact target
in D.2.54 gives \textup{(D.2.57b)}, while retaining
all its source-changing ambiguity.
\end{proof}

\begin{claim}{D.2.58}{Proved}{The bottom Adams class remains nonzero on the full target}
On the actual space $L=Y\langle7\rangle$ there is a class
$\ell_{14}\in H^{14}(L;\FF_2)$ satisfying
\[
 \operatorname{Sq}^1\ell_{14}
       =\operatorname{Sq}^2z_{13},\qquad
 z_{13}=\operatorname{Sq}^4\operatorname{Sq}^2a_7.
 \tag{D.2.58a}
\]
The image of the sphere class $h_0h_3$ is nonzero in
\[
 \operatorname{Ext}_{\mathcal A}^{2,16}
                 (H^*L,\FF_2).
 \tag{D.2.58b}
\]
It supports no Adams differential.  Unlike the finite-skeleton
case of D.2.57, it can still be hit by
\[
 d_2:\operatorname{Ext}_{\mathcal A}^{0,15}(H^*L,\FF_2)
       \longrightarrow
       \operatorname{Ext}_{\mathcal A}^{2,16}(H^*L,\FF_2).
 \tag{D.2.58c}
\]
This is the only possible incoming Adams differential, and it is
not evaluated here.  In particular \textup{(D.2.58b)} is an
algebraic nonboundary statement, not a claim of full-target
Adams permanence or a value of $h_Ml_4$.
\end{claim}
\begin{proof}
\textbf{The actual degree-fourteen Bockstein correction.}
Continue the notation $a_7,s_9,n_{11},z_{13}$ of D.2.57, and put
\[
 v_{15}=\operatorname{Sq}^2z_{13}.
 \tag{D.2.58h1}
\]
There is an actual class $\ell_{14}\in H^{14}(L;\FF_2)$ such that
\[
             v_{15}=\operatorname{Sq}^1\ell_{14}.
 \tag{D.2.58h2}
\]
We prove this as an identity of classes on $L$, rather than as a
formal relation in a candidate quotient.

Write $x\in H_6B$ for the label dual to $v_B^2$; thus
$\beta x=d+e$, where $d,e$ are the degree-five labels dual to
$w_B,u_Bv_B$. On $H_{12}T$, let $\ell$ be the coefficient functional
of the ordinary Browder word $\{d,x\}$. Its restriction to $H_{12}E$
annihilates the complete double-suspension kernel in D.2.56:
the five Pontryagin-product directions have no indecomposable
component, and $\{u,D\},\{u,E_6\}$ have ordinary Lie length at
least three, whereas $\{d,x\}$ has length two. Therefore this
functional is the double cohomology suspension of an actual
$\ell_{14}$ on $L$.

One can also see its nonvanishing directly. The class
\[
             r_{12}=\{d,x\}+b\{b,d\}
             \tag{D.2.58h3}
\]
is primitive: the two reduced diagonals are both
$\{b,d\}\otimes b+b\otimes\{b,d\}$ and cancel.
Its image in $B$ is zero, so the actual product splitting puts it
in $H_{12}F$. Every base-generator functional for
$F\to K(\FF_2,4)$ vanishes on it: a composite of homology squares
on its ordinary Lie and product terms cannot produce the bottom
Pontryagin square $a^2$. Primitivity also annihilates all positive
cohomology products. The exact D.2.56 edge consequently places
$r_{12}$ in $H_{12}E$, and $\ell(r_{12})=1$.

The following is an exact identity of functionals on the whole
ambient group $H_{13}T$:
\[
       \operatorname{Sq}^1(\{d,x\}^{*})=(\xi x)^{*}.
       \tag{D.2.58h4}
\]
Indeed the top Bockstein formula gives
$\beta\xi x=\{x,d+e\}$. No ordinary Hall word can contribute a
$\{d,x\}$ coefficient under $\beta$. Lie length is preserved;
in length two, no label has Bockstein $x$, and the only label with
a $d$ component in its Bockstein is $x$ itself. The latter would
require the zero word $\{x,x\}$. Pontryagin products cannot
contribute an indecomposable word. The other restricted generators
in degree thirteen are $\xi f$ and $\xi\{a,b\}$, both with zero
Bockstein. This proves \textup{(D.2.58h4)} using the exact CLM top
formula \cite[Chapter III, Theorem 1.3(3), p.~218]
{I_CohenLadaMay1976}.

Likewise, the even top Nishida formula identifies the restriction
of $(\xi x)^*$ to $H_{13}E$ with the double suspension of
$\operatorname{Sq}^2z_{13}$: the only degree-six input whose
Bockstein has a $d$ component is $x$.
Thus \textup{(D.2.58h4)} proves the desired equality after double
cohomology suspension. It remains to justify injectivity in this
one additional degree.

The two actual bar constructions give that injectivity without
assuming a new Eilenberg--Moore edge range. In total degree fourteen
of the first bar construction, the only possible filtration-two
cycles beyond filtration one are
\[
 u|S+S|u,\qquad u|R+R|u,\qquad D|E_6+E_6|D.
 \tag{D.2.58h5}
\]
They are the leading filtration terms of actual ordered products
in $H_{14}(\Omega L)$ and hence are killed by the next homology
suspension. In total degree fifteen of the second bar construction,
the four possible filtration-two terms have factors of degrees
six and seven. Their differentials are precisely the four
independent ordered products proved in D.2.56. Hence they supply
no additional class. Filtration three starts above both ranges,
so no higher differential changes these conclusions. Therefore
\[
       H_{13}E\longrightarrow H_{15}L
       \quad\hbox{is surjective}.
       \tag{D.2.58h6}
\]
Dualizing proves the required injectivity and establishes
\textup{(D.2.58h2)} on the actual space $L$.

In particular every reduction of an integral homology class in
degree fifteen annihilates $v_{15}$. This is the ordinary
Bockstein pairing consequence of \textup{(D.2.58h2)}; it does not
assert that every mod-two Bockstein cycle has an integral lift.


\medskip\noindent
\textbf{Extending the complete finite certificate.}
Use the normalized left-comodule cobar convention of D.2.57.
Regard its functional $\Phi$ as a functional on
$C^{2,16}(H_*L)$ by using the same actual cohomology
coordinates $a_7,s_9,z_{13}$ on the module factor and
setting all unlisted values to zero.  Define
\[
 \widetilde\Phi([u|v]m)=\Phi([u|v]m)
   +\mathbf1_{u=\xi_1}\mathbf1_{v=\xi_1}
                   \langle\ell_{14},m\rangle.
 \tag{D.2.58d}
\]
Here the last pairing is zero unless $m$ has degree
fourteen.  We check every possible incoming cobar column,
including all newly available degree-fifteen module classes.

For a module class of degree at most thirteen, its entire
cobar differential has module degree at most thirteen.
The additional term in \textup{(D.2.58d)} therefore vanishes,
and the complete certificate of D.2.57 applies.  For a
degree-fourteen module class the only algebra monomial in
$C^{1,16}$ is $\xi_1^2$.  Its reduced coproduct is zero,
so the new functional again sees no differential term;
the old functional is zero by the same truncated-comodule
calculation.

Finally, every degree-fifteen incoming column is
$[\xi_1]m_{15}$.  The internal differential is zero.
Of the eleven old supported pairs in D.2.57, precisely two
begin with $\xi_1$.  They give
\[
 \Phi\bigl(d[\xi_1]m_{15}\bigr)
   =\langle\operatorname{Sq}^8a_7+
                \operatorname{Sq}^2z_{13},m_{15}\rangle
   =\langle\operatorname{Sq}^2z_{13},m_{15}\rangle.
 \tag{D.2.58e}
\]
The last equality is instability on the degree-seven
class.  The additional supported pair in
\textup{(D.2.58d)} contributes
$\langle\operatorname{Sq}^1\ell_{14},m_{15}\rangle$,
which cancels \textup{(D.2.58e)} by
\textup{(D.2.58a)}.  There are no other columns:
a positive algebra factor in internal degree sixteen
requires module degree at most fifteen.  Thus
\[
             \widetilde\Phi\,d=0
                       \quad\hbox{on all }C^{1,16}(H_*L).
             \tag{D.2.58f}
\]

The bottom cocycle remains
$t=[\xi_1|\xi_1^8]A$, with $A\in H_7L$ dual to $a_7$.
The added term does not change its pairing, so
$\widetilde\Phi(t)=1$.  This proves
\textup{(D.2.58b)} by a complete symbolic certificate,
without requiring a basis for all of $H_{15}L$.

As an independent finite algebra check, the authored
Milnor-cobar engine also verifies the same functional on
the supplied degree-fifteen operation matrices.  The
geometrically specified $\ell_{14}$ uses the coordinate
\texttt{L14\_22}; the exact twelve-term support and all
166 incoming column counts are saved in
\texttt{code/outputs/EH5\_L15\_ell\_functional\_certificate.json}.
That larger check is unnecessary for the symbolic proof
\textup{(D.2.58d)--(D.2.58f)}.

\medskip\noindent
\textbf{The remaining Adams differential.}
The actual bottom sphere map sends the permanent sphere
class $h_0h_3$ to the nonzero class just proved.
Adams naturality, with the same conventions and primary
sources as D.2.57, prevents any outgoing differential
while the image exists.  An incoming $d_r$ to bidegree
$(2,16)$ would originate at
$(2-r,17-r)$.  For $r>2$ this has negative filtration.
For $r=2$ the source is exactly \textup{(D.2.58c)}.
The full space can have degree-fifteen homology, so the
finite-skeleton vanishing argument cannot remove this
source.  None of the preceding calculations determines
its differential or the final obstruction in D.2.57.
\end{proof}

\begin{claim}{D.2.59}{Proved}{Integral lifts of all seven primary classes}
Put $L=Y\langle7\rangle$ and $E=\Omega^2L$.  With mod-two
coefficients, $\dim H_{13}E=47$, $\dim H_{15}L=33$, and the
subspace
\[
 \mathcal P_{15}=\bigcap_{i>0}\ker\bigl(
      \operatorname{Sq}_*^i:H_{15}L\longrightarrow H_{15-i}L\bigr)
 \tag{D.2.59a}
\]
has dimension seven.  Every element of $\mathcal P_{15}$ is the
reduction of an integral $2$-local homology class.

There is an explicit basis $K_1,\ldots,K_7$ specified below for which
\[
 \begin{gathered}
 \mathcal U=\langle K_1+K_2,K_3,K_4,K_5,K_6,K_7\rangle
   =\mathcal P_{15}\cap\ker\theta,\\
 \langle\theta,K_1\rangle=\langle\theta,K_2\rangle=1,
 \qquad \langle\theta,K_i\rangle=0\quad(3\leq i\leq7).
 \end{gathered}
 \tag{D.2.59b}
\]
Every class in $\mathcal U$ is the reduction of an integral double
homology suspension from $E$.  Neither $K_1$ nor $K_2$ is such an
integral double suspension.  This is an integral-source statement;
it does not, by itself, supply stable spherical representatives
or compute the incoming Adams differential of D.2.58.
\end{claim}
\begin{proof}
All unnamed integral groups in this proof are localized at $2$.
We retain the labels $a,b,c,d,e,h$ dual respectively to
$u_B,v_B,u_B^2,w_B,u_Bv_B,u_Bw_B$, and the classes
$u,D,E_6,S,R,J_1,J_2,P_r,B_1,\ldots,B_7$ of D.2.22 and
D.2.56.  Thus their respective label degrees are
$2,3,4,5,5,7$; the $J_i$ have degree eight and the $B_i,P_r$
have degree ten.

\textbf{The next edge and the double-bar quotient.}
In the Hopf-algebra argument of D.2.56, the new primitive directions
in degree fourteen of $H^*K(\FF_2,4)$ are exactly
\[
 (\operatorname{Sq}^3x_4)^2,\quad
 (\operatorname{Sq}^2\operatorname{Sq}^1x_4)^2,\quad
 \operatorname{Sq}^6\operatorname{Sq}^3\operatorname{Sq}^1x_4.
 \tag{D.2.59c}
\]
Here the first two are the squares of the two degree-seven
generators, and the last is the new Serre polynomial generator.
Let $r:T\to F$ be the actual projection from the space splitting
$T\simeq B\times F$.  On the three actual homology classes
\[
 r_*([v_B^2][u_B^4]),\qquad
 r_*([u_B^2][w_B^2]),\qquad r_*([w_B][t_B]),
\]
the pullbacks of the three classes in \textup{(D.2.59c)} have
pairing matrix
\[
 \begin{pmatrix}1&1&0\\0&1&1\\0&0&1\end{pmatrix}.
 \tag{D.2.59d}
\]
For clarity, this calculation uses $r_*(bc)=P$, $r_*(ad)=Q$,
and $r_*(g)=r_*(h)=r_*(ae)=r_*(a^2b)=0$ in degree seven.
The first two columns are the diagonal pairings with
$(V^*+P^*)^2$ and $(P^*+Q^*)^2$.  In the last column, the reverse
chain $\operatorname{Sq}_*^1\operatorname{Sq}_*^3
\operatorname{Sq}_*^6$ takes the three representatives respectively
through $c^2\mapsto0$, $ax\mapsto ab\mapsto a^2$, and
$bd\mapsto ab\mapsto a^2$.  Thus no multiplicativity of $r$ is
assumed.  This proves injectivity of $p^*$ through degree fourteen.

The first possible kernel degree is now fifteen.  The bounded
free-resolution argument of D.2.56 therefore gives
$t\geq15+4(s-1)$ for positive Tor in homological degree $s$.
Its total degree is at least $14+3(s-1)$, so the actual
Eilenberg--Moore edge remains the quotient edge through degree
thirteen.  Its Hilbert coefficient gives $\dim H_{13}E=47$.

The two ordinary bar filtrations give
\[
 \begin{gathered}
 H_{15}L\cong H_{13}E/\mathcal J,\\
 \begin{split}
 \mathcal J=\langle{}&uH_8E,DS,DR,E_6S,E_6R,\\
       &\{u,S\},\{u,R\},\{D,E_6\},\xi D,\xi E_6\rangle .
 \end{split}
 \end{gathered}
 \tag{D.2.59e}
\]
Here $uH_8E$ contributes five dimensions.  The first nine product
directions are independent by their reduced diagonals in
bidegrees $(5,8)$ and $(6,7)$.  In the first bar, degree fourteen
has their $38$-dimensional quotient and the three additional
cycles $u|S+S|u$, $u|R+R|u$, $D|E_6+E_6|D$.
In the second bar, the four possible length-two terms in total
degree fifteen have independent ordered-product differentials,
as in D.2.56.  The eight products removed from $H_{14}\Omega L$
have three independent leading terms, the three cycles just listed.
Their five remaining terms are exactly the last five directions
of \textup{(D.2.59e)}: the reversed-product differences give the
three Browder brackets, and the two squares give the top
operations \cite[III, Theorem 1.4]{I_CohenLadaMay1976}.
The brackets have distinct ordinary-Lie leading label multisets
$(a,a,g)/(a,a,h)$, $(a,a,a,c)$, and $(a,b,f)$; the two top
operations have distinct outer restricted terms $\xi\{a,b\}$
and $\xi f$.  They are independent modulo the nine products.
Thus $\dim\mathcal J=14$ and $\dim H_{15}L=33$.

\textbf{The actual normalized Browder map.}
We use a reduced ordinary map, including its corrections on
nonprimitive inputs.  A loop-group model of $T$ and the fundamental
circle of its binary little-cubes operation give a loop of maps
$T\times T\to T$ based at multiplication.  Multiply pointwise
by the inverse of that multiplication.  On each axis the resulting
loop comes from unary little-cubes operations, a contractible
space.  Contract these loops and use the homotopy extension
property of the wedge inclusion.  The result is constant on the
axes and at the circle basepoint, and hence gives an actual map
\[
 \mathcal N:\Sigma(T\wedge T)\longrightarrow T.
 \tag{D.2.59f}
\]
Equivalently, it is the double adjoint of the generalized Whitehead
product of the two evaluation maps $\Sigma^2T\to Y$.
The defining circle action is the Browder operation.  The diagonal
and inverse give the precise mod-two formula
\[
 \mathcal N_*(\sigma(x\otimes y))
   =\sum\{x',y'\}\,S_T(x''y''),
 \quad
 \Delta x=\sum x'\otimes x'',\quad
 \Delta y=\sum y'\otimes y''.
 \tag{D.2.59g}
\]
The inverse factor has no circle dependence, so this formula
follows directly by placing the circle fundamental class in the
action factor.  $S_T$ is the antipode of $H_*T$, not that of
$H_*B$.  For example, if $y$ is primitive then
$\mathcal N_*(\sigma(c\otimes y))=\{c,y\}+a\{a,y\}$.

Restrict \textup{(D.2.59f)} along $B\wedge E$ or $B\wedge B$.
The suspended sources are respectively seven- and four-connected;
their maps to $T$ lift uniquely to the four-connected cover $E$.
Double suspension and evaluation consequently give actual maps
\[
 \Sigma^3(B\wedge E)\longrightarrow L,
 \qquad \Sigma^3(B\wedge B)\longrightarrow L.
 \tag{D.2.59h}
\]
Write $n(x,y)$ for the induced degree-fifteen class when
$|x|+|y|=12$.

\textbf{The seven primary classes.}
An explicit choice of basis for \textup{(D.2.59a)} is
\[
 \begin{aligned}
 K_1&=n(a,B_7),&K_2&=n(a,B_3),\\
 K_3&=n(a,B_2+B_6),&K_4&=n(a,B_2)+n(c,J_2),\\
 K_5&=n(a,P_r),\\
 K_6&=n(a,B_2+B_4)+n(c,J_1)+n(b,K_{\rm sum}),\\
 K_7&=n(e,h),&K_{\rm sum}&=K_2^{\rm old}+K_3^{\rm old}+K_4^{\rm old}.
 \end{aligned}
 \tag{D.2.59i}
\]
The superscript distinguishes the earlier degree-nine cover
classes from the present degree-fifteen basis.
We specify the finite rank verification to make completeness
checkable.  In the ordered quotient basis of
\textup{(D.2.59e)}, the seven vectors have supports
\[
 \begin{array}{c|l}
 K_1&05\\K_2&08\\K_3&06,10\\K_4&06,12\\K_5&15\\
 K_6&06,07,11,13,18,19\\K_7&30.
 \end{array}
 \tag{D.2.59j}
\]
These are the indices $\texttt{L15\_01},\ldots,\texttt{L15\_33}$
of the complete $33$-vector quotient basis in the certificate.
Applying Cartan, the Serre label operations, and the CLM
restricted-operation and Browder formulas gives rank $26$ for
the combined map $(\operatorname{Sq}_*^1,
\operatorname{Sq}_*^2,\operatorname{Sq}_*^4)$.
Its kernel consists exactly of \textup{(D.2.59j)}.
$\operatorname{Sq}_*^8$ is zero by instability on degree-seven
cohomology, and the squares of power-of-two degree generate the
Steenrod algebra.  This proves the asserted primary dimension.
The complete basis, operation matrices, and exact row reductions
are in \texttt{I\_E2\_primary\_15.py},
\texttt{I\_L15\_primary.json}, and
\texttt{I\_L15\_primary\_cycles.json}; the normalized image
calculation is \texttt{I\_integral\_browder\_sources.py}.
These use the actual edge and quotient proved above.  They also
check membership in $H_*E$ before taking the bar quotient,
and all Cartan, coproduct, and Adem compatibilities in this range.

\textbf{Prescribed integral tensors and Tor cycles.}
The classes $a,c,d+e,h$ in $B$ are reductions of integral cycles;
the first two have orders $2,4$, and the last two have order two.
For example $d+e$ and $h$ are reductions of the integral
Bocksteins of $[v_B^2]$ and $[v_Bw_B]$.
The classes $J_1,J_2$ in $E$ also lift integrally:
$H_7(E;\mathbb Z_{(2)})=\mathbb Z/2$, so reduction is injective
on the possible integral connecting image, and their first
Bocksteins vanish.  In degree ten the needed integral reductions
are $P_r,B_2,B_1+B_3,B_5+B_6$.  Indeed
\[
 H_9(E;\mathbb Z_{(2)})=
       \mathbb Z/8\{N'\}\oplus(\mathbb Z/2)^3.
 \tag{D.2.59k}
\]
A first-Bockstein cycle can therefore have integral connecting
image only $0$ or $4N'$.  The latter is detected by the third
Bockstein.  The complete D.2.22 table has
$\beta_3B_5=\beta_3B_6=N$ and zero on
$P_r,B_2,B_1+B_3$, proving the assertion for the four combinations.
Thus no general inference from a first-Bockstein cycle to an
integral lift is used.

Choose integral chains with the following boundaries and the
indicated mod-two reductions:
\[
 \begin{gathered}
 d\widetilde b=2\widetilde a,
 \quad d\widetilde d=4\widetilde c,
 \quad d\widetilde l=2\widetilde h,\\
 d\widetilde B_5=8\widetilde N,
 \qquad d\widetilde B_4=2\widetilde K_{\rm sum}.
 \end{gathered}
 \tag{D.2.59l}
\]
Here $l=[v_Bw_B]$.  These are the established bottom Moore,
second-Bockstein, and D.2.22 Bockstein representatives, with
orientations chosen to give the displayed signs.  The following
literal tensor chains are cycles:
\[
 \begin{aligned}
 \tau(a,N)&=4\widetilde b\otimes\widetilde N
                 -\widetilde a\otimes\widetilde B_5,
       &\rho_2\tau(a,N)&=a\otimes B_5,\\
 \tau(a,K_{\rm sum})&=\widetilde b\otimes\widetilde K_{\rm sum}
                 -\widetilde a\otimes\widetilde B_4,
       &\rho_2\tau(a,K_{\rm sum})&=b\otimes K_{\rm sum}+a\otimes B_4,\\
 \tau(c,h)&=\widetilde d\otimes\widetilde h
                 -2\widetilde c\otimes\widetilde l,
       &\rho_2\tau(c,h)&=d\otimes h.
 \end{aligned}
 \tag{D.2.59m}
\]
The boundaries cancel as $8aN-8aN$, $2aK_{\rm sum}-2aK_{\rm sum}$,
and $4ch-4ch$.  This proves the prescribed reductions directly,
without selecting an abstract K\"unneth splitting.

The actual maps \textup{(D.2.59h)} applied to the following
integral tensor/Tor cycles give six independent directions:
\[
 \begin{array}{c|c}
 \text{integral source cycle}&\text{mod-two image in }H_{15}L\\ \hline
 a\otimes P_r&K_5\\
 a\otimes(B_1+B_3+B_2)&K_1+K_2\\
 a\otimes B_2+c\otimes J_2&K_4\\
 a\otimes(B_2+B_5+B_6)+\tau(a,N)&K_3\\
 a\otimes B_2+c\otimes J_1+\tau(a,K_{\rm sum})&K_6\\
 (d+e)\otimes h+\tau(c,h)&K_7.
 \end{array}
 \tag{D.2.59n}
\]
Integral representatives are understood in the first column.
The first five rows use $B\wedge E$ and the last $B\wedge B$.
The fourth reduction is $a\otimes(B_2+B_6)$ and the last is
$e\otimes h$.  Formula \textup{(D.2.59g)} verifies every image;
in particular its correction on $c$ is essential in the third
and fifth rows.  Independence follows at once from
\textup{(D.2.59j)}.

\textbf{The integral double-suspension restriction.}
By D.2.50, $\beta_{\mathbb Z}\theta=b_8^2$.
Cohomology suspension kills products, so $\sigma^2\theta$ has
an integral cohomology lift on $E$.  The space $E$ is rationally
acyclic of finite type, hence $H_{13}(E;\mathbb Z_{(2)})$ is
torsion.  An integral cohomology class pairs trivially with this
group.  Therefore $\theta$ vanishes on every integral double
homology suspension from $E$, in particular on all six rows of
\textup{(D.2.59n)}.

D.2.52 supplies an integral class $\omega$ whose reduction lies
in $\mathcal P_{15}$ and pairs to one with $\theta$.
Its reduction is outside the six-dimensional subspace
$\mathcal U$.  Thus the reductions of integral classes span all
seven dimensions of $\mathcal P_{15}$.  The functional $\theta$
has kernel exactly $\mathcal U$, proving \textup{(D.2.59b)}.
The nonzero evaluations on $K_1,K_2$ also prove that neither is
an integral double suspension, despite mod-two double-suspension
surjectivity.  The argument imposes no stable-sphericity claim.
\end{proof}

\begin{claim}{D.2.60}{Proved}{The auxiliary James--Hopf class survives all twelve-cells}
Put $B=K(\mathbb F_2,2)$, $X_0=\Sigma B$, $Y=\Sigma X_0$, and
\[
 Z=\Sigma(X_0\wedge X_0)=\Sigma^3(B\wedge B).
 \tag{D.2.60a}
\]
Work at the prime two and use the actual CW model specified below.
If $i_{7,12}:S^7\to Z^{(12)}$ is its bottom map, then
\[
 0\ne i_{7,12}\sigma'\in\pi_{14}Z^{(12)},\qquad
                         2i_{7,12}\sigma'=0.
 \tag{D.2.60b}
\]
Its image in $\pi_{14}Z$ is the actual second James--Hopf image
of $\alpha\sigma'$, where $\alpha=u_4\nu_4\in\pi_7Y$.
Thus none of the ten-, eleven-, or twelve-dimensional attaching
maps kills the indicated auxiliary class.  This assertion does
not determine its image after the cells in dimensions thirteen
through fifteen, or the value of $\alpha\sigma'$ or $h_Ml_4$ in
the full target.
\end{claim}
\begin{proof}
All free groups and cellular degrees below are two-local.
Write $M(2,j)=M(\mathbb Z/2,j)$ for a Moore space with bottom
cell in degree $j$, and put
\[
 C_1^j=\operatorname{Cof}
 \bigl((2,\eta_{j-2}):S^{j-2}\vee S^{j-1}\longrightarrow S^{j-2}\bigr).
 \tag{D.2.60c}
\]
We retain actual attaching maps when their composites leave
the stable range.

\textbf{The bottom class and the borderline nine-stage.}
Naturality of the second James--Hopf map with respect to
$\Sigma u_3:S^4\to Y$, together with the sphere calculation
$H(\nu_4\sigma')=\sigma'$ in
\cite[Proposition~13, p.~9]{FGHStricklandToda}, gives
\[
 H_2(\alpha\sigma')=i_7\sigma',\qquad
                  i_7=\Sigma(u_3\wedge u_3):S^7\to Z.
 \tag{D.2.60d}
\]
This uses the stated sphere calculation, not a general
distributivity formula for a Hopf invariant of a composite.
The low tensor cells, equivalently the twice-suspended Ganea
calculation in D.2.49, give
\[
                 Z^{(9)}\simeq C_1^9\vee M(2,8)\vee S^9.
 \tag{D.2.60e}
\]
The map $M(2,7)\to C_1^9$ is injective on $\pi_{14}$.
Here an explicit unstable argument is necessary.  For
$f=i_7\eta_7:S^8\to M(2,7)$, the fiber of the pinch
$C_1^9\to S^9$ is the relative James construction.  Its
second stage attaches the cone on $M(2,7)\wedge S^7$ by
$[\mathrm{id},f]$; see
\cite[Lemmas~2.1 and 2.3--2.5, pp.~4--5]{FGHJinZhu2024}.
Its first new cell has dimension fifteen and attaches by
\[
 [i_7,i_7\eta_7]=i_7[\iota_7,\iota_7]\eta_{13}=0,
 \tag{D.2.60f}
\]
since $S^7$ is an $H$-space.  The next cell has dimension
sixteen; later word lengths start in dimension twenty-three.
Thus $M(2,7)$ maps isomorphically to the fiber on $\pi_{14}$.
The remaining boundary is zero: suspension
$\pi_{14}S^8\to\pi_{15}S^9$ is onto, and its image under the
fiber boundary is $i_7\eta_7\nu_8^2=0$.
Indeed $\eta_7\nu_8^2$ stabilizes to zero, while
$\pi_{14}S^7=\mathbb Z/8\{\sigma'\}$ stabilizes injectively
by $\sigma'\mapsto2\sigma$.  This verifies the hypotheses
of the cited relative-fiber exact sequence in the required
degree.  The exact Moore relation
$i_7\sigma'=2l_4\ne0$ has order two
\cite[Lemma~2.8(ii)]{EH5MorisugiMukai2000}.
Consequently the bottom class is nonzero of order two in
\textup{(D.2.60e)}.

\textbf{Literal integral cells and complete attaching matrices.}
Use the fixed mod-two labels
$a_2,b_3,c_4,d_5,e_5,f_6,x_6$ of $B$, where $d,e,f,x$
are dual to $w,uv,u^3,v^2$.  After the parity-preserving
integral normalization used in D.2.18 and D.2.22, choose chains
\[
 \partial b=2a,\quad \partial c=0,\quad
 \partial d=\partial e=4c,\quad
 \partial x=2(e-d),\quad \partial f=0.
 \tag{D.2.60g}
\]
The last cycle can also be chosen from the spherical class
of D.2.17.  Ordered tensor labels below suppress the shift
by three.  In the nine-stage set
\[
 B_8=ab,\quad e_8=ba-ab,\quad
 D_9=ac,\quad r_9=ca-ac,\quad m_9=bb-ac.
 \tag{D.2.60h}
\]
These give respectively the Chang, Moore and free factors
in \textup{(D.2.60e)}.  Write $\gamma,\delta\in\pi_9C_1^9$
for the free class with Hurewicz image $2D$ and the order-two
class detected by $\operatorname{Sq}^2$ on $B_8$.
The other possible torsion component is $e_8\eta_8$.

The literal ten-cycles are
$AD=ad-2bc$, $AE=ae-2bc$, $DA=da-2cb$, $EA=ea-2cb$;
the remaining cells are $BC=bc$ and $CB=cb$.
Their boundaries and second-square detectors determine all
components of their attaching maps:
\[
\begin{array}{c|rrrrrr}
\text{cell}&AD&AE&DA&EA&BC&CB\\ \hline
\text{attachment}&\delta&0&\delta+e_8\eta_8&0&
 \gamma+\delta+e_8\eta_8&\gamma+2r+\delta
\end{array}
\tag{D.2.60i}
\]
For example $\operatorname{Sq}_*^2(ad)=ab$ and
$\operatorname{Sq}_*^2(da)=ba$; the corresponding values on
$bc,cb$ are $ba,ab$.  There is no further component in
$\pi_9$ of the wedge.  The explicit cell changes
\[
 Q=AD,\quad V=DA-AD,\quad G=BC-DA,\quad
 H=CB-BC+DA-AD
\]
isolate $\delta,e_8\eta_8,\gamma,2r$.  Define
$\mathrm{Cb}_j=C_1^j\cup_\delta e^{j+1}\cup_\gamma e^{j+1}$.
We obtain an actual splitting
\[
 Z^{(10)}\simeq\mathrm{Cb}_9\vee C_1^{10}\vee M(2,9)
                         \vee S^{10}_{AE}\vee S^{10}_{EA}.
 \tag{D.2.60j}
\]
The complete degree-ten groups of these factors have free
generators $\kappa,\gamma',AE,EA$, with
$h(\kappa)=2Q$, $h(\gamma')=2V$, and three order-two
generators $v=i_7\nu_7,\delta',r\eta_9$.
Their independent torsion detectors are respectively
$\operatorname{Sq}^4$ on $aa$ and $\operatorname{Sq}^2$
on $m,r$.  Normalize $\kappa$ to have zero fourth-square
detector; its second-square detector is $D$.
These follow from the same Chang cofiber sequences and
Adem calculation as in D.2.38, shifted once.

Here is the entire signed eleven-cell matrix.  It records
the literal cellular boundary as well as the actual map;
thus no higher-exponent conclusion is being inferred from
the first Bockstein alone.
\[
\begin{array}{c|c|c}
\text{cell}&\partial&\text{attaching map}\\ \hline
AF&0&0\\
AX&2(AE-AD)&-\kappa+2AE\\
FA&0&0\\
XA&2(EA-DA)&-\kappa-\gamma'+2EA+r\eta_9\\
BD&2AD&\kappa+\delta'\\
BE&2AE&2AE\\
DB&-2DA&-\kappa-\gamma'+\delta'\\
EB&-2EA&-2EA\\
CC&0&v+r\eta_9
\end{array}
\tag{D.2.60k}
\]
In this order the second-square vectors are
$0,D,0,D+r,D+m,0,D+m,0,r$; only $CC$ has a nonzero
fourth square, namely $aa$.  Together with the integral
boundaries these are complete detectors in $\pi_{10}$.
First isolate $2AE,2EA$; then $DB+BD$ isolates $-\gamma'$,
and $AX+BD-BE$ isolates $\delta'$.  The remaining rows
isolate $\kappa,r\eta_9,v$ and two zero maps.  Consequently
\[
\begin{gathered}
 Z^{(11)}\simeq J\vee\mathrm{Cb}_{10}\vee C_1^{11}
                 \vee M(2,10)^{\vee2}\vee(S^{11})^{\vee2},\\
 J=\mathrm{Cb}_9\cup_\kappa e^{11}_K\cup_v e^{11}_N.
\end{gathered}
\tag{D.2.60l}
\]
These ordinary cell slides take place within the stable
range through dimension eleven and fix the bottom $S^7$.
In particular there is an actual retraction onto $J$.

\textbf{Every degree-eleven class of the bottom block.}
The Chang cofiber sequence gives
\[
\begin{gathered}
 \pi_{10}\mathrm{Cb}_9=\mathbb Z_{(2)}\{\kappa\}
                                  \oplus\mathbb Z/2\{v\},\\
 0\longrightarrow\mathbb Z/2\{\omega\}
 \longrightarrow\pi_{11}\mathrm{Cb}_9
 \longrightarrow\mathbb Z/2\{G\eta_{10}\}\longrightarrow0,
 \qquad\kappa\eta_{10}=\omega.
\end{gathered}
\tag{D.2.60m}
\]
For completeness, before the two ten-cells the relevant
Chang group is $\mathbb Z/4\{t\}$, with pinch coordinates
$(2\nu_8,\eta_9^2)$ and $\delta\eta_9^2=2t$.
The $\delta$ cell leaves its order-two image $\omega$;
the $\gamma$ cell supplies the displayed $G\eta_{10}$
quotient.  This is the cofiber calculation of D.2.38,
in degrees shifted once.  Attaching $\kappa$ kills
$\omega$ and supplies no free lift, because $\kappa$ has
infinite order.  Its remaining group is
$\mathbb Z/2\{\zeta\}$.  Finally $v$ has order two
and $v\eta_{10}=0$, so its eleven-cell gives
\[
 \pi_{11}J=\mathbb Z_{(2)}\{p\}\oplus\mathbb Z/2\{\zeta\},
 \qquad h(p)=2N.
 \tag{D.2.60n}
\]
The second-square detector of $\zeta$ on $G_{10}$ is one.

\textbf{A quotient detects the torsion composite.}
Collapsing the bottom gives the actual equivalence
\[
 \mathrm{Cb}_9/S^7\simeq C(\eta_8)\vee M(2,9).
 \tag{D.2.60o}
\]
Their degree-eleven groups are respectively
$\mathbb Z/4\{i_8\nu_8\}$ and
$\mathbb Z/4\{\widetilde\eta_9\}$.
The Moore coordinate of $\zeta$ is odd because its pinch
on $G$ is $\eta_{10}$.  The connecting map to $S^8$
sends the first generator to $2\nu_8$ and the second
to an odd multiple of $2\nu_8$, by the ordinary bracket
$\{\eta_8,2,\eta_9\}=\{2\nu_8,6\nu_8\}$.
Exactness forces its $C(\eta_8)$ coordinate to be odd too.
Therefore $\zeta\nu_{11}$ has the nonzero coordinate
$i_8\nu_8^2$ in that factor.  The fourth stable stem is
zero, so the $\eta$ cell does not kill this bottom
$\nu^2$; the Moore coordinate has zero composite by
the fifth stable stem of the mod-two Moore spectrum.
Adding the two eleven-cells to obtain $J/S^7$ still
preserves this coordinate: these quotient spaces are
seven-connected, and the possible kernel on $\pi_{14}$
comes from $\pi_{14}(S^{10}\vee S^{10})=0$.
In particular
\[
 \zeta\nu_{11}\text{ has nonzero image in }\pi_{14}(J/S^7).
 \tag{D.2.60p}
\]

\textbf{The free composite and the full twelve-cell kernel.}
There is an actual map $j:C(\nu_7)\to J$ on the bottom
and the $N$ cell.  The unique $v_8\in\pi_{11}C(\nu_7)$
of top degree eight exists by its stable-range cofiber
sequence: $\pi_{11}S^7=0$ and $\nu_7$ has order eight.
Both $jv_8$ and $4p$ have integral Hurewicz image $8N$.
Their difference could only be $\zeta$, but both have
zero second-square detector on $G_{10}$.  For $4p$ use
the degree-four cone map; for $jv_8$ use naturality and
$H^{10}C(\nu_7)=0$.  Hence $4p=jv_8$ exactly.
The ordinary coextension formula, with its full
indeterminacy, now gives
\[
 4p\nu_{11}=i_7\{\nu_7,8\iota_{10},\nu_{10}\}
            =i_7(4\sigma')=0.
 \tag{D.2.60q}
\]
The singleton bracket is the one proved in D.2.55 from
the ordinary $\sigma'''$ bracket and its two suspensions;
see also \cite[Section~2.5 and accompanying program]
{FGHStricklandToda}.  Its indeterminacy is
$\nu_7\pi_{14}S^{10}+\pi_{11}S^7\nu_{11}=0$.
The last equality in \textup{(D.2.60q)} is actual:
$2i_7=0$ and $[i_7,i_7]=0$ because $S^7$ is an
$H$-space, so the cyclic Hilton correction vanishes.

If $(mp+\epsilon\zeta)\nu_{11}$ lies in the old bottom
subgroup, pinching $N$ gives $2m\nu_{11}=0$ in
$\pi_{14}S^{11}=\mathbb Z/8\{\nu_{11}\}$.
Thus $4\mid m$, and \textup{(D.2.60q)} kills that part.
Equation \textup{(D.2.60p)} then forces $\epsilon=0$.
We have proved the complete intersection statement
\[
 \bigl(\pi_{11}J\circ\nu_{11}\bigr)
                    \cap\langle i_7\sigma'\rangle=0.
 \tag{D.2.60r}
\]
For the actual attachment of all twelve-cells, the
Blakers--Massey comparison from the wedge of attaching
$S^{11}$ spheres to the homotopy fiber is
sixteen-connected.  Thus its entire kernel on old
$\pi_{14}$ consists of the $\nu_{11}$ composites of
the attaching maps.  Retracting \textup{(D.2.60l)} to
$J$ and applying \textup{(D.2.60r)} excludes a bottom
relation.  The same comparison is fourteen-connected
for the ten-cells and fifteen-connected for the
eleven-cells; their possible contributions vanish by
$\pi_{14}S^9=\pi_5^s\mathbb S=0$ and
$\pi_{14}S^{10}=\pi_4^s\mathbb S=0$.
Together with the nine-stage argument, this proves
\textup{(D.2.60b)}.  No comparison in this proof
replaces the actual degree-fourteen homotopy group by
its stabilization, and no higher-cell conclusion follows.
\end{proof}

\begin{claim}{D.2.61}{Proved}{A secondary operation kills the remaining Moore obstruction}
Put $L=Y\langle7\rangle$, $E=\Omega^2L$, and let
$\widetilde\alpha:S^7\to L$ be the lift of $\alpha=u_4\nu_4$.
The bottom image of $h_0h_3$ in
$\operatorname{Ext}_{\mathcal A}^{2,16}(H^*L,\FF_2)$ is
a nonzero permanent cycle.  In particular
\[
                 2\widetilde\alpha\sigma\ne0
                       \quad\hbox{in }\pi_{14}^sL.
                 \tag{D.2.61a}
\]
The remaining stable Atiyah--Hirzebruch differential of
D.2.54 is zero on its whole source:
\[
 d_8:E^8_{15,0}\longrightarrow E^8_{7,7}=\ZZ/4
       \quad\hbox{is zero},\qquad E^\infty_{7,7}=\ZZ/4.
       \tag{D.2.61b}
\]
Consequently the actual Moore composite and the actual
sphere composite have the values
\[
                h_Ml_4=0,\qquad
                q=\widetilde\alpha\sigma'
                      \text{ has exact order two}.
                \tag{D.2.61c}
\]
The integer $n$ in D.2.53 is even, and the detecting
integral class of D.2.54 is a stable Hurewicz class.
These statements determine the specified Moore obstruction.
They do not construct the full fourth-suspension extension
or evaluate the nullification in D.2.4, which remains open.
\end{claim}
\begin{proof}
All unmarked cohomology in this proof has coefficients
$\FF_2$.  Write $a=a_7\in H^7L$ for the bottom class.
D.2.57 and instability give
\[
       \operatorname{Sq}^1a=\operatorname{Sq}^6a
                         =\operatorname{Sq}^8a=0.
       \tag{D.2.61d}
\]

\medskip\noindent
\textbf{The relation and a complete four-term functional.}
The Adem relations give
\[
 R=\operatorname{Sq}^1\operatorname{Sq}^8
   +\operatorname{Sq}^2\operatorname{Sq}^1\operatorname{Sq}^6
   +\operatorname{Sq}^8\operatorname{Sq}^1=0.
   \tag{D.2.61e}
\]
Indeed the first term is $\operatorname{Sq}^9$, and the
middle term is
$\operatorname{Sq}^2\operatorname{Sq}^7
 =\operatorname{Sq}^9+\operatorname{Sq}^8\operatorname{Sq}^1$.
In the normalized bar complex dual to the left-comodule
cobar convention of D.2.57, consider
\[
 r=[\operatorname{Sq}^1|\operatorname{Sq}^8]a
  +[\operatorname{Sq}^2\operatorname{Sq}^1|
                                  \operatorname{Sq}^6]a
  +[\operatorname{Sq}^8|\operatorname{Sq}^1]a.
  \tag{D.2.61f}
\]
The bar differential is
$d([b|c]m)=[bc]m+[b](cm)$.
Thus \textup{(D.2.61d)--(D.2.61e)} make $r$ a cycle.
Its pairing defines a functional $\lambda_R$ on
$\operatorname{Ext}_{\mathcal A}^{2,16}(H^*L,\FF_2)$.
In Milnor coordinates its entire support is
\[
 [\xi_1|\xi_1^8]A,\quad
 [\xi_1^3|\xi_1^6]A,\quad
 [\xi_2|\xi_1^6]A,\quad
 [\xi_1^8|\xi_1]A,
 \tag{D.2.61g}
\]
where $A$ is dual to $a$.  In particular
\[
              \lambda_R([\xi_1|\xi_1^8]A)=1.
              \tag{D.2.61h}
\]
This is a symbolic certificate on the full module.
Independently, the authored script
\texttt{code/EH5\_secondary\_relation\_certificate.py}
checks these four terms against all 166 incoming columns
of the actual degree-fifteen Milnor input and records
every supported term in
\texttt{code/outputs/EH5\_secondary\_relation\_certificate.json}.
No secondary topological assertion is inferred from that
finite algebra check alone.

\medskip\noindent
\textbf{An actual universal choice of the secondary operation.}
Let $H=H\FF_2$ and let $\mathcal P$ be the fiber of the
stable map
\[
 (\operatorname{Sq}^1,\operatorname{Sq}^6,\operatorname{Sq}^8):
 \Sigma^7H\longrightarrow
           \Sigma^8H\oplus\Sigma^{13}H\oplus\Sigma^{15}H.
 \tag{D.2.61i}
\]
The row
$(\operatorname{Sq}^8,\operatorname{Sq}^2\operatorname{Sq}^1,
\operatorname{Sq}^1)$ lands in $\Sigma^{16}H$.
A fixed nullhomotopy of its composite, supplied by $R$,
defines a secondary map $\mathcal P\to\Sigma^{15}H$.
Its restriction to the fiber of
$\mathcal P\to\Sigma^7H$ is the loop of that row.
This is the universal construction of secondary operations;
the precise fiber restriction, existence, and uniqueness
up to a primary operation are given by
\cite[Section~3.6, especially Theorems~3.6.1--3.6.2 and
Lemma~3.6.3, pp.~68--74]{EH5Adams1960Secondary}.

Use the actual space
\[
 \mathcal F=\operatorname{hofib}
  \bigl(\operatorname{Sq}^6\rho:
       K(\ZZ/4,7)\longrightarrow K(\FF_2,13)\bigr).
 \tag{D.2.61j}
\]
Here $\rho$ denotes reduction modulo two.
On the base $K(\ZZ/4,7)$ choose based nullhomotopies
of $\operatorname{Sq}^1\rho$ and
$\operatorname{Sq}^8\rho$.
They exist respectively from the mod-four lift and
instability.  Together with the tautological fiber path
for $\operatorname{Sq}^6\rho$, they give a map
$\mathcal F\to\Omega^\infty\mathcal P$ and hence a
particular class $\phi\in H^{15}\mathcal F$.
On its fiber $K(\FF_2,12)$ the two base-chosen
nullhomotopies are constant.  Therefore
\[
             \phi|_{K(\FF_2,12)}
                    =\operatorname{Sq}^2\operatorname{Sq}^1
                                                \iota_{12}.
             \tag{D.2.61k}
\]
This fixes an actual representative; it does not replace
that restriction by a coset allowing an arbitrary
$\operatorname{Sq}^3\iota_{12}$ term.

Choose a mod-four lift $\widehat a\in H^7(L;\ZZ/4)$
of $a$.  Such a lift exists by
$\operatorname{Sq}^1a=0$.  Since
$\operatorname{Sq}^6a=0$, it lifts to an actual map
$f:L\to\mathcal F$.
The class $f^*\phi$ is an admissible value of the
degree-eight operation $\Phi_R(a)$.
Its indeterminacy is
\[
 \operatorname{Sq}^8H^7L+
 \operatorname{Sq}^2\operatorname{Sq}^1H^{12}L+
 \operatorname{Sq}^1H^{14}L.
 \tag{D.2.61l}
\]
Every class annihilated by positive homology Steenrod
operations annihilates this entire indeterminacy.

Double looping \textup{(D.2.61j)} gives an actual splitting
\[
 \Omega^2\mathcal F\simeq
            K(\ZZ/4,5)\times K(\FF_2,10)
 \tag{D.2.61m}
\]
over $K(\ZZ/4,5)$: the defining map is now
$\operatorname{Sq}^6$ on a degree-five class and is null.
Choose one such splitting.  There are no mixed positive
products in degree thirteen, since $5+10>13$.
Writing $z_{10}$ for the second factor's fundamental
class, \textup{(D.2.61k)} implies
\[
 \sigma^2\phi=\operatorname{Sq}^2\operatorname{Sq}^1z_{10}
                       +\eta,
       \qquad \eta\in H^{13}K(\ZZ/4,5).
 \tag{D.2.61n}
\]
Serre's computation
\cite[Section~2, no.~11, Theorem~4 and its proof,
pp.~204--205]{I_Serre1953}
says every polynomial generator of this last algebra
is a Steenrod operation on the fundamental reduction
$x_5$ or on $y_6=\rho\beta_4\widehat x_5$.
In degree thirteen, $\eta$ is therefore a sum of
positive-primary images of $x_5,y_6$ and decomposables.
The higher coefficient connecting map here is the one
for $0\to\FF_2\to\ZZ/8\to\ZZ/4\to0$;
it is $\rho\beta_4$ and commutes with suspension.
Under $\Omega^2f$, both $x_5$ and $y_6$ are double
cohomology suspensions of the actual classes
$a$ and $\rho\beta_4\widehat a$ on $L$.

\medskip\noindent
\textbf{Evaluation on every possible Adams source.}
Let $x\in H_{15}(L;\FF_2)$ be annihilated by every
positive homology Steenrod operation.
D.2.58 proves that $\sigma^2:H_{13}E\to H_{15}L$
is onto.  We can choose a primitive preimage $z$.
Indeed $E$ is four-connected, so its degree-thirteen
reduced diagonal has only bidegrees $(5,8),(6,7)$
and their reverses.  All factors in those degrees are
primitive.  The nine products
$uH_8E$ and $H_6E\,H_7E$ have independent symmetric
diagonals and realize all these tensors.  Subtracting
a suitable sum of them makes a preimage primitive;
these products have zero double homology suspension.

Naturality makes $\operatorname{Sq}^2_*z$ primitive
and places it in the complete D.2.56 kernel
\[
 \ker(\sigma^2:H_{11}E\to H_{13}L)
                    =\langle uD,uE_6,\xi u\rangle.
 \tag{D.2.61o}
\]
The first two products have independent nonzero
reduced diagonals, so only $\xi u$ could remain.
Its coefficient is zero already in the ambient free
$E_2$ homology of $T=\Omega^2\Sigma^2B$.
The only degree-six restricted-Lie inputs for a
degree-thirteen restricted generator are $f,x,\{a,b\}$;
their Bocksteins are $0,d+e,0$.
Thus the top Nishida formula cannot produce the outer
restricted generator $\xi u$, where $u=\xi a$.
Ordinary labels, Lie words and products cannot produce
that outer restricted generator either.
These are the free-$E_2$ basis and top-square formulas
already used in D.2.56--D.2.58
\cite[Chapter~III, Theorem~1.3(3), pp.~217--218]{I_CohenLadaMay1976}.
The established injection $H_{11}E\to H_{11}T$ gives
\[
                         \operatorname{Sq}^2_*z=0.
                         \tag{D.2.61p}
\]
This argument uses no new degree-thirteen cover-edge
calculation and does not assert $\operatorname{Sq}^1_*z=0$.

Pair \textup{(D.2.61n)} with $(\Omega^2f)_*z$.
The decomposable part of $\eta$ pairs zero by primitivity.
Every positive-primary part is a double suspension of
a positive operation on one of the two stated classes
of $L$, so it pairs zero with the primary-annihilated
image $x$.  The fiber term pairs as
$\operatorname{Sq}^1_*\operatorname{Sq}^2_*z=0$.
Consequently
\[
             \langle\Phi_R(a),x\rangle=0
            \quad\text{for every such }x.
             \tag{D.2.61q}
\]
The controlled universal representative is essential:
arbitrary secondary indeterminacy after looping could
include $\operatorname{Sq}^1H^{12}E$, which need not
pair zero with the chosen $z$.

\medskip\noindent
\textbf{The precise secondary-operation/Adams comparison.}
We verify the needed comparison rather than identifying
operations merely by their degree.  Put $M=H^*L$.
Use the actual GEM Adams coresolution of
$X=\Sigma^\infty L$, identified with the nonminimal
free bar resolution
\[
 P_0=\mathcal A\otimes M,\qquad
 P_1=\mathcal A\otimes\overline{\mathcal A}\otimes M,
 \qquad
 P_2=\mathcal A\otimes\overline{\mathcal A}^{\otimes2}
                                                   \otimes M.
 \tag{D.2.61r}
\]
Its coherent zero tracks come from the actual Adams
fiber tower, as in
\cite[Theorem~6.9 and Section~7]{EH5BauesJibladze2004};
they are not assigned solely from a primary chain complex.
Here the indicated free resolution can be realized explicitly.
Choose $g_0:X\to A_0$ to induce its augmentation
$P_0\to M$.  Its cofiber $C_0$ has cohomology
$\ker(P_0\to M)$, because the augmentation is onto
in every degree.  Realize the free surjection from $P_1$
onto this kernel by $C_0\to A_1$, and repeat with its
cofiber and the free surjection from $P_2$.
The resulting cofiber triangles give the required zero
tracks and their coherence.  There are only finitely many
generators in each bounded degree, so the connective GEM
products and these cohomology identifications are legitimate.
Write the corresponding GEM maps as
$X\xrightarrow{g_0}A_0\xrightarrow{v_0}A_1
\xrightarrow{v_1}A_2$, with null tracks
$\eta_0$ on $v_0g_0$ and $\eta_1$ on $v_1v_0$.
Coherence says $\eta_1g_0=v_1\eta_0$.

In the free resolution,
\[
 d_P([b|c]m)=b[c]m+[bc]m+[b](cm).
 \tag{D.2.61s}
\]
For the lift of $r$ in \textup{(D.2.61f)}, the last
two sums vanish by \textup{(D.2.61d)--(D.2.61e)}.
Its coordinate of $v_1$ is therefore exactly the row
$(\operatorname{Sq}^8,
\operatorname{Sq}^2\operatorname{Sq}^1,
\operatorname{Sq}^1)$ on the three coordinates of $A_1$
in degrees $8,13,15$.
The expression $r$ is a nonzero scalar combination of
three free $P_2$ generators, so an ordinary change of
free basis makes it a single coordinate as used here.
Let $H_R$ be the chosen relation nullhomotopy on $A_0$
factoring through its bottom $\Sigma^7H$ coordinate.
The difference of tracks
$\eta_{1,r}-H_R$ is a stable map
$q_0:A_0\to\Sigma^{15}H$.
Coherence gives
\[
                       g_0^*q_0=\Phi_R(a)
                         \quad\text{modulo primary indeterminacy}.
                       \tag{D.2.61t}
\]

A class in $\operatorname{Ext}^{0,15}_{\mathcal A}(M,\FF_2)$
is exactly a primary-annihilated $x$ as above.
Represent it by $c:S^{15}\to A_0$, chosen with nonzero
components only in degree-fifteen GEM coordinates, and
choose a null track $\gamma$ for $v_0c$.
The actual secondary differential is represented by
the glued track $\eta_1c+v_1\gamma$.
This is the construction in
\cite[equation~(4.1), Theorem~4.2, and Theorem~7.3 with
its proof, pp.~13--15 and 19--21]{EH5BauesJibladze2004},
which identifies it with the Adams $d_2$.
On the $r$ coordinate, the relevant components of
$\gamma$ map to
$\Omega\Sigma^8H$, $\Omega\Sigma^{13}H$ and
$\Omega\Sigma^{15}H$; their groups of maps from
$S^{15}$ are all zero.  The track $H_Rc$ is constant,
since the chosen $c$ has zero bottom $\Sigma^7H$
component.  Thus only $q_0c$ remains, giving the exact
pairing formula
\[
       \lambda_R(d_2x)
          =\langle g_0^*q_0,x\rangle
          =\langle\Phi_R(a),x\rangle.
       \tag{D.2.61u}
\]
Changing the three primary nullhomotopies adds precisely
\textup{(D.2.61l)}.  Changing $H_R$ adds a degree-eight
primary operation on $a$, by the same universal
construction or Adams's uniqueness theorem.
Every such change pairs zero with $x$.
This accounts for both the source-dependent tracks
and the full indeterminacy in \textup{(D.2.61u)}.

\medskip\noindent
\textbf{Permanence and the actual conclusion.}
Equations \textup{(D.2.61h)}, \textup{(D.2.61q)} and
\textup{(D.2.61u)} show that no incoming Adams $d_2$
can hit the bottom $h_0h_3$ image.
There is no later incoming differential to filtration
two, and sphere naturality prevents outgoing
differentials, as in D.2.58.
The connective finite-type Adams convergence theorem
used in D.2.57 therefore detects a nonzero image of the
actual bottom $2\sigma$ in the two-completed stable
homotopy of $L$
\cite[Part~III, Theorem~15.1 and its construction,
pp.~370--377]{EH5Adams1974AHSS}.
A zero actual stable class would have zero completed
image, so \textup{(D.2.61a)} follows without requiring
a separate completion-isomorphism assertion.

The bottom subgroup $E^8_{7,7}$ is $\ZZ/4$ by D.2.57.
Its twice-generator is the actual class detected in
\textup{(D.2.61a)}.  Any nonzero subgroup of this
$\ZZ/4$ contains that element.  Thus a nonzero incoming
$d_8$ would kill the actual bottom stable class.
There is no smaller cell filtration, since $L$ is
six-connected; hence it cannot disappear into a lower
filtration term.  This proves $d_8=0$ on its whole
source.  No later incoming differential is possible
because the sphere spectrum is connective, and no
outgoing differential from the bottom filtration is
possible.  This proves \textup{(D.2.61b)}.

The detecting class $\omega$ of D.2.54 now survives.
It has no later possible outgoing differential and no
incoming differential in the $q=0$ row.  Finite-range
cellular convergence makes it the integral Hurewicz
image of a stable map.  Its $\theta$ pairing is one,
so the exact D.2.51 criterion gives $h_Ml_4=0$.
Equivalently, the exact D.2.54 equation gives
$2n\bar\alpha=0$ in $\ZZ/4$, so $n$ is even and the
D.2.54 source-changing criterion holds with zero change.
Finally D.2.53 gives $2q=2h=0$, whereas the stabilization
of $q=\widetilde\alpha\sigma'$ is the nonzero class
$2\widetilde\alpha\sigma$ of \textup{(D.2.61a)}.
Thus $q$ has exact order two.
None of these implications supplies an extension over
all the other cells of the fourth-suspension source;
the main D.2.4 test is unchanged.
\end{proof}


\begin{claim}{D.3.1}{Proved}{There is no largest non-total reversible localizer}
Two distinct primes already give reversible localizers $W_{\mathbb F_p},W_{\mathbb F_q}$ having no common non-total reversible upper bound. Both lie between $W_{\mathbb Z}$ and $W_0$. Here non-total means strictly smaller than all functors; it is distinct from the technical properness condition of Part E.
\end{claim}
\begin{proof}
The coefficient-change theorem gives $W_{\mathbb Z}\subseteq W_{\mathbb F_p}$, and the actual map on degree-zero homology over a field detects a bijection of components, giving $W_{\mathbb F_p}\subseteq W_0$. Reversibility follows from the cofiber criterion for ordinary field homology. The same statements hold for $q$.

Let $Y=S^1[1/p]$ be the telescope of degree-$p$ circle maps, and let $i:S^1\to Y$ be its initial-stage map. The space $Y$ is a $K(\mathbb Z[1/p],1)$ and has reduced integral homology only in degree one, where it is $\mathbb Z[1/p]$. Consequently it is $\mathbb F_p$-acyclic. Its homotopy cofiber $C(i)$ is connected and has reduced integral homology only in degree one, where it is $\mathbb Z[1/p]/\mathbb Z$. Multiplication by $q\neq p$ is an automorphism of that group. The integral universal coefficient sequence therefore makes $C(i)$ $\mathbb F_q$-acyclic.

In any reversible upper bound, both $Y$ and $C(i)$ would be aspherical. Cofiber two out of three would make $S^1$ aspherical. Desuspension would make $S^0$ aspherical, and pointed cell constructions would then make every pointed space aspherical. The empty-to-point map has cofiber $S^0$, so the cofiber criterion includes it as well. Thus the localizer would be all functors. This proves the assertion without inferring totality merely from a full stable part, which would be invalid.
\end{proof}

\begin{claim}{D.3.2}{Open}{The fixed stable part and its projective-space interval}
For $\mathcal L=\loc(\mathbb S/2)$ the lower endpoint is the accessible $W_{\mathrm{rev}}$ of B.4 and the upper endpoint is $R(\mathcal L)=W_{\mathbb Z[1/2]}$. They are distinct by A.4--A.5. The existence of a strict intermediate reversible localizer, and the requested set-size assertion for an arbitrary fixed stable part, remain active questions here.
\end{claim}

\paragraph{Literature update (15 November 2026).}
\citet[Theorem~1.4]{EH5Putzstuck2026} proves that the smash-homology
Bousfield lattice has cardinality $2^{2^{\aleph_0}}$, including an antichain
of that size, and that every spectrum is Bousfield equivalent to a sum of
countable spectra. Here a Bousfield class is the kernel of $E\wedge-$.
The theorem counts these kernels; it does not bound all localizing
subcategories or the reversible localizers with a fixed $T(W)$.
It therefore does not settle D.3.2. The same paper constructs arbitrary
finite tensor-nilpotence heights and refutes the Hovey--Palmieri retract
conjecture; that conjecture concerns the Bousfield lattice, separately
from the $RT$-fixedness studied here.
\citet[Remarks~2.3, 2.9 and 2.10]{EH5BCHSZ2026} likewise distinguishes
this cardinality theorem from the general set-size question for
cohomological Bousfield classes, and records that the conjecture that
every localizing ideal of spectra is a homological Bousfield class
remains open. Its chromatic localization results do not resolve these
fixed-$T$ questions.

\begin{claim}{D.3.3}{Proved}{Every nonzero finite primary addition reaches the upper endpoint}
Fix a prime $p$. Let $V$ be a finite connected pointed complex with
nonzero reduced integral homology, all of it $p$-primary torsion.
Then
\[
 \operatorname{Rev}(B\ZZ/p\vee V)
       =W_{\ZZ[1/p]}=R(\operatorname{loc}(\mathbb S/p)).
 \tag{D.3.3.1}
\]
Consequently a proper intermediate localizer in the projective-space
interval cannot make any nonzero finite $p$-primary complex aspherical.
In particular adding a finite complex of type two, or of any higher
finite chromatic type, cannot produce such an intermediate.
\end{claim}
\begin{proof}
Write $W=\operatorname{Rev}(B\ZZ/p\vee V)$. Its two generators
are $\ZZ[1/p]$-acyclic, so $W\subseteq W_{\ZZ[1/p]}$.
Replacing $V$ by a sufficiently high suspension does not change
$W$, by reversibility. We may therefore suppose that $V$ is a
simply connected suspension. The nonzero finite spectrum
$\Sigma^\infty V$ has some finite $p$-type $n+1\geq1$.
Let $d$ be the lowest degree with $\pi_dV\ne0$, and choose an
odd integer $m\geq3$ with $m-1\geq d$.

Distinguish the ordinary nullification $P_V$ from the Goodwillie
stages $G_s=P_s^{\mathrm{Id}}(S^m_{(p)})$.
The precise periodicity input is
\cite[Theorem~4.1]{EH5AroneMahowald1999}: the natural map
\[
                 S^m_{(p)}\longrightarrow G_{p^n}
 \tag{D.3.3.2}
\]
is a $v_j$-periodic equivalence for every $0\leq j\leq n$,
where $j=0$ means a rational homotopy equivalence.
The nullification comparison is
\cite[Corollary~4.8]{EH5Bousfield2001}, in the arbitrary finite
suspension form explicitly stated in
\cite[Theorem~3.1]{EH5BarthelHeutsMeier2021}: a map of
$d$-connected $p$-local spaces is a $P_V$-equivalence if and
only if it is a $v_j$-equivalence for $0\leq j\leq n$.
All the spaces $G_s$ are $(m-1)$-connected. To check the bound,
the derivative formula in
\cite[Theorem~2.1]{EH5AroneMahowald1999} gives the spectrum of
the $s$th layer as
\[
 E_s=\bigl(\operatorname{Map}_*(\Sigma K_s,
          \Sigma^\infty(S^m_{(p)})^{\wedge s})\bigr)_{h\Sigma_s}.
\]
The finite complex $\Sigma K_s$ has dimension $s-1$, so $E_s$
has no homotopy below $sm-s+1\geq m$. Homotopy orbits preserve
this lower bound. Starting with $G_1=QS^m_{(p)}$ and using
the fiber sequences $\Omega^\infty E_s\to G_s\to G_{s-1}$
proves the assertion by induction. These spaces are $p$-local
as well, by the same induction. Thus (D.3.3.2) is a
$P_V$-equivalence and hence belongs to $W$.

We next remove the finite Goodwillie tower by $B\ZZ/p$
nullification. For $s>1$, the spectrum $E_s$ is rationally
zero, by \cite[Proposition~3.1]{EH5AroneMahowald1999}.
Since it is $p$-local, $E_s[1/p]=0$; and the preceding
connectivity estimate makes it a connected spectrum.
\cite[Theorem~3.2]{EH5Levy2025} therefore gives
\[
                    P_{B\ZZ/p}\Omega^\infty E_s\simeq *.
 \tag{D.3.3.3}
\]
Every projection $G_s\to G_{s-1}$ is consequently a
$P_{B\ZZ/p}$-equivalence. Here only the following elementary
nullification property is used: a fibration with connected
$A$-acyclic fibers is a $P_A$-equivalence. Indeed, mapping its
total space into an $A$-null target is the homotopy limit over
the base of the fiber mapping spaces; the natural maps from
the target to those mapping spaces are equivalences. This
identifies the resulting homotopy limit with the mapping space
from the base. No preservation of periodic homotopy by
$P_{B\ZZ/p}$ is asserted.

The map
\[
                 QS^m_{(p)}\longrightarrow K(\QQ,m)
 \tag{D.3.3.4}
\]
is also a $P_{B\ZZ/p}$-equivalence. Its fiber is the infinite
loop space of
\[
     F=\operatorname{fib}(\Sigma^m\mathbb S_{(p)}
                            \longrightarrow\Sigma^mH\QQ).
\]
The spectrum $F$ satisfies $F[1/p]=0$, and its lowest possible
nonzero homotopy is in degree $m-1\geq2$. Apply Levy's theorem
again. The target $K(\QQ,m)$ is $B\ZZ/p$-null, as is also
immediate from the rational cohomology of $B\ZZ/p$.

Both ordinary nullification-equivalence classes used here are
contained in $W$, since $W$ contains their generating maps.
Combining (D.3.3.2)--(D.3.3.4), we have a natural $W$-equivalence
\[
                         r:S^m_{(p)}\longrightarrow K(\QQ,m).
\]
Naturality for a map $p:S^m_{(p)}\to S^m_{(p)}$ of degree $p$
gives $r p\simeq p r$, with the lower $p$ the multiplication
automorphism of $K(\QQ,m)$. Two out of three therefore puts
the degree-$p$ map of $S^m_{(p)}$ in $W$. Its homotopy cofiber
is the usual finite Moore space $M(\ZZ/p,m)$: the comparison
from the integral degree-$p$ cofiber is an integral homology
equivalence of simply connected spaces. The reversible cofiber
criterion makes this Moore space $W$-aspherical.

Suspension and desuspension now put $M(\ZZ/p,2)$ in the
aspherical class. Its ordinary nullification kills every
simply connected space with $p$-primary homotopy groups, by
\cite[Theorem~7.5]{EH5Bousfield1997}. The Serre-class theorem
identifies these with the simply connected spaces having
$p$-primary reduced integral homology. If $X$ is any connected
space with such homology, its suspension is of this form;
reversibility makes $X$ aspherical as well. These are exactly
the aspherical spaces of $W_{\ZZ[1/p]}$. Applying the cofiber
criterion gives $W_{\ZZ[1/p]}\subseteq W$ and proves (D.3.3.1).
\end{proof}

\begin{claim}{D.3.4}{Proved}{Large target families and the collapse of Moore wedges}
Let $\kappa$ be an infinite cardinal and $d\geq3$, and put
\[
 Y_\kappa^d=\bigvee_{\alpha<\kappa}P^d(2),\qquad
 \mathcal N_\kappa^d=
       \{X:\Map_*(X,Y_\kappa^d)\simeq *\}.
\]
Then $\mathcal N_\kappa^d$ is closed under desuspension, for
arbitrary pointed source spaces. In fact all its cardinal
parameters collapse:
\[
 \mathcal N_\kappa^d
   =\{X:\Map_*(X,P^d(2))\simeq *\}
   =\{X:\Map_*(X,P^3(2))\simeq *\}.
 \tag{D.3.4.0}
\]
Consequently its cofiber test,
intersected with $R(\loc(\mathbb S/2))$, is a reversible
fundamental localizer $W_\kappa^d$ satisfying
\[
 W_{\mathrm{rev}}\subseteq W_\kappa^d
       \subsetneq R(\loc(\mathbb S/2)),\qquad
 T(W_\kappa^d)=\loc(\mathbb S/2).
 \tag{D.3.4.1}
\]
There is also such a construction using all simply connected
finite-dimensional CW targets with at most $\kappa$ cells.
The infinite Moore wedges give one localizer, independently of
cardinality and dimension. D.3.5 below shows that the broader
target families also give one localizer, namely
$W_S\cap R(\loc(\mathbb S/2))$, independently of cardinality.
Whether either construction differs strictly from
$W_{\mathrm{rev}}$ remains unresolved here.
\end{claim}
\begin{proof}
Fix a source $X$ and again put
$\mathcal R_X=\{Z:\Map_*(X,Z)\simeq *\}$.
We first prove the large-wedge assertion, including its
bottom-dimensional case without a localization convention for
$\RP^2$.

\medskip\noindent
\textbf{A three-factor finite retraction.}
For $A=P^{d-1}(2)$, let $a=d-2\geq1$ and $s=2a+1$.
There is an actual map
\[
 \alpha:S^{2a+1}\longrightarrow A\wedge A,
 \qquad \alpha_*(\iota)=uv+vu\quad\hbox{over }\FF_2.
\]
Here is an elementary verification valid also for $a=1$.
The bottom skeleton of $A\wedge A$ has one cell in degree
$2a$ and two cells in degree $2a+1$. The two attaching degrees
are $2$ and $(-1)^a2$. A unimodular change of the two attaching
spheres changes this row to $(2,0)$. The skeleton is therefore
homotopy equivalent to $P^{2a+1}(2)\vee S^{2a+1}$.
The sphere summand is represented by the integral cycle
$v\otimes u-(-1)^a u\otimes v$; its image modulo two is
$uv+vu$. Include this sphere into the full smash product to
obtain $\alpha$. This is the Moore-space instance of
\cite[Proposition~2.1]{I_ChenWu2013}.

The suspended space $\Sigma A^{\wedge3}$ is simply connected
and has finite $2$-primary integral homology, even when
$A=\RP^2$. The complementary Dynkin--Specht--Wever telescopes
of D.1.8 can therefore be formed on this suspended space:
multiplication by $3$ is a homotopy equivalence and can be
inverted. The same eight-word calculation (D.1.8.4) gives
\[
 \Sigma^{s+1}A=P^{3d-3}(2)
       \quad\hbox{as a homotopy retract of}\quad
       \Sigma A^{\wedge3}.
 \tag{D.3.4.2}
\]
Explicitly, project to the $\beta_3$ telescope and use the
equivalence from $\Sigma^{s+1}A$ induced by
$\Sigma(\id_A\wedge\alpha)$; its two homology images are
$uuv+vuu$ and $uvv+vvu$. The inclusion of the telescope and
the inverse equivalence supply the two retraction maps.
No assertion that $\RP^2$ itself is $2$-local is used.

\medskip\noindent
\textbf{A retraction for the entire infinite wedge.}
Partition the index set into three subsets, each of cardinality
$\kappa$, and write
\[
 Y_\kappa^d\simeq\Sigma U_1\vee\Sigma U_2\vee\Sigma U_3,
 \qquad U_i=\bigvee_{\alpha<\kappa}A.
\]
Use Hilton--Milnor with these three connected spaces as its
three letters. The factor for the basic word $[[1,2],3]$ makes
$\Omega\Sigma(U_1\wedge U_2\wedge U_3)$ a homotopy retract of
$\Omega Y_\kappa^d$. The theorem uses finitely many letters;
there is no finiteness assumption on each letter's CW complex.
Its factors' connectivities tend to infinity with word length,
so the factor inclusions and projections are legitimate.

The smash product here is the wedge of copies of $A^{\wedge3}$
indexed by $\kappa^3$. Choose a diagonal set of $\kappa$ triples,
include those summands, and collapse the others. Wedge together
the actual finite retractions (D.3.4.2) over this diagonal set.
After looping, their composite with the Hilton--Milnor
retraction gives
\[
             \Omega Y_\kappa^{3d-3}
              \quad\hbox{as a retract of}\quad
             \Omega Y_\kappa^d .                   \tag{D.3.4.3}
\]
This construction includes $d=3$: in that case (D.3.4.2) is
the retract $P^6(2)$ of $\Sigma(\RP^2)^{\wedge3}$.
The use of three whole wedges, followed by an actual wedge
projection, avoids replacing an infinite collection of
same-connectivity Hilton--Milnor factors by a product.

Suppose $\Omega Y_\kappa^d\in\mathcal R_X$, and set
$D=3d-3>d$ and $K=Y_\kappa^{D-1}$.
Equation (D.3.4.3) gives $\Omega\Sigma K\in\mathcal R_X$.
Since $\kappa+\kappa=\kappa$, the space $\Sigma K$ is
homotopy equivalent to $\Sigma K\vee\Sigma K$.
The two-letter Hilton--Milnor theorem now puts
\[
                     \Omega\Sigma K^{\wedge q}\in\mathcal R_X
                     \qquad(q\geq1).
\]
Apply the two-copy argument proved in D.1.8, treating the
whole simply connected space $K$ as one letter. For every
finite wedge $V$ of copies of $K$, its iterated James
splittings have finite multiplicities at each smash weight.
The tails become arbitrarily highly connected, regardless of
the cardinality of each weight. Thus all the positive iterates
$(\Omega\Sigma)^{r+1}V$ belong to $\mathcal R_X$, and the
convergent simply connected resolution gives
$V\in\mathcal R_X$.
Barratt's finite cone cubes lower this whole finite-wedge
family from $Y_\kappa^{D-1}$ to $Y_\kappa^d$.
We have proved
\[
 \Map_*(X,\Omega Y_\kappa^d)\simeq *
       \quad\Longrightarrow\quad
 \Map_*(X,Y_\kappa^d)\simeq *.
\]
All inverse limits are in the targets. This proves the
arbitrary-source desuspension assertion.

\medskip\noindent
\textbf{Why the Moore cardinal parameter collapses.}
We strengthen the result just proved. Suppose first that $X$
is countable and $P^d(2)\in\mathcal R_X$. The proof of D.1.8,
including its iteration and finite cone cubes, puts every
finite bouquet of copies of $P^d(2)$ in $\mathcal R_X$.
Write $A=P^{d-1}(2)$ and $F_q=\Sigma A^{\wedge q}$ for
$q\geq1$. The finite-letter Hilton--Milnor theorem therefore
gives $\Omega F_q\in\mathcal R_X$ for every $q$.
Each $F_q$ is a finite simply connected complex with finite
$2$-primary integral homology. Its homotopy groups are finite
$2$-groups, by the Serre-class version of the finiteness
theorem. Also its connectivity tends to infinity with $q$.

Choose an exhaustion $X=\bigcup_i X_i$ by finite pointed
subcomplexes. It is not necessary that $X$ have finite skeleta.
For fixed $q$, the groups
\[
                 [X_i,\Omega F_q]_*
\]
are finite. The Milnor sequence and the vanishing of $\lim^1$
for towers of finite groups identify their inverse limit with
$[X,\Omega F_q]_*=0$. A tower of finite groups with zero
inverse limit is pro-zero: at any fixed stage its descending
images stabilize, and the stabilized image is the image of
the inverse limit, so is trivial. The same assertion holds
with $X_i$ replaced by any fixed suspension $\Sigma^rX_i$.

The Hilton--Milnor decomposition for a countable bouquet,
obtained compatibly from its finite subbouquets, is the weak
product of copies of $\Omega F_q$, indexed by basic words of
length $q$. For any finite $X_i$, maps to this weak product
factor through a finite subproduct. Thus it identifies the
restriction tower $[X_i,\Omega Y_{\aleph_0}^d]_*$, as a tower
of pointed sets, with the weak products of the factor towers.
For fixed $i$, only finitely many values of $q$ can be visible,
by the dimension of $X_i$ and the connectivity of $F_q$.
Choose $j\geq i$ that kills the image of every one of those
finitely many factor types. It then kills every copy of each
type, regardless of its word label. Consequently
\[
 [X_j,\Omega Y_{\aleph_0}^d]_*\longrightarrow
 [X_i,\Omega Y_{\aleph_0}^d]_*
             \quad\text{is the zero map}.
 \tag{D.3.4.4}
\]
This step uses the weak product itself, not an identification
of it with the full product. The same proof applies to all
the suspensions $\Sigma^rX_i$. These pro-zero towers have
zero inverse limit and zero $\lim^1$, so their Milnor
sequences give
\[
             \Map_*(X,\Omega Y_{\aleph_0}^d)\simeq *.
\]
The desuspension result already proved makes
$Y_{\aleph_0}^d\in\mathcal R_X$. Every map from a suspension
of countable $X$ to $Y_\kappa^d$ factors through a countable
subbouquet: the image of each closed source cell lies in a
finite target subcomplex. That subbouquet is a retract of
$Y_{\aleph_0}^d$. Hence $Y_\kappa^d\in\mathcal R_X$ for all
$\kappa$. This is the uniform finite-factor-type argument
behind \cite[Proposition~6]{I_Strom2012}, with a finite-cell
exhaustion and finite homotopy groups replacing its
finite-type and rational hypotheses.

To pass from countable to arbitrary $X$, apply the bounded
extension lemma of G.7 to a finite simplicial Moore model
$K$ and its Bousfield--Kan $\FF_2$-resolution. Exactly the
compact model construction in G.7 makes
\[
                  Z=\operatorname{Tot}(R^{\bullet+1}K)
\]
a Kan complex with compact metrizable levels. It represents
the $2$-completion of $P^d(2)$. Since the latter space is
simply connected and all its homotopy groups are finite
$2$-groups, the completion map is a weak equivalence. Apply
the lemma simultaneously to $Z,\Omega Z,\Omega^2 Z,\ldots$.
It says that every countable simplicial subcomplex of a
$P^d(2)$-null $X$ is contained in a countable $P^d(2)$-null
subcomplex. These countable null subcomplexes form a filtered
system with homotopy colimit $X$. Each is $Y_\kappa^d$-null
by the countable argument. Nullity is closed under pointed
homotopy colimits, so $X$ is $Y_\kappa^d$-null as well.
The reverse implication uses the retraction onto one target
summand. Equality of the dimensions follows from D.1.8.
This proves (D.3.4.0) for arbitrary sources.

\medskip\noindent
\textbf{The family with a cell-cardinality bound.}
Let $\mathcal F_\kappa$ consist of representatives of the simply
connected finite-dimensional CW complexes with at most
$\kappa$ cells, and put
\[
 \mathcal N(\mathcal F_\kappa)=
       \{X:\Map_*(X,Z)\simeq *\text{ for all }Z\in\mathcal F_\kappa\}.
\]
Suspension, finite wedges and finite smash products preserve
this target family up to weak equivalence, since $\kappa$ is
infinite. Suppose $\Sigma X\in\mathcal N(\mathcal F_\kappa)$.
Then $\Omega Z\in\mathcal R_X$ for every
$Z\in\mathcal F_\kappa$. For fixed $Z$ in the family, each
finite weight truncation of
\[
 \Sigma(\Omega\Sigma)^rZ
     \simeq\bigvee_{m\geq1}\bigvee^{c_{r,m}}\Sigma Z^{\wedge m}
\]
is finite-dimensional and has at most $\kappa$ cells. Its loop
space is therefore in $\mathcal R_X$. Increasing connectivity
of the tails puts $(\Omega\Sigma)^{r+1}Z$ in that class; the
convergent resolution puts $Z$ there. This holds for every
$Z\in\mathcal F_\kappa$, proving desuspension closure of
$\mathcal N(\mathcal F_\kappa)$.

\medskip\noindent
\textbf{The stable part and the precise remaining question.}
Both null classes are closed under pointed homotopy colimits.
Desuspension and the mapping-space fibre sequence give all
three cofiber implications, as in A.3. Their cofiber tests
therefore define reversible fundamental localizers. Intersect
each with $R(\mathcal L)$, where $\mathcal L=\loc(\mathbb S/2)$.
Miller's theorem \cite[Theorem~A]{Miller1984} puts
$B\ZZ/2$ in each null class: all the targets used are
finite-dimensional, including the possibly enormous Moore
wedge. Its reduced suspension spectrum belongs to $\mathcal L$.
Thus the intersection contains $W_{\mathrm{rev}}$, and its
stable part contains $T(W_{\mathrm{rev}})=\mathcal L$.
The opposite containment follows from the intersection with
$R(\mathcal L)$, proving equality of stable parts.

The finite space $P^d(2)$ is in the upper endpoint's aspherical
class and is outside both new null classes. For the single
large wedge use its inclusion as one wedge summand, followed
by the retraction onto that summand; for
$\mathcal F_\kappa$ use $P^d(2)$ itself as target. These maps
are not null, proving the strict upper inclusion in (D.3.4.1).

For the family $\mathcal F_\kappa$, larger cardinal bounds
initially yield smaller or equal null classes. The complete
countable-source and compactness argument in D.3.5 shows that
all these interval localizers are equal to
$W_S\cap R(\loc(\mathbb S/2))$. Thus cardinality cannot
supply a separator within this family. Its distinctness
from $W_{\mathrm{rev}}$ remains the relevant unresolved issue.
\end{proof}

\begin{claim}{D.3.5}{Proved}{Cardinality bounds on finite-dimensional targets give the sphere test}
Let $\mathcal F_\kappa$ be the family from D.3.4, where
$\kappa$ is any infinite cardinal. Write
\[
 \mathcal P_2=\{X:\widetilde H_*(X;\mathbb Z[1/2])=0\}.
\]
Then
\[
 \begin{aligned}
 \mathcal N(\mathcal F_\kappa)\cap\mathcal P_2
   &=\Nul\cap\mathcal P_2\\
   &=\{X\in\mathcal P_2:
       \Map_*(X,Z)\simeq *\text{ for every simply connected}\\[-2pt]
   &\hspace{42mm}\text{finite-dimensional CW complex }Z\}.
 \end{aligned}
 \tag{D.3.5.1}
\]
In particular every localizer constructed from the families
$\mathcal F_\kappa$ in D.3.4 is precisely
$W_S\cap R(\loc(\mathbb S/2))$.
Increasing the cardinal bound produces no new element of that
interval. This equality does not decide whether this common
element differs from $W_{\mathrm{rev}}$.
\end{claim}
\begin{proof}
We first prove a countable-source extension of the sphere
criterion. Suppose $X$ has countably many cells and
\[
             \Map_*(X,S^m)\simeq *\quad(m\geq2).
 \tag{D.3.5.2}
\]
Choose an increasing exhaustion $X=\bigcup_{i\geq0}X_i$
by finite pointed subcomplexes. These are finite cell
subcomplexes, not necessarily skeleta of $X$.
For each $a\geq0$ and $m\geq2$ put
\[
       G_i(a,m)=[\Sigma^aX_i,\Omega S^m]_*.
\]
These are countable groups: the source is finite and the
homotopy groups of a sphere are countable, so the finite
obstruction calculation of the mapping set is countable.
The exact sequence for maps out of a sequential homotopy
colimit, applied to $\Sigma^aX$ and $\Sigma^{a+1}X$, gives
\[
       \lim_iG_i(a,m)=0,\qquad
       \lim_i^1G_i(a,m)=*.
\]
By \cite[Proposition~4(b)]{I_Strom2012}, this tower of
countable groups is Mittag--Leffler because its $\lim^1$ vanishes.
Its stable images are the images of its inverse limit:
after passing to stable images the bonding maps are
surjective, so every element of a stable image extends to a
compatible sequence. Since the limit is zero, every stable
image is zero. Thus each tower is \emph{pro-zero}:
for every $i,a,m$ there is $j\geq i$ for which
$G_j(a,m)\to G_i(a,m)$ is zero.

Let $W=\bigvee_{\alpha\in I}S^{n_\alpha}$ be any simply
connected wedge of spheres, of any cardinality. The
Hilton--Milnor decomposition identifies $\Omega W$ with
the weak product of spaces $\Omega S^{d(w)}$, indexed by
basic words $w$ in the sphere labels. This is a weak
product, not the categorical product. For a finite source
$\Sigma^aX_i$, maps and homotopies factor through finite
subproducts, giving a natural bijection of pointed sets
\[
 [\Sigma^aX_i,\Omega W]_*
   \cong\prod_w' [\Sigma^aX_i,\Omega S^{d(w)}]_*.
 \tag{D.3.5.3}
\]
One may verify the assertion first on finite subwedges by
Hilton--Milnor and then pass to the filtered union. Compactness
of the finite source, and of the source of a homotopy, gives
the displayed assertion for the entire wedge. The same
weak-product decomposition is the one used in
\cite[proof of Proposition~6]{I_Strom2012}.

For fixed $i,a$, only finitely many \emph{sphere dimensions}
$d(w)$ can contribute to the right-hand side of
(D.3.5.3): if $d(w)>\dim(\Sigma^aX_i)+1$, the factor is
trivial by connectivity. Choose $j$ large enough that all
the finitely many relevant sphere towers have zero
restriction to stage $i$. This one $j$ works for every
labelled copy of each dimension, regardless of the size of
$I$. Naturality of (D.3.5.3) therefore gives
\[
 [\Sigma^aX_j,\Omega W]_*
       \longrightarrow[\Sigma^aX_i,\Omega W]_*
       \quad\text{equal to zero}.
\]
The tower is pro-zero for every $a$. The sequential
homotopy-colimit exact sequence now gives
$\Map_*(X,\Omega W)\simeq *$ and, in particular,
\[
                \lim_i^1[\Sigma X_i,W]_* = *.
 \tag{D.3.5.4}
\]

It remains to check the unsuspended components of
$\Map_*(X,W)$; loop-space nullity alone would not do so.
By \cite[Corollary~12(a)]{I_Strom2012}, condition
(D.3.5.2) makes every finite subbouquet $V\subset W$
$X$-null. Given $f:X\to W$ and a finite $X_i$, its image
lies in some such $V$. The composite of $f$ with the
retraction $W\to V$ is nullhomotopic. Its restriction to
$X_i$, followed by $V\to W$, is $f|_{X_i}$, so the latter
is nullhomotopic. Thus $f$ represents the zero element of
$\lim_i[X_i,W]_*$, and (D.3.5.4) makes $f$ nullhomotopic.
Apply the same argument to every suspension of $X$.
It follows that $\Map_*(X,W)\simeq *$ for every simply
connected wedge of spheres.

The class $\{Z:\Map_*(X,Z)\simeq *\}$ is a resolving
kernel. \cite[Corollary~12(b)]{I_Strom2012} now puts in
it every simply connected finite-dimensional space.
This proves the countable-source assertion without a
finite-type hypothesis on $X$.

\medskip\noindent
\textbf{Comparison of integral and $2$-local sphere tests.}
For every $X\in\mathcal P_2$ and every simply connected
nilpotent space $Y$, the natural map induces
\[
        \Map_*(X,Y)\simeq\Map_*(X,Y_{(2)}).
 \tag{D.3.5.5}
\]
Here is an arithmetic-square verification, including the
possible infinite size of $X$. Such an $X$ is connected,
rationally acyclic, and mod-$q$ acyclic for every odd
prime $q$. Apply the general arithmetic fracture theorem
\cite[Theorem~13.1.4]{FGHMayPonto2012} first with the set
of all primes and then with the singleton set $\{2\}$.
The completion vertex in the first square is
$\prod_qY_q^{\wedge}$. Each factor for $q\ne2$ has
contractible pointed mapping space from $X$, since it is
mod-$q$-homology-local by
\cite[Theorems~11.1.1--11.1.2]{FGHMayPonto2012}.
Both rational vertices have contractible mapping space
from $X$, since rational nilpotent spaces are rational
homology-local. Mapping into the first square consequently
gives $\Map_*(X,Y)\simeq\Map_*(X,Y_2^{\wedge})$.
The second square gives the same equivalence with
$Y_{(2)}$ in place of $Y$. Naturality proves (D.3.5.5).
These fracture theorems require no finite-generation
hypothesis. Only their special case $Y=S^m$ is needed here.

\medskip\noindent
\textbf{Reduction of arbitrary sources to countable ones.}
Suppose now that $X\in\Nul\cap\mathcal P_2$, with no
cardinality assumption. G.7 supplies the countable
enlargement property for $\Nul$. The class $\mathcal P_2$
has the same property inside $X$: for a countable
simplicial subcomplex, adjoin the finite supports of
bounding chains in $X$ for all its reduced
$\mathbb Z[1/2]$-cycles, and iterate. There are only
countably many such finite chains, since
$\mathbb Z[1/2]$ is countable. Every cycle in the resulting
union bounds at the next stage. This also applies in
reduced degree zero and ensures connectedness.

Alternate the two countable enlargement constructions.
Each of the two null classes is closed under pointed
homotopy colimits, so the sequential union lies in both
classes. Hence every countable collection of simplices
of $X$ lies in a countable subcomplex $X'$ with
$X'\in\Nul\cap\mathcal P_2$. These subcomplexes form
a filtered system with homotopy colimit $X$. The use of
simplicial models and transfer to CW homotopy types is
exactly as in G.7; it makes no assumption that the
singular complex of a countable CW complex is countable.

By (D.3.5.5), each $X'$ satisfies the integral sphere
condition (D.3.5.2). The countable-source assertion
therefore makes $X'$ null into every simply connected
finite-dimensional target. Mapping the filtered homotopy
colimit into a fixed target gives a homotopy limit of
contractible spaces. Thus the same nullity holds for $X$.

Conversely, $\mathcal F_\kappa$ contains every simply
connected sphere. If $X\in\mathcal P_2$ is
$\mathcal F_\kappa$-null, (D.3.5.5) makes it null into
every $2$-local sphere, hence $X\in\Nul$. This proves
all equalities in (D.3.5.1). Finally the aspherical class
of $R(\loc(\mathbb S/2))$ is $\mathcal P_2$, and the
reversible cofiber criterion turns equality of the
displayed null classes into the stated equality of
fundamental localizers.
\end{proof}

\begin{claim}{D.3.6}{Proved}{The Moore target gives the sphere test in the primary interval}
Fix a prime $p$ and $d\geq3$. For every pointed source $X$,
\[
 \Map_*(X,P^d(p))\simeq *
 \quad\Longrightarrow\quad
 \Map_*(X,(S^n)^{\wedge}_p)\simeq *
                       \quad(n\geq2).
 \tag{D.3.6.1}
\]
The reverse implication holds if
$\widetilde H_*(X;\QQ)=0$. In particular, if
$\widetilde H_*(X;\ZZ[1/p])=0$, the single Moore target and
the whole family of simply connected $p$-completed spheres
have the same nullity test. For $p=2$, the localizer furnished
by the Moore wedges of D.3.4 is therefore exactly
\[
                 W_S\cap R(\loc(\mathbb S/2)).
 \tag{D.3.6.2}
\]
This identification does not establish that it is strictly
larger than $W_{\mathrm{rev}}$.
\end{claim}
\begin{proof}
Fix $X$ satisfying the left side of (D.3.6.1) and put
$\mathcal R_X=\{Z:\Map_*(X,Z)\simeq *\}$.
We give both the higher-exponent construction and the
unstable inverse-limit justification explicitly.

\medskip\noindent
\textbf{A smash-closed resolving family.}
Put $M=P^3(p)$ and
\[
 \mathcal E=\{\Sigma^i M^{\wedge q}:i\geq0,\ q\geq1\}.
\]
This family is closed under suspension and smash product,
and its members are simply connected. We claim that every
member of $\mathcal E$ belongs to $\mathcal R_X$.

The constructions of D.1.7--D.1.8 supply arbitrarily large
integers $D$ for which
$\Omega(P^D(p)\vee P^D(p))\in\mathcal R_X$.
For odd $p$, use the finite-bouquet conclusion of D.1.7.
For $p=2$, use the iterated two-copy retraction of D.1.8;
its upper dimensions increase without bound.
Set $K=P^{D-1}(p)=\Sigma^{D-4}M$. The two-letter
Hilton--Milnor argument then gives
$\Omega\Sigma K^{\wedge m}\in\mathcal R_X$ for all $m\geq1$.
For fixed $q$, apply Hilton--Milnor to
$\Sigma K^{\wedge q}\vee\Sigma K^{\wedge q}$.
Its factors have the form $\Omega\Sigma K^{\wedge qj}$,
and their connectivities tend to infinity, so its loop space
also belongs to $\mathcal R_X$. The two-copy
James-resolution argument of D.1.8, with $K^{\wedge q}$ as
one letter, now puts every finite bouquet of $K^{\wedge q}$
in $\mathcal R_X$.

Since
$K^{\wedge q}=\Sigma^{q(D-4)}M^{\wedge q}$,
the finite cone cubes of \cite[Theorem~8]{I_Strom2012}
descend this entire finite-bouquet family to
$\Sigma^iM^{\wedge q}$ whenever $0\leq i\leq q(D-4)$.
There are arbitrarily large choices of $D$, so the asserted
membership holds for every $i,q$.
Theorem~7 of the same paper now gives
$\Sigma\mathcal E^\vee\subseteq\mathcal R_X$, where the
superscript denotes finite-type wedges. Because
$\mathcal R_X$ is a strong resolving class, its Corollary~11
puts in $\mathcal R_X$ every simply connected space with
finite $\mathcal E^\vee$-cone length. No claim of automatic
closure under arbitrary cofiber sequences is being made.

\medskip\noindent
\textbf{The exponent can be increased.}
For $n\geq3$ and $r\geq1$, consider the actual composite
\[
 P^n(p)\xrightarrow{q}S^n
       \xrightarrow{\iota}P^{n+1}(p^r),
 \tag{D.3.6.3}
\]
where $q$ pinches the bottom cell and $\iota$ includes the
bottom sphere. Its cofiber has two cells in degree $n$ and
two in degree $n+1$, with integral cellular boundary, after
choosing orientations,
\[
                    \begin{pmatrix}p^r&1\\0&p\end{pmatrix}.
\]
Changing the bases of the two wedges of $n$-spheres and
cancelling the degree-one cell pair leaves the attaching
degree $p^{r+1}$. This is an unstable cellular argument:
for $n\geq2$, maps from an $n$-sphere to a wedge of
$n$-spheres are classified by their degrees. Hence the
cofiber of (D.3.6.3) is $P^{n+1}(p^{r+1})$.

Starting from the cofiber of $P^n(p)\to *$, this constructs
$P^{n+1}(p^r)$ with finite $\mathcal E$-cone length.
The same holds for every finite bouquet of these spaces.
The preceding resolving-class result makes all such bouquets
null. The Barratt cone cubes lower the bouquets with top
dimension four to top dimension three, if necessary.
Consequently
\[
               P^{n+1}(p^r)\in\mathcal R_X
                         \quad(n\geq2,\ r\geq1).
 \tag{D.3.6.4}
\]

\medskip\noindent
\textbf{The unstable Moore tower.}
For $n\geq2$, put $M_r=P^{n+1}(p^r)$ and use the transition
$M_{r+1}\to M_r$ induced by the commuting square of sphere
attaching maps whose top degree is $p$ and whose bottom
degree is $1$. The bottom-sphere inclusions are compatible.
We prove
\[
                   \holim_r M_r\simeq(S^n)^{\wedge}_p.
 \tag{D.3.6.5}
\]
The assertion concerns spaces, not the inverse limit of
Moore spectra.

Here are the precise pro-space tools. In
\cite[Theorems~3.3, 6.3, 6.7 and~7.2; Proposition~7.3]
{I_Isaksen2005}, a pro-space is fibrant for the
$\FF_p$-cohomology model structure if it is strictly fibrant
and essentially levelwise
$\{K(\FF_p,j):j\geq0\}$-nilpotent. Proposition~3.5 there
identifies this last condition with a finite tower of
principal fibrations with those fibres. The
$\FF_p$-homology structure has fewer acyclic cofibrations,
so an object fibrant for the cohomology structure is
also fibrant for the homology structure. Proposition~7.3
identifies a homological fibrant replacement of a connected
space with the Postnikov-truncated Bousfield--Kan tower.
The cited numbering is that of the opened author preprint.

Apply these facts to the pro-space
\[
                    Y_{r,N}=P_NM_r\qquad(r,N\geq1).
\]
Each of its levels is simply connected and has only
finitely many nonzero homotopy groups, all finite
$p$-groups. Its Postnikov tower is a finite principal
tower with fibres $K(A,j)$ for finite abelian $p$-groups
$A$. Refining each $A$ by a finite filtration with
successive quotient $\FF_p$ refines this to a finite
principal tower with fibres $K(\FF_p,j)$. Thus every level
is nilpotent in the precise sense required above.
A strict fibrant replacement $Y'$ remains essentially
levelwise of this type, and so is fibrant in both of
Isaksen's model structures. Equivalently, choose a
cofinal diagonal tower and its ordinary fibration
replacement; each replacement level is weakly equivalent
to the original finite Postnikov level.

The compatible map $cS^n\to Y$ is an isomorphism on
pro-homology with coefficients in $\FF_p$. Indeed for
each fixed degree, sufficiently high Postnikov truncation
does not change that homology. The Moore space has one
reduced mod-$p$ class in degree $n$ and one in degree
$n+1$. The first has identity transition, while the second
has zero transition because the top cellular degree is $p$.
The second pro-group is therefore zero. This verifies the
homological weak-equivalence hypothesis, including
coefficients in every $\FF_p$-module.

It follows that $Y'$ is a homological fibrant model of the
completion of $S^n$. By Proposition~7.3 and Theorem~7.2
of the cited source, it is weakly equivalent between
fibrant objects to a strict fibrant model of
$\{P_NR_NS^n\}_N$, where $R_NS^n$ is the $N$th
Bousfield--Kan $\FF_p$-tower stage.
Taking the derived inverse limit preserves that
equivalence. To check this model-categorically, the
constant-pro-space functor preserves cofibrations and
ordinary weak equivalences, and hence is left Quillen
from spaces to the homological pro-space structure;
its right adjoint is inverse limit.
The derived limit of the latter tower is
$\holim_NR_NS^n=(S^n)^{\wedge}_p$: removing the
Postnikov truncations does not change the limit, since
in each fixed homotopy degree the truncation is eventually
an isomorphism, also one degree above it for the
$\lim^1$ term.
On the other hand,
\[
 \holim_{r,N}P_NM_r
       \simeq\holim_r\bigl(\holim_NP_NM_r\bigr)
       \simeq\holim_rM_r.
 \]
This proves (D.3.6.5). Combining it with (D.3.6.4) and
closure under target homotopy limits proves (D.3.6.1)
for arbitrary sources.

\medskip\noindent
\textbf{The reverse test on rationally acyclic sources.}
Suppose now that $X$ is rationally acyclic and null for
all $(S^n)^{\wedge}_p$, $n\geq2$. The arithmetic square
for each finite sphere identifies its mapping space
from $X$ to $S^n_{(p)}$ with that to $(S^n)^{\wedge}_p$:
the other two vertices are rational and have contractible
mapping spaces from $X$. This is the same direct
arithmetic-square argument as in C.8 and G.6, with $p$
in place of $2$. Thus $S^n_{(p)}\in\mathcal R_X$ for
all $n\geq2$.

The family of simply connected $p$-local spheres is
closed under smash products and suspension. Theorems~7
and~10 and Corollary~11 of \cite{I_Strom2012} therefore
put in $\mathcal R_X$ every sufficiently suspended
$p$-local finite simply connected complex, and every
finite bouquet of its sufficiently suspended copies:
its finite skeletal cone decomposition has attaching
spaces that are finite wedges of such spheres.
Localization preserves these cofiber, suspension and
wedge comparisons by
\cite[Propositions~6.4.1--6.4.6]{FGHMayPonto2012}.
In particular take finite bouquets of $\Sigma^2P^d(p)$.
The finite cone cubes descend them twice and give
$P^d(p)\in\mathcal R_X$.
This proves the reverse implication under the stated
rational hypothesis.

Finally every cofiber in $R(\loc(\mathbb S/2))$ is
rationally acyclic. On such cofibers G.6 identifies the
completed-sphere test with the localized-sphere test
defining $W_S$. The cofiber characterizations of the
localizers now give (D.3.6.2).
\end{proof}

\subsection*{A countable list for reversible classes generated by polyGEMs}

\begin{claim}{D.3.7}{Proved}{Abelian groups do not retain arbitrary-cardinal distinctions}
\label{D37GEMnormalform}
Put $C_{p^\infty}=\mathbb Z[1/p]/\mathbb Z$.  There is a countable
list of countable connected spaces
\[
 \mathcal E=\{B\mathbb Q\}\ \cup
 \{K(\mathbb Z_{(p)},n),\ K(\mathbb F_p,n),\
                  K(C_{p^\infty},n):p\text{ prime},\ n\geq1\}
 \tag{D.3.7a}
\]
with the following property.  For every abelian group $A$ and $n\geq1$,
the smallest reversible fundamental localizer containing
$K(A,n)\to *$ is generated reversibly by a subset of
\[
 \{B\mathbb Q\}\cup
 \{K(\mathbb Z_{(p)},n),K(\mathbb F_p,n),
                          K(C_{p^\infty},n):p\text{ prime}\}.
 \tag{D.3.7b}
\]
The same conclusion with the universal list $\mathcal E$ holds for
any class of connected polyGEM generators, in particular for any
class of connected GEMs or nilpotent finite Postnikov pieces.
Consequently these localizers
are accessible, each has a countable CW complex as a single reversible
aspherical generator, and there are at most $2^{\aleph_0}$ of them.
These conclusions hold in the fixed small universe and use no large
cardinal hypothesis.

More precisely, if $A\ne0$ is torsion-free and
$J=\{p:pA=A\}$, then
\[
 \min\mathrm{Rev}(K(A,n))
    =\min\mathrm{Rev}(K(\mathbb Z[J^{-1}],n)).
 \tag{D.3.7c}
\]
For a nonzero $p$-primary torsion group $A$, its ordinary nullity,
and hence its reversible closure, is that of $K(C_{p^\infty},n)$
if $A$ is divisible, and that of $K(\mathbb F_p,n)$ otherwise.
The list is a list of sufficient normal forms; no assertion that all
its subsets generate distinct localizers is intended.
\end{claim}

\begin{proof}
We write $\mathcal Q=R(\loc(H\mathbb Q))$ and write
$\operatorname{Rev}(X)$ for the smallest reversible localizer with
$X$ aspherical.  Ordinary $X$-acyclicity implies asphericity for
$\operatorname{Rev}(X)$.  We use ordinary nullity only where explicitly
stated.  Its acyclic objects are closed under retracts, pointed
homotopy colimits and fibre extensions
\citep[beginning of Section~5]{EH5Bousfield1997}; in particular they
are closed under finite products.

\medskip\noindent
\emph{Rational changes of abelian groups.}
First, $\operatorname{Rev}(K(\mathbb Q,n))=\mathcal Q$ for every
$n\geq1$: its suspension spectrum is nonzero and rational, so D.2.5
applies to its reversible closure.  If $A\otimes\mathbb Q\ne0$,
the telescope of multiplication by $2,3,4,\ldots$ on $K(A,n)$ is
$K(A\otimes\mathbb Q,n)$.  The latter has $K(\mathbb Q,n)$ as a
retract.  Thus
\[
 A\otimes\mathbb Q\ne0\quad\Longrightarrow\quad
                \mathcal Q\subseteq\operatorname{Rev}(K(A,n)).
 \tag{D.3.7d}
\]

If $0\to U\to V\to D\to0$ is exact and $D$ is a rational vector
space, then $K(U,n)\to K(V,n)$ is an $H\mathbb F_\ell$-equivalence
for every prime $\ell$.  Indeed, in
\[
 K(U,n)\longrightarrow K(V,n)\longrightarrow K(D,n)
\]
the base is mod-$\ell$ acyclic and its action on the homology of the
fibre is trivial.  For $n=1$ the latter assertion follows from
commutativity of the abelian groups; for $n\geq2$ the base is simply
connected.  The Serre spectral sequence proves the assertion also at
$n=1$, without a claim that a disconnected homotopy fibre is
aspherical.  A rational kernel instead gives a rational fibre and the
same homology conclusion.  A map between connected spaces which is a
mod-$\ell$ homology equivalence for every $\ell$ has a homotopy
cofiber with uniquely divisible integral homology, by the integral
universal coefficient sequences.  Its suspension spectrum is
rational: the fibre of its rationalization is bounded below and
integrally homology-acyclic, and hence zero by the stable Hurewicz
theorem.  Such a map is therefore in $\mathcal Q$.

Consequently kernels or cokernels which are rational vector spaces
may be discarded when comparing reversible aspherical classes that
already contain $\mathcal Q$.  All the comparisons below involving
large products use this fact, not closure under arbitrary products.

\medskip\noindent
\emph{Torsion groups.}
Let $0\ne A$ be $p$-primary.  If $A$ is divisible, the elementary
structure theorem for divisible groups expresses it as a nonempty
direct sum of copies of $C_{p^\infty}$.  Direct-summand retracts and
filtered colimits of finite products give equality of ordinary
nullity with $K(C_{p^\infty},n)$.

If $A$ is not divisible, then $A/pA\ne0$, since all other primes
already act invertibly on $A$.  Thus
\[
 [K(A,n),K(\mathbb F_p,n)]_*
   =H^n(K(A,n);\mathbb F_p)
   =\operatorname{Hom}(A,\mathbb F_p)\ne0.
\]
The exact dichotomy in
\citet[Lemma~9.13]{EH5Bousfield1997} says that
$K(\mathbb F_p,n)$ is either local or acyclic for any localization.
It is not $K(A,n)$-null, so it is $K(A,n)$-acyclic.  Conversely,
$A$ is the filtered union of its finite subgroups.  A finite abelian
$p$-group has a finite filtration with quotients $\mathbb F_p$.
The associated fibre extensions of Eilenberg--Mac Lane spaces,
followed by the filtered homotopy colimit, make $K(A,n)$
$K(\mathbb F_p,n)$-acyclic.  This proof works for every $n\geq1$.

For a general torsion group $A=\bigoplus_p A_p$, each $K(A_p,n)$
is a retract of $K(A,n)$, and the latter is the filtered homotopy
colimit of the finite products of these factors.  This proves the
torsion part of the claim using only the two stated types at each
prime.

\medskip\noindent
\emph{Two module facts, including the required realization.}
We record the algebraic and homotopical details needed for a mixed
group.  Every $\mathbb Z_p$-module $M$, viewed as an abstract abelian
group, satisfies
\[
              K(M,n)\text{ is }K(\mathbb Z_p,n)\text{-acyclic}.
 \tag{D.3.7e}
\]
The discrete valuation ring $\mathbb Z_p$ is a principal ideal
domain, so $M$ has a free resolution
$0\to F_1\to F_0\to M\to0$ of length one.  Each $K(F_i,n)$ is
a filtered homotopy colimit of finite products of
$K(\mathbb Z_p,n)$.  In simplicial abelian group models, realize the
bar construction for the additive action of $K(F_1,n)$ on
$K(F_0,n)$.  Its terms are
$K(F_0,n)\times K(F_1,n)^j$ and are all
$K(\mathbb Z_p,n)$-acyclic.  Its realization is $K(M,n)$: the
homotopy quotient has fibre sequence
\[
 K(F_0,n)\longrightarrow
 K(F_0,n)\mathbin{\!//\!}K(F_1,n)
                  \longrightarrow K(F_1,n+1),
\]
and the homotopy exact sequence, with injective connecting map
$F_1\to F_0$, leaves just $M$ in degree $n$.  Homotopy-colimit
closure proves (D.3.7e).  This is a genuine bar quotient; it does not
use a two-out-of-three rule for arbitrary homotopy fibres.

Next let $F$ be a torsion-free abelian group, uniquely divisible by
every prime other than $p$.  Its $p$-adic completion
\[
                 \widehat F_p=\lim_k F/p^kF
\]
is a torsion-free, separated and complete $\mathbb Z_p$-module, and
$\widehat F_p/p\widehat F_p=F/pF$.  For completeness, torsion-freeness
of $F$ makes division by $p$ unique modulo $p^k$ whenever an element
vanishes modulo $p$.  Applied to compatible residue classes, this
identifies the kernel of
$\widehat F_p\to F/pF$ with $p\widehat F_p$ and shows that
$p$-torsion in the limit is zero.  The same argument with $p^k$
gives all the displayed completeness assertions.

Every torsion-free separated complete $\mathbb Z_p$-module $M$ is
the completion of a free $\mathbb Z_p$-module.  Choose lifts of a
basis of $M/pM$ and map the completion of the free module on this
basis to $M$.  Successive finite expressions modulo $p$ give
surjectivity by $p$-adic convergence.  An element of the kernel is
divisible by $p$ in the source; torsion-freeness of the target puts
its quotient back in the kernel.  Iteration and separatedness make
the kernel zero.  If $M/pM\ne0$, a chosen coordinate of this
completed free module exhibits $\mathbb Z_p$ as a retract of $M$.
The completed free module need not be free as an abstract module.
Only the preceding length-one resolution uses abstract freeness.

\medskip\noindent
\emph{The completion groups of a reduced abelian group.}
Split $A=D\oplus A_0$, with $D$ its maximal divisible subgroup and
$A_0$ reduced.  Divisible groups are injective, and
\[
 D\cong\mathbb Q^{(\kappa_0)}
              \oplus\bigoplus_p C_{p^\infty}^{(\kappa_p)}.
\]
Its contribution is already on the stated list.  We may suppose
$A\otimes\mathbb Q\ne0$, since torsion groups were dealt with
above, and hence retain $\mathcal Q$ among the reversible classes
being compared.  Put
\[
                   C_p=\operatorname{Ext}^1_{\mathbb Z}
                                  (C_{p^\infty},A_0).
\]
Applying $\operatorname{Hom}(-,A_0)$ to
$0\to\mathbb Z\to\mathbb Q\to\mathbb Q/\mathbb Z\to0$
gives
\[
 0\longrightarrow A_0\longrightarrow\prod_p C_p
       \longrightarrow\operatorname{Ext}^1_{\mathbb Z}
                                 (\mathbb Q,A_0)\longrightarrow0.
 \tag{D.3.7f}
\]
Here $\operatorname{Hom}(\mathbb Q,A_0)=0$ because the image of
a divisible group is divisible, and the final Ext group is a
rational vector space by the action of $\mathbb Q$ on its first
argument.  Thus (D.3.7f) identifies $K(A_0,n)$ and
$K(\prod_p C_p,n)$ modulo $\mathcal Q$.

Each $C_p$ is a $\mathbb Z_p$-module by the endomorphism action on
$C_{p^\infty}$.  In particular, multiplication by any prime
$q\ne p$ is invertible.  Also
\[
       \operatorname{Hom}(\mathbb Q,C_p)=0,
       \qquad \operatorname{Ext}^1(\mathbb Q,C_p)=0.
 \tag{D.3.7g}
\]
Here is a direct justification rather than an assumption of ordinary
separated completeness.  Since $A_0$ is reduced,
$\operatorname{Hom}(C_{p^\infty},A_0)=0$.  Since $\mathbb Z$ has
global dimension one, the derived group
$\mathbf R\operatorname{Hom}(C_{p^\infty},A_0)$ has its only
cohomology $C_p$ in degree one.  Derived tensor--Hom adjunction gives
\[
 \mathbf R\operatorname{Hom}\bigl(\mathbb Q,
       \mathbf R\operatorname{Hom}(C_{p^\infty},A_0)\bigr)
  \simeq
 \mathbf R\operatorname{Hom}
       (\mathbb Q\otimes^{\mathbf L}C_{p^\infty},A_0)=0,
\]
which proves (D.3.7g).  This is the Ext-completeness condition of
\citet[Section~5.6]{EH5Bousfield1997}.  It does not assert that the
possibly torsion group $C_p$ is separated for its ordinary $p$-adic
topology.

The quotient
\[
                     (\prod_p C_p)/(\bigoplus_p C_p)
 \tag{D.3.7h}
\]
is a rational vector space.  Indeed, for each prime $\ell$, the
kernel and cokernel of multiplication by $\ell$ on either the sum
or the product are exactly those on its $\ell$-coordinate.  The
maps on these kernels and cokernels induced by the inclusion are
isomorphisms.  The snake lemma makes multiplication by $\ell$ an
isomorphism on (D.3.7h).  Hence modulo $\mathcal Q$ the contribution
of the product is the join of the individual $K(C_p,n)$
contributions: the sum is a filtered homotopy colimit of finite
products, and each coordinate is a retract of the product.

\medskip\noindent
\emph{Reduction at one prime.}
Fix $p$, set $M=C_p$, let $T$ be its torsion subgroup, and put
$F=M/T$.  The torsion in $M$ is $p$-primary.  Both $M$ and $F$ are
uniquely divisible by all other primes, and $F$ is torsion-free.
The long Ext sequence and (D.3.7g) give
$\operatorname{Ext}^1(\mathbb Q,F)=0$, since
$\operatorname{Ext}^2_{\mathbb Z}=0$.

If $F/pF=0$, then $F$ is a rational vector space.  Thus $K(T,n)$
and $K(M,n)$ agree modulo $\mathcal Q$.  The group $T$ is reduced:
a nonzero divisible subgroup would contain a copy of
$C_{p^\infty}$ and receive a nonzero map from $\mathbb Q$,
contradicting (D.3.7g).  When $T\ne0$ its ordinary nullity is
therefore that of $K(\mathbb F_p,n)$ by the torsion argument.
When $T=0$ the group $M=F$ is rational and (D.3.7g) makes $M=0$.

Suppose instead that $F/pF\ne0$.  The same exact sequence used in
(D.3.7f), now applied to the torsion-free group $F$, gives
\[
 0\longrightarrow D_F\longrightarrow F
                    \longrightarrow\widehat F_p\longrightarrow0,
 \tag{D.3.7i}
\]
where $D_F$ is a rational vector space.  The other prime completions
vanish, and the cokernel is
$\operatorname{Ext}^1(\mathbb Q,F)=0$.  We used here the precise
completion formula
\[
 \operatorname{Ext}^1(C_{p^\infty},F)=\lim_k F/p^kF
\]
for torsion-free $F$; more generally it holds for bounded
$p$-torsion, as stated in
\citet[Section~4, p.~41]{Bousfield1996ULP}.  The subgroup $D_F$
is injective, so (D.3.7i) splits.  The completed-free-module
argument above supplies a surjection
$\widehat F_p\to\mathbb Z_p$.  Its composite gives a surjection
$h:M\to\mathbb Z_p$.

Every additive map $h:M\to\mathbb Z_p$ is
$\mathbb Z_p$-linear.  For $r\in\mathbb Z_p$, choose integers
$a_k\equiv r\pmod {p^k}$; then
$h(rx)-a_kh(x)\in p^k\mathbb Z_p$ for all $k$.
Separatedness of the target gives $h(rx)=rh(x)$.  Choosing
$x\in M$ with $h(x)=1$ therefore splits $h$ by $r\mapsto rx$.
Thus $K(\mathbb Z_p,n)$ is a retract of $K(M,n)$, while (D.3.7e)
gives the reverse ordinary acyclicity.  In this case $K(M,n)$ and
$K(\mathbb Z_p,n)$ have the same ordinary nullity.

Finally $\mathbb Z_p/\mathbb Z_{(p)}$ is a rational vector space.
Multiplication by $p$ is injective on the quotient because
$\mathbb Q\cap\mathbb Z_p=\mathbb Z_{(p)}$, and is surjective by
subtracting an integer with the prescribed residue modulo $p$;
every other prime is already invertible on both groups.  Consequently
$K(\mathbb Z_p,n)$ can be replaced by the countable
$K(\mathbb Z_{(p)},n)$ in the presence of $\mathcal Q$.
This completes the reduction of every mixed abelian group to
(D.3.7b).

\medskip\noindent
\emph{The torsion-free formula, arbitrary GEMs and polyGEMs.}
For a torsion-free group $A$, the argument before reduction at one
prime has
$C_p=\widehat A_p$; it is zero exactly when $A/pA=0$.
Every nonzero $C_p$ is torsion-free and has a $\mathbb Z_p$
retract, so its contribution is precisely the
$K(\mathbb Z_{(p)},n)$ contribution modulo $\mathcal Q$.
This profile is identical for $A$ and
$\mathbb Z[J^{-1}]$, proving (D.3.7c).

A connected GEM $G=\prod_{n\geq1}K(A_n,n)$ is weakly equivalent
to the homotopy colimit of its finite partial products: at any fixed
homotopy degree all sufficiently late new factors are more highly
connected.  Each factor is a retract.  Thus its reversible closure
is the join of the closures of its Eilenberg--Mac Lane factors.
For any class of such GEM generators, apply (D.3.7b) to all these
factors and take the union of the resulting subsets of the single
countable list $\mathcal E$.  The wedge of that subset is a
countable CW complex: each listed abelian group has a countable
simplicial Eilenberg--Mac Lane model.  Its reversible closure is
the required localizer, accessible by B.4.  There are only
$2^{\aleph_0}$ subsets of the list, proving the cardinal bound.

The extension to polyGEMs uses a precise ordinary-acyclicity theorem,
not closure under inverse limits.  In the terminology of
\citet{FGHCFFS2015}, a GEM is a $1$-polyGEM, and an $s$-polyGEM
is a homotopy retract of the fibre of a map from an
$(s-1)$-polyGEM to a GEM.  Their Proposition~6.2(4) identifies the ordinary
acyclic classes generated by a polyGEM $X$ and by its free abelian
group construction $\mathbb ZX$.  For connected $X$, Dold--Thom
identifies $\mathbb ZX$ with the GEM
$\prod_{n\geq1}K(\widetilde H_n(X;\mathbb Z),n)$.
Mutual ordinary acyclicity gives equality of their reversible
closures, so the preceding countable-list reduction applies to
every class of connected polyGEMs.  Finally,
\citep[Proposition~5.2]{FGHCFFS2015} makes every connected
nilpotent finite Postnikov piece a polyGEM.  These two cited
statements impose no cardinality or finite-generation bound.
\end{proof}

The ordinary nullity classes obtained from rigid abelian groups in
\citet[Theorems~8.2 and~8.5]{FGHStanley2006} therefore cannot supply
arbitrarily many reversible classes.  Their unsuspended
$K(A,n)$ separators do not survive this reduction.  D.3.7 proves a
set-size statement for all classes generated by connected polyGEMs,
including their intersection with any specified stable-part fibre.
It does not assert that an arbitrary reversible localizer is
generated by polyGEMs, and does not settle the general set-size question
in D.3.2.

\medskip\noindent
\emph{Consequence for the projective-space interval.}
Let $p$ be any prime.  A reversible localizer generated by connected
infinite-loop spaces and satisfying
\[
 \operatorname{Rev}(B\mathbb Z/p)\ \subseteq\ W
          \ \subseteq\ R(\loc(\mathbb S/p))
\]
must equal $\operatorname{Rev}(B\mathbb Z/p)$.  Indeed, each
aspherical generator $X$ has trivial reduced
$\mathbb Z[1/p]$-homology.  A connected infinite-loop space is
nilpotent, so the nilpotent homological localization criterion
\citep[Theorems~6.1.1--6.1.2]{FGHMayPonto2012} gives
$\pi_i(X)[1/p]=0$ for every $i\geq1$.  Choose its connected
connective spectrum $E$, so $X\simeq\Omega^\infty E$ and
$E[1/p]=0$.  By \citet[Theorem~3.2]{EH5Levy2025}, $X$ is already
ordinarily $B\mathbb Z/p$-acyclic.  Thus every generator belongs
to the lower endpoint and $W$ equals that endpoint.  The same
conclusion holds for connected polyGEM generators, by replacing
each by the ordinarily equivalent GEM just used in the proof.
Hence such generators cannot produce a strict intermediate in
the interval of D.3.2.

\begin{claim}{D.3.8}{Proved}{A countable bound for primary infinite-loop presentations of a fixed stable part}
Fix a prime $p$. Let $\{X_i\}_{i\in I}$ be a set or class of
connected infinite-loop spaces with $p$-primary torsion homotopy
groups, at least one of which is noncontractible. Let $n\geq1$
be the least degree in which some $\pi_nX_i$ is nonzero, and put
\[
 \mathcal L=\loc\{\Sigma^\infty X_i:i\in I\},\qquad
 W=\min\mathrm{Rev}\{X_i\to *:i\in I\}.
\]
Then $T(W)=\mathcal L$. If some $\pi_nX_i$ is not divisible,
equivalently some $H_n(X_i;\mathbb F_p)\ne0$, then
\[
 W=\min\mathrm{Rev}
   \{\Omega^\infty E\to *:
                         E\in\mathcal L\cap\Ho(\Sp)_{\geq n}\}.
 \tag{D.3.8a}
\]
If every $\pi_nX_i$ is divisible, then
\[
 W=\min\mathrm{Rev}\left(
   \{\Omega^\infty E\to *:
                   E\in\mathcal L\cap\Ho(\Sp)_{\geq n+1}\}
       \cup\{K(C_{p^\infty},n)\to *\}\right).
 \tag{D.3.8b}
\]
Thus $W$ is determined by $\mathcal L$, the least degree $n$,
and one divisibility flag. For fixed $p$ and $\mathcal L$ there
are at most countably many reversible localizers admitting such
a presentation. If $\mathcal L$ is generated by a small set of
spectra, each is accessible even when the original family $I$
is a class. A family of contractible generators gives the
smallest reversible localizer $W_{\mathbb Z}$ separately.
\end{claim}

\begin{proof}
Choose for each $i$ a connective spectrum $E_i$ giving its indicated
infinite-loop structure, with $X_i\simeq\Omega^\infty E_i$.
Its connective choice has $E_i\in\Ho(\Sp)_{\geq n}$.
The homotopy groups are $p$-primary, so $E_i[1/p]=0$.
The corresponding reduced homology of $X_i$ is $p$-primary as
well, either by the nilpotent localization criterion or by its
Postnikov tower degree by degree.

The exact space--spectrum interface is
\citet[Theorem~2.10 and Proposition~3.1]{Bousfield1996ULP}:
for a space $A$ and a connective spectrum $E$,
\[
 P_A\Omega^\infty E\simeq
          \Omega^\infty P_{\Sigma^\infty A}^{\Sp}E,
 \tag{D.3.8c}
\]
and $A$ kills $\Omega^\infty E$ in ordinary nullification if and
only if $\Sigma^\infty A$ kills $E$ in spectral nullification.
In particular, since $X_i$ kills itself, $E_i$ belongs to the
spectral nullification kernel generated by $\Sigma^\infty X_i$,
and hence to $\mathcal L$. Each $X_i$ is therefore among the
generators on the right of (D.3.8a). The equality
$T(W)=\mathcal L$ also follows at once: $R(\mathcal L)$ contains
every specified aspherical generator, so
$W\subseteq R(\mathcal L)$ and $T(W)\subseteq\mathcal L$;
the suspension spectra of the original generators give the
reverse inclusion.

\medskip\noindent
\emph{The nondivisible case.}
First suppose $I$ is small and let $A=\bigvee_{i\in I}X_i$.
Then $C=\Sigma^\infty A$ is $n$-connective, has $p$-primary
integral homology, and has $H_n(C;\mathbb F_p)\ne0$.
Also $\loc(C)=\mathcal L$.
The spectral conclusion of D.2.9 gives
$P_C^{\Sp}E=0$ for every
$E\in\mathcal L\cap\Ho(\Sp)_{\geq n}$.
Equation (D.3.8c) then makes $\Omega^\infty E$ ordinarily
$A$-acyclic, and consequently $W$-aspherical. This proves
(D.3.8a). The use of ordinary nullification here is justified
by (D.3.8c); it is not an inference from stable asphericity alone.

\medskip\noindent
\emph{The divisible case and the larger GEM-error bound.}
Still suppose $I$ small, and let $A$ and $C$ have the same meaning.
Now $H_n(A;\mathbb Z)=\bigoplus_i\pi_nX_i$ is a nonzero
divisible $p$-primary group. By
\citet[Corollary~3.2]{Bousfield1996ULP}, $A$ kills $SP^\infty A$.
The latter has $K(H_n(A;\mathbb Z),n)$ as a retract, and that
space has $K(C_{p^\infty},n)$ as a retract. Thus this last
space is ordinarily $A$-acyclic. The nonsplit extension
\[
 0\longrightarrow\mathbb F_p\longrightarrow C_{p^\infty}
       \xrightarrow{p}C_{p^\infty}\longrightarrow0
\]
gives a nonzero class in
$H^{n+1}(K(C_{p^\infty},n);\mathbb F_p)$ by the universal
coefficient theorem. Lemma~9.13 of
\citet{EH5Bousfield1997} therefore makes
$K(\mathbb F_p,n+1)$ ordinarily
$K(C_{p^\infty},n)$-acyclic, hence $A$-acyclic.
The precise comparison (D.3.8c), or Proposition~3.1 of its source,
shows that $C$ spectrally kills $\Sigma^{n+1}H\mathbb F_p$.

We claim that
\[
 P_C^{\Sp}E=0
       \quad\text{for every }E\in\mathcal L\cap\Ho(\Sp)_{\geq n+1}.
 \tag{D.3.8d}
\]
Set $C_i=\Sigma^{-i}C$ and $Y_i=P_{C_i}^{\Sp}E$.
The general part of
\citet[Theorem~2.6]{Bousfield1996ULP}, without a support assumption,
says that $F_i=\fib(Y_i\to Y_{i+1})$ is a stable GEM.
Both endpoints are $C_i$-null, so $F_i$ is $C_i$-null.
All $Y_i$ belong to $\mathcal L$, because their nullification
cones are built from shifts of $C$. In particular $F_i$ is
$p$-primary: $\mathcal L$ is contained in the kernel of
inverting $p$.

Since $C_i$ kills $\Sigma^{n+1-i}H\mathbb F_p$, $F_i$ is
null for this latter spectrum. If the stable GEM $F_i$ had a
nonzero homotopy group $G$ in degree $k\geq n+1-i$, its retract
$\Sigma^kHG$ would receive the nonzero group
$\operatorname{Hom}(\mathbb F_p,G)$ in mapping-space degree
$k-(n+1-i)\geq0$. Every nonzero $p$-primary group has an
element of order $p$, so this is impossible. Therefore
\[
 F_i\in\Ho(\Sp)_{\leq n-i},\qquad
 \fib(Y_i\to Y_\infty)\in\Ho(\Sp)_{\leq n-i},
       \quad Y_\infty=\colim_iY_i.
 \tag{D.3.8e}
\]
The second bound follows by extensions and exact filtered
colimits, exactly as in D.2.9.

For a fixed $i$ compare $Y_{i+2}$ with $Y_\infty$.
Their fibre lies in degree at most $n-i-2$, and its suspension
in degree at most $n-i-1$. Since $C_i$ is $(n-i)$-connective,
$t$-structure orthogonality gives
\[
 \Map_{\Sp}(C_i,Y_{i+2})
                 \simeq\Map_{\Sp}(C_i,Y_\infty).
\]
The left side is contractible, since $Y_{i+2}$ is
$C_{i+2}$-null and $C_i=\Sigma^2C_{i+2}$.
Thus $Y_\infty$ is null for every shift of $C$, while it belongs
to $\mathcal L=\loc(C)$, so $Y_\infty=0$.
Equation (D.3.8e) yields $P_C^{\Sp}E\in\Ho(\Sp)_{\leq n}$.
On the other hand nullification by $n$-connective $C$ preserves
$(n+1)$-connectivity of an input which already has it:
the cofiber of any map $\Sigma^kC\to Z$, with $k\geq0$ and
$Z\in\Ho(\Sp)_{\geq n+1}$, still lies in
$\Ho(\Sp)_{\geq n+1}$ by its homotopy exact sequence.
The cone-attachment construction and its transfinite colimits
preserve this bound. Hence $P_C^{\Sp}E=0$, proving (D.3.8d).
The two-step comparison with $Y_{i+2}$, rather than $Y_{i+1}$,
is what controls components with this enlarged error bound.

By (D.3.8c), every higher infinite-loop generator on the right
of (D.3.8b) is ordinarily $A$-acyclic; the Prüfer-space
generator was already shown to be so. This gives one inclusion
in (D.3.8b). For the other write $D_i=\pi_nE_i$.
Each $D_i$ is a direct sum of copies of $C_{p^\infty}$,
possibly zero. Since $K(C_{p^\infty},n)$ is $W$-aspherical,
its suspension spectrum belongs to $\mathcal L$.
Its ordinary acyclic class contains $K(D_i,n)$, and
Proposition~3.1 then puts $\Sigma^nHD_i$ in $\mathcal L$ as well.
Consequently $E_i\in\mathcal L$ implies
$\tau_{\geq n+1}E_i\in\mathcal L$.
In the Postnikov fibre sequence
\[
 \Omega^\infty\tau_{\geq n+1}E_i
          \longrightarrow X_i\longrightarrow K(D_i,n)
\]
both the fibre and the base are aspherical for the right side of
(D.3.8b). The ordinary acyclic class generated by their wedge
is closed under fibre extensions and is contained in that
reversible class. Thus $X_i$ is aspherical there. This uses
fibre and base to obtain the total space; it makes no inference
of fibre acyclicity from the other two terms.

\medskip\noindent
\emph{Class families and the cardinal bound.}
For a class $I$, every individual member of
$\loc\{\Sigma^\infty X_i:i\in I\}$ lies in the localizing
subcategory generated by a small subfamily. Indeed the union
of those subcategories over all small subfamilies is itself
closed under shifts, cofibers and small coproducts: a small
collection of small subfamilies is contained in their small
union. This union contains all the generators and therefore
is $\mathcal L$. Given an individual $E$ in the required
connective range, choose such a small subfamily and adjoin one
source $X_j$ witnessing the least degree and its indicated
nondivisibility or divisibility case. The corresponding
preceding argument applied to its wedge makes
$\Omega^\infty E$ ordinarily acyclic for that wedge, hence
$W$-aspherical. The Prüfer generator and the reverse Postnikov
argument in the divisible case use the same witness. Both
formulas therefore hold in the class case. There are just two
displayed possibilities at each positive integer $n$, proving
the countable bound for fixed $p$ and $\mathcal L$.

Finally suppose $\mathcal L=\loc\{F_a:a\in J\}$ for a small
set $J$. Apply the small-subfamily observation to each $F_a$
and take the union of the resulting source subfamilies,
together with the indicated witness $X_j$. This is a small
subfamily with the same stable part, least degree and
divisibility flag. Its reversible closure equals $W$ by the
appropriate formula, and B.4 makes it accessible.
\end{proof}

This fixes the possible number of primary infinite-loop
presentations of a given stable part. It neither determines
which of the countably many displayed localizers coincide nor
decides the adjacent strictness test in D.2.9. In particular,
it does not replace the nonlinear suspension-nullification
test needed for the zero-finite-part separator.

